\section{Scalar three-cocycles and fusion-tensor gauge freedom}
\label{sec:mpug_scalar_gauge}

Let $G$ be a group and write $\C^\times$ for the multiplicative group of
nonzero complex numbers. The fusion tensors have the scalar gauge freedom
\begin{align}
    F^>_{g,h} & \longmapsto \beta_{g,h}F^>_{g,h},      \\
    F^<_{g,h} & \longmapsto \beta_{g,h}^{-1}F^<_{g,h}.
    \label{eq:mpug_fusion_gauge}
\end{align}
% Source anchor: arXiv:2502.20257, eq:scalar_fus_ten.
This specifies the fusion-tensor gauge freedom of \cite{FrancoRubio2025MPUGauging}
for the scalar functions $\beta$, $\omega$, $\gamma$, $L$, $\sigma$, and $\Xi$,
with representation-level data as in
Definition~\ref{def:mpug_group_representation}.

\begin{lemma}[Scalar factors of an identity Kronecker product]
    \label{thm:mpug_kronecker_identity_factors}
    \lean{Matrix.exists_eq_smul_one_of_kronecker_eq_one}
    \leanok
    Let $M$ and $N$ be nonempty sets with decidable equality, and let
    $A\in\operatorname{Mat}_M(\C)$ and $B\in\operatorname{Mat}_N(\C)$. If
    $A\otimes B=\mathbf 1$, then there is a nonzero $c\in\C$ such that
    \begin{align}
        A & =c\mathbf 1, & B & =c^{-1}\mathbf 1.
        \label{eq:mpug_kronecker_identity_factors}
    \end{align}
\end{lemma}

\begin{proof}\leanok
    For any $i_0\in M$ and $j_0\in N$, the
    $((i_0,j_0),(i_0,j_0))$ entry gives
    $A_{i_0i_0}B_{j_0j_0}=1$. Thus $c=A_{i_0i_0}$ is nonzero and
    $B_{j_0j_0}=c^{-1}$. Comparing every
    $((i,j_0),(i',j_0))$ entry gives $A_{ii'}=c\delta_{ii'}$; comparing every
    $((i_0,j),(i_0,j'))$ entry gives $B_{jj'}=c^{-1}\delta_{jj'}$.
\end{proof}

\begin{lemma}[Unitarity from an adjoint Kronecker product]
    \label{thm:matrix_unitary_adjoint_kronecker}
    \lean{Matrix.isUnitaryBetween_of_conjTranspose_kronecker}
    \leanok
    Let $M,N$ be nonempty finite sets, and let
    $K\colon\C^N\to\C^M$ be a complex matrix. If
    $K^\dagger\otimes K\colon\C^{M\times N}\to\C^{N\times M}$
    is unitary, then $K$ is unitary between its coordinate spaces:
    \begin{align}
        K^\dagger K & =I_N, \\
        KK^\dagger  & =I_M.
    \end{align}
    In the matrix implication used in
    \cite[lines~5444--5487]{FrancoRubio2025MPUGauging}, no equality of
    $|M|$ and $|N|$ is assumed; the tensor-network identity to which it is
    applied is a separate result.
\end{lemma}

\begin{proof}\leanok
    \uses{thm:mpug_kronecker_identity_factors}
    The isometry identity for $K^\dagger\otimes K$ gives
    \begin{align}
        (KK^\dagger)\otimes(K^\dagger K) & =I_{M\times N}.
    \end{align}
    By Lemma~\ref{thm:mpug_kronecker_identity_factors}, there is a
    nonzero $c\in\C$ with $KK^\dagger=cI_M$ and
    $K^\dagger K=c^{-1}I_N$. Consequently,
    \begin{align}
        (KK^\dagger)^2
         & =K(K^\dagger K)K^\dagger
        =c^{-1}KK^\dagger=I_M.
    \end{align}
    For any $i\in M$, the $(i,i)$ entry gives $c^2=1$, while positivity
    of the Gram matrix gives $c=(KK^\dagger)_{ii}\geq0$.
    Thus $c=1$, and the two scalar Gram identities become
    $K^\dagger K=I_N$ and $KK^\dagger=I_M$.
\end{proof}
\begin{lemma}[Unitary factors of a unitary Kronecker product]
    \label{thm:mpug_unitary_kronecker_factors}
    \lean{Matrix.exists_pos_real_rescaling_mem_unitaryGroup_of_kronecker_mem_unitaryGroup,
        Matrix.exists_pos_real_smul_unitary_factors_of_kronecker_mem_unitaryGroup}
    \leanok
    Let $M$ and $N$ be nonempty finite sets. If
    $K_1\in\operatorname{Mat}_M(\C)$ and
    $K_2\in\operatorname{Mat}_N(\C)$ satisfy that $K_1\otimes K_2$ is
    unitary, then there are a positive real number $\delta$ and unitary
    matrices $W_1,W_2$ such that
    \begin{align}
        K_1 & =\delta W_1, & K_2 & =\delta^{-1}W_2.
        \label{eq:mpug_unitary_kronecker_factors}
    \end{align}
    Equivalently, $\delta^{-1}K_1$ and $\delta K_2$ are unitary.
    % Source anchor: arXiv:2502.20257, Lemma lem:deco, lines 1052--1066.
    This is only the tensor-factor step in the decomposition result
    of~\cite{FrancoRubio2025MPUGauging}; it does not construct or compare the
    three-legged decompositions occurring there.
\end{lemma}

\begin{proof}\leanok
    \uses{thm:mpug_kronecker_identity_factors}
    Unitarity gives
    \begin{align}
        (K_1^\dagger K_1)\otimes(K_2^\dagger K_2)=\mathbf 1.
    \end{align}
    Lemma~\ref{thm:mpug_kronecker_identity_factors} gives a nonzero
    $c\in\C$ such that
    $K_1^\dagger K_1=c\mathbf 1$ and
    $K_2^\dagger K_2=c^{-1}\mathbf 1$. Positivity of the first Gram matrix
    makes $c$ a positive real number. Taking $\delta=\sqrt c$ shows directly
    that $\delta^{-1}K_1$ and $\delta K_2$ have identity Gram matrices; these
    are $W_1$ and $W_2$ in
    (\ref{eq:mpug_unitary_kronecker_factors}).
\end{proof}

\begin{lemma}[Cross-product rigidity]
    \label{thm:mpug_cross_product_rigidity}
    \lean{Kraus.IsNormal.eq_smul_of_cross_mul_eq,
        Kraus.IsInjective.eq_smul_of_cross_mul_eq}
    \leanok
    \uses{def:injective, def:normal}
    Let $A_i$ and $B_i$ be two families of square matrices of the same size,
    indexed by the same finite alphabet, and suppose that, for every $i,j$,
    \begin{align}
        B_iA_j & =A_iB_j.
        \label{eq:mpug_cross_product_rigidity}
    \end{align}
    If $A$ is injective, or more generally normal, then there is a scalar
    $c\in\C$ with $B_i=cA_i$ for every $i$;
    % Source anchor: arXiv:2502.20257, prop:MPUsplus, lines 1255--1325.
    in~\cite[lines~1255--1325]{FrancoRubio2025MPUGauging}, this
    commuting-tensor step shows that the tensor commuting with a normal
    matrix product unitary tensor is a multiple of it.
\end{lemma}

\begin{proof}\leanok
    \uses{def:decomposition_map, thm:normal_tensor_gram_rigidity}
    Let $N\geq1$ be a blocking length at which the length-$N$ words of $A$
    span the full matrix algebra, and decompose the identity as
    $\mathbf 1=\sum_\sigma c_\sigma A^{\sigma_1}\cdots A^{\sigma_N}$. The
    cross relation moves the marked letter successively through the word,
    \begin{align}
        B_iA^{\sigma_1}\cdots A^{\sigma_N}
         & =A_iA^{\sigma_1}\cdots A^{\sigma_{N-1}}B^{\sigma_N}.
        \notag
    \end{align}
    When $N=1$, the middle word on the right is absent. Thus $B_i=A_iC$ with
    $C=\sum_\sigma c_\sigma
        A^{\sigma_1}\cdots A^{\sigma_{N-1}}B^{\sigma_N}$.
    Substituting this identity into
    (\ref{eq:mpug_cross_product_rigidity}) gives
    $A_i(CA_j-A_jC)=0$ for all $i$ and $j$. Multiplying on the left by a
    length-$N$ word, whose last letter absorbs one of the $A_i$, and
    recombining the words into the identity gives $CA_j=A_jC$. A matrix
    commuting with every letter of a normal family is a multiple of the
    identity, so $C=c\mathbf 1$ and $B_i=cA_i$. The injective case is the
    case of blocking length one.
\end{proof}

\begin{definition}[Canonical boundaries under physical adjunction]
    \label{def:mpug_adjoint_canonical_boundaries}
    \lean{MPOTensor.IsMPUCanonicalFormII.physicalAdjointTensor}
    \leanok
    \uses{def:mpu_canonical_form_ii, def:mpdo_physical_adjoint}
    Let $\mathcal U$ be an MPU with a canonical-form-II presentation and
    recorded positive diagonal fixed matrix $\rho$, with
    $\operatorname{tr}\rho=1$. Its physical adjoint has letters
    $(\mathcal U^\dagger)^{ij}=\overline{\mathcal U^{ji}}$; the virtual
    indices are not transposed. Give its normalized flattening the
    conjugated canonical blocks with the physical indices exchanged, and
    retain the same full support and the same fixed matrix $\rho$.
    As in the reflected boundaries of
    \cite[lines~1164--1254]{FrancoRubio2025MPUGauging}, this defines
    canonical-form-II data for $\mathcal U^\dagger$ with the same left and
    right boundaries $I,\rho$.
    % This construction does not assert preservation of simplicity.
\end{definition}

\begin{proof}\leanok
    \uses{thm:mpu_physical_adjoint_identities,
        def:mpu_physical_adjoint_cfii_data,
        thm:mpu_physical_adjoint_normalized_flattening,
        thm:cpsv_map_star_transfer_fixed_point}
    Physical adjunction preserves periodic unitarity and conjugates the
    canonical blocks without changing their dimensions. Writing
    $\mathcal E_{\mathcal U}$ for the transfer map of $\mathcal U$, the
    conjugated fixed-point equation is
    \begin{align}
        \mathcal E_{\mathcal U^\dagger}(\rho)
         & =\overline{\mathcal E_{\mathcal U}(\overline\rho)}
        =\overline{\mathcal E_{\mathcal U}(\rho)}
        =\overline\rho=\rho,
    \end{align}
    using $\overline\rho=\rho$ and $\mathcal E_{\mathcal U}(\rho)=\rho$. The transfer map is invariant under the physical-index permutation, so the
    conjugated data share the same fixed matrix.
\end{proof}

\begin{lemma}[Canonical boundary contractions]
    \label{thm:mpug_simple_canonical_boundary_contractions}
    \lean{MPOTensor.IsMPUCanonicalFormII.simple2_recorded_fixed_pair,
        MPOTensor.IsMPUCanonicalFormII.doubleLayer_boundary_contractions}
    \leanok
    \discussion{7317}
    \uses{def:mpu_canonical_form_ii, def:mpu_simple, def:mpu_double_layer}
    Let $\mathcal U$ be a simple MPU with a canonical-form-II presentation
    and recorded fixed matrix $\rho$. Let $W$ be its unnormalized double layer,
    put $(\Phi|_{(a,b)}=\delta_{a,b}$ and
    $|\rho)_{(a,b)}=\rho_{ba}$. For all physical indices $i,j,k,l$, the
    boundary removals used in the six-tensor argument of
    \cite[lines~1164--1254]{FrancoRubio2025MPUGauging} hold:
    \begin{align}
        W^{ij}W^{kl}       & =W^{ij}|\rho)(\Phi|W^{kl},
        \label{eq:mpug_simple_canonical_split}          \\
        (\Phi|W^{ij}W^{kl} & =\delta_{ij}(\Phi|W^{kl},
        \label{eq:mpug_simple_canonical_bdry_left}      \\
        W^{ij}W^{kl}|\rho) & =\delta_{kl}W^{ij}|\rho).
        \label{eq:mpug_simple_canonical_bdry_right}
    \end{align}
\end{lemma}

\begin{proof}\leanok
    \uses{thm:mpu_canonical_form_ii_transfer_stabilization,
        thm:mpu_supplied_projector_simple2,
        thm:mpu_all_later_simple_blockings}
    Aligning the simplicity fixed vectors with the canonical fixed pair using
    the stabilized normalized transfer power gives~(\ref{eq:mpug_simple_canonical_split}) and
    $(\Phi|W^{ij}|\rho)=\delta_{ij}$. Contracting~(\ref{eq:mpug_simple_canonical_split}) on
    the left with $(\Phi|$ yields~(\ref{eq:mpug_simple_canonical_bdry_left}); contracting on the right with $|\rho)$ yields~(\ref{eq:mpug_simple_canonical_bdry_right}).
\end{proof}

\begin{lemma}[Canonical boundary contractions for a simple physical adjoint]
    \label{thm:mpug_adjoint_canonical_boundary_contractions}
    \lean{MPOTensor.IsMPUCanonicalFormII.physicalAdjoint_doubleLayer_boundary_contractions}
    \leanok
    \uses{def:mpu_canonical_form_ii, def:mpu_simple,
        def:mpdo_physical_adjoint, def:mpu_double_layer}
    Let $\mathcal U$ be an MPU with a canonical-form-II presentation and
    recorded fixed matrix $\rho$. Assume that $\mathcal U^\dagger$ is
    simple. Let $W_\dagger$ be the unnormalized double layer of
    $\mathcal U^\dagger$, put $(\Phi|_{(a,b)}=\delta_{a,b}$ and
    $|\rho)_{(a,b)}=\rho_{ba}$. For all physical indices $i,j,k,l$,
    \begin{align}
        (\Phi|W_\dagger^{ij}W_\dagger^{kl}
         & =\delta_{ij}(\Phi|W_\dagger^{kl}, \\
        W_\dagger^{ij}W_\dagger^{kl}|\rho)
         & =\delta_{kl}W_\dagger^{ij}|\rho).
    \end{align}
    The simplicity assumption is precisely the additional hypothesis used
    for the reflected reduction in
    \cite[lines~1221--1254]{FrancoRubio2025MPUGauging}.
\end{lemma}

\begin{proof}\leanok
    \uses{def:mpug_adjoint_canonical_boundaries,
        thm:mpug_simple_canonical_boundary_contractions}
    Passing the canonical presentation to $\mathcal U^\dagger$ with
    the same $\rho$ and applying
    Lemma~\ref{thm:mpug_simple_canonical_boundary_contractions}
    to the simple tensor $\mathcal U^\dagger$ yields both boundary
    contractions.
\end{proof}

\begin{definition}[The three-tensor contraction $s$]
    \label{def:mpug_adjoint_simple_contraction}
    \lean{MPOTensor.adjointSimpleContraction}
    \leanok
    Let $\mathcal U$ be an MPO tensor of physical dimension $d$ and bond
    dimension $D$, and let $\rho$ be a complex $D\times D$ matrix.
    Define an MPO tensor $s$ of the same dimensions by
    \begin{align}
        s^{ij}_{ab}
         & =\sum_{p,q=0}^{d-1}\sum_{x,r,t=0}^{D-1}
        \mathcal U^{ip}_{xb}\,
        \overline{\mathcal U^{qp}_{xr}}\,
        \mathcal U^{qj}_{at}\,\rho_{tr}.
        \label{eq:mpug_adjoint_simple_contraction}
    \end{align}
    For the tensor marked $s$ in
    \cite[lines~1190--1218]{FrancoRubio2025MPUGauging}, the left open bond
    is on the bottom tensor, the right open bond is on the top tensor,
    only $t=r$ contributes when $\rho$ is diagonal, and no factor $1/d$
    occurs because the letters in this contraction are unnormalized.
\end{definition}

\begin{lemma}[The two reductions of the six-tensor network]
    \label{thm:mpug_adjoint_simple_cross_product}
    \lean{MPOTensor.IsMPUCanonicalFormII.adjointSimpleContraction_cross_mul}
    \leanok
    \discussion{7317}
    \uses{def:mpu_canonical_form_ii, def:mpu_simple,
        def:mpdo_physical_adjoint, def:mpug_adjoint_simple_contraction}
    Let $\mathcal U$ have a canonical-form-II presentation with recorded
    fixed matrix $\rho$, and assume that both $\mathcal U$ and
    $\mathcal U^\dagger$ are simple. Construct $s$ from this same $\rho$
    by (\ref{eq:mpug_adjoint_simple_contraction}). For all physical
    indices $i,j,k,l$, the two reductions in
    \cite[lines~1164--1254]{FrancoRubio2025MPUGauging} give
    \begin{align}
        s^{ij}\mathcal U^{kl} & =\mathcal U^{ij}s^{kl}.
        \label{eq:mpug_adjoint_simple_cross_product}
    \end{align}
\end{lemma}

\begin{proof}\leanok
    \uses{thm:mpug_simple_canonical_boundary_contractions,
        thm:mpug_adjoint_canonical_boundary_contractions}
    For fixed open indices $i,j,k,l,a,b$, the six-tensor network expands as
    \begin{align}
        \mathcal N^{ij,kl}_{ab}
         & =\sum_{p,q,v,w}\sum_{x,m,n,z,r,t}
        \mathcal U^{ip}_{xm}\mathcal U^{kq}_{mb}
        \overline{\mathcal U^{vp}_{xn}}
        \overline{\mathcal U^{wq}_{nr}}
        \mathcal U^{vj}_{az}\mathcal U^{wl}_{zt}\rho_{tr}.
    \end{align}
    The first sum ranges over physical indices and the second over virtual
    indices. Applying the right-boundary contraction to the lower two
    layers gives $\mathcal N^{ij,kl}_{ab}=(s^{ij}\mathcal U^{kl})_{ab}$.
    Regrouping the same finite sum and applying the left-boundary
    contraction for the physical adjoint to the upper two layers gives
    $\mathcal N^{ij,kl}_{ab}=(\mathcal U^{ij}s^{kl})_{ab}$.
    Both expressions are therefore equal.
\end{proof}

\begin{lemma}[The contraction $s$ is a scalar multiple of the MPU tensor]
    \label{thm:mpug_adjoint_simple_scalar}
    \lean{MPOTensor.IsMPUCanonicalFormII.adjointSimpleContraction_eq_smul}
    \leanok
    \discussion{7317}
    \uses{def:mpu_canonical_form_ii, def:mpu_simple,
        def:mpdo_physical_adjoint, def:mpug_adjoint_simple_contraction}
    Let $\mathcal U$ have a canonical-form-II presentation with recorded
    fixed matrix $\rho$, and suppose that both $\mathcal U$ and
    $\mathcal U^\dagger$ are simple. Then, as in the proportionality
    conclusion of \cite[lines~1255--1325]{FrancoRubio2025MPUGauging}, the
    tensor $s$ constructed from this $\rho$ satisfies
    \begin{align}
        s^{ij} & =\delta\mathcal U^{ij}
    \end{align}
    for a single scalar $\delta\in\C$ and every physical pair $i,j$.
    % This theorem proves only proportionality. Normalization requires
    % the separate source-factor and trace arguments below; the raw
    % proportionality theorem itself does not evaluate its scalar.
\end{lemma}

\begin{proof}\leanok
    \uses{thm:mpug_adjoint_simple_cross_product,
        thm:mpu_canonical_form_ii_normal, thm:mpug_cross_product_rigidity}
    With $\widetilde{\mathcal U}^{ij}=d^{-1/2}\mathcal U^{ij}$,
    canonical form II gives $d,D>0$ and normality of
    $\widetilde{\mathcal U}$. Scaling
    (\ref{eq:mpug_adjoint_simple_cross_product}) gives
    $s^{ij}\widetilde{\mathcal U}^{kl}
        =\widetilde{\mathcal U}^{ij}s^{kl}$.
    Cross-product rigidity, using blocking to an injective tensor, gives
    $s^{ij}=c\widetilde{\mathcal U}^{ij}$ for a single scalar $c$, so
    $\delta=c/\sqrt d$ satisfies $s^{ij}=\delta\mathcal U^{ij}$.
\end{proof}

\subsection{The anomaly three-cocycle from fusion tensors}

For matrix product operator tensors $M$ and $N$ of bond dimensions $D_M$
and $D_N$, write $MN$ for the stacked product tensor, whose periodic operators
are the products of those of $M$ and $N$, and write $\mathsf a$ for the
bond associator of Definition~\ref{def:mpdo_operator_product_associator},
the permutation from $\C^{D_1}\otimes(\C^{D_2}\otimes\C^{D_3})$, the bond
space of $M(NP)$, to $(\C^{D_1}\otimes\C^{D_2})\otimes\C^{D_3}$, the bond
space of $(MN)P$.

\begin{lemma}[Composition of reductions]
    \label{thm:mpug_reduction_composition}
    \lean{MPSTensor.IsReduction.trans,
        MPSTensor.IsReduction.mulTensor_kronId,
        MPSTensor.IsReduction.mulTensor_idKron,
        MPSTensor.IsReduction.mulTensor_assoc_left,
        MPSTensor.IsReduction.mulTensor_assoc_right}
    \leanok
    \uses{def:mps_rectangular_reduction, def:mpdo_operator_product,
        def:mpdo_operator_product_associator}
    Reductions compose: if $(V_1,W_1)$ reduces $C$ onto $B$ and
    $(V_2,W_2)$ reduces $B$ onto $A$, then $(V_2V_1,W_1W_2)$ reduces $C$
    onto $A$. If $(V,W)$ reduces the doubled-index tensor of $X$ onto that
    of $A$, then $(V\otimes\mathbf 1,W\otimes\mathbf 1)$ reduces $XP$ onto
    $AP$ and $(\mathbf 1\otimes V,\mathbf 1\otimes W)$ reduces $PX$ onto
    $PA$. A reduction $(V,W)$ of $M(NP)$ gives the reduction
    $(V\mathsf a^{-1},\mathsf aW)$ of $(MN)P$, and conversely.
\end{lemma}

\begin{proof}\leanok
    \uses{lem:mpdo_operator_product_eval_word}
    The first claim is immediate from the intertwining equations. For the
    second, Lemma~\ref{lem:mpdo_operator_product_eval_word} expands a word of
    $XP$ along a configuration pair $(\sigma,\tau)$ as a sum over the
    contracted middle configurations $\rho$, and $V\otimes\mathbf 1$ and
    $W\otimes\mathbf 1$ act on the first factor only:
    \begin{align}
        (V\otimes\mathbf 1)\,(XP)^{\sigma\tau}\,(W\otimes\mathbf 1)
         & =\sum_\rho VX^{\sigma\rho}W\otimes P^{\rho\tau}
        =\sum_\rho A^{\sigma\rho}\otimes P^{\rho\tau}
        =(AP)^{\sigma\tau}.
        \notag
    \end{align}
    The case $PX$ is the same with the factors exchanged. For the third, the
    letters satisfy
    $((MN)P)^{i}\,\mathsf a=\mathsf a\,(M(NP))^{i}$, so the words do too.
\end{proof}

\begin{lemma}[Scalar comparison of two reductions]
    \label{thm:mpug_dressed_proportionality}
    \lean{MPSTensor.IsReduction.exists_isDressedProportional,
        MPSTensor.IsDressedProportional.eq_of_forall_exists_ne_zero,
        MPSTensor.IsDressedProportional.pullback,
        MPSTensor.IsReduction.exists_mul_evalWord_ne_zero}
    \leanok
    \uses{thm:mps_reduction_boundary_dressed_uniqueness}
    Let $A$ be normal with positive bond dimension, and let $(V,W)$ and
    $(V',W')$ reduce $B$ onto $A$, where $B$ and $A$ generate the same
    matrix product vectors on nonempty chains. Then there is a unique
    $z\in\C^\times$ such that $VB^{\mathbf w}=z\,V'B^{\mathbf w}$ for all
    sufficiently long words $\mathbf w$. If
    $XA^{\mathbf w}=z\,YA^{\mathbf w}$ for all long words, then
    $XVB^{\mathbf w}=z\,YVB^{\mathbf w}$ for all long words.
\end{lemma}

\begin{proof}\leanok
    Existence is Lemma~\ref{thm:mps_reduction_boundary_dressed_uniqueness}.
    Uniqueness holds because $VB^{\mathbf w}W=A^{\mathbf w}$ and a normal
    tensor has nonzero words of every multiple of its injectivity length.
    For the last claim, once both exterior parts of a long word
    $\mathbf w=\mathbf p\mathbf c\mathbf q$ exceed the residual nilpotency
    length, $VB^{\mathbf w}=A^{\mathbf p\mathbf c}VB^{\mathbf q}$.
\end{proof}

\begin{definition}[The anomaly three-cochain]
    \label{def:mpug_anomaly_three_cochain}
    \lean{MPOTensor.GroupFamily.IsNormalRepresentation,
        MPOTensor.GroupFamily.FusionData,
        MPOTensor.GroupFamily.FusionData.leftV,
        MPOTensor.GroupFamily.FusionData.rightV,
        MPOTensor.GroupFamily.FusionData.IsAssociator,
        MPOTensor.GroupFamily.FusionData.omega}
    \leanok
    \uses{def:mpug_group_representation, def:mpdo_operator_product,
        def:mpo_to_mps, def:mps_rectangular_reduction,
        def:mpdo_operator_product_associator, def:normal}
    Let $\{O_g\}_{g\in G}$ be matrix product operator tensors of positive
    bond dimension whose doubled-index tensors are normal and whose periodic
    operators satisfy $O_g O_h=O_{gh}$ on every nonempty chain. A choice of
    fusion tensors is, for every pair $g,h$, a reduction
    $(F^<_{g,h},F^>_{g,h})$ of $O_gO_h$ onto $O_{gh}$. The triple product
    $T=(O_gO_h)O_k$ reduces onto $O_{ghk}$ along the two fusion trees with
    left boundaries
    \begin{align}
        V^{L} & =F^<_{gh,k}\,(F^<_{g,h}\otimes\mathbf 1),                  &
        V^{R} & =F^<_{g,hk}\,(\mathbf 1\otimes F^<_{h,k})\,\mathsf a^{-1},
    \end{align}
    where the target bond space of $F^<_{g,hk}$, that of $O_{g(hk)}$, is
    identified with that of $O_{(gh)k}$ through the equality
    $g(hk)=(gh)k$ in $G$; the same identification enters the right
    boundaries $W^L=(F^>_{g,h}\otimes\mathbf 1)F^>_{gh,k}$ and
    $W^R=\mathsf a(\mathbf 1\otimes F^>_{h,k})F^>_{g,hk}$.
    The anomaly three-cochain value $\omega(g,h,k)\in\C^\times$ is the scalar with
    $V^LT^{\mathbf w}=\omega(g,h,k)\,V^RT^{\mathbf w}$ for all long words
    $\mathbf w$. Equivalently $T^{\mathbf w}W^R=\omega(g,h,k)\,T^{\mathbf w}W^L$
    for the right boundaries, which is the defining display of
    \cite{FrancoRubio2025MPUGauging} before its three-cocycle equation.
    No normalization of $\omega$ is imposed.
\end{definition}

\begin{center}
    % Formula: V^LT^{\mathbf w}=\omega(g,h,k)\,V^RT^{\mathbf w} with
    % V^L=F^<_{gh,k}(F^<_{g,h}\otimes\mathbf 1) and
    % V^R=F^<_{g,hk}(\mathbf 1\otimes F^<_{h,k})\mathsf a^{-1}; source:
    % arXiv:2502.20257, the display before eq:3-cocycle, lines 1505--1535
    % (drawn there with right boundaries).
    % Ink-to-index map: the bracketed three-row column stands for the word
    % \mathbf w of T=(O_gO_h)O_k, rows O_g, O_h, O_k from top to bottom;
    % each pill has its single side west.  The inner pill fuses two rows,
    % and the outer pill fuses its output with the remaining row.  The
    % associator \mathsf a only regroups the Kronecker factors of the triple
    % bond and is not drawn.
    % Contraction set: the middle physical indices, the bonds between the
    % pills, and the bonds from the pills into the word; the three east
    % virtual legs, the single west leg, and the outer physical legs remain
    % open.
    % Boundary signature: (virtual-west=1 | virtual-east=3 |
    % physical-up=1 | physical-down=1) on both sides.
    \tenkzkernel
    \begin{tenkz}[rows={op,op,op}, cols=3, bonds=none]
        \tnfuse[at=(1,1), name=Fo, skin=pill, wires=3]{F^<_{gh,k}}
        \tnfuse[at=(1,2), name=Fi, skin=pill]{F^<_{g,h}}
        \tn[at=(3,2), void=sealed]{}
        \tn[at=(1,3), name=Tg, skin=mpo, ports={90:physical, 270:physical, 180:virtual, 0:virtual}]{O_g}
        \tn[at=(2,3), name=Th, skin=mpo, ports={90:physical, 270:physical, 180:virtual, 0:virtual}]{O_h}
        \tn[at=(3,3), name=Tk, skin=mpo, ports={90:physical, 270:physical, 180:virtual, 0:virtual}]{O_k}
        \tnwire{Tg.270}{Th.90}
        \tnwire{Th.270}{Tk.90}
        \tnwire{open w}{Fo.180}
        \tnwire{Fo.0@1}{Fi.180}
        \tnwire{Fo.0@3}{Tk.180}
        \tnwire{Fi.0@1}{Tg.180}
        \tnwire{Fi.0@2}{Th.180}
        \tnwire{Tg.0}{open e}
        \tnwire{Th.0}{open e}
        \tnwire{Tk.0}{open e}
        \tnmark[form=bracket, label pos=90]{Tg}{\mathbf w}
    \end{tenkz}
    \quad $=\omega(g,h,k)$ \quad
    \begin{tenkz}[rows={op,op,op}, cols=3, bonds=none]
        \tnfuse[at=(1,1), name=Fo, skin=pill, wires=3]{F^<_{g,hk}}
        \tnfuse[at=(2,2), name=Fi, skin=pill]{F^<_{h,k}}
        \tn[at=(1,2), void=sealed]{}
        \tn[at=(1,3), name=Tg, skin=mpo, ports={90:physical, 270:physical, 180:virtual, 0:virtual}]{O_g}
        \tn[at=(2,3), name=Th, skin=mpo, ports={90:physical, 270:physical, 180:virtual, 0:virtual}]{O_h}
        \tn[at=(3,3), name=Tk, skin=mpo, ports={90:physical, 270:physical, 180:virtual, 0:virtual}]{O_k}
        \tnwire{Tg.270}{Th.90}
        \tnwire{Th.270}{Tk.90}
        \tnwire{open w}{Fo.180}
        \tnwire{Fo.0@1}{Tg.180}
        \tnwire{Fo.0@2}{Fi.180}
        \tnwire{Fi.0@1}{Th.180}
        \tnwire{Fi.0@2}{Tk.180}
        \tnwire{Tg.0}{open e}
        \tnwire{Th.0}{open e}
        \tnwire{Tk.0}{open e}
        \tnmark[form=bracket, label pos=90]{Tg}{\mathbf w}
    \end{tenkz}
\end{center}

\begin{lemma}[The two fusion trees are reductions]
    \label{thm:mpug_fusion_trees_reductions}
    \lean{MPOTensor.GroupFamily.FusionData.isReduction_left,
        MPOTensor.GroupFamily.FusionData.isReduction_right}
    \leanok
    \uses{def:mpug_anomaly_three_cochain, thm:mpug_reduction_composition}
    The pairs $(V^L,W^L)$ and $(V^R,W^R)$ of
    Definition~\ref{def:mpug_anomaly_three_cochain} are reductions of the
    triple product $(O_gO_h)O_k$ onto $O_{ghk}$.
\end{lemma}

\begin{proof}\leanok
    Tensoring the reduction $(F^<_{g,h},F^>_{g,h})$ with the
    identity on the third factor and composing with $(F^<_{gh,k},F^>_{gh,k})$
    gives $(V^L,W^L)$. Similarly, tensoring $(F^<_{h,k},F^>_{h,k})$ with the
    identity on the first factor, composing with $(F^<_{g,hk},F^>_{g,hk})$,
    and conjugating by the associator $\mathsf a$ gives $(V^R,W^R)$; all
    three steps follow from Lemma~\ref{thm:mpug_reduction_composition}.
\end{proof}

\begin{lemma}[Characterization of the anomaly three-cochain]
    \label{thm:mpug_anomaly_three_cochain_spec}
    \lean{MPOTensor.GroupFamily.FusionData.isAssociator_omega,
        MPOTensor.GroupFamily.FusionData.eq_omega_of_isAssociator}
    \leanok
    \uses{def:mpug_anomaly_three_cochain, thm:mpug_fusion_trees_reductions,
        thm:mpug_dressed_proportionality}
    The scalar $\omega(g,h,k)$ exists, is nonzero, and is the only scalar
    with the dressed identity of
    Definition~\ref{def:mpug_anomaly_three_cochain}.
\end{lemma}

\begin{proof}\leanok
    By Lemma~\ref{thm:mpug_fusion_trees_reductions} both fusion trees are
    reductions of $T$ onto the normal tensor $O_{ghk}$, and $T$ and
    $O_{ghk}$ have the same periodic operators on nonempty chains, so
    Lemma~\ref{thm:mpug_dressed_proportionality} yields existence and
    uniqueness.
\end{proof}

\begin{theorem}[The anomaly three-cocycle]
    \label{thm:mpug_anomaly_three_cocycle}
    \lean{MPOTensor.GroupFamily.FusionData.isCocycle_omega,
        MPOTensor.GroupFamily.FusionData.isCocycle_omega_of_isRepresentation}
    \leanok
    \uses{thm:mpug_anomaly_three_cochain_spec, thm:mpug_reduction_composition,
        thm:mpug_dressed_proportionality}
    For every $g,h,k,l\in G$,
    \begin{align}
        \omega(gh,k,l)\,\omega(g,h,kl)
         & =\omega(g,h,k)\,\omega(g,hk,l)\,\omega(h,k,l).
        \label{eq:mpug_anomaly_three_cocycle}
    \end{align}
\end{theorem}

\begin{proof}\leanok
    \uses{thm:mpug_anomaly_three_cochain_spec, thm:mpug_reduction_composition,
        thm:mpug_dressed_proportionality}
    The five fusion trees of the fourfold product $((O_gO_h)O_k)O_l$ are
    reductions onto $O_{ghkl}$. Two trees that differ by one change of
    bracketing are related against long words by one value of $\omega$:
    either directly, through Lemma~\ref{thm:mpug_reduction_composition}
    with an identity factor, or after pulling back along a fusion tensor
    by Lemma~\ref{thm:mpug_dressed_proportionality}. Write
    $\mathsf a_{p,q,r}$ for $\mathsf a$ on $\C^p\otimes\C^q\otimes\C^r$.
    Identifying the boundaries uses the naturality of $\mathsf a$, for
    $X\colon\C^n\to\C^m$,
    \begin{align}
        \mathsf a^{-1}_{m,q,r}\,(X\otimes\mathbf 1_q\otimes\mathbf 1_r)
         & =(X\otimes\mathbf 1_{qr})\,\mathsf a^{-1}_{n,q,r},
        \notag                                                                  \\
        \mathsf a^{-1}_{p,m,r}\,(\mathbf 1_p\otimes X\otimes\mathbf 1_r)
         & =(\mathbf 1_p\otimes(X\otimes\mathbf 1_r))\,\mathsf a^{-1}_{p,n,r},
        \notag                                                                  \\
        \mathsf a^{-1}_{p,q,m}\,(\mathbf 1_{pq}\otimes X)
         & =(\mathbf 1_p\otimes(\mathbf 1_q\otimes X))\,\mathsf a^{-1}_{p,q,n},
        \notag
    \end{align}
    and the pentagon identity for the reassociation permutations,
    \begin{align}
        (\mathbf 1_p\otimes\mathsf a^{-1}_{q,r,s})\,\mathsf a^{-1}_{p,qr,s}\,
        (\mathsf a^{-1}_{p,q,r}\otimes\mathbf 1_s)
         & =\mathsf a^{-1}_{p,q,rs}\,\mathsf a^{-1}_{pq,r,s}.
        \notag
    \end{align}
    The two paths around the pentagon
    compare the same pair of trees, and uniqueness of the scalar on a
    nonvanishing word gives (\ref{eq:mpug_anomaly_three_cocycle}).
\end{proof}

\begin{lemma}[Independence of the fusion tensors]
    \label{thm:mpug_anomaly_three_cocycle_gauge}
    \lean{MPOTensor.GroupFamily.FusionData.relativeScalar,
        MPOTensor.GroupFamily.FusionData.isDressedProportional_relativeScalar,
        MPOTensor.GroupFamily.FusionData.omega_eq_fusionGauge,
        MPOTensor.GroupFamily.FusionData.omega_cohomologousTo}
    \leanok
    \uses{thm:mpug_anomaly_three_cochain_spec,
        thm:mpug_dressed_proportionality}
    Let $\omega$ and $\omega'$ come from two choices of fusion tensors, and
    let $\beta_{g,h}\in\C^\times$ be the scalar with
    $F^<_{g,h}B^{\mathbf w}=\beta_{g,h}F'^<_{g,h}B^{\mathbf w}$ for all long
    words of $B=O_gO_h$. Then
    \begin{align}
        \omega'(g,h,k)
         & =\frac{\beta_{g,hk}\,\beta_{h,k}}{\beta_{g,h}\,\beta_{gh,k}}
        \,\omega(g,h,k),
    \end{align}
    so the cohomology class of $\omega$ does not depend on the fusion
    tensors, as in (\ref{eq:mpug_fusion_gauge}).
\end{lemma}

\begin{proof}\leanok
    \uses{thm:mpug_anomaly_three_cochain_spec,
        thm:mpug_dressed_proportionality}
    By the two parts of Lemma~\ref{thm:mpug_dressed_proportionality}, the
    left boundaries of the two choices of the tree fusing $g$ with $h$
    first differ by $\beta_{g,h}\beta_{gh,k}$, and those of the tree fusing
    $h$ with $k$ first differ by $\beta_{h,k}\beta_{g,hk}$; comparing with
    Lemma~\ref{thm:mpug_anomaly_three_cochain_spec} gives the ratio for
    $\omega'$.
\end{proof}

\begin{definition}[The anomaly class]
    \label{def:mpug_anomaly_class}
    \lean{MPOTensor.GroupFamily.IsNormalRepresentation.anomalyClass}
    \leanok
    \uses{thm:mpug_anomaly_three_cocycle, thm:mpug_h3_group_cohomology}
    Let the doubled-index tensors be normal. The anomaly class in
    $H^3(G,\C^\times_{\mathrm{triv}})$ is the class of the anomaly
    three-cocycle $\omega$, for any choice of fusion tensors.
\end{definition}

\begin{corollary}[The anomaly class is independent of the fusion tensors]
    \label{cor:mpug_anomaly_class}
    \lean{MPOTensor.GroupFamily.FusionData.anomalyClass_omega_eq,
        MPOTensor.GroupFamily.FusionData.anomalyClass_omega,
        MPOTensor.GroupFamily.FusionData.isTrivialGaugeClass_omega_iff}
    \leanok
    \uses{def:mpug_anomaly_class, thm:mpug_anomaly_three_cocycle,
        thm:mpug_anomaly_three_cocycle_gauge, thm:mpug_h3_group_cohomology}
    The class $[\omega]\in H^3(G,\C^\times_{\mathrm{triv}})$ of the anomaly
    three-cocycle is the same for every choice of fusion tensors. It vanishes
    exactly when $\omega$ has trivial gauge class.
\end{corollary}

\begin{proof}\leanok
    \uses{thm:mpug_anomaly_three_cocycle, thm:mpug_anomaly_three_cocycle_gauge,
        thm:mpug_h3_group_cohomology}
    Two choices of fusion tensors give cohomologous three-cocycles by
    Lemma~\ref{thm:mpug_anomaly_three_cocycle_gauge}, hence equal classes
    by (\ref{eq:mpug_h3_cohomologous}). The second statement is
    (\ref{eq:mpug_h3_trivial}).
\end{proof}

\subsection{Reductions through action tensors}

\begin{definition}[The action tensor]
    \label{def:mpug_action_tensor}
    \lean{MPOTensor.actTensor}
    \leanok
    \uses{def:mpo_tensor, def:mps_tensor}
    For a matrix product operator tensor $T$ of bond dimension $D_1$ and a
    matrix product state tensor $B$ of bond dimension $D_2$, the action
    tensor $T\cdot B$ has bond dimension $D_1D_2$ and letters
    \begin{align}
        (T\cdot B)^i
         & =\sum_jT^{ij}\otimes B^j.
        \label{eq:mpug_action_tensor}
    \end{align}
    It is the tensor whose local decomposition onto an image block gives the
    action tensors of the periodic-boundary case of
    \cite[``Periodic boundary condition case'']{GarreRubio2023MPOSym}.
\end{definition}

\begin{lemma}[The action tensor generates the acted state]
    \label{thm:mpug_action_tensor_periodic}
    \lean{MPOTensor.evalWord_actTensor, MPOTensor.mpv_actTensor,
        MPOTensor.mpo_mulVec_mpv}
    \leanok
    \uses{def:mpug_action_tensor, def:mpo_family, def:mpv_state}
    For every word $\sigma$ of length $N$,
    \begin{align}
        (T\cdot B)^{\sigma}
         & =\sum_{\tau}T^{\sigma\tau}\otimes B^{\tau},
        \label{eq:mpug_action_tensor_word}
    \end{align}
    the sum running over words $\tau$ of length $N$. Consequently
    $\ket{V^{(N)}(T\cdot B)}=\rho^{(N)}(T)\ket{V^{(N)}(B)}$ for every $N$.
\end{lemma}

\begin{center}
    % Formula: \ket{V^{(N)}(T\cdot B)}=\rho^{(N)}(T)\ket{V^{(N)}(B)}, the
    % consequence of (\ref{eq:mpug_action_tensor_word}).
    % Ink-to-index map: the left ring is the periodic vector of the action
    % tensor T\cdot B; on the right the upper row is T and the lower row B,
    % and the vertical leg at each site is the summed index j of
    % (\ref{eq:mpug_action_tensor}).
    % Contraction set: the periodic virtual bonds of each ring and the
    % summed indices j; the physical legs \sigma_n remain open.
    % Boundary signature: (virtual-west=0 | virtual-east=0 |
    % physical-up=3 | physical-down=0) on both sides.
    \tenkzkernel
    \begin{tenkz}[rows={ket}, cols=4, boundary=periodic]
        \tn[ports={90:physical:$\sigma_1$}]{T\cdot B} &
        \tn[ports={90:physical:$\sigma_2$}]{T\cdot B} &
        \tn[skin=dots, ports={180:virtual, 0:virtual}]{} &
        \tn[ports={90:physical:$\sigma_N$}]{T\cdot B}
    \end{tenkz}
    \quad $=$ \quad
    \begin{tenkz}[rows={op,ket}, cols=4, boundary=periodic]
        \tn[skin=mpo, ports={90:physical:$\sigma_1$}]{T} &
        \tn[skin=mpo, ports={90:physical:$\sigma_2$}]{T} &
        \tn[skin=dots, ports={180:virtual, 0:virtual}]{} &
        \tn[skin=mpo, ports={90:physical:$\sigma_N$}]{T} \\
        \tn[label pos=s, ports={180:virtual, 0:virtual}]{B} &
        \tn[label pos=s, ports={180:virtual, 0:virtual}]{B} &
        \tn[skin=dots, ports={180:virtual, 0:virtual}]{} &
        \tn[label pos=s, ports={180:virtual, 0:virtual}]{B}
    \end{tenkz}
\end{center}

\begin{proof}\leanok
    Induction on $N$: the empty word gives $\mathbf 1\otimes\mathbf 1$, and
    one more letter multiplies both sides by $(T\cdot B)^i$, since
    $(T^{\sigma\tau}\otimes B^{\tau})(T^{ij}\otimes B^j)
        =T^{\sigma\tau}T^{ij}\otimes B^{\tau}B^j$. Taking traces and using
    $\tr(X\otimes Y)=\tr X\,\tr Y$ gives
    $V^{(N)}(T\cdot B)_\sigma
        =\sum_\tau\tr(T^{\sigma\tau})\tr(B^\tau)
        =(\rho^{(N)}(T)V^{(N)}(B))_\sigma$.
\end{proof}

\begin{lemma}[Reductions through action tensors]
    \label{thm:mpug_action_tensor_reductions}
    \lean{MPSTensor.IsReduction.actTensor_idKron,
        MPSTensor.IsReduction.actTensor_kronId,
        MPOTensor.actTensor_mulTensor,
        MPOTensor.actTensor_mulTensor_mul_assocMatrix,
        MPOTensor.assocInvMatrix_mul_actTensor_mulTensor,
        MPSTensor.IsReduction.actTensor_assoc_left,
        MPSTensor.IsReduction.actTensor_assoc_right,
        MPOTensor.actTensor_mul_kronId_of_intertwine,
        MPSTensor.IsDressedProportional.actTensor_idKron,
        MPSTensor.IsDressedProportional.actTensor_kronId}
    \leanok
    \uses{def:mpug_action_tensor, def:mpdo_operator_product,
        def:mpdo_operator_product_associator, def:mpo_to_mps,
        def:mps_rectangular_reduction}
    If $(V,W)$ reduces $B$ onto $A$, then
    $(\mathbf 1\otimes V,\mathbf 1\otimes W)$ reduces $T\cdot B$ onto
    $T\cdot A$. If $(V,W)$ reduces the doubled-index tensor of $X$ onto that
    of $Y$, then $(V\otimes\mathbf 1,W\otimes\mathbf 1)$ reduces $X\cdot B$
    onto $Y\cdot B$. The letters satisfy
    $((MN)\cdot A)^i\,\mathsf a=\mathsf a\,(M\cdot(N\cdot A))^i$, so a
    reduction $(V,W)$ of $M\cdot(N\cdot A)$ onto $Z$ gives the reduction
    $(V\mathsf a^{-1},\mathsf aW)$ of $(MN)\cdot A$ onto $Z$, and a reduction
    $(V,W)$ of $(MN)\cdot A$ onto $Z$ gives the reduction
    $(V\mathsf a,\mathsf a^{-1}W)$ of $M\cdot(N\cdot A)$ onto $Z$, as in
    Lemma~\ref{thm:mpug_reduction_composition}. If $P$ intertwines the
    letters of two operator tensors, $X'^{ij}P=PX^{ij}$ for all $i,j$, then
    $(X'\cdot B)^i\,(P\otimes\mathbf 1)=(P\otimes\mathbf 1)\,(X\cdot B)^i$.
    Finally, a relation
    $VB^{\mathbf w}=z\,V'B^{\mathbf w}$ for all long words lifts to
    $(\mathbf 1\otimes V)(T\cdot B)^{\mathbf w}
        =z\,(\mathbf 1\otimes V')(T\cdot B)^{\mathbf w}$, and the analogous
    relation for the doubled-index tensor of $X$ lifts through
    $V\otimes\mathbf 1$ to $X\cdot B$.
\end{lemma}

\begin{proof}\leanok
    \uses{thm:mpug_action_tensor_periodic, thm:mpug_reduction_composition}
    By Lemma~\ref{thm:mpug_action_tensor_periodic}, every word of
    $T\cdot B$ is the sum~(\ref{eq:mpug_action_tensor_word}) of Kronecker
    products of a word of $T$ and a word of $B$ of the same length. A
    boundary of the form $\mathbf 1\otimes V$ passes through the first
    factor of each summand:
    \begin{align}
        (\mathbf 1\otimes V)\,(T\cdot B)^{\sigma}\,(\mathbf 1\otimes W)
         & =\sum_\tau T^{\sigma\tau}\otimes VB^{\tau}W
        =\sum_\tau T^{\sigma\tau}\otimes A^{\tau}
        =(T\cdot A)^{\sigma},
        \notag
    \end{align}
    which is the reduction of the state. For the operator, the word
    $X^{\sigma\tau}$ is the word of the doubled-index tensor of $X$ on the
    pair configuration $(\sigma,\tau)$, so
    $(V\otimes\mathbf 1)(X\cdot B)^\sigma(W\otimes\mathbf 1)
        =\sum_\tau VX^{\sigma\tau}W\otimes B^\tau
        =(Y\cdot B)^\sigma$. For a dressed proportionality, if
    $VB^{\tau}=z\,V'B^{\tau}$ for all words $\tau$ of length at least $N_0$,
    then for every word $\sigma$ of such a length
    \begin{align}
        (\mathbf 1\otimes V)\,(T\cdot B)^{\sigma}
         & =\sum_\tau T^{\sigma\tau}\otimes VB^{\tau}
        =z\sum_\tau T^{\sigma\tau}\otimes V'B^{\tau}
        =z\,(\mathbf 1\otimes V')\,(T\cdot B)^{\sigma},
        \notag
    \end{align}
    because every $\tau$ in the sum has the length of $\sigma$; the same
    substitution on the pair words of $X$ gives the lift through
    $V\otimes\mathbf 1$. For the intertwiner,
    $(X'^{ij}\otimes B^j)(P\otimes\mathbf 1)=X'^{ij}P\otimes B^j
        =PX^{ij}\otimes B^j=(P\otimes\mathbf 1)(X^{ij}\otimes B^j)$.
    The letter identity is the reassociation of Kronecker products:
    \begin{align}
        ((MN)\cdot A)^i
         & =\sum_{j,k}(M^{ij}\otimes N^{jk})\otimes A^k
        =\mathsf a\,\Bigl(\sum_{j,k}M^{ij}\otimes(N^{jk}\otimes A^k)\Bigr)
        \,\mathsf a^{-1}
        =\mathsf a\,(M\cdot(N\cdot A))^i\,\mathsf a^{-1}.
        \notag
    \end{align}
\end{proof}

\subsection{An anomaly obstructs invariant normal states}

\begin{definition}[L-symbol of an invariant normal state]
    \label{def:mpug_single_block_l_symbol}
    \lean{MPOTensor.GroupFamily.FixesMPV,
        MPOTensor.GroupFamily.ActionData,
        MPOTensor.GroupFamily.ActionData.actV,
        MPOTensor.GroupFamily.ActionData.fuseV,
        MPOTensor.GroupFamily.ActionData.lSymbol}
    \leanok
    \uses{def:mpug_anomaly_three_cochain, def:mpug_l_symbols, def:mpv_state}
    Let $\{O_g\}$ be as in Definition~\ref{def:mpug_anomaly_three_cochain},
    and let $A$ be a matrix product state tensor with
    $O_g\ket{V^{(N)}(A)}=\ket{V^{(N)}(A)}$ for every $g$ and every $N\geq1$.
    Action tensors for $A$ are, for every $g$, a reduction $(V_g,W_g)$ of
    $O_g\cdot A$ onto $A$. The stacked action $(O_gO_h)\cdot A$ then has the
    two left boundaries
    \begin{align}
        \alpha_{g,h}  & =V_g(\mathbf 1\otimes V_h)\,\mathsf a^{-1}, &
        \varphi_{g,h} & =V_{gh}(F^<_{g,h}\otimes\mathbf 1),
        \notag
    \end{align}
    and the $L$-symbol of a single invariant block in
    \cite{FrancoRubio2025MPUGauging} is a nonzero scalar $L_{g,h}=z$ with
    $\alpha_{g,h}B^{\mathbf w}=z\,\varphi_{g,h}B^{\mathbf w}$ for all long
    words $\mathbf w$ of $B=(O_gO_h)\cdot A$, or $1$ if there is none.
\end{definition}

\begin{center}
    % Formula: \alpha_{g,h}B^{\mathbf w}=L_{g,h}\,\varphi_{g,h}B^{\mathbf w}
    % with \alpha_{g,h}=V_g(\mathbf 1\otimes V_h)\mathsf a^{-1},
    % \varphi_{g,h}=V_{gh}(F^<_{g,h}\otimes\mathbf 1), and
    % B=(O_gO_h)\cdot A; source: arXiv:2502.20257, eq:defL, lines 1875--1914
    % (drawn there with right boundaries).
    % Ink-to-index map: the bracketed column stands for the word \mathbf w
    % of B, with rows O_g, O_h and the state tensor A from top to bottom;
    % each pill has its single side west.  On the left, V_h reduces
    % O_h\cdot A and V_g then reduces O_g\cdot A; on the right, F^<_{g,h}
    % fuses O_gO_h and V_{gh} reduces O_{gh}\cdot A.  The associator is
    % not drawn.
    % Contraction set: the physical indices between the rows, the bonds
    % between the pills, and the bonds into the word; the three east virtual
    % legs, the west leg, and the upper physical leg remain open.
    % Boundary signature: (virtual-west=1 | virtual-east=3 |
    % physical-up=1 | physical-down=0) on both sides.
    \tenkzkernel
    \begin{tenkz}[rows={op,op,ket}, cols=3, bonds=none]
        \tnfuse[at=(1,1), name=Fo, skin=pill, wires=3]{V_g}
        \tnfuse[at=(2,2), name=Fi, skin=pill]{V_h}
        \tn[at=(1,2), void=sealed]{}
        \tn[at=(1,3), name=Tg, skin=mpo, ports={90:physical, 270:physical, 180:virtual, 0:virtual}]{O_g}
        \tn[at=(2,3), name=Th, skin=mpo, ports={90:physical, 270:physical, 180:virtual, 0:virtual}]{O_h}
        \tn[at=(3,3), name=A, label pos=s, ports={90:physical, 180:virtual, 0:virtual}]{A}
        \tnwire{Tg.270}{Th.90}
        \tnwire{Th.270}{A.90}
        \tnwire{open w}{Fo.180}
        \tnwire{Fo.0@1}{Tg.180}
        \tnwire{Fo.0@2}{Fi.180}
        \tnwire{Fi.0@1}{Th.180}
        \tnwire{Fi.0@2}{A.180}
        \tnwire{Tg.0}{open e}
        \tnwire{Th.0}{open e}
        \tnwire{A.0}{open e}
        \tnmark[form=bracket, label pos=90]{Tg}{\mathbf w}
    \end{tenkz}
    \quad $=L_{g,h}$ \quad
    \begin{tenkz}[rows={op,op,ket}, cols=3, bonds=none]
        \tnfuse[at=(1,1), name=Fo, skin=pill, wires=3]{V_{gh}}
        \tnfuse[at=(1,2), name=Fi, skin=pill]{F^<_{g,h}}
        \tn[at=(3,2), void=sealed]{}
        \tn[at=(1,3), name=Tg, skin=mpo, ports={90:physical, 270:physical, 180:virtual, 0:virtual}]{O_g}
        \tn[at=(2,3), name=Th, skin=mpo, ports={90:physical, 270:physical, 180:virtual, 0:virtual}]{O_h}
        \tn[at=(3,3), name=A, label pos=s, ports={90:physical, 180:virtual, 0:virtual}]{A}
        \tnwire{Tg.270}{Th.90}
        \tnwire{Th.270}{A.90}
        \tnwire{open w}{Fo.180}
        \tnwire{Fo.0@1}{Fi.180}
        \tnwire{Fo.0@3}{A.180}
        \tnwire{Fi.0@1}{Tg.180}
        \tnwire{Fi.0@2}{Th.180}
        \tnwire{Tg.0}{open e}
        \tnwire{Th.0}{open e}
        \tnwire{A.0}{open e}
        \tnmark[form=bracket, label pos=90]{Tg}{\mathbf w}
    \end{tenkz}
\end{center}

\begin{lemma}[Characterization of the single-block L-symbol]
    \label{thm:mpug_single_block_l_symbol_spec}
    \lean{MPOTensor.GroupFamily.nonempty_actionData,
        MPOTensor.GroupFamily.ActionData.isDressedProportional_lSymbol}
    \leanok
    \uses{def:mpug_single_block_l_symbol}
    If $A$ in Definition~\ref{def:mpug_single_block_l_symbol} is normal
    of positive bond dimension, then action tensors for $A$ exist
    \cite[lines~1064--1090]{GarreRubio2023MPOSym}, and for every choice of
    them $\alpha_{g,h}B^{\mathbf w}=L_{g,h}\,\varphi_{g,h}B^{\mathbf w}$ for
    all long words $\mathbf w$ of $B=(O_gO_h)\cdot A$.
\end{lemma}

\begin{proof}\leanok
    \uses{cor:asym_action_tensors, thm:mpug_action_tensor_reductions,
        thm:mpug_reduction_composition, thm:mpug_dressed_proportionality}
    Existence is Corollary~\ref{cor:asym_action_tensors}. By
    Lemmas~\ref{thm:mpug_action_tensor_reductions}
    and~\ref{thm:mpug_reduction_composition}, both $\alpha_{g,h}$ and
    $\varphi_{g,h}$ are left boundaries of reductions of $B$ onto $A$, and
    $B$ generates the same vectors as $A$, so
    Lemma~\ref{thm:mpug_dressed_proportionality} gives a nonzero scalar
    relating them against long words.
\end{proof}

\begin{lemma}[Compatibility of the L-symbol with the anomaly]
    \label{thm:mpug_single_block_l_compatibility}
    \lean{MPOTensor.GroupFamily.ActionData.isCompatible_lSymbol}
    \leanok
    \uses{def:mpug_single_block_l_symbol, thm:mpug_single_block_l_symbol_spec,
        def:mpug_anomaly_three_cochain, def:mpug_l_symbols}
    In the setting of Lemma~\ref{thm:mpug_single_block_l_symbol_spec},
    (\ref{eq:mpug_l_compatibility}) for a one-point set of blocks holds for
    all $g,h,k\in G$:
    \begin{align}
        L_{g,hk}L_{h,k}
         & =\omega(g,h,k)\,L_{g,h}L_{gh,k}.
    \end{align}
\end{lemma}

\begin{proof}\leanok
    \uses{thm:mpug_single_block_l_symbol_spec,
        thm:mpug_action_tensor_reductions, thm:mpug_reduction_composition,
        thm:mpug_anomaly_three_cochain_spec, thm:mpug_dressed_proportionality}
    For the triple action $B=((O_gO_h)O_k)\cdot A$ and the two fusion trees
    $V^L$, $V^R$ of $(g,h,k)$ in
    Definition~\ref{def:mpug_anomaly_three_cochain}, the naturality and
    pentagon identities of $\mathsf a$ give five identities between left
    boundaries of reductions of $B$ onto $A$:
    \begin{align}
        \alpha_{g,h}(\mathbf 1\otimes V_k)\,\mathsf a^{-1}
         & =V_g(\mathbf 1\otimes\alpha_{h,k})\,\mathsf a^{-1}
        (\mathsf a^{-1}\otimes\mathbf 1),
        \label{eq:mpug_l_compat_first_boundary}               \\
        V_g(\mathbf 1\otimes\varphi_{h,k})\,\mathsf a^{-1}
        (\mathsf a^{-1}\otimes\mathbf 1)
         & =\alpha_{g,hk}
        \bigl((\mathbf 1\otimes F^<_{h,k})\,\mathsf a^{-1}\otimes\mathbf 1\bigr),
        \notag                                                \\
        \varphi_{g,h}(\mathbf 1\otimes V_k)\,\mathsf a^{-1}
         & =\alpha_{gh,k}
        \bigl((F^<_{g,h}\otimes\mathbf 1)\otimes\mathbf 1\bigr),
        \notag                                                \\
        \varphi_{gh,k}\bigl((F^<_{g,h}\otimes\mathbf 1)\otimes\mathbf 1\bigr)
         & =V_{ghk}(V^L\otimes\mathbf 1),
        \notag                                                \\
        \varphi_{g,hk}
        \bigl((\mathbf 1\otimes F^<_{h,k})\,\mathsf a^{-1}\otimes\mathbf 1\bigr)
         & =V_{ghk}(V^R\otimes\mathbf 1).
        \notag
    \end{align}
    Starting from $\alpha_{g,h}(\mathbf 1\otimes V_k)\mathsf a^{-1}B^{\mathbf w}$,
    replacing $\alpha_{g,h}$ by $L_{g,h}\varphi_{g,h}$, $\alpha_{gh,k}$ by
    $L_{gh,k}\varphi_{gh,k}$, and $V^L$ by $\omega(g,h,k)V^R$ along one
    path, or applying (\ref{eq:mpug_l_compat_first_boundary}) and replacing
    $\alpha_{h,k}$ by $L_{h,k}\varphi_{h,k}$ inside the action of $g$ and
    $\alpha_{g,hk}$ by $L_{g,hk}\varphi_{g,hk}$ along the other, uses only
    replacements that hold against long words by
    Lemma~\ref{thm:mpug_single_block_l_symbol_spec}, lifted through the
    action tensors by Lemma~\ref{thm:mpug_action_tensor_reductions} and
    pulled back along the reductions by
    Lemma~\ref{thm:mpug_dressed_proportionality}. For all long words,
    \begin{align}
        \alpha_{g,h}(\mathbf 1\otimes V_k)\,\mathsf a^{-1}B^{\mathbf w}
         & =L_{g,h}L_{gh,k}\,\omega(g,h,k)\,V_{ghk}(V^R\otimes\mathbf 1)
        B^{\mathbf w}
        =L_{h,k}L_{g,hk}\,V_{ghk}(V^R\otimes\mathbf 1)B^{\mathbf w}.
        \notag
    \end{align}
    The boundary $V_{ghk}(V^R\otimes\mathbf 1)$ belongs to a reduction onto
    the normal tensor $A$, so it is nonzero against arbitrarily long words,
    and the two scalars agree.
\end{proof}

\begin{theorem}[An anomaly obstructs invariant normal states]
    \label{thm:mpug_anomaly_obstruction}
    \lean{MPOTensor.GroupFamily.IsNormalRepresentation.isTrivialGaugeClass_omega_of_invariant,
        MPOTensor.GroupFamily.IsNormalRepresentation.not_exists_invariant_of_not_isTrivialGaugeClass}
    \leanok
    \uses{def:mpug_anomaly_three_cochain, def:mpug_trivial_gauge_class,
        def:mpv_state}
    Let $\{O_g\}$ be matrix product operator tensors with normal
    doubled-index tensors and $O_gO_h=O_{gh}$ on every nonempty chain. If a
    normal matrix product state tensor of positive bond dimension satisfies
    $O_g\ket{V^{(N)}(A)}=\ket{V^{(N)}(A)}$ for every $g$ and every $N\geq1$, then
    the anomaly three-cocycle has trivial gauge class. Equivalently, as in
    \cite[line~1545]{FrancoRubio2025MPUGauging} and
    \cite[line~738]{GarreRubio2023MPOSym} with normal in place of injective
    and without unitarity or simplicity of the operators, a group with
    nontrivial anomaly has no invariant normal matrix product state.
\end{theorem}

\begin{proof}\leanok
    \uses{thm:mpug_single_block_l_symbol_spec,
        thm:mpug_single_block_l_compatibility,
        thm:mpug_fixed_block_trivializes}
    Action tensors exist by
    Lemma~\ref{thm:mpug_single_block_l_symbol_spec}, so
    Lemma~\ref{thm:mpug_single_block_l_compatibility} applies, and Lemma~\ref{thm:mpug_fixed_block_trivializes} for the
    single block gives $\omega=dL$.
\end{proof}

\subsection{Blocks permuted by the group}

Let $\{O_g\}$ be as in Definition~\ref{def:mpug_anomaly_three_cochain}, and
let $G$ act on a set $\mathsf X$. Following
\cite[line~1064]{GarreRubio2023MPOSym}, consider normal matrix product
state tensors $A_x$, $x\in\mathsf X$, of positive bond dimensions, possibly
different for different $x$, with
$O_g\ket{V^{(N)}(A_x)}=\ket{V^{(N)}(A_{gx})}$ for every $g$, $x$ and
$N\geq1$.

\begin{definition}[L-symbols of permuted blocks]
    \label{def:mpug_permuted_block_l_symbol}
    \lean{MPOTensor.GroupFamily.CarriesMPV,
        MPOTensor.GroupFamily.castIndex,
        MPOTensor.GroupFamily.BlockActionData,
        MPOTensor.GroupFamily.BlockActionData.actV,
        MPOTensor.GroupFamily.BlockActionData.fuseV,
        MPOTensor.GroupFamily.BlockActionData.lSymbol}
    \leanok
    \uses{def:mpug_anomaly_three_cochain, def:mpug_l_symbols,
        def:mpug_single_block_l_symbol}
    Action tensors for the family $(A_x)$ are, for every $g$ and $x$, a
    reduction $(V_{g,x},W_{g,x})$ of $O_g\cdot A_x$ onto $A_{gx}$. The stacked
    action $B=(O_gO_h)\cdot A_x$ then reduces onto $A_{g(hx)}=A_{(gh)x}$ along
    two trees, with left boundaries
    \begin{align}
        V_{g,hx}(\mathbf 1\otimes V_{h,x})\,\mathsf a^{-1}, &  &
        V_{gh,x}(F^<_{g,h}\otimes\mathbf 1),
        \notag
    \end{align}
    and $L^x_{g,h}$ is a nonzero scalar relating them against all long
    words of $B$, or $1$ if there is none, as in
    Definition~\ref{def:mpug_single_block_l_symbol}.
\end{definition}

\begin{lemma}[Characterization of the L-symbols of permuted blocks]
    \label{thm:mpug_permuted_block_l_symbol_spec}
    \lean{MPOTensor.GroupFamily.nonempty_blockActionData,
        MPOTensor.GroupFamily.BlockActionData.isDressedProportional_lSymbol,
        MPOTensor.GroupFamily.actTensor_idTensor_apply,
        MPOTensor.GroupFamily.isReduction_actTensor_idTensor}
    \leanok
    \uses{def:mpug_permuted_block_l_symbol}
    Action tensors for the family $(A_x)$ of
    Definition~\ref{def:mpug_permuted_block_l_symbol} exist; for the bond-one
    identity tensor they are the identification of the bond spaces. For
    every choice of them, $L^x_{g,h}$ relates the two left boundaries against
    all long words of $(O_gO_h)\cdot A_x$.
\end{lemma}

\begin{proof}\leanok
    \uses{cor:asym_action_tensors, thm:mpug_action_tensor_reductions,
        thm:mpug_reduction_composition, thm:mpug_dressed_proportionality}
    As in Lemma~\ref{thm:mpug_single_block_l_symbol_spec}: existence is
    Corollary~\ref{cor:asym_action_tensors} applied to $O_g\cdot A_x$ and
    $A_{gx}$, which generate the same vectors. By
    Lemmas~\ref{thm:mpug_action_tensor_reductions}
    and~\ref{thm:mpug_reduction_composition} both boundaries are left
    boundaries of reductions of $(O_gO_h)\cdot A_x$ onto the normal tensor
    $A_{(gh)x}$, so Lemma~\ref{thm:mpug_dressed_proportionality} gives a
    nonzero scalar relating them against long words.
\end{proof}

\begin{lemma}[Compatibility for permuted blocks]
    \label{thm:mpug_permuted_block_l_compatibility}
    \lean{MPOTensor.GroupFamily.BlockActionData.isCompatible_lSymbol,
        MPOTensor.GroupFamily.IsNormalRepresentation.isTrivialGaugeClass_comap_omega_of_fixed}
    \leanok
    \uses{def:mpug_permuted_block_l_symbol, thm:mpug_anomaly_three_cochain_spec,
        thm:mpug_permuted_block_l_symbol_spec,
        thm:mpug_dressed_proportionality, thm:mpug_fixed_block_trivializes}
    For all $x\in\mathsf X$ and $g,h,k\in G$,
    \begin{align}
        L^x_{g,hk}L^x_{h,k}
         & =\omega(g,h,k)\,L^{kx}_{g,h}L^x_{gh,k},
    \end{align}
    which is (\ref{eq:mpug_l_compatibility}). Consequently, as in
    \cite[lines~740--745]{GarreRubio2023MPOSym}, the restriction of
    $\omega$ to any subgroup fixing one block has trivial gauge class.
\end{lemma}

\begin{proof}\leanok
    \uses{thm:mpug_permuted_block_l_symbol_spec,
        thm:mpug_anomaly_three_cochain_spec,
        thm:mpug_fixed_block_trivializes}
    Write
    \begin{align}
        \alpha^x_{g,h}  & =V_{g,hx}(\mathbf 1\otimes V_{h,x})\,\mathsf a^{-1}, &
        \varphi^x_{g,h} & =V_{gh,x}(F^<_{g,h}\otimes\mathbf 1)
        \notag
    \end{align}
    for the two left boundaries of Definition~\ref{def:mpug_permuted_block_l_symbol},
    so that $\alpha^x_{g,h}=L^x_{g,h}\varphi^x_{g,h}$ against long words of
    $(O_gO_h)\cdot A_x$ by Lemma~\ref{thm:mpug_permuted_block_l_symbol_spec}.
    The block labels $h(kx)$ and $(hk)x$ are equal, and so are $g(h(kx))$ and
    $(ghk)x$, and the bond spaces of blocks with equal labels are identified,
    so that $V_{g,h(kx)}=V_{g,(hk)x}$. With this identification, the
    naturality and
    pentagon identities of $\mathsf a$ give, for $B=((O_gO_h)O_k)\cdot A_x$
    and the fusion trees $V^L$, $V^R$ of $(g,h,k)$, the five identities
    \begin{align}
        \alpha^{kx}_{g,h}(\mathbf 1\otimes V_{k,x})\,\mathsf a^{-1}
         & =V_{g,(hk)x}(\mathbf 1\otimes\alpha^x_{h,k})\,\mathsf a^{-1}
        (\mathsf a^{-1}\otimes\mathbf 1),
        \label{eq:mpug_permuted_l_compat_first_boundary}                \\
        V_{g,(hk)x}(\mathbf 1\otimes\varphi^x_{h,k})\,\mathsf a^{-1}
        (\mathsf a^{-1}\otimes\mathbf 1)
         & =\alpha^x_{g,hk}
        \bigl((\mathbf 1\otimes F^<_{h,k})\,\mathsf a^{-1}\otimes\mathbf 1\bigr),
        \notag                                                          \\
        \varphi^{kx}_{g,h}(\mathbf 1\otimes V_{k,x})\,\mathsf a^{-1}
         & =\alpha^x_{gh,k}
        \bigl((F^<_{g,h}\otimes\mathbf 1)\otimes\mathbf 1\bigr),
        \notag                                                          \\
        \varphi^x_{gh,k}\bigl((F^<_{g,h}\otimes\mathbf 1)\otimes\mathbf 1\bigr)
         & =V_{ghk,x}(V^L\otimes\mathbf 1),
        \notag                                                          \\
        \varphi^x_{g,hk}
        \bigl((\mathbf 1\otimes F^<_{h,k})\,\mathsf a^{-1}\otimes\mathbf 1\bigr)
         & =V_{ghk,x}(V^R\otimes\mathbf 1).
        \notag
    \end{align}
    The first factor $\alpha^{kx}_{g,h}$ acts on the block $kx$ reached by
    $O_k$, which is why $L^{kx}_{g,h}$ appears. Replacing, as in the proof of
    Lemma~\ref{thm:mpug_single_block_l_compatibility}, $\alpha^{kx}_{g,h}$,
    $\alpha^x_{gh,k}$ and $V^L$ along one route, and
    (\ref{eq:mpug_permuted_l_compat_first_boundary}), $\alpha^x_{h,k}$ and
    $\alpha^x_{g,hk}$ along the other, gives for all long words
    \begin{align}
        \alpha^{kx}_{g,h}(\mathbf 1\otimes V_{k,x})\,\mathsf a^{-1}B^{\mathbf w}
         & =L^{kx}_{g,h}L^x_{gh,k}\,\omega(g,h,k)\,
        V_{ghk,x}(V^R\otimes\mathbf 1)B^{\mathbf w}
        \notag                                                                \\
         & =L^x_{h,k}L^x_{g,hk}\,V_{ghk,x}(V^R\otimes\mathbf 1)B^{\mathbf w}.
        \notag
    \end{align}
    The boundary $V_{ghk,x}(V^R\otimes\mathbf 1)$ belongs to a reduction onto
    the normal tensor $A_{(ghk)x}$, so it is nonzero against arbitrarily long
    words, and the two scalars agree. The second statement is
    Lemma~\ref{thm:mpug_fixed_block_trivializes}.
\end{proof}

\begin{lemma}[The anomaly at an involution]
    \label{thm:mpug_involution_anomaly_ratio}
    \lean{TNLean.Algebra.LSymbol.eq_div_of_isCompatible_of_mul_self_eq_one}
    \leanok
    \uses{def:mpug_l_symbols}
    If $L$ is compatible with $\omega$ and $g^2=e$, then for every
    $x\in\mathsf X$ the formula of
    \cite[line~5121]{FrancoRubio2025MPUGauging} holds:
    \begin{align}
        \omega(g,g,g)
         & =\frac{L^x_{g,g}L^x_{g,e}}{L^{gx}_{g,g}L^x_{e,g}}.
    \end{align}
\end{lemma}

\begin{proof}\leanok
    Setting $h=k=g$ in (\ref{eq:mpug_l_compatibility}) and using $gg=e$
    gives the ratio.
\end{proof}

\begin{lemma}[The involution ratio from the self-pair and identity L-symbols]
    \label{thm:mpug_involution_anomaly_identity_ratio}
    \lean{TNLean.Algebra.LSymbol.div_eq_mul_of_isCompatible_of_mul_self_eq_one}
    \leanok
    \uses{def:mpug_l_symbols}
    If $L$ is compatible with $\omega$ and $g^2=e$, then for every
    $x\in\mathsf X$
    \begin{align}
        \frac{L^x_{g,g}L^x_{e,e}}{L^{gx}_{g,g}L^{gx}_{e,e}}
         & =\omega(g,g,g)\,\omega(g,e,g).
        \notag
    \end{align}
    For L-symbols and cocycles that are trivial when an argument is the
    identity, this recovers the relation $L_A/L_B=\omega$ of
    \cite[Section~III.C]{GarreRubioSchuch2024DomainWall} in the form
    $L^x_{g,g}/L^{gx}_{g,g}=\omega(g,g,g)$.
\end{lemma}

\begin{proof}\leanok
    \uses{thm:mpug_involution_anomaly_ratio}
    Compatibility at $(g,e,g)$ gives $L^x_{e,g}=\omega(g,e,g)L^{gx}_{g,e}$,
    and at $(g,e,e)$ it gives $L^z_{e,e}=\omega(g,e,e)L^z_{g,e}$ for $z=x$ and
    $z=gx$. Substituting into the formula of
    Lemma~\ref{thm:mpug_involution_anomaly_ratio} gives the claim.
\end{proof}

\subsection{Domain walls between permuted blocks}

Keep the setting of the previous subsection. A domain wall between the
blocks $x$ and $y$ is a family $e=(e^i)_{i}$ of $D_x\times D_y$ matrices,
placed between a region of $A_x$ and a region of $A_y$
\cite[Section~III.B]{GarreRubioSchuch2024DomainWall}. For an operator tensor
$T$ write $(T\cdot e)^i=\sum_j T^{ij}\otimes e^j$.

\begin{definition}[Local action on domain walls]
    \label{def:mpug_domain_wall_action}
    \lean{MPOTensor.actRect,
        MPOTensor.physAct,
        MPOTensor.wallChain,
        MPOTensor.GroupFamily.BlockActionData.IsDomainWallAction}
    \leanok
    \uses{def:mpug_permuted_block_l_symbol}
    Let $e$ be a domain wall between $x$ and $y$, and $e'$ a domain wall
    between $gx$ and $gy$. The element $g$ carries $e$ to $e'$ with phase $c$
    if $e\neq0$, $e'\neq0$, $c\neq0$, and there is $N$ such that
    \begin{align}
        V_{g,x}\,(O_g\cdot A_x)^{\mathbf u}\,(O_g\cdot e)^i\,
        (O_g\cdot A_y)^{\mathbf v}\,W_{g,y}
         & =c\,A_{gx}^{\mathbf u}\,{e'}^{i}\,A_{gy}^{\mathbf v}
        \notag
    \end{align}
    for every letter $i$ and all words $\mathbf u,\mathbf v$ of length at
    least $N$. This is the local relation of
    \cite[Section~III.B]{GarreRubioSchuch2024DomainWall}, drawn there at the
    renormalization fixed point with one site on each side of the wall; the
    words $\mathbf u,\mathbf v$ are the blocked sites of the discussion away
    from the fixed point \cite{gap:gs24_domain_wall_nondegenerate}.
\end{definition}

\begin{lemma}[Fractionalization of the symmetry on domain walls]
    \label{thm:mpug_domain_wall_pentagon}
    \lean{MPOTensor.GroupFamily.BlockActionData.IsDomainWallAction.mul,
        MPOTensor.GroupFamily.BlockActionData.IsDomainWallAction.mul_eq,
        MPOTensor.GroupFamily.BlockActionData.IsDomainWallAction.phase_eq,
        MPOTensor.GroupFamily.BlockActionData.actW,
        MPOTensor.GroupFamily.BlockActionData.fuseW,
        MPOTensor.GroupFamily.BlockActionData.exists_evalWord_mul_actW_eq,
        MPOTensor.physAct_mulTensor,
        MPOTensor.physAct_kronId,
        MPOTensor.physAct_wallChain,
        MPOTensor.exists_evalWord_mul_mul_evalWord_ne_zero}
    \leanok
    \uses{def:mpug_domain_wall_action, def:mpug_permuted_block_l_symbol}
    Suppose $h$ carries the domain wall $e_1$ between $x$ and $y$ to $e_2$
    with phase $B^h_{x,y}$ and $g$ carries $e_2$ to $e_3$ with phase
    $B^g_{hx,hy}$. Then $gh$ carries $e_1$ to $e_3$ with phase
    $(L^y_{g,h}/L^x_{g,h})\,B^g_{hx,hy}B^h_{x,y}$. Consequently, as in
    \cite[Section~IV.B]{GarreRubioSchuch2024DomainWall}, if $gh$
    carries $e_1$ to $e_3$ with phase $B^{gh}_{x,y}$, then
    \begin{align}
        B^g_{hx,hy}\,B^h_{x,y}
         & =\frac{L^x_{g,h}}{L^y_{g,h}}\,B^{gh}_{x,y}.
        \notag
    \end{align}
\end{lemma}

\begin{proof}\leanok
    \uses{thm:mpug_permuted_block_l_symbol_spec, thm:mpug_dressed_proportionality}
    Acting with $O_g$ on the local relation of $h$, which holds for every
    configuration of a chain with long regions, passes the boundary matrices
    $V_{h,x}$, $W_{h,y}$ through as $\mathbf 1\otimes V_{h,x}$ and
    $\mathbf 1\otimes W_{h,y}$, while the action of $O_g$ on the action of
    $O_h$ is the action of $O_gO_h$ up to the bond associator. The local
    relation of $g$ then gives
    \begin{align}
        V_{g,hx}(\mathbf 1\otimes V_{h,x})\,\mathsf a^{-1}\,
        (O_gO_h\cdot\mathsf C)\,\mathsf a\,(\mathbf 1\otimes W_{h,y})\,W_{g,hy}
         & =B^g_{hx,hy}B^h_{x,y}\,\mathsf C_3
        \notag
    \end{align}
    for the chain $\mathsf C=A_x^{\mathbf u}e_1^iA_y^{\mathbf v}$ and the chain
    $\mathsf C_3$ with $e_3$. Against long words the left boundary equals
    $L^x_{g,h}$ times the left boundary through the fusion tensor
    $F^<_{g,h}$; by the two-sided comparison of reductions onto a normal
    tensor, the right boundary equals $(L^y_{g,h})^{-1}$ times the right
    boundary through $F^>_{g,h}$. The fusion tensors reduce
    $O_gO_h\cdot\mathsf C$ to $O_{gh}\cdot\mathsf C$, which gives the local
    action of $gh$ with the stated phase. Since $A_{ghx}$ and $A_{ghy}$ are
    normal and $e_3\neq0$, some chain $\mathsf C_3$ with long regions is
    nonzero, so the phase of a local action is unique, and the two phases of
    $gh$ agree.
\end{proof}

\begin{lemma}[The symmetry on states with two domain walls]
    \label{thm:mpug_domain_wall_state}
    \lean{MPOTensor.twoWallChain,
        MPOTensor.twoWallMPV,
        MPOTensor.mpo_mulVec_trace,
        MPOTensor.GroupFamily.BlockActionData.IsDomainWallAction.mpo_mulVec_twoWallMPV}
    \leanok
    \uses{def:mpug_domain_wall_action}
    Let $\ket{\psi(x\text{-}y\text{-}x)}$ be the periodic state with an $A_x$
    region, a domain wall $e$ between $x$ and $y$, an $A_y$ region, a domain
    wall $f$ between $y$ and $x$, and a second $A_x$ region. If $g$ carries
    $e$ to $e'$ with phase $c$ and $f$ to $f'$ with phase $c'$, then
    \begin{align}
        O_g\ket{\psi(x\text{-}y\text{-}x)}
         & =c\,c'\,\ket{\psi(gx\text{-}gy\text{-}gx)}
        \notag
    \end{align}
    whenever the first $A_x$ region and the $A_y$ region are long enough.
    For $\Z_2$, as in \cite[Section~III.B]{GarreRubioSchuch2024DomainWall},
    this reduces to $U\ket{\psi(A\text{-}B\text{-}A)}
        =c_{AB}c_{BA}\ket{\psi(B\text{-}A\text{-}B)}$.
\end{lemma}

\begin{proof}\leanok
    \uses{thm:mps_reduction_residual_dimension_bound,
        thm:mps_reduction_exterior_buffer}
    The coefficient of $O_g\ket{\psi}$ is the trace of the acted chain.
    Inserting the reduced block $W_{g,z}A_{gz}^{\mathbf c}V_{g,z}$ into each
    of the two acted regions leaves long words unchanged by the residual
    nilpotency of a reduction \cite{Molnar2018SemiInjective}. By cyclicity
    of the trace the chain splits into the two local relations, with phases
    $c$ and $c'$, and the remaining trace is the coefficient of
    $\ket{\psi(gx\text{-}gy\text{-}gx)}$.
\end{proof}

\begin{lemma}[A symmetry string over a pair of domain walls]
    \label{thm:mpug_domain_wall_pair}
    \lean{MPOTensor.GroupFamily.BlockActionData.IsDomainWallAction.pair}
    \leanok
    \uses{def:mpug_domain_wall_action}
    If $g$ carries the domain wall $e$ between $x$ and $y$ to $e'$ with phase
    $c$, and the domain wall $f$ between $y$ and $z$ to $f'$ with phase $c'$,
    then for all words $\mathbf u,\mathbf v,\mathbf w$ longer than a fixed
    buffer
    \begin{align}
        V_{g,x}\,(O_g\cdot A_x)^{\mathbf u}(O_g\cdot e)^i
        (O_g\cdot A_y)^{\mathbf v}(O_g\cdot f)^j(O_g\cdot A_z)^{\mathbf w}\,
        W_{g,z}
         & =c\,c'\,A_{gx}^{\mathbf u}{e'}^{i}A_{gy}^{\mathbf v}{f'}^{j}
        A_{gz}^{\mathbf w}.
        \notag
    \end{align}
    In the comparison of two overlapping strings in
    \cite[Section~III.F]{GarreRubioSchuch2024DomainWall}, this is the
    tensor-level phase acquired by a symmetry string passing over a pair of
    domain walls.
\end{lemma}

\begin{proof}\leanok
    \uses{thm:mps_reduction_residual_dimension_bound,
        thm:mps_reduction_exterior_buffer}
    Splitting the middle word as $\mathbf p\mathbf c\mathbf q$ with long
    outer parts and inserting the reduced block
    $W_{g,y}A_{gy}^{\mathbf c}V_{g,y}$ leaves the word unchanged by the
    residual nilpotency of a reduction \cite{Molnar2018SemiInjective}. The
    two resulting factors are the local relations of $e$ and $f$.
\end{proof}

\begin{definition}[Domain-wall string operators]
    \label{def:mpug_domain_wall_string}
    \lean{MPOTensor.physPairing,
        MPOTensor.leftEndpoint,
        MPOTensor.rightEndpoint,
        MPOTensor.leftAct,
        MPOTensor.rightAct,
        MPOTensor.wallSplit,
        MPOTensor.wallConfig,
        MPOTensor.wallConfig_comp_cast,
        MPOTensor.ofFn_wallConfig,
        MPOTensor.stringKernel,
        MPOTensor.stringOperator,
        MPOTensor.GroupFamily.IsSeparatingLeftInverse,
        MPOTensor.GroupFamily.BlockActionData.wallLeftEndpoint,
        MPOTensor.GroupFamily.BlockActionData.wallRightEndpoint,
        MPOTensor.GroupFamily.BlockActionData.wallString}
    \leanok
    \uses{def:mpug_domain_wall_action}
    Let $g$ exchange the blocks $A=A_x$ and $B=A_y$, with left and right action
    tensors $V_B=V_{g,x}$, $W_B=W_{g,x}$ on $A$ and $V_A=V_{g,y}$,
    $W_A=W_{g,y}$ on $B$, and let $e_{AB}$ and $e_{BA}$ be domain walls. Let
    $\hat A$ and $\hat B$ be left inverses of $A$ and $B$ with separated
    supports,
    \begin{align}
        \sum_\sigma \hat A^\sigma_{\alpha\beta}A^\sigma_{\gamma\delta}
         & =\delta_{\alpha\gamma}\delta_{\beta\delta},
         &
        \sum_\sigma \hat A^\sigma_{\alpha\beta}B^\sigma_{\gamma\delta}
         & =0,
        \notag
    \end{align}
    and likewise for $\hat B$. The left endpoint tensor is
    $V_Be_{AB}\hat A+V_Ae_{BA}\hat B$: its input leg is mapped by the left
    inverse to a pair of virtual indices, which are closed by the domain wall
    and the left action tensor, leaving the output leg and one bond of the
    operator string. The right endpoint tensor is
    $W_Be_{BA}\hat A+W_Ae_{AB}\hat B$. Following
    \cite[Section~III.E]{GarreRubioSchuch2024DomainWall}, the string
    operator $O^{[i,j]}$ on the periodic chain acts by the identity outside
    the sites $i,\dots,j$, by the two endpoint tensors on the sites $i$ and
    $j$, and by the tensor of $O_g$ on the sites between them.
\end{definition}

\begin{lemma}[Domain-wall strings create pairs of domain walls]
    \label{thm:mpug_domain_wall_string_creation}
    \lean{MPOTensor.stringOperator_mulVec_wallConfig,
        MPOTensor.stringOperator_mulVec_trace,
        MPOTensor.mpv_wallConfig,
        MPOTensor.leftAct_mul_physAct_mul_rightAct,
        MPOTensor.leftAct_leftEndpoint,
        MPOTensor.leftAct_leftEndpoint_eq_zero,
        MPOTensor.rightAct_rightEndpoint,
        MPOTensor.rightAct_rightEndpoint_eq_zero,
        MPOTensor.GroupFamily.BlockActionData.wallString_mulVec_mpv_left,
        MPOTensor.GroupFamily.BlockActionData.wallString_mulVec_mpv_right}
    \leanok
    \uses{def:mpug_domain_wall_string, thm:mpug_domain_wall_state}
    On the periodic chain of every length $L\geq2$ and for all sites $i<j$,
    \begin{align}
        O^{[i,j]}\ket{\psi_A}
         & =\ket{\psi(A\text{-}B\text{-}A)},
         &
        O^{[i,j]}\ket{\psi_B}
         & =\ket{\psi(B\text{-}A\text{-}B)},
        \notag
    \end{align}
    with, as in \cite[Section~III.E]{GarreRubioSchuch2024DomainWall}, the
    domain walls $e_{AB}$ at $i$ and $e_{BA}$ at $j$ in the first state and
    $e_{BA}$ at $i$ and $e_{AB}$ at $j$ in the second.
\end{lemma}

\begin{proof}\leanok
    On a state given by the trace of a chain, the string operator acts by
    its endpoint tensors on the letters of the two endpoint sites and by the
    physical action of $O_g$ on the chain between them. On $A$ the left
    inverse $\hat A$ turns the left endpoint into $e_{AB}V_B$ and the right
    endpoint into $W_Be_{BA}$, while the terms with $\hat B$ vanish; the
    action tensors reduce the string of $O_g\cdot A$ between them to $B$.
    On $B$ the terms with $\hat B$ act in the same way.
\end{proof}

\begin{lemma}[Exchange of two domain-wall strings]
    \label{thm:mpug_domain_wall_string_exchange}
    \lean{MPOTensor.wallConfig_resplit,
        MPOTensor.twoWallChain_append,
        MPOTensor.wallConfig_eq_wallConfig,
        MPOTensor.twoWallMPV_wallConfig_comp_cast,
        MPOTensor.fourWallCoeff,
        MPOTensor.exists_wallConfig_wallSplit,
        MPOTensor.GroupFamily.BlockActionData.wallString_outer_mul_inner_mulVec_mpv_apply,
        MPOTensor.GroupFamily.BlockActionData.wallString_inner_mul_outer_mulVec_mpv_apply,
        MPOTensor.GroupFamily.BlockActionData.wallString_mul_wallString_mulVec_mpv}
    \leanok
    \uses{thm:mpug_domain_wall_string_creation, thm:mpug_domain_wall_pair}
    Suppose $g$ carries $e_{AB}$ to $e_{BA}$ with phase $c_{AB}$ and $e_{BA}$
    to $e_{AB}$ with phase $c_{BA}$. For sites $i_2<i_1<j_1<j_2$ whose three
    inner regions are longer than a fixed buffer,
    \begin{align}
        O^{[i_2,j_2]}O^{[i_1,j_1]}\ket{\psi_A}
         & =c_{AB}c_{BA}\,O^{[i_1,j_1]}O^{[i_2,j_2]}\ket{\psi_A}.
        \notag
    \end{align}
    The phase stands on the side where it is acquired; the relation printed
    in \cite[Section~III.F]{GarreRubioSchuch2024DomainWall} has it on the
    other side. The buffer is the blocked form of the local action
    \cite{gap:gs24_domain_wall_nondegenerate}.
\end{lemma}

\begin{proof}\leanok
    Both sides are multiples of the state with the domain walls $e_{AB}$,
    $e_{BA}$, $e_{AB}$, $e_{BA}$ at $i_2,i_1,j_1,j_2$. If the inner string
    acts first, the outer string passes over its two domain walls, and
    Lemma~\ref{thm:mpug_domain_wall_pair} gives the phase $c_{AB}c_{BA}$.
    If the outer string acts first, the inner string acts on the $B$ region
    it created, and Lemma~\ref{thm:mpug_domain_wall_string_creation} gives
    no phase.
\end{proof}

\begin{lemma}[The square of the exchange phase]
    \label{lem:mpug_domain_wall_exchange_phase_square}
    \lean{MPOTensor.exists_twoWallMPV_ne_zero,
        MPOTensor.GroupFamily.BlockActionData.IsDomainWallAction.mul_sq_eq_one_of_mpo_mul_self_eq_one}
    \leanok
    \uses{thm:mpug_domain_wall_state}
    Let $g$ carry $e_{AB}$ to $e_{BA}$ with phase $c_{AB}$ and $e_{BA}$ to
    $e_{AB}$ with phase $c_{BA}$, where $A$ and $B$ are normal, and suppose
    that $U=O_g$ satisfies $U^2=\mathbf 1$ on the periodic chain of every
    length. Then, as in \cite[Section~III.B]{GarreRubioSchuch2024DomainWall},
    $(c_{AB}c_{BA})^2=1$.
\end{lemma}

\begin{proof}\leanok
    By Lemma~\ref{thm:mpug_domain_wall_state},
    $U\ket{\psi(A\text{-}B\text{-}A)}=c_{AB}c_{BA}\ket{\psi(B\text{-}A\text{-}B)}$
    and $U\ket{\psi(B\text{-}A\text{-}B)}=c_{BA}c_{AB}\ket{\psi(A\text{-}B\text{-}A)}$
    whenever the regions are long, so $U^2=\mathbf 1$ gives
    $\ket{\psi(A\text{-}B\text{-}A)}=(c_{AB}c_{BA})^2\ket{\psi(A\text{-}B\text{-}A)}$.
    The state is nonzero for suitable long regions: the words of a fixed
    long length span the full matrix algebras of $A$ and $B$, so if every
    coefficient vanished, then $\operatorname{tr}(Xe_{AB}^iYe_{BA}^j)=0$ for
    all matrices $X$ and $Y$, and choosing matrix units for $X$ and $Y$
    gives $e_{AB}^i=0$ or $e_{BA}^j=0$ for every pair of letters.
\end{proof}

\begin{corollary}[Exchange of two domain-wall strings, printed orientation]
    \label{cor:mpug_domain_wall_string_exchange_printed}
    \lean{MPOTensor.isInjective_of_physPairing_eq_one,
        MPOTensor.GroupFamily.BlockActionData.wallString_mul_wallString_mulVec_mpv_of_mpo_mul_self_eq_one}
    \leanok
    \uses{thm:mpug_domain_wall_string_exchange,
        lem:mpug_domain_wall_exchange_phase_square}
    In the setting of Lemma~\ref{thm:mpug_domain_wall_string_exchange},
    suppose that $U=O_g$ satisfies $U^2=\mathbf 1$ on the periodic chain of
    every length. Then the relation of
    \cite[Section~III.F]{GarreRubioSchuch2024DomainWall} holds:
    \begin{align}
        O^{[i_1,j_1]}O^{[i_2,j_2]}\ket{\psi_A}
         & =c_{AB}c_{BA}\,O^{[i_2,j_2]}O^{[i_1,j_1]}\ket{\psi_A}.
        \notag
    \end{align}
\end{corollary}

\begin{proof}\leanok
    The tensors $A$ and $B$ have left inverses, so their matrices span the
    full matrix algebras and they are normal. Lemma~\ref{lem:mpug_domain_wall_exchange_phase_square}
    gives $(c_{AB}c_{BA})^2=1$. Multiplying the identity of
    Lemma~\ref{thm:mpug_domain_wall_string_exchange} by $c_{AB}c_{BA}$
    gives the result.
\end{proof}

\begin{lemma}[Exchange of two left endpoints of domain-wall strings]
    \label{thm:mpug_domain_wall_endpoint_exchange}
    \lean{MPOTensor.GroupFamily.BlockActionData.wallString_cross_over_right_mulVec_mpv_apply,
        MPOTensor.GroupFamily.BlockActionData.wallString_cross_over_left_mulVec_mpv_apply,
        MPOTensor.GroupFamily.BlockActionData.wallString_mul_wallString_mulVec_mpv_swap_left_far,
        MPOTensor.GroupFamily.BlockActionData.wallString_mul_wallString_mulVec_mpv_swap_left_near}
    \leanok
    \uses{thm:mpug_domain_wall_string_exchange}
    In the setting of Lemma~\ref{thm:mpug_domain_wall_string_exchange},
    let $j<i$ and let $r_1\neq r_2$ be two sites to the right of $i$. If the
    regions between the four sites $j,i,r_1,r_2$ are longer than a fixed
    buffer, then
    \begin{align}
        O^{[j,r_1]}O^{[i,r_2]}\ket{\psi_A}
         & =c_{AB}\,O^{[i,r_1]}O^{[j,r_2]}\ket{\psi_A}.
        \notag
    \end{align}
    This is the exchange of two left endpoints
    $O^{[j]}_xO^{[i]}_y\ket{\psi_A}=c_{AB}\,O^{[i]}_xO^{[j]}_y\ket{\psi_A}$
    of \cite[Section~III.F]{GarreRubioSchuch2024DomainWall}, with the
    half-chain operators glued with the right endpoints of their strings at
    $r_1$ and $r_2$; the identity of half-chain objects with open legs is not
    formalized \cite{gap:gs24_domain_wall_half_string_gluing}.
\end{lemma}

\begin{proof}\leanok
    Every side is a multiple of the state with the domain walls $e_{AB}$,
    $e_{BA}$, $e_{AB}$, $e_{BA}$ at the four sites in increasing order. For
    $r_2<r_1$ the left side is a nested pair whose outer string acts second,
    with phase $c_{AB}c_{BA}$ by
    Lemma~\ref{thm:mpug_domain_wall_string_exchange}; on the right side the
    string on $[i,r_1]$ starts in the $B$ region, passes over the domain wall
    $e_{BA}$ at $r_2$ and ends in the $A$ region, so the local action gives
    $c_{BA}$. For $r_1<r_2$ the string on $[j,r_1]$ on the left side passes
    over the domain wall $e_{AB}$ at $i$, with phase $c_{AB}$, and the right
    side is a nested pair whose inner string acts second, with no phase.
\end{proof}

\begin{lemma}[Rescaling domain walls]
    \label{thm:mpug_domain_wall_rescaling}
    \lean{MPOTensor.actRect_smul,
        MPOTensor.GroupFamily.BlockActionData.IsDomainWallAction.smul_source,
        MPOTensor.GroupFamily.BlockActionData.IsDomainWallAction.smul_target,
        MPOTensor.GroupFamily.BlockActionData.IsDomainWallAction.exists_eq_of_mul_self}
    \leanok
    \uses{def:mpug_domain_wall_action}
    If $g$ carries $e$ to $e'$ with phase $c$ and $a\neq0$, then it carries $a\,e$
    to $e'$ with phase $ac$, and $e$ to $a\,e'$ with phase $c/a$. In
    particular, as in \cite[Section~III.F]{GarreRubioSchuch2024DomainWall},
    if an involution $g$ exchanges $e_{AB}$ and $e_{BA}$ with phases $c_{AB}$
    and $c_{BA}$, then for every $s$ with $s^2=c_{AB}c_{BA}$ the walls
    $e_{AB}$ and $(c_{AB}/s)\,e_{BA}$ are exchanged with both phases equal
    to $s$.
\end{lemma}

\begin{proof}\leanok
    The rectangular action is linear in the domain wall, and rescaling
    $e_{BA}$ by $c_{AB}/s$ gives the last claim.
\end{proof}

\begin{definition}[A family of domain walls permuted by the group]
    \label{def:mpug_domain_wall_family}
    \lean{MPOTensor.GroupFamily.BlockActionData.IsDomainWallFamily,
        MPOTensor.GroupFamily.BlockActionData.IsDomainWallAction.congr_target}
    \leanok
    \uses{def:mpug_domain_wall_action}
    Following \cite[Section~IV.B]{GarreRubioSchuch2024DomainWall}, with
    nonzero walls and phases by convention
    \cite{gap:gs24_domain_wall_nondegenerate}, a family of domain walls
    permuted by the group is a nonzero domain wall $e_{yz}$ between $y$ and
    $z$ for every ordered pair of blocks, together with nonzero phases
    $B^g_{y,z}$ such that every $g$ carries $e_{yz}$ to $e_{gy,gz}$ with
    phase $B^g_{y,z}$.
\end{definition}

\begin{lemma}[The interchange phase of a family of domain walls]
    \label{thm:mpug_domain_wall_interchange}
    \lean{MPOTensor.GroupFamily.BlockActionData.IsDomainWallFamily.mul_eq,
        MPOTensor.GroupFamily.BlockActionData.IsDomainWallFamily.prod_eq_prod_lSymbol,
        MPOTensor.GroupFamily.BlockActionData.IsDomainWallFamily.prod_eq_inv_cyclicInvariant,
        MPOTensor.GroupFamily.BlockActionData.IsDomainWallFamily.prod_eq_domainWallPhase,
        TNLean.Algebra.LSymbol.prod_div_eq_inv_cyclicInvariant,
        TNLean.Algebra.LSymbol.prod_Ico_div_eq_inv_cyclicInvariant}
    \leanok
    \uses{def:mpug_domain_wall_family, def:mpug_permuted_block_l_symbol,
        def:mpug_anomaly_three_cochain}
    As in \cite[Section~IV.C]{GarreRubioSchuch2024DomainWall}, for a family
    of domain walls permuted by the group, every $g$ with $g^n=e$, and every
    block $x$, the interchange phase depends only on the anomaly:
    \begin{align}
        \prod_{k=0}^{n-1}B^g_{g^k(gx),\,g^kx}
         & =\prod_{k=0}^{n-1}\frac{L^{gx}_{g,g^k}}{L^x_{g,g^k}}
        =\Bigl(\prod_{k=0}^{n-1}\omega(g,g^k,g)\Bigr)^{-1}.
        \notag
    \end{align}
    The middle product includes the factor $k=0$, $L^{gx}_{g,e}/L^x_{g,e}$,
    which is one for normalized L-symbols; the source's middle product runs
    over $1\leq k<n$, and for normalized L-symbols that product is again
    $\bigl(\prod_{k=0}^{n-1}\omega(g,g^k,g)\bigr)^{-1}$.
\end{lemma}

\begin{proof}\leanok
    \uses{thm:mpug_domain_wall_pentagon, thm:mpug_permuted_block_l_compatibility}
    With $\beta_k=B^{g^k}_{gx,x}$,
    Lemma~\ref{thm:mpug_domain_wall_pentagon} for $(g,g^k)$ gives
    $B^g_{g^k(gx),g^kx}\,\beta_k=(L^{gx}_{g,g^k}/L^x_{g,g^k})\,\beta_{k+1}$.
    Multiplying over $k<n$ and using $\beta_n=\beta_0$ gives the first
    equality. Compatibility at $(g,g^k,g)$ gives
    $L^{gx}_{g,g^k}/L^x_{g,g^k}=\omega(g,g^k,g)^{-1}
        (L^x_{g,g^{k+1}}/L^x_{g,g^k})(L^x_{g^{k+1},g}/L^x_{g^k,g})^{-1}$,
    and both quotients telescope to one over $k<n$.
\end{proof}

\begin{lemma}[The stabilizer of two blocks acts projectively]
    \label{thm:mpug_domain_wall_projective}
    \lean{TNLean.Algebra.LSymbol.ratioCocycle,
        TNLean.Algebra.LSymbol.ratioCocycle_isMulCocycle₂,
        MPOTensor.GroupFamily.BlockActionData.IsDomainWallFamily.mul_eq_of_fixed,
        MPOTensor.GroupFamily.BlockActionData.IsDomainWallFamily.mul_eq_of_fixed_self}
    \leanok
    \uses{def:mpug_domain_wall_family, def:mpug_permuted_block_l_symbol}
    Let $H_{y,z}$ be the subgroup of elements fixing the blocks $y$ and $z$.
    For $a,b\in H_{y,z}$,
    \begin{align}
        B^a_{y,z}B^b_{y,z}
         & =\alpha(a,b)\,B^{ab}_{y,z},
         &
        \alpha(a,b)
         & =\frac{L^y_{a,b}}{L^z_{a,b}},
        \notag
    \end{align}
    and $\alpha$ is a two-cocycle on $H_{y,z}$; for $y=z$ the phases form a
    linear representation. The source states this for the subgroup fixing
    every block \cite[Section~IV.C]{GarreRubioSchuch2024DomainWall}, which is
    contained in $H_{y,z}$; see \cite{gap:gs24_unbroken_subgroup_stabilizer}.
\end{lemma}

\begin{proof}\leanok
    \uses{thm:mpug_domain_wall_pentagon, thm:mpug_permuted_block_l_compatibility}
    The first identity is Lemma~\ref{thm:mpug_domain_wall_pentagon} with
    $by=y$ and $bz=z$. Dividing the compatibility relations at $y$ and at $z$
    for $(a,b,c)$ with $c\in H_{y,z}$ cancels $\omega(a,b,c)$ and gives the
    two-cocycle equation.
\end{proof}

\begin{lemma}[Domain walls of an involution detect the anomaly]
    \label{thm:mpug_domain_wall_involution}
    \lean{MPOTensor.GroupFamily.BlockActionData.IsDomainWallAction.mul_eq_lSymbol_of_mul_self_eq_one,
        MPOTensor.GroupFamily.BlockActionData.IsDomainWallAction.mul_eq_omega_of_mul_self_eq_one,
        MPOTensor.GroupFamily.BlockActionData.IsDomainWallAction.mul_eq_mul_of_mul_self_eq_one,
        MPOTensor.GroupFamily.BlockActionData.IsDomainWallAction.mpo_mulVec_twoWallMPV_eq_omega,
        MPOTensor.GroupFamily.BlockActionData.IsDomainWallAction.congr_elem}
    \leanok
    \uses{def:mpug_domain_wall_action, def:mpug_permuted_block_l_symbol,
        def:mpug_anomaly_three_cochain}
    Let $g^2=e$ exchange the blocks $x$ and $y$, let $g$ carry the domain wall
    $e_{AB}$ between $x$ and $y$ to $e_{BA}$ with phase $c_{AB}$ and
    $e_{BA}$ to $e_{AB}$ with phase $c_{BA}$. Then
    \begin{align}
        c_{AB}\,c_{BA}
         & =\frac{L^x_{g,g}L^x_{e,e}}{L^y_{g,g}L^y_{e,e}}
        =\omega(g,g,g)\,\omega(g,e,g),
        \notag
    \end{align}
    so $c_{AB}c_{BA}$ does not depend on the action tensors, the fusion
    tensors or the domain walls, and
    $U\ket{\psi(A\text{-}B\text{-}A)}=\omega(g,g,g)\,\omega(g,e,g)
        \ket{\psi(B\text{-}A\text{-}B)}$. When the tensors of the identity are
    trivial, as in \cite[Section~III.C]{GarreRubioSchuch2024DomainWall},
    $L_{e,e}=1$ and $\omega(g,e,g)=1$, which yields
    $c_{AB}c_{BA}=L_A/L_B=\omega$.
\end{lemma}

\begin{proof}\leanok
    \uses{thm:mpug_domain_wall_pentagon, thm:mpug_domain_wall_state,
        thm:mpug_permuted_block_l_compatibility,
        thm:mpug_involution_anomaly_identity_ratio}
    Applying $g$ twice, Lemma~\ref{thm:mpug_domain_wall_pentagon} shows that
    $g^2=e$ carries $e_{AB}$ to itself with phase
    $b=(L^y_{g,g}/L^x_{g,g})\,c_{BA}c_{AB}\neq0$. The same theorem for $(e,e)$
    gives $b^2=(L^x_{e,e}/L^y_{e,e})\,b$, so $b=L^x_{e,e}/L^y_{e,e}$, which is
    the first equality. The second equality is
    Lemma~\ref{thm:mpug_involution_anomaly_identity_ratio} for the compatible
    L-symbols of Lemma~\ref{thm:mpug_permuted_block_l_compatibility}. Its
    right-hand side involves only the fusion tensors, and it is invariant
    under their gauge action, so the product is independent of all choices.
    Finally, by Lemma~\ref{thm:mpug_domain_wall_state},
    $U\ket{\psi(A\text{-}B\text{-}A)}=c_{AB}c_{BA}\ket{\psi(B\text{-}A\text{-}B)}
        =\omega(g,g,g)\,\omega(g,e,g)\ket{\psi(B\text{-}A\text{-}B)}$.
\end{proof}

\subsection{Detecting the anomaly with truncated symmetries}

\begin{lemma}[The commutator of two unitaries on a vector]
    \label{lem:mpug_truncation_commutator}
    \lean{Matrix.groupCommutator,
        Matrix.groupCommutator_mulVec_eq_smul,
        Matrix.star_dotProduct_groupCommutator_mulVec}
    \leanok
    Let $U_1$ and $U_2$ be unitary, let $\psi$ be a vector and $c$ a scalar
    with $U_2U_1\psi=c\,U_1U_2\psi$. Then
    \begin{align}
        U_2^\dagger U_1^\dagger U_2U_1\psi
         & =c\,\psi,
         &
        \langle\psi|\,U_2^\dagger U_1^\dagger U_2U_1\,|\psi\rangle
         & =c\,\langle\psi|\psi\rangle.
        \notag
    \end{align}
    For two truncations $U_1=U^{[i_1,j_1]}$ and $U_2=U^{[i_2,j_2]}$ of a
    symmetry whose exchange relation on a ground state has the phase
    $c=\omega$, this is the last step of the detection of the anomaly in
    \cite[Section~V]{GarreRubioSchuch2024DomainWall}: the expectation value of
    the commutator measures the exchange phase.
\end{lemma}

\begin{proof}\leanok
    Substituting the exchange relation and canceling $U_1^\dagger U_1$ and
    $U_2^\dagger U_2$ gives both identities.
\end{proof}

\begin{lemma}[Truncations of an on-site symmetry]
    \label{thm:mpug_onsite_truncation_commutator}
    \lean{Matrix.onSiteTruncation,
        Matrix.onSiteTruncation_univ,
        Matrix.onSiteTruncation_conjTranspose_mul_self,
        Matrix.groupCommutator_onSiteTruncation,
        Matrix.star_dotProduct_groupCommutator_onSiteTruncation_mulVec}
    \leanok
    Let $u$ satisfy $u^\dagger u=\mathbf 1$ and, for a set $S$ of sites of a
    chain, let $U^S=\bigotimes_{i\in S}u^{[i]}$ act by $u$ on the sites of $S$
    and by the identity elsewhere. For any two sets $S_1$ and $S_2$,
    \begin{align}
        (U^{S_2})^\dagger(U^{S_1})^\dagger U^{S_2}U^{S_1}
         & =\mathbf 1,
        \notag
    \end{align}
    so, as in \cite[Section~V]{GarreRubioSchuch2024DomainWall} for the
    trivial anomaly of on-site symmetries, the expectation value of the
    commutator on every normalized state is $1$.
\end{lemma}

\begin{proof}\leanok
    The four operators are tensor products over the sites. On a site in both
    sets the product of the four factors is
    $u^\dagger u^\dagger uu=\mathbf 1$, on a site in one set it is
    $u^\dagger u=\mathbf 1$, and elsewhere all four factors are the
    identity.
\end{proof}

\begin{lemma}[The detector depends on the endpoints of the truncations]
    \label{thm:mpug_truncation_detector_endpoint_dependence}
    \lean{Matrix.finKronecker_ite_smul_one,
        Matrix.dressed_onSiteTruncation_conjTranspose_mul_self,
        Matrix.groupCommutator_dressed_onSiteTruncation_pauliX,
        Matrix.star_dotProduct_groupCommutator_dressed_onSiteTruncation_pauliX}
    \leanok
    Let $X$ and $Z$ be the Pauli matrices, and for sets of sites $S_1$, $S_2$
    of a chain and sites $j_1\in S_2$, $j_2\notin S_1$ with $j_1\neq j_2$ put
    $U_1=Z_{j_1}X^{S_1}$ and $U_2=Z_{j_2}X^{S_2}$. Then $U_1$ and $U_2$ are
    unitary and
    \begin{align}
        U_2^\dagger U_1^\dagger U_2U_1
         & =-\mathbf 1,
        \notag
    \end{align}
    so the expectation value of the commutator on every normalized state is
    $-1$, although the on-site symmetry $X^{\otimes N}$ has trivial anomaly.
\end{lemma}

\begin{proof}\leanok
    The four operators are tensor products over the sites. At $j_1$ the
    product of the four factors is $X(XZ)X(ZX)=ZXZX=-\mathbf 1$ if $j_1\in S_1$,
    and $XZXZ=-\mathbf 1$ otherwise. At $j_2$ only $U_2$ acts, and the product is
    $(ZX)^\dagger(ZX)=\mathbf 1$ or $Z^\dagger Z=\mathbf 1$. At every other site
    the factors are $X$ or the identity, and they cancel.
\end{proof}

\begin{remark}[Circuit truncations with a nontrivial detector]
    \label{rem:mpug_truncation_detector_circuit}
    \uses{thm:mpug_truncation_detector_endpoint_dependence}
    The symmetry $X^{\otimes N}$ on a periodic chain of even length is the
    depth-two circuit whose first layer applies $X\otimes ZX$ to the pairs
    $(2k,2k+1)$ and whose second layer applies $Z\otimes\mathbf 1$ to the pairs
    $(2k+1,2k+2)$, since $Z\cdot ZX=X$. For $i$ even and $j$ odd, dropping the
    gates not contained in $[i,j]$ leaves both factors on every odd site of
    $[i,j]$ except $j$, which keeps only $ZX$, so the truncation is
    $Z_jX^{\otimes[i,j]}$. For $i_2<i_1<j_1<j_2$ of these parities,
    Lemma~\ref{thm:mpug_truncation_detector_endpoint_dependence} applies.
    The expectation value of the commutator of two circuit truncations on a
    symmetry-broken ground state is therefore not determined by the anomaly,
    and the detection asserted in
    \cite[Section~V]{GarreRubioSchuch2024DomainWall} fails for an arbitrary
    circuit truncation \cite{gap:gs24_truncation_detector_endpoint_dependence}.
    The identification of the circuit truncation with $Z_jX^{\otimes[i,j]}$
    is not formalized.
\end{remark}

\begin{lemma}[Detection for truncations acting as domain-wall strings]
    \label{thm:mpug_truncation_detector_wall_strings}
    \lean{MPOTensor.GroupFamily.BlockActionData.star_dotProduct_groupCommutator_mulVec_mpv_of_wallString,
        MPOTensor.GroupFamily.BlockActionData.star_dotProduct_groupCommutator_mulVec_mpv_eq_omega}
    \leanok
    \uses{lem:mpug_truncation_commutator, thm:mpug_domain_wall_string_exchange,
        thm:mpug_domain_wall_involution}
    In the setting of Lemma~\ref{thm:mpug_domain_wall_string_exchange}, let
    $O_1=O^{[i_1,j_1]}$ and $O_2=O^{[i_2,j_2]}$ be domain-wall strings on
    sites $i_2<i_1<j_1<j_2$ whose three inner regions are longer than a fixed
    buffer, and let $U_1$ and $U_2$ be unitary operators with
    \begin{align}
        U_2U_1\ket{\psi_A} & =O_2O_1\ket{\psi_A}, &
        U_1U_2\ket{\psi_A} & =O_1O_2\ket{\psi_A}.
        \notag
    \end{align}
    Then
    \begin{align}
        \bra{\psi_A}U_2^\dagger U_1^\dagger U_2U_1\ket{\psi_A}
         & =c_{AB}c_{BA}\braket{\psi_A}{\psi_A}.
        \notag
    \end{align}
    If in addition $g$ is an involution of a normal representation, then
    $c_{AB}c_{BA}=\omega(g,g,g)\,\omega(g,e,g)$. This is the detection of the
    anomaly of \cite[Section~V]{GarreRubioSchuch2024DomainWall} restricted to
    unitaries whose two ordered products on $\ket{\psi_A}$ agree with the
    corresponding products of the domain-wall strings. These identities are
    expected for truncations whose walls on $\ket{\psi_B}$ are the images under
    $U$ of their walls on $\ket{\psi_A}$, the property the source's argument
    uses; that implication is not formalized. By
    Lemma~\ref{thm:mpug_truncation_detector_endpoint_dependence} some
    restriction of this kind is needed
    \cite{gap:gs24_truncation_detector_endpoint_dependence}.
\end{lemma}

\begin{proof}\leanok
    \uses{lem:mpug_truncation_commutator, thm:mpug_domain_wall_string_exchange,
        thm:mpug_domain_wall_involution}
    By Lemma~\ref{thm:mpug_domain_wall_string_exchange} and the hypotheses,
    $U_2U_1\ket{\psi_A}=c_{AB}c_{BA}\,U_1U_2\ket{\psi_A}$, and
    Lemma~\ref{lem:mpug_truncation_commutator} gives the first identity. The
    second is Lemma~\ref{thm:mpug_domain_wall_involution}.
\end{proof}

\subsection{Normalization of the first source factor}

\begin{lemma}[The unsimplified first-cut contraction]
    \label{thm:mpug_raw_first_cut_contraction}
    \lean{MPOTensor.sourceCutM₁_adjointSimpleContraction_raw}
    \leanok
    \uses{def:mpu_source_matrix_one, def:mpu_source_weight,
        def:mpug_adjoint_simple_contraction}
    For an arbitrary tensor $A$ of physical dimension $d$ and bond
    dimension $D$, and an arbitrary $D\times D$ matrix $\rho$, put
    $M=M_1(A)$ and $W=I_d\otimes\rho$. Then
    \begin{align}
        M_1(s(A,\rho)) & =MM^\dagger W^{\mathsf T}M.
        \label{eq:mpug_raw_first_cut_contraction}
    \end{align}
    No positivity, canonical-form, or simplicity hypothesis is needed.
\end{lemma}
\begin{proof}\leanok
    \uses{def:mpu_source_matrix_one, def:mpu_source_weight,
        def:mpug_adjoint_simple_contraction}
    In the first-cut coordinates, multiplication and the physical
    Kronecker delta in $W^{\mathsf T}$ give
    \begin{align}
        (MM^\dagger W^{\mathsf T}M)_{(i,b),(a,j)}
         & =\sum_{p,q=0}^{d-1}\sum_{x,r,t=0}^{D-1}
        A^{i,p}_{xb}\,\overline{A^{q,p}_{xr}}\,
        A^{q,j}_{at}\,\rho_{tr}                    \\
         & =s(A,\rho)^{i,j}_{ab}.
    \end{align}
    The transpose preserves the boundary entry $\rho_{tr}$ in the
    defining contraction.
\end{proof}

\begin{lemma}[The first cut of the adjoint-simple contraction]
    \label{thm:mpug_first_cut_adjoint_contraction}
    \lean{MPOTensor.IsMPUCanonicalFormII.sourceCutM₁_adjointSimpleContraction}
    \leanok
    \discussion{7317}
    \uses{def:mpu_canonical_form_ii, def:mpu_weighted_source_factors,
        def:mpu_source_matrix_one, def:mpug_adjoint_simple_contraction}
    Let $\mathcal U$ be an MPU in canonical form II with recorded fixed
    matrix $\rho$, and let $X_1,Y_1$ be its chosen source factors for this
    same $\rho$. Construct $s$ from $\mathcal U,\rho$ as in
    Definition~\ref{def:mpug_adjoint_simple_contraction}. With the first
    cut convention $(M_1(A))_{(i,b),(a,j)}=A^{i,j}_{ab}$,
    \begin{align}
        M_1(s) & =X_1(Y_1Y_1^\dagger)Y_1.
        \label{eq:mpug_first_cut_adjoint_contraction}
    \end{align}
    In the source-factor expansion of
    \cite[lines~1326--1374]{FrancoRubio2025MPUGauging}, this factorization
    requires no simplicity assumption.
\end{lemma}

\begin{proof}\leanok
    \uses{thm:mpug_raw_first_cut_contraction,
        thm:mpu_source_factorizations, thm:mpu_source_factor_normalizations,
        def:mpu_weighted_source_factors, def:mpu_source_weight}
    With $W=I_d\otimes\rho$, diagonality gives $W^{\mathsf T}=W$.
    The cut factorization and weighted isometry then yield
    \begin{align}
        M_1(s)
         & = (X_1Y_1)(X_1Y_1)^\dagger W(X_1Y_1)       \\
         & =X_1(Y_1Y_1^\dagger)(X_1^\dagger W X_1)Y_1 \\
         & =X_1(Y_1Y_1^\dagger)Y_1.
    \end{align}
\end{proof}

\begin{lemma}[The first source Gram matrix is scalar]
    \label{thm:mpug_source_y_one_scalar_gram}
    \lean{MPOTensor.IsMPUCanonicalFormII.sourceY₁_mul_conjTranspose_eq_smul}
    \leanok
    \uses{def:mpu_canonical_form_ii, def:mpu_simple,
        def:mpdo_physical_adjoint, def:mpu_weighted_source_factors}
    Let $\mathcal U$ be an MPU in canonical form II, and suppose that
    both $\mathcal U$ and $\mathcal U^\dagger$ are simple. Construct the
    source factors from $\mathcal U$ and its recorded fixed matrix
    $\rho$. If $d_r$ is the first source-cut rank, then, as in the scalar
    Gram identity of \cite[lines~1375--1397]{FrancoRubio2025MPUGauging},
    there is a scalar $\delta\in\C$ such that
    \begin{align}
        Y_1Y_1^\dagger & =\delta I_{d_r}.
    \end{align}
\end{lemma}

\begin{proof}\leanok
    \uses{thm:mpug_adjoint_simple_scalar,
        thm:mpug_first_cut_adjoint_contraction,
        thm:mpu_source_factor_normalizations,
        thm:mpu_source_factor_right_inverses,
        thm:mpu_source_factorizations}
    By Lemma~\ref{thm:mpug_adjoint_simple_scalar}, $s=\delta\mathcal U$.
    Applying the first cut and
    Lemma~\ref{thm:mpug_first_cut_adjoint_contraction} gives
    \begin{align}
        X_1(Y_1Y_1^\dagger)Y_1 & =\delta X_1Y_1.
    \end{align}
    Multiplying on the left by $X_1^\dagger(I_d\otimes\rho)$ and on the
    right by $Z_1$ cancels the two outer factors via
    $X_1^\dagger(I_d\otimes\rho)X_1=I_{d_r}$ and $Y_1Z_1=I_{d_r}$.
\end{proof}

\begin{lemma}[Trace of the first source Gram matrix]
    \label{thm:mpug_source_y_one_gram_trace}
    \lean{MPOTensor.IsMPUCanonicalFormII.trace_sourceY₁_mul_conjTranspose}
    \leanok
    \uses{def:mpu_canonical_form_ii, def:mpu_simple,
        def:mpu_weighted_source_factors}
    Let $\mathcal U$ be a simple MPU in canonical form II, with physical
    dimension $d$, bond dimension $D$, first source-cut rank $d_r$, and
    recorded fixed matrix $\rho$. As in the physical-trace calculation of
    \cite[line~1398]{FrancoRubio2025MPUGauging}, which does not require
    simplicity of $\mathcal U^\dagger$, its chosen first source factor
    satisfies
    \begin{align}
        \operatorname{tr}(Y_1Y_1^\dagger) & =d.
    \end{align}
\end{lemma}

\begin{proof}\leanok
    \uses{thm:mpug_source_y_physical_contractions}
    Cyclic invariance gives
    $\operatorname{tr}(Y_1Y_1^\dagger)=\operatorname{tr}(Y_1^\dagger Y_1)$.
    Grouping the latter sum by its physical index gives
    \begin{align}
        \operatorname{tr}(Y_1^\dagger Y_1)
         & =\sum_{p=0}^{d-1}\sum_{r=0}^{d_r-1}\sum_{\alpha=0}^{D-1}
        \overline{(Y_1)_{r,(\alpha,p)}}(Y_1)_{r,(\alpha,p)}
        =\sum_{p=0}^{d-1}1=d
    \end{align}
    by Lemma~\ref{thm:mpug_source_y_physical_contractions}.
\end{proof}

\begin{lemma}[Normalization of the first source Gram matrix]
    \label{thm:mpug_source_y_one_normalization}
    \lean{MPOTensor.IsMPUCanonicalFormII.sourceY₁_mul_conjTranspose}
    \leanok
    \discussion{7317}
    \uses{def:mpu_canonical_form_ii, def:mpu_simple,
        def:mpdo_physical_adjoint, def:mpu_weighted_source_factors}
    Let $\mathcal U$ be an MPU in canonical form II of physical dimension
    $d$, and suppose that both $\mathcal U$ and $\mathcal U^\dagger$ are
    simple. Let $Y_1$ be the first source factor constructed from
    $\mathcal U$ and its recorded fixed matrix $\rho$, and let $d_r$ be
    the first source-cut rank. Then, as in the source-factor normalization
    proposition of \cite[lines~1064--1119]{FrancoRubio2025MPUGauging},
    \begin{align}
        Y_1Y_1^\dagger & =\frac{d}{d_r}I_{d_r}.
    \end{align}
    % Source: first identity of eq:MPUnice3. The other three constants
    % and the full prop:MPUsplus are not established by this entry.
\end{lemma}

\begin{proof}\leanok
    \uses{thm:mpug_source_y_one_scalar_gram,
        thm:mpug_source_y_one_gram_trace}
    Writing $Y_1Y_1^\dagger=\delta I_{d_r}$ by
    Lemma~\ref{thm:mpug_source_y_one_scalar_gram} and taking traces via
    Lemma~\ref{thm:mpug_source_y_one_gram_trace} gives $\delta d_r=d$.
    Since canonical form II gives $d>0$ and hence $d_r\ne0$, dividing as in
    \cite[line~1398]{FrancoRubio2025MPUGauging} yields $\delta=d/d_r$.
\end{proof}

\subsection{Normalization of the weighted second source factor}

\begin{lemma}[The first cut of the reflected contraction]
    \label{thm:mpug_adjoint_cut_second_gram}
    \lean{MPOTensor.IsMPUCanonicalFormII.sourceCutM₁_adjointSimpleContraction_physicalAdjoint}
    \leanok
    \discussion{7317}
    \uses{def:mpu_canonical_form_ii, def:mpu_weighted_source_factors,
        def:mpu_source_matrix_one, def:mpug_adjoint_simple_contraction,
        def:mpdo_physical_adjoint, def:mpu_source_weight}
    Let $\mathcal U$ be an MPU in canonical form II with recorded fixed
    matrix $\rho$, and let $X_2,Y_2$ be its chosen second source factors.
    Set $W=I_d\otimes\rho$ and let $s_\dagger$ be the contraction of
    Definition~\ref{def:mpug_adjoint_simple_contraction}, applied to
    $\mathcal U^\dagger$ and this same $\rho$. Then, in the reflected
    factorization for $Y_2(I_d\otimes\rho)Y_2^\dagger$ in the normalization
    proposition of \cite[lines~1064--1119 and 1398]{FrancoRubio2025MPUGauging},
    \begin{align}
        M_1(s_\dagger) & =Y_2^\dagger(Y_2WY_2^\dagger)X_2^\dagger,
    \end{align}
    without any simplicity hypothesis.
\end{lemma}

\begin{proof}\leanok
    \uses{thm:mpug_raw_first_cut_contraction,
        thm:mpu_physical_adjoint_source_cuts,
        thm:mpu_source_factorizations, thm:mpu_source_factor_normalizations,
        def:mpu_weighted_source_factors, def:mpu_source_weight}
    Lemma~\ref{thm:mpug_raw_first_cut_contraction} gives
    \begin{align}
        M_1(s(A,\rho)) & =M_1(A)M_1(A)^\dagger W^{\mathsf T}M_1(A).
    \end{align}
    The transpose records the boundary coefficient $\rho_{tr}$. Here
    $\rho$ is diagonal, so $W^{\mathsf T}=W$. Substituting
    $M_1(\mathcal U^\dagger)=M_2(\mathcal U)^\dagger=Y_2^\dagger X_2^\dagger$
    and canceling $X_2^\dagger X_2=I$ gives the factorization. All factors
    are those of $\mathcal U$, not newly chosen factors of its adjoint.
\end{proof}

\begin{lemma}[The weighted second source Gram matrix is scalar]
    \label{thm:mpug_source_y_two_scalar_gram}
    \lean{MPOTensor.IsMPUCanonicalFormII.sourceY₂_weighted_mul_conjTranspose_eq_smul}
    \leanok
    \uses{def:mpu_canonical_form_ii, def:mpu_simple,
        def:mpdo_physical_adjoint, def:mpu_weighted_source_factors,
        def:mpu_source_weight}
    Let $\mathcal U$ be an MPU in canonical form II with recorded fixed
    matrix $\rho$, and assume that both $\mathcal U$ and
    $\mathcal U^\dagger$ are simple. Let $Y_2$ be its chosen second
    source factor and $d_l$ its second source-cut rank. There is a scalar
    $\varepsilon\in\C$ such that
    \begin{align}
        Y_2(I_d\otimes\rho)Y_2^\dagger & =\varepsilon I_{d_l}.
    \end{align}
\end{lemma}

\begin{proof}\leanok
    \uses{def:mpug_adjoint_canonical_boundaries,
        thm:mpug_adjoint_simple_scalar, thm:mpu_physical_adjoint_identities,
        thm:mpug_adjoint_cut_second_gram,
        thm:mpu_physical_adjoint_source_cuts,
        thm:mpu_source_factorizations, thm:mpu_source_factor_normalizations,
        thm:mpu_source_factor_right_inverses}
    Passing canonical form II to $\mathcal U^\dagger$ with the same
    $\rho$, Lemma~\ref{thm:mpug_adjoint_simple_scalar} gives
    $s_\dagger=\varepsilon\mathcal U^\dagger$, since physical adjunction
    is involutive and both required simplicity hypotheses hold.
    Taking the first cut and applying
    Lemma~\ref{thm:mpug_adjoint_cut_second_gram} gives
    \begin{align}
        Y_2^\dagger(Y_2(I_d\otimes\rho)Y_2^\dagger)X_2^\dagger
         & =\varepsilon Y_2^\dagger X_2^\dagger.
    \end{align}
    Multiplying on the left by $Z_2^\dagger$ and on the right by $X_2$
    cancels the outer factors via $Z_2^\dagger Y_2^\dagger=I$ and
    $X_2^\dagger X_2=I$. The scalar comes from the adjoint contraction; it
    has not been identified with the scalar for the original tensor.
\end{proof}

\begin{lemma}[Trace of the weighted second source Gram matrix]
    \label{thm:mpug_source_y_two_gram_trace}
    \lean{MPOTensor.IsMPUCanonicalFormII.trace_sourceY₂_weighted_mul_conjTranspose}
    \leanok
    \uses{def:mpu_canonical_form_ii, def:mpu_simple,
        def:mpu_weighted_source_factors, def:mpu_source_weight}
    Let $\mathcal U$ be a simple MPU in canonical form II of physical
    dimension $d$, with recorded fixed matrix $\rho$ and chosen second
    source factor $Y_2$. Then, as in the physical-trace calculation for
    $Y_2$ in \cite[line~1398]{FrancoRubio2025MPUGauging}, which does not
    require simplicity of $\mathcal U^\dagger$,
    \begin{align}
        \operatorname{tr}\bigl(Y_2(I_d\otimes\rho)Y_2^\dagger\bigr) & =d.
    \end{align}
\end{lemma}

\begin{proof}\leanok
    \uses{thm:mpug_source_y_physical_contractions, def:mpu_source_weight}
    Cyclic invariance and expanding $I_d\otimes\rho$ give
    \begin{align}
        \operatorname{tr}\bigl(Y_2(I_d\otimes\rho)Y_2^\dagger\bigr)
         & =\sum_{p=0}^{d-1}\sum_{x,y=0}^{D-1}\sum_{l=0}^{d_l-1}
        \overline{(Y_2)_{l,(p,y)}}(Y_2)_{l,(p,x)}\rho_{xy}
        =\sum_{p=0}^{d-1}1=d
    \end{align}
    by Lemma~\ref{thm:mpug_source_y_physical_contractions} at $q=p$.
\end{proof}

\begin{lemma}[Normalization of the weighted second source Gram matrix]
    \label{thm:mpug_source_y_two_normalization}
    \lean{MPOTensor.IsMPUCanonicalFormII.sourceY₂_weighted_mul_conjTranspose}
    \leanok
    \discussion{7317}
    \uses{def:mpu_canonical_form_ii, def:mpu_simple,
        def:mpdo_physical_adjoint, def:mpu_weighted_source_factors,
        def:mpu_source_weight}
    Let $\mathcal U$ be an MPU in canonical form II of physical dimension
    $d$, and suppose that both $\mathcal U$ and $\mathcal U^\dagger$ are
    simple. Let $Y_2$ be the second source factor constructed from
    $\mathcal U$ and its recorded fixed matrix $\rho$, and let $d_l$ be
    the second source-cut rank. Then, as in the source-factor normalization
    proposition of \cite[lines~1064--1119]{FrancoRubio2025MPUGauging},
    \begin{align}
        Y_2(I_d\otimes\rho)Y_2^\dagger & =\frac{d}{d_l}I_{d_l}.
    \end{align}
    % Source: second identity of eq:MPUnice3. Neither physical X
    % normalization nor the full prop:MPUsplus follows from this entry alone.
\end{lemma}

\begin{proof}\leanok
    \uses{thm:mpug_source_y_two_scalar_gram,
        thm:mpug_source_y_two_gram_trace}
    Writing the Gram matrix as $\varepsilon I_{d_l}$ by
    Lemma~\ref{thm:mpug_source_y_two_scalar_gram}, its trace is $d$ by
    Lemma~\ref{thm:mpug_source_y_two_gram_trace}, so $\varepsilon d_l=d$.
    As $d>0$, the rank $d_l$ is nonzero, and the trace argument of
    \cite[line~1398]{FrancoRubio2025MPUGauging} gives $\varepsilon=d/d_l$.
\end{proof}

\subsection{Physical normalizations of the two source X factors}

\begin{lemma}[The canonical one-letter output contraction]
    \label{thm:mpug_one_letter_output_contraction}
    \lean{MPOTensor.IsMPUCanonicalFormII.oneLetter_output_contraction}
    \leanok
    \discussion{7317}
    \uses{def:mpu_canonical_form_ii, def:mpu_simple,
        def:mpdo_physical_adjoint}
    Let $\mathcal U$ be an MPU in canonical form II of physical dimension
    $d$ and bond dimension $D$, with recorded fixed matrix $\rho$.
    Suppose that $\mathcal U^\dagger$ is simple. For $0\leq i,j<d$,
    \begin{align}
        \sum_{x,\alpha,\beta=0}^{D-1}\sum_{p=0}^{d-1}
        \mathcal U^{i,p}_{x\alpha}
        \overline{\mathcal U^{j,p}_{x\beta}}\rho_{\beta\alpha}
         & =\delta_{i,j}.
    \end{align}
    In the reflected one-letter contraction used for the physical
    source-factor normalizations in
    \cite[lines~1120--1162 and 1398]{FrancoRubio2025MPUGauging}, the open
    indices are the physical output indices and simplicity of $\mathcal U$
    is not required.
\end{lemma}

\begin{proof}\leanok
    \uses{def:mpug_adjoint_canonical_boundaries,
        thm:mpug_one_letter_physical_contraction, def:mpdo_physical_adjoint}
    Applying Lemma~\ref{thm:mpug_one_letter_physical_contraction} to
    $\mathcal U^\dagger$ with its canonical presentation carrying the
    same $\rho$, where
    $(\mathcal U^\dagger)^{p,i}_{x\alpha}=\overline{\mathcal U^{i,p}_{x\alpha}}$,
    and taking double conjugates gives the displayed sum.
\end{proof}

\begin{lemma}[Physical normalization of the first source X factor]
    \label{thm:mpug_source_x_one_physical_normalization}
    \lean{MPOTensor.IsMPUCanonicalFormII.sourceX₁_physical_contraction}
    \leanok
    \discussion{7317}
    \uses{def:mpu_canonical_form_ii, def:mpu_simple,
        def:mpdo_physical_adjoint, def:mpu_weighted_source_factors}
    Let $\mathcal U$ be an MPU in canonical form II of physical dimension
    $d$ and bond dimension $D$. Assume that both $\mathcal U$ and
    $\mathcal U^\dagger$ are simple. Let $X_1$ be the first source factor
    constructed from $\mathcal U$ and its recorded fixed matrix $\rho$,
    and let $d_r$ be the first source-cut rank. For $0\leq i,j<d$,
    \begin{align}
        \sum_{r=0}^{d_r-1}\sum_{\alpha,\beta=0}^{D-1}
        (X_1)_{(i,\alpha),r}\rho_{\beta\alpha}
        \overline{(X_1)_{(j,\beta),r}}
         & =\frac{d_r}{d}\delta_{i,j}.
    \end{align}
    In this first physical source-factor normalization of
    \cite[lines~1120--1162]{FrancoRubio2025MPUGauging}, the matrix $\rho$
    acts on the virtual legs and the identity acts on the physical space.
\end{lemma}

\begin{proof}\leanok
    \uses{thm:mpug_source_y_one_normalization,
        thm:mpug_source_y_one_gram_trace, thm:mpug_one_letter_output_contraction,
        thm:mpu_source_factorizations, def:mpu_source_matrix_one}
    Factoring $M_1(\mathcal U)=X_1Y_1$ and applying
    Lemma~\ref{thm:mpug_source_y_one_normalization} yields
    \begin{align}
        M_1(\mathcal U)M_1(\mathcal U)^\dagger
         & =\frac{d}{d_r}X_1X_1^\dagger.
    \end{align}
    Multiplying the $((i,\alpha),(j,\beta))$ entry by
    $\rho_{\beta\alpha}$ and summing over $\alpha,\beta$ makes the left-hand
    side $\delta_{i,j}$ by Lemma~\ref{thm:mpug_one_letter_output_contraction}.
    The trace identity in Lemma~\ref{thm:mpug_source_y_one_gram_trace}
    gives $d_r\ne0$, and canonical form II gives $d>0$. Multiplication by
    $d_r/d$ gives the required physical contraction.
\end{proof}

\begin{lemma}[Physical normalization of the second source X factor]
    \label{thm:mpug_source_x_two_physical_normalization}
    \lean{MPOTensor.IsMPUCanonicalFormII.sourceX₂_physical_contraction}
    \leanok
    \discussion{7317}
    \uses{def:mpu_canonical_form_ii, def:mpu_simple,
        def:mpdo_physical_adjoint, def:mpu_weighted_source_factors}
    Let $\mathcal U$ be an MPU in canonical form II of physical dimension
    $d$ and bond dimension $D$. Assume that both $\mathcal U$ and
    $\mathcal U^\dagger$ are simple. Let $X_2$ be its chosen second source
    factor and let $d_l$ be the second source-cut rank. For $0\leq i,j<d$,
    the second physical source-factor normalization in
    \cite[lines~1120--1162]{FrancoRubio2025MPUGauging} holds:
    \begin{align}
        \sum_{l=0}^{d_l-1}\sum_{x=0}^{D-1}
        (X_2)_{(x,i),l}\overline{(X_2)_{(x,j),l}}
         & =\frac{d_l}{d}\delta_{i,j}.
    \end{align}
\end{lemma}

\begin{proof}\leanok
    \uses{thm:mpug_source_y_two_normalization,
        thm:mpug_source_y_two_gram_trace, thm:mpug_one_letter_output_contraction,
        thm:mpu_source_factorizations, def:mpu_source_matrix_two,
        def:mpu_source_weight}
    With $W=I_d\otimes\rho$, the factorization $M_2(\mathcal U)=X_2Y_2$
    and Lemma~\ref{thm:mpug_source_y_two_normalization} give
    \begin{align}
        M_2(\mathcal U)WM_2(\mathcal U)^\dagger
         & =\frac{d}{d_l}X_2X_2^\dagger.
    \end{align}
    Summing the $((x,i),(x,j))$ entry over $x$ and expanding the left-hand
    side with $\rho_{\alpha\beta}=\rho_{\beta\alpha}$ gives
    $\delta_{i,j}$ by Lemma~\ref{thm:mpug_one_letter_output_contraction}.
    The trace identity in Lemma~\ref{thm:mpug_source_y_two_gram_trace}
    gives $d_l\ne0$ with $d>0$, so multiplying by $d_l/d$ yields the result.
\end{proof}

Lemmas~\ref{thm:mpug_source_y_one_normalization},
\ref{thm:mpug_source_y_two_normalization},
\ref{thm:mpug_source_x_one_physical_normalization}, and
\ref{thm:mpug_source_x_two_physical_normalization} establish all four
identities of the source-factor normalization proposition in
\cite[lines~1064--1162]{FrancoRubio2025MPUGauging} for the compact source
factors under canonical form~II and simplicity of $\mathcal U$ and
$\mathcal U^\dagger$; no additional boundary or Gram identities are assumed.
% This establishes prop:MPUsplus for the constructed source factors.
% It does not identify arbitrary separately supplied factors or modify
% the scope of the earlier raw proportionality theorem.

\begin{lemma}[Comparison of factorizations from one-sided inverses]
    \label{thm:mpug_factorization_comparison}
    \lean{Matrix.factorization_comparison_of_one_sided_inverses}
    \leanok
    Let $A,B,\widetilde B,C$ be finite index sets, possibly empty, and let
    $\mathbb S$ be a semiring. Suppose that
    \begin{align}
        X            & \in\operatorname{Mat}_{A\times B}(\mathbb S),           \\
        Y            & \in\operatorname{Mat}_{B\times C}(\mathbb S),           \\
        \widetilde X & \in\operatorname{Mat}_{A\times\widetilde B}(\mathbb S), \\
        \widetilde Y & \in\operatorname{Mat}_{\widetilde B\times C}(\mathbb S)
    \end{align}
    satisfy $XY=\widetilde X\widetilde Y$. If
    $L\in\operatorname{Mat}_{\widetilde B\times A}(\mathbb S)$ and
    $R\in\operatorname{Mat}_{C\times B}(\mathbb S)$ satisfy
    $L\widetilde X=I_{\widetilde B}$ and $YR=I_B$, then the rectangular
    comparison matrix $K=LX\in\operatorname{Mat}_{\widetilde B\times B}(\mathbb S)$
    obeys
    \begin{align}
        K            & =\widetilde YR, \\
        X            & =\widetilde XK, \\
        \widetilde Y & =KY.
    \end{align}
    In these one-sided-inverse identities from the proof of
    \cite[lines~1052--1066]{FrancoRubio2025MPUGauging}, the intermediate
    index sets $B$ and $\widetilde B$ are independent, and no equality of
    their cardinalities is assumed or concluded.
\end{lemma}
\begin{proof}\leanok
    Associativity and the two inverse identities give
    \begin{align}
        LX & =(LX)(YR)=L(XY)R
        =L(\widetilde X\widetilde Y)R=\widetilde YR.
    \end{align}
    Multiplying $XY=\widetilde X\widetilde Y$ on the right by $R$ gives
    $X=\widetilde XK$. Multiplying it on the left by $L$ gives
    $\widetilde Y=KY$.
\end{proof}

\begin{lemma}[Reciprocal unitary gauges from a Kronecker comparison]
    \label{thm:mpug_reciprocal_unitary_gauges}
    \lean{Matrix.exists_reciprocal_unitary_gauges_of_comparison}
    \leanok
    For $i=1,2$, suppose that two compatible matrix factorizations are related
    by square comparison matrices $K_i$ through
    \begin{align}
        X_i & =\widetilde X_iK_i, & \widetilde Y_i & =K_iY_i.
        \label{eq:mpug_comparison_matrices}
    \end{align}
    If $K_1\otimes K_2$ is unitary, then there are a positive real number
    $\delta$ and unitary matrices $W_1,W_2$ such that
    \begin{align}
        \widetilde X_1 & =\delta^{-1}X_1W_1,          &
        \widetilde Y_1 & =\delta W_1^\dagger Y_1,       \\
        \widetilde X_2 & =\delta X_2W_2,              &
        \widetilde Y_2 & =\delta^{-1}W_2^\dagger Y_2.
        \label{eq:mpug_reciprocal_unitary_gauges}
    \end{align}
    Thus the positive scalar factors in the orientation of
    \cite[Lemma at lines~1052--1066]{FrancoRubio2025MPUGauging} are
    $\delta_1=\delta^{-1}$ and $\delta_2=\delta$, and hence
    $\delta_1\delta_2=1$.
    % Source anchor: arXiv:2502.20257, Lemma lem:deco, lines 1052--1066.
\end{lemma}

\begin{proof}\leanok
    \uses{thm:mpug_unitary_kronecker_factors}
    Applying Lemma~\ref{thm:mpug_unitary_kronecker_factors} to
    $K_1\otimes K_2$ gives $\delta>0$ and unitary matrices $U_1,U_2$ with
    $K_1=\delta U_1$ and $K_2=\delta^{-1}U_2$. Setting
    $W_i=U_i^\dagger$ and substituting into
    (\ref{eq:mpug_comparison_matrices}) gives
    (\ref{eq:mpug_reciprocal_unitary_gauges}), while positivity of $\delta$
    gives $\delta^{-1}\delta=1$.
\end{proof}

\subsection{Uniqueness of the source decompositions}

\begin{lemma}[Transformation of the source contraction]
    \label{thm:mpug_source_contraction_transport}
    \lean{MPOTensor.sourceU_contraction_transport}
    \leanok
    \discussion{7315}
    \uses{def:mpu_source_matrix_one, def:mpu_source_matrix_two}
    Let $U$ have physical dimension $d$ and bond dimension $D$, and write
    $r=\operatorname{rank}M_1(U)$ and $\ell=\operatorname{rank}M_2(U)$.
    Let $Y_1,\widetilde Y_1\in\operatorname{Mat}_{r\times(D\times d)}(\C)$
    and $Y_2,\widetilde Y_2\in\operatorname{Mat}_{\ell\times(d\times D)}(\C)$.
    For square matrices $K_1\in\operatorname{Mat}_r(\C)$ and
    $K_2\in\operatorname{Mat}_\ell(\C)$, suppose that
    $\widetilde Y_i=K_iY_i$ for $i=1,2$. Define
    $u,\widetilde u\colon\C^d\otimes\C^d\to\C^\ell\otimes\C^r$ by
    \begin{align}
        u_{(l,r'),(p,q)}
         & =\sum_{\beta=0}^{D-1}(Y_2)_{l,(p,\beta)}(Y_1)_{r',(\beta,q)}, \\
        \widetilde u_{(l,r'),(p,q)}
         & =\sum_{\beta=0}^{D-1}(\widetilde Y_2)_{l,(p,\beta)}
        (\widetilde Y_1)_{r',(\beta,q)}.
        \label{eq:mpug_raw_source_contractions}
    \end{align}
    Here $0\leq l<\ell$, $0\leq r'<r$, and $0\leq p,q<d$.
    Then $\widetilde u=(K_2\otimes K_1)u$.
\end{lemma}

\begin{proof}\leanok
    Substituting $\widetilde Y_i=K_iY_i$ into
    (\ref{eq:mpug_raw_source_contractions}) and interchanging the three finite
    sums isolates the coefficient of $u_{(l',r''),(p,q)}$ as
    $(K_2)_{l,l'}(K_1)_{r',r''}$, the corresponding entry of $K_2\otimes K_1$.
\end{proof}

\begin{lemma}[Reciprocal unitary gauges for the source decompositions]
    \label{thm:mpug_source_decomposition_gauges}
    \lean{MPOTensor.exists_reciprocal_unitary_source_gauges}
    \leanok
    \discussion{7315}
    \uses{def:mpu_source_matrix_one, def:mpu_source_matrix_two}
    Let $U$ have physical dimension $d$ and bond dimension $D$, with
    $r=\operatorname{rank}M_1(U)$ and $\ell=\operatorname{rank}M_2(U)$.
    Suppose that
    \begin{align}
        X_1,\widetilde X_1 & \in\operatorname{Mat}_{(d\times D)\times r}(\C),   &
        Y_1,\widetilde Y_1 & \in\operatorname{Mat}_{r\times(D\times d)}(\C),      \\
        X_2,\widetilde X_2 & \in\operatorname{Mat}_{(D\times d)\times\ell}(\C), &
        Y_2,\widetilde Y_2 & \in\operatorname{Mat}_{\ell\times(d\times D)}(\C)
    \end{align}
    satisfy $X_iY_i=\widetilde X_i\widetilde Y_i=M_i(U)$ for $i=1,2$.
    Suppose there are matrices $L_i,R_i$ with
    $L_i\widetilde X_i=I$ and $Y_iR_i=I$.
    Assume that the maps $u,\widetilde u\colon\C^d\otimes\C^d\to
        \C^\ell\otimes\C^r$ with coordinates
    \begin{align}
        u_{(l,r'),(p,q)}
         & =\sum_{\beta=0}^{D-1}(Y_2)_{l,(p,\beta)}(Y_1)_{r',(\beta,q)}, \\
        \widetilde u_{(l,r'),(p,q)}
         & =\sum_{\beta=0}^{D-1}(\widetilde Y_2)_{l,(p,\beta)}
        (\widetilde Y_1)_{r',(\beta,q)}
    \end{align}
    are unitary. The indices range over $0\leq l<\ell$, $0\leq r'<r$,
    and $0\leq p,q<d$. Then there exist $\delta>0$ and unitaries
    $W_1\in\operatorname{Mat}_r(\C)$, $W_2\in\operatorname{Mat}_\ell(\C)$ such that
    \begin{align}
        \widetilde X_1 & =\delta^{-1}X_1W_1,          &
        \widetilde Y_1 & =\delta W_1^\dagger Y_1,       \\
        \widetilde X_2 & =\delta X_2W_2,              &
        \widetilde Y_2 & =\delta^{-1}W_2^\dagger Y_2.
        \label{eq:mpug_source_decomposition_gauges}
    \end{align}
    Thus $\delta_1=\delta^{-1}$ and $\delta_2=\delta$ satisfy
    $\delta_1\delta_2=1$, establishing the decomposition lemma of
    \cite[lines~1052--1066]{FrancoRubio2025MPUGauging} from the cut factorizations,
    one-sided inverses, and unitarity of $u,\widetilde u$ alone, without requiring
    unitarity of the second contraction $v$.
    % Source lem:deco, issue #7315. No clause (c), weighted normalization,
    % canonical-form-II, density, or physical-adjoint hypothesis is used.
    % This result does not assert completion of the separate #7313 assembly.
\end{lemma}

\begin{proof}\leanok
    \uses{thm:mpug_factorization_comparison,
        thm:mpug_source_contraction_transport, thm:mpug_reciprocal_unitary_gauges}
    If $d=0$, every factor equality has an empty physical row or column
    index. Thus $\delta=1$ and identity gauges satisfy all four
    equalities, and $\delta^{-1}\delta=1$.
    Suppose now that $d>0$. Since $u$ is an isometry,
    $\ell r\geq d^2>0$; hence both
    source ranks are positive. Lemma~\ref{thm:mpug_factorization_comparison}
    gives $K_i=L_iX_i=\widetilde Y_iR_i$ with
    $X_i=\widetilde X_iK_i$ and $\widetilde Y_i=K_iY_i$.
    Lemma~\ref{thm:mpug_source_contraction_transport} yields
    $\widetilde u=(K_2\otimes K_1)u$. Therefore
    $K_2\otimes K_1=\widetilde u u^\dagger$ is unitary.
    Interchanging the two tensor factors makes $K_1\otimes K_2$ unitary, and
    Lemma~\ref{thm:mpug_reciprocal_unitary_gauges} then gives
    (\ref{eq:mpug_source_decomposition_gauges}). In particular, if
    $K_1=\delta V_1$ and $K_2=\delta^{-1}V_2$ with $V_i$ unitary,
    the required orientation is $W_i=V_i^\dagger$.
\end{proof}

\begin{corollary}[Unitary gauges for normalized source decompositions]
    \label{cor:mpug_normalized_source_decomposition_gauges}
    \lean{MPOTensor.exists_unitary_source_gauges_of_isometric_second_factors}
    \leanok
    \uses{thm:mpug_source_decomposition_gauges}
    Under the hypotheses of Lemma~\ref{thm:mpug_source_decomposition_gauges},
    suppose in addition that
    \begin{align}
        X_2^\dagger X_2 & =I_\ell,
                        & \widetilde X_2^\dagger\widetilde X_2 & =I_\ell.
    \end{align}
    In this isometrically normalized specialization of
    \cite[lines~1052--1066]{FrancoRubio2025MPUGauging}, the additional
    isometry hypotheses remove the scalar freedom: there exist unitaries
    $W_1\in\operatorname{Mat}_r(\C)$ and
    $W_2\in\operatorname{Mat}_\ell(\C)$ such that, for $i=1,2$,
    \begin{align}
        \widetilde X_i & =X_iW_i,
                       & \widetilde Y_i & =W_i^\dagger Y_i.
    \end{align}
\end{corollary}

\begin{proof}\leanok
    \uses{thm:mpug_source_decomposition_gauges}
    If $d=0$, all four equalities have an empty physical index and hold
    with identity gauges. If $d>0$, unitarity of $u$ gives $r\ell=d^2>0$,
    so $\ell>0$. Lemma~\ref{thm:mpug_source_decomposition_gauges} gives
    $\widetilde X_2=\delta X_2W_2$ with $\delta>0$. The two isometry
    identities and unitarity of $W_2$ imply
    $I_\ell=\delta^2I_\ell$. Since $\ell>0$, a diagonal entry gives
    $\delta^2=1$, hence $\delta=1$. Substitution into the four comparison
    identities proves the claim.
\end{proof}

\begin{definition}[Unitary factorization comparison]
    \label{def:mpug_unitary_factorization_comparison}
    \lean{Matrix.unitaryFactorizationComparison,
        Matrix.unitaryFactorizationComparison_coe}
    \leanok
    Let $I$ be a set and let $R_{a,b}$ be a family of unitary matrices of a
    fixed size, indexed by $(a,b)\in I^2$. Following the orientation in the
    corollary after Proposition~3
    of~\cite[lines~2782--2821]{FrancoRubio2025MPUGauging}, the unitary
    comparison from $(a,b)$ to the reference pair $(c,d)$ is
    \begin{align}
        \lambda_{a,b}^{c,d}
         & =(R_{c,d})^{-1}R_{a,b}
        =R_{c,d}^\dagger R_{a,b}.
        \label{eq:mpug_unitary_factorization_comparison}
    \end{align}
    % Source anchor: arXiv:2502.20257, Corollary after Proposition 3,
    % lines 2782--2821, especially the defining formula at line 2820.
    In that corollary, $I=G$ and comparisons are used when $ab=cd$; this
    equality selects the factorizations and is not needed in the unitary
    calculation.
\end{definition}

\begin{lemma}[Identity, inverse, composition, and right-reference laws]
    \label{thm:mpug_unitary_factorization_comparison_laws}
    \lean{Matrix.unitaryFactorizationComparison_self,
        Matrix.unitaryFactorizationComparison_inv,
        Matrix.unitaryFactorizationComparison_trans,
        Matrix.unitaryFactorizationComparison_of_reference_eq_one,
        Matrix.unitaryFactorizationComparison_right_reference}
    \leanok
    \uses{def:mpug_unitary_factorization_comparison}
    For all indices $a,b,c,d,f,g\in I$,
    \begin{align}
        \lambda_{a,b}^{a,b} & =\mathbf 1.
    \end{align}
    Reversing the comparison gives
    \begin{align}
        \left(\lambda_{a,b}^{c,d}\right)^{-1}
         & =\lambda_{c,d}^{a,b}.
    \end{align}
    Comparisons compose as
    \begin{align}
        \lambda_{a,b}^{f,g}
         & =\lambda_{c,d}^{f,g}\lambda_{a,b}^{c,d}.
        \label{eq:mpug_unitary_factorization_comparison_laws}
    \end{align}
    If $R_{c,d}=\mathbf 1$, then
    $\lambda_{a,b}^{c,d}=R_{a,b}$. In particular, if $I=G$ is a group, the
    normalization $R_{e,x}=\mathbf 1$ gives the right-reference specialization
    $\lambda^R_{g,h}=\lambda_{g,h}^{e,gh}$ in the cited corollary,
    \begin{align}
        R_{g,h} & =\lambda_{g,h}^{e,gh}.
        \label{eq:mpug_unitary_factorization_right_reference}
    \end{align}
    % Source anchor: arXiv:2502.20257, Corollary after Proposition 3,
    % line 2817 (right reference and composition).
\end{lemma}

\begin{proof}\leanok
    \uses{def:mpug_unitary_factorization_comparison}
    The identity, inverse, and right-reference formulas follow from
    $R_{a,b}^{-1}R_{a,b}=\mathbf 1$, inversion of a product, and
    $\mathbf 1^{-1}R_{a,b}=R_{a,b}$. For composition, inserting the identity at
    the intermediate pair gives
    \begin{align}
        (R_{f,g})^{-1}R_{a,b}
         & =(R_{f,g})^{-1}R_{c,d}(R_{c,d})^{-1}R_{a,b}.
    \end{align}
\end{proof}

\begin{definition}[Gauss action on the gauge legs]
    \label{def:mpug_gauss_leg_action}
    \lean{TNLean.Algebra.gaussLegAction,
        TNLean.Algebra.gaussLegAction_apply}
    \leanok
    \discussion{7503}
    Let $G$ be a group. Following the two gauge-leg labels in the local Gauss
    operator of \cite[lines~3380--3387]{FrancoRubio2025MPUGauging}, for
    $g\in G$ define the permutation $T_g\colon G^2\to G^2$ by
    \begin{align}
        T_g(a,b) & =\left(ag^{-1},gb\right).
        \label{eq:mpug_gauss_leg_action}
    \end{align}
\end{definition}

\begin{lemma}[Composition of the Gauss leg action]
    \label{thm:mpug_gauss_leg_action_mul}
    \lean{TNLean.Algebra.gaussLegAction_one,
        TNLean.Algebra.gaussLegAction_mul}
    \leanok
    \discussion{7503}
    \uses{def:mpug_gauss_leg_action}
    For all $g,h\in G$,
    \begin{align}
        T_e          & =\id,    \\
        T_g\circ T_h & =T_{gh}.
        \label{eq:mpug_gauss_leg_action_mul}
    \end{align}
\end{lemma}

\begin{proof}\leanok
    \uses{def:mpug_gauss_leg_action}
    The two coordinates give
    \begin{align}
        T_g(T_h(a,b))
         & =\left(ah^{-1}g^{-1},ghb\right)
        =\left(a(gh)^{-1},(gh)b\right).
    \end{align}
\end{proof}

\begin{definition}[Local Gauss operator]
    \label{def:mpug_gauss_operator}
    \lean{TNLean.Algebra.gaussOperator,
        TNLean.Algebra.gaussOperator_apply_target,
        TNLean.Algebra.gaussOperator_apply_of_ne}
    \leanok
    \discussion{7503}
    \uses{def:mpug_unitary_factorization_comparison,
        def:mpug_gauss_leg_action}
    Let $G$ be finite, let $R_{a,b}\in\mathrm U(n)$ for $a,b\in G$, and take
    the basis $\ket{i,a,b}$ of $\C^n\otimes\C[G]^{\otimes2}$. As in
    \cite[lines~3380--3387]{FrancoRubio2025MPUGauging}, define the local Gauss
    operator
    \begin{align}
        \mathcal G_g
         & =\sum_{a,b\in G}
        \lambda_{a,b}^{ag^{-1},gb}
        \otimes\ket{ag^{-1},gb}\!\bra{a,b}.
        \label{eq:mpug_gauss_operator}
    \end{align}
    Equivalently, the block from $(a,b)$ to $T_g(a,b)$ is
    $(R_{T_g(a,b)})^{-1}R_{a,b}$ and every other block in that column
    vanishes.
\end{definition}

\begin{lemma}[Gauss operators form a unitary representation]
    \label{thm:mpug_gauss_representation}
    \lean{TNLean.Algebra.gaussOperator_one,
        TNLean.Algebra.gaussOperator_mul,
        TNLean.Algebra.star_gaussOperator,
        TNLean.Algebra.gaussOperator_mem_unitaryGroup,
        TNLean.Algebra.gaussRepresentation,
        TNLean.Algebra.gaussRepresentation_coe}
    \leanok
    \discussion{7503}
    \uses{def:mpug_gauss_operator}
    The local Gauss operators satisfy
    \begin{align}
        \mathcal G_e & =\id.
    \end{align}
    Their product law is
    \begin{align}
        \mathcal G_g\mathcal G_h & =\mathcal G_{gh}.
        \label{eq:mpug_gauss_representation}
    \end{align}
    Finally,
    \begin{align}
        \mathcal G_g^\dagger & =\mathcal G_{g^{-1}}.
    \end{align}
    Hence $g\mapsto\mathcal G_g$ is a unitary representation, establishing the
    representation part of the local Gauss-law proposition in
    \cite[lines~3380--3387]{FrancoRubio2025MPUGauging}; invariance of the
    gauged tensor is a separate assertion and is not included here.
\end{lemma}

\begin{proof}\leanok
    \uses{thm:mpug_unitary_factorization_comparison_laws,
        thm:mpug_gauss_leg_action_mul,
        def:mpug_gauss_operator}
    In a column labelled by $(a,b)$, multiplication has only the intermediate
    gauge label $T_h(a,b)$. The surviving matter block is
    \begin{align}
        \lambda_{T_h(a,b)}^{T_g(T_h(a,b))}
        \lambda_{a,b}^{T_h(a,b)}
         & =\lambda_{a,b}^{T_{gh}(a,b)},
    \end{align}
    by the comparison composition law and
    $T_gT_h=T_{gh}$. At $e$, every surviving comparison is the identity.
    Reversing a comparison gives
    $(\lambda_{a,b}^{c,d})^\dagger=\lambda_{c,d}^{a,b}$, so taking the adjoint
    gives $\mathcal G_g^\dagger=\mathcal G_{g^{-1}}$. The product with this
    adjoint is $\mathcal G_e=\id$, which proves unitarity.
\end{proof}

\begin{definition}[Unplaced local Gauss projector kernel]
    \label{def:mpug_gauss_projector}
    \lean{TNLean.Algebra.gaussProjector,
        TNLean.Algebra.gaussProjector_eq_average}
    \leanok
    \discussion{7578}
    \uses{def:mpug_finite_group_unitary_average,
        thm:mpug_gauss_representation}
    Let $G$ be finite and let $\mathcal G_g$ act on
    $\mathcal H=\C^n\otimes\C[G]^{\otimes2}$. Define the unplaced local
    projector kernel
    \begin{align}
        P_G(R) & =\frac{1}{|G|}\sum_{g\in G}\mathcal G_g.
        \label{eq:mpug_gauss_projector}
    \end{align}
    This is the local average introduced in
    \cite[lines~3494--3499]{FrancoRubio2025MPUGauging}, before a position on
    the chain is chosen. This definition neither places the kernel on
    $(j,j+1)$ nor asserts invariance of the gauged matrix product state.
\end{definition}

\begin{lemma}[The local Gauss average is a projection]
    \label{thm:mpug_gauss_projector_projection}
    \lean{TNLean.Algebra.isStarProjection_gaussProjector}
    \leanok
    \discussion{7578}
    \uses{def:mpug_gauss_projector}
    The kernel in~(\ref{eq:mpug_gauss_projector}) satisfies
    \begin{align}
        P_G(R)^2       & =P_G(R), \\
        P_G(R)^\dagger & =P_G(R).
    \end{align}
\end{lemma}

\begin{proof}\leanok
    \uses{thm:mpug_finite_group_unitary_average_projection,
        thm:mpug_gauss_representation}
    Identifying the matrix average with the averaging operator of the unitary
    representation $g\mapsto\mathcal G_g$ makes this operator idempotent by
    the invariant-space averaging theorem. For self-adjointness,
    $\mathcal G_g^\dagger=\mathcal G_{g^{-1}}$ and reindexing the finite sum by
    the bijection $g\mapsto g^{-1}$ yield $P_G(R)^\dagger=P_G(R)$.
\end{proof}

\begin{definition}[Gauss projector on a periodic chain]
    \label{def:mpug_placed_gauss_projector}
    \lean{MPOTensor.gaussLocalCoordinateEquiv,
        MPOTensor.gaussLocalCoordinateEquiv_apply_zero,
        MPOTensor.gaussLocalCoordinateEquiv_apply_one,
        MPOTensor.gaussWindowProjector,
        MPOTensor.placedGaussProjector}
    \leanok
    \discussion{7605}
    \uses{def:mpug_gauss_projector}
    Let the matter degree of freedom at one site have dimension $d$.  The
    local kernel has row and column coordinates
    \begin{align}
        \bigl((i_0,i_1),(a,b)\bigr)
         & \in (\{0,\ldots,d-1\}^{\{0,1\}})\times G^2.
    \end{align}
    Identifying these with two site coordinates by the bijection
    \begin{align}
        \kappa\bigl((i_0,i_1),(a,b)\bigr)
         & =\bigl((i_0,a),(i_1,b)\bigr),
    \end{align}
    followed at each site by a fixed enumeration of
    $\{0,\ldots,d-1\}\times G$, let $\widetilde P_G(R)$ denote the matrix
    obtained from $P_G(R)$ by this simultaneous row and column relabelling.
    Following the position-dependent projector $P_j$
    in~\cite[lines~3494--3499]{FrancoRubio2025MPUGauging}, define its placement
    on a periodic chain of length $N\geq2$ by
    \begin{align}
        P_j(R) & =\operatorname{Embed}_{j,j+1}\bigl(\widetilde P_G(R)\bigr),
        \label{eq:mpug_placed_gauss_projector}
    \end{align}
    where the second site is $j+1$ modulo $N$.
    % Source anchor: arXiv:2502.20257, eq:gauge_projs.
\end{definition}

\begin{lemma}[The placed Gauss projector is a projection]
    \label{thm:mpug_placed_gauss_projector_projection}
    \lean{MPOTensor.isStarProjection_gaussWindowProjector,
        MPOTensor.isStarProjection_placedGaussProjector}
    \leanok
    \discussion{7605}
    \uses{def:mpug_placed_gauss_projector}
    For every $N\geq2$ and every position $j$, the operator
    in~(\ref{eq:mpug_placed_gauss_projector}) satisfies
    \begin{align}
        P_j(R)^2       & =P_j(R), \\
        P_j(R)^\dagger & =P_j(R).
    \end{align}
\end{lemma}

\begin{proof}\leanok
    \uses{thm:mpug_gauss_projector_projection,
        thm:mpdo_cyclic_local_embedding_reindex}
    Let $U_\kappa$ be the permutation unitary determined by the coordinate
    bijection in Definition~\ref{def:mpug_placed_gauss_projector}, with the
    orientation
    \begin{align}
        U_\kappa\ket{((i_0,i_1),(a,b))}
         & =\ket{\kappa((i_0,i_1),(a,b))}.
        \notag
    \end{align}
    Simultaneous row and column reindexing is therefore the conjugation
    \begin{align}
        \widetilde P_G(R)
         & =U_\kappa P_G(R)U_\kappa^{-1}.
        \label{eq:mpug_gauss_reindex_conjugation}
    \end{align}
    Since $U_\kappa^{-1}=U_\kappa^\dagger$, Lemma~\ref{thm:mpug_gauss_projector_projection}
    and~(\ref{eq:mpug_gauss_reindex_conjugation}) give
    \begin{align}
        \widetilde P_G(R)^2
         & =U_\kappa P_G(R)^2U_\kappa^{-1}
        =\widetilde P_G(R),
        \label{eq:mpug_gauss_reindexed_idempotent} \\
        \widetilde P_G(R)^\dagger
         & =U_\kappa P_G(R)^\dagger U_\kappa^{-1}
        =\widetilde P_G(R).
        \label{eq:mpug_gauss_reindexed_adjoint}
    \end{align}
    Let $V_j$ be the permutation unitary that sends a chain configuration to
    its restriction to the cyclic window $(j,j+1)$ followed by its restriction
    to the complement. By Lemma~\ref{thm:mpdo_cyclic_local_embedding_reindex},
    \begin{align}
        V_jP_j(R)V_j^{-1}
         & =\widetilde P_G(R)\otimes\Id_{\mathrm{comp}}.
        \label{eq:mpug_gauss_window_tensor_identity}
    \end{align}
    Using $V_j^{-1}=V_j^\dagger$ and
    (\ref{eq:mpug_gauss_reindexed_idempotent})--(\ref{eq:mpug_gauss_reindexed_adjoint})
    in~(\ref{eq:mpug_gauss_window_tensor_identity}), we obtain
    \begin{align}
        V_jP_j(R)^2V_j^{-1}
         & =\bigl(\widetilde P_G(R)\otimes\Id_{\mathrm{comp}}\bigr)^2
        =\widetilde P_G(R)^2\otimes\Id_{\mathrm{comp}}
        =V_jP_j(R)V_j^{-1},                                                 \\
        V_jP_j(R)^\dagger V_j^{-1}
         & =\bigl(\widetilde P_G(R)\otimes\Id_{\mathrm{comp}}\bigr)^\dagger
        =\widetilde P_G(R)^\dagger\otimes\Id_{\mathrm{comp}}
        =V_jP_j(R)V_j^{-1}.
    \end{align}
    Conjugating by $V_j^{-1}$ gives $P_j(R)^2=P_j(R)$ and
    $P_j(R)^\dagger=P_j(R)$.
\end{proof}

\begin{definition}[Common fixed subspace]
    \label{def:mpug_common_fixed_subspace}
    \lean{LinearMap.commonFixedSubmodule}
    \leanok
    Let $\mathcal H$ be a complex vector space and let $(P_j)_{j\in J}$ be any
    family of endomorphisms of $\mathcal H$. Define
    their common fixed subspace by
    \begin{align}
        \mathcal V(P)
         & =\bigcap_{j\in J}\ker(P_j-\Id_{\mathcal H}).
    \end{align}
    No commutativity of the operators is required. For the averaged Gauss
    projectors, this is the common $+1$ eigenspace called the gauge-invariant
    subspace in
    \cite[lines~4327--4335 and line~5198]{FrancoRubio2025MPUGauging}.
    Here $P_j$ is an averaged projector, not an individual Gauss operator
    $\mathcal G_g$.
    % Source anchors: arXiv:2502.20257, lines 4327--4335 and line 5198.
\end{definition}

\begin{lemma}[Membership in the common fixed subspace]
    \label{thm:mpug_mem_common_fixed_subspace}
    \lean{LinearMap.mem_commonFixedSubmodule_iff}
    \leanok
    \uses{def:mpug_common_fixed_subspace}
    A vector $\psi\in\mathcal H$ belongs to $\mathcal V(P)$ if and only if
    $P_j\psi=\psi$ for every $j\in J$.
\end{lemma}

\begin{proof}\leanok
    \uses{def:mpug_common_fixed_subspace}
    Membership in an indexed intersection is equivalent to membership in
    every factor, and membership in $\ker(P_j-\Id_{\mathcal H})$ is equivalent
    to $P_j\psi=\psi$.
\end{proof}

\begin{lemma}[Nonzero-vector lower bound]
    \label{thm:mpug_common_fixed_subspace_finrank}
    \lean{LinearMap.one_le_finrank_commonFixedSubmodule}
    \leanok
    \uses{def:mpug_common_fixed_subspace}
    If $\mathcal H$ is finite-dimensional, $\psi\in\mathcal V(P)$, and
    $\psi\neq0$, then $1\leq\dim_{\C}\mathcal V(P)$.
\end{lemma}

\begin{proof}\leanok
    The common fixed subspace is nonzero because it contains $\psi$, so its
    dimension is at least one.
\end{proof}

\begin{definition}[Prescribed defect maps]
    \label{def:mpug_defect_maps}
    \lean{TNLean.Algebra.DefectMaps}
    \leanok
    \discussion{7569}
    Let $\mathcal H_{\mathrm m}$ be the matter space of two sites. Prescribed
    defect maps assign to every ordered pair $(a,b)\in G^2$ a defect subspace
    \begin{align}
        \mathcal K_{a,b} & \subseteq\mathcal H_{\mathrm m}
    \end{align}
    together with a map $D_{a,b}$ of $\mathcal H_{\mathrm m}$ whose
    restriction to that subspace is an isometry: for every
    $\ket{\xi}\in\mathcal K_{a,b}$,
    \begin{align}
        D_{a,b}^\dagger D_{a,b}\ket{\xi}
         & =\ket{\xi}.
        \label{eq:mpug_defect_isometry}
    \end{align}
    When the fusion scalars are not block independent, the state-level
    construction of~\cite[lines~4215--4326]{FrancoRubio2025MPUGauging}
    prescribes the modified fusion operators only on such subspaces.
    % Source anchor: arXiv:2502.20257, lemma:modif and the sentence closing
    % its proof.
\end{definition}

\begin{definition}[Full unitary completion class]
    \label{def:mpug_completion_class}
    \lean{TNLean.Algebra.DefectMaps.IsCompletion,
        TNLean.Algebra.DefectMaps.completionClass,
        TNLean.Algebra.DefectMaps.mem_completionClass_iff}
    \leanok
    \discussion{7569}
    \uses{def:mpug_defect_maps}
    The full completion class of prescribed defect maps
    $D=(D_{a,b})_{a,b\in G}$ is
    \begin{align}
        \mathfrak C(D)
         & =\left\{R=(U_{a,b})_{a,b\in G}:
        U_{a,b}\in\mathrm U(\mathcal H_{\mathrm m}),\
        U_{a,b}|_{\mathcal K_{a,b}}=D_{a,b}|_{\mathcal K_{a,b}}\right\}.
        \label{eq:mpug_completion_class}
    \end{align}
    As noted in~\cite[lines~4215--4326]{FrancoRubio2025MPUGauging}, when
    $\mathcal K_{a,b}$ is a proper subspace of the matter space, the unitary
    extension of $D_{a,b}$ is non-unique.
\end{definition}

\begin{lemma}[Existence of full unitary completions]
    \label{thm:mpug_completion_class_nonempty}
    \lean{TNLean.Algebra.DefectMaps.completionClass_nonempty,
        TNLean.Algebra.DefectMaps.exists_completion}
    \leanok
    \discussion{7702}
    \uses{def:mpug_defect_maps, def:mpug_completion_class}
    Let $\mathcal H_{\mathrm m}$ be finite-dimensional. Every family of
    prescribed defect maps satisfying (\ref{eq:mpug_defect_isometry}) admits a
    simultaneous unitary extension:
    \begin{align}
        \mathfrak C(D) & \neq\varnothing.
        \label{eq:mpug_completion_class_nonempty}
    \end{align}
    Equivalently, there is a family
    $R=(U_{a,b})_{(a,b)\in G^2}$ of unitary operators such that, for every
    $a,b\in G$ and every $\ket{\xi}\in\mathcal K_{a,b}$,
    \begin{align}
        U_{a,b}\ket{\xi}
         & =D_{a,b}\ket{\xi}.
        \label{eq:mpug_completion_exists}
    \end{align}
    This is the finite-dimensional extension step implicit in the freedom of
    the modified fusion operators at
    \cite[lines~4215--4326]{FrancoRubio2025MPUGauging}, where the subspaces
    $\mathcal K_{a,b}$ and their prescribed isometries enter as hypotheses.
    % Source anchor: arXiv:2502.20257, lemma:modif and its proof,
    % lines 4215--4326. The theorem does not assert that the defect data arise
    % from locally orthogonal MPS blocks or action tensors; their construction
    % remains in issues #7349, #7350, and #7569.
\end{lemma}

\begin{proof}\leanok
    \uses{def:mpug_defect_maps, def:mpug_completion_class}
    For each pair $(a,b)$, the restriction of $D_{a,b}$ to
    $\mathcal K_{a,b}$ is a linear isometry and therefore extends to a unitary
    operator on the whole finite-dimensional matter space. Choosing one such
    extension for every pair gives a family $R$ satisfying
    (\ref{eq:mpug_completion_exists}), hence a member of $\mathfrak C(D)$.
\end{proof}

\begin{lemma}[Comparison identity for a completion]
    \label{thm:mpug_completion_transport}
    \lean{TNLean.Algebra.DefectMaps.IsCompletion.mulVec_eq_prescribed,
        TNLean.Algebra.DefectMaps.IsCompletion.mulVec_eq_mulVec,
        TNLean.Algebra.DefectMaps.IsCompletion.unitaryFactorizationComparison_mulVec,
        TNLean.Algebra.DefectMaps.IsCompletion.unitaryFactorizationComparison_gaussLegAction_mulVec}
    \leanok
    \discussion{7569}
    \uses{def:mpug_completion_class,
        def:mpug_unitary_factorization_comparison,
        def:mpug_gauss_leg_action}
    Let $R=(U_{a,b})\in\mathfrak C(D)$, let
    $\ket{\xi}\in\mathcal K_{a,b}$ and $\ket{\eta}\in\mathcal K_{c,d}$, and
    assume the exact covariance
    \begin{align}
        D_{a,b}\ket{\xi} & =D_{c,d}\ket{\eta}.
        \label{eq:mpug_prescribed_covariance}
    \end{align}
    Then, as required for the modified Gauss laws
    of~\cite[lines~4325--4335]{FrancoRubio2025MPUGauging} when $(c,d)=T_g(a,b)$,
    \begin{align}
        U_{a,b}\ket{\xi}             & =D_{a,b}\ket{\xi},
        \label{eq:mpug_completion_agrees}                  \\
        U_{a,b}\ket{\xi}             & =U_{c,d}\ket{\eta},
        \label{eq:mpug_completion_covariance}              \\
        \lambda_{a,b}^{c,d}\ket{\xi} & =\ket{\eta}.
        \label{eq:mpug_completion_transport}
    \end{align}
\end{lemma}

\begin{proof}\leanok
    \uses{def:mpug_completion_class,
        thm:mpug_unitary_factorization_comparison_laws}
    Identity~(\ref{eq:mpug_completion_agrees}) is the defining property of a
    completion on the defect subspace. Applying it at $(a,b)$ and at $(c,d)$
    and inserting~(\ref{eq:mpug_prescribed_covariance})
    gives~(\ref{eq:mpug_completion_covariance}), an equality of vectors in the
    whole matter space. Multiplying that equality by $U_{c,d}^\dagger$ and
    using $U_{c,d}^\dagger U_{c,d}=\Id$ gives
    \begin{align}
        \lambda_{a,b}^{c,d}\ket{\xi}
         & =U_{c,d}^\dagger U_{a,b}\ket{\xi}
        =U_{c,d}^\dagger U_{c,d}\ket{\eta}
        =\ket{\eta}.
    \end{align}
    The adjoint is applied only after the equality in the whole matter space
    is available; no step asserts that $U_{c,d}^\dagger$ restricts to
    $D_{c,d}^\dagger$ on a defect subspace.
\end{proof}

\begin{definition}[Label action at one window]
    \label{def:mpug_chain_label_action}
    \lean{MPOTensor.chainLabelAction,
        MPOTensor.extractWindow_chainLabelAction_zero,
        MPOTensor.extractWindow_chainLabelAction_one}
    \leanok
    \discussion{7569}
    \uses{def:mpug_gauss_leg_action}
    Let $\alpha=(a_1,\ldots,a_N)$ be a configuration of gauge labels on a
    periodic chain of length $N\geq2$. For a window $(j,j+1)$ and $g\in G$,
    the configuration $T_{j,g}\alpha$ carries the labels
    \begin{align}
        \left(a_jg^{-1},ga_{j+1}\right)
    \end{align}
    on that window and agrees with $\alpha$ at every other site.
\end{definition}

\begin{lemma}[Composition of the label action]
    \label{thm:mpug_chain_label_action_mul}
    \lean{MPOTensor.chainLabelAction_one,
        MPOTensor.chainLabelAction_chainLabelAction,
        MPOTensor.chainLabelActionEquiv,
        MPOTensor.chainLabelActionEquiv_apply}
    \leanok
    \discussion{7569}
    \uses{def:mpug_chain_label_action}
    For all $g,h\in G$,
    \begin{align}
        T_{j,e}        & =\id,      \\
        T_{j,g}T_{j,h} & =T_{j,gh}.
        \label{eq:mpug_chain_label_action_mul}
    \end{align}
    In particular $T_{j,g}$ is a bijection of gauge-label configurations with
    inverse $T_{j,g^{-1}}$.
\end{lemma}

\begin{proof}\leanok
    \uses{def:mpug_chain_label_action}
    Outside the window both sides leave $\alpha$ unchanged. On the window,
    replacing the labels twice keeps only the second replacement, and the two
    labels are
    \begin{align}
        \left(a_jh^{-1}g^{-1},gha_{j+1}\right)
         & =\left(a_j(gh)^{-1},(gh)a_{j+1}\right).
    \end{align}
    At $g=e$ the labels are unchanged, so $T_{j,e}=\id$, and the composition
    law with $h=g^{-1}$ gives the inverse.
\end{proof}

\begin{definition}[Placed Gauss operator of a completion]
    \label{def:mpug_placed_gauss_operator}
    \lean{MPOTensor.gaussWindowOperator,
        MPOTensor.gaussWindowOperator_apply,
        MPOTensor.placedGaussOperator}
    \leanok
    \discussion{7569}
    \uses{def:mpug_gauss_operator,
        def:mpug_placed_gauss_projector,
        def:mpug_completion_class}
    Let $R\in\mathfrak C(D)$ and let a chain site carry one matter index and
    one gauge label. Relabelling the row and column coordinates of
    $\mathcal G_g$ by the bijection $\kappa$ of
    Definition~\ref{def:mpug_placed_gauss_projector} and embedding the result
    on the periodic window $(j,j+1)$ defines
    \begin{align}
        \widetilde{\mathcal G}^{j,j+1}_g(R)
         & =\operatorname{Embed}_{j,j+1}\bigl(\widetilde{\mathcal G}_g(R)\bigr),
        \label{eq:mpug_placed_gauss_operator}
    \end{align}
    the modified Gauss law
    of~\cite[lines~4325--4335]{FrancoRubio2025MPUGauging}, whose block from
    $(a,b)$ to $T_g(a,b)$ is $U_{T_g(a,b)}^\dagger U_{a,b}$.
\end{definition}

\begin{lemma}[The placed projector is the average of the placed Gauss operators]
    \label{thm:mpug_placed_gauss_projector_average}
    \lean{MPOTensor.placedGaussProjector_eq_average}
    \leanok
    \discussion{7569}
    \uses{def:mpug_placed_gauss_operator}
    For every window $(j,j+1)$,
    \begin{align}
        \widetilde P_j(R)
         & =\frac{1}{|G|}\sum_{g\in G}\widetilde{\mathcal G}^{j,j+1}_g(R).
        \label{eq:mpug_placed_gauss_projector_average}
    \end{align}
\end{lemma}

\begin{proof}\leanok
    \uses{def:mpug_gauss_projector, def:mpug_placed_gauss_operator}
    Comparing matrix entries, the relabelling is entrywise and the embedding
    multiplies each entry by the indicator that the two chain configurations
    agree outside the window. Both operations commute with the normalized sum
    over $G$ defining the local projector kernel
    in~(\ref{eq:mpug_gauss_projector}).
\end{proof}

\begin{definition}[Gauged vector of matter defect states]
    \label{def:mpug_gauged_vector}
    \lean{MPOTensor.gaugedMatter,
        MPOTensor.gaugedLabel,
        MPOTensor.gaugedVector,
        MPOTensor.gaugedVector_apply}
    \leanok
    \discussion{7569}
    Every chain configuration splits into its matter part and its gauge-label
    part. Given matter defect states $\ket{\Psi_\alpha}$ indexed by the
    gauge-label configurations $\alpha$ of a chain of length $N$, the gauged
    vector of \cite[lines~4325--4335]{FrancoRubio2025MPUGauging} is
    \begin{align}
        \ket{\Psi_{\mathrm{gauged}}}
         & =\sum_\alpha\ket{\Psi_\alpha}\otimes\ket{\alpha},
        \label{eq:mpug_gauged_vector}
    \end{align}
    whose component at a chain configuration is the component of
    $\ket{\Psi_\alpha}$ at its matter part, $\alpha$ being its gauge-label
    part.
\end{definition}

\begin{definition}[Local defect support and exact defect covariance]
    \label{def:mpug_defect_state_hypotheses}
    \lean{MPOTensor.HasDefectSupport,
        MPOTensor.HasPrescribedDefectCovariance}
    \leanok
    \discussion{7569}
    \uses{def:mpug_defect_maps, def:mpug_chain_label_action,
        def:mpug_gauged_vector}
    The matter defect states have local defect support when, for every
    $\alpha$ and every window $(j,j+1)$ whose adjacent labels are
    $(a_j,a_{j+1})$, the matter support of $\ket{\Psi_\alpha}$ lies in
    \begin{align}
        \mathcal K_{a_j,a_{j+1}}\otimes\mathcal H_{\mathrm{rest}}.
    \end{align}
    They satisfy exact defect covariance when
    \begin{align}
        \left(D_{a_j,a_{j+1}}\otimes\id_{\mathrm{rest}}\right)
        \ket{\Psi_\alpha}
         & =\left(D_{T_g(a_j,a_{j+1})}\otimes\id_{\mathrm{rest}}\right)
        \ket{\Psi_{T_{j,g}\alpha}}
        \label{eq:mpug_exact_defect_covariance}
    \end{align}
    for every $\alpha$, every window $(j,j+1)$, and every $g\in G$.
\end{definition}

\begin{lemma}[Universal invariant vector]
    \label{thm:mpug_universal_invariant_vector}
    \lean{MPOTensor.embedLocalOperator_mulVec_replaceWindow,
        MPOTensor.placedGaussOperator_mulVec_gaugedVector,
        MPOTensor.placedGaussProjector_mulVec_gaugedVector}
    \leanok
    \discussion{7569}
    \uses{def:mpug_placed_gauss_operator,
        def:mpug_placed_gauss_projector,
        def:mpug_completion_class,
        def:mpug_gauged_vector,
        def:mpug_defect_state_hypotheses}
    Assume local defect support and exact defect
    covariance~(\ref{eq:mpug_exact_defect_covariance}), which as recorded
    in~\cite{gap:fbc25_state_level_gauging_covariance} enter here as hypotheses
    rather than consequences of locally orthogonal MPS blocks. For every
    $R\in\mathfrak C(D)$, every window $(j,j+1)$, and every $g\in G$,
    \begin{align}
        \widetilde{\mathcal G}^{j,j+1}_g(R)\ket{\Psi_{\mathrm{gauged}}}
         & =\ket{\Psi_{\mathrm{gauged}}},
        \label{eq:mpug_universal_invariant_vector} \\
        \widetilde P_j(R)\ket{\Psi_{\mathrm{gauged}}}
         & =\ket{\Psi_{\mathrm{gauged}}}.
        \label{eq:mpug_universal_invariant_proj}
    \end{align}
\end{lemma}

\begin{center}
    % Formula: \widetilde{\mathcal G}^{j,j+1}_g(R)\ket{\Psi_{\mathrm{gauged}}}
    % =\ket{\Psi_{\mathrm{gauged}}}, equation
    % (\ref{eq:mpug_universal_invariant_vector}); source: arXiv:2502.20257,
    % lines 4325--4335.
    % Ink-to-index map: the wide lower box is the gauged vector, one leg per
    % chain site carrying the matter index and the gauge label of that site;
    % sites 1, j, j+1, N are drawn.  The upper box is the placed Gauss
    % operator acting on the window (j,j+1) and as the identity elsewhere.
    % Contraction set: the two window legs between the operator and the
    % vector; the output legs remain open.
    % Boundary signature: (virtual-west=0 | virtual-east=0 |
    % physical-up=4 | physical-down=0) on both sides.
    \tenkzkernel
    \begin{tenkz}[rows={op,ket}, cols=4, bonds=none]
        \tn[at=(1,2), name=G, skin=box, wide=2, ports={90@1:physical:$j$, 90@2:physical:$j{+}1$, 270@1:physical, 270@2:physical}]{\widetilde{\mathcal G}^{j,j+1}_g(R)}
        \tn[at=(2,1), name=P, skin=box, wide=4, ports={90@1:physical:$1$, 90@2:physical, 90@3:physical, 90@4:physical:$N$}]{\Psi_{\mathrm{gauged}}}
        \tn[at=(1,1), void=sealed]{}
        \tn[at=(1,4), void=sealed]{}
        \tnwire{G.270@1}{P.90@2}
        \tnwire{G.270@2}{P.90@3}
    \end{tenkz}
    \quad $=$ \quad
    \begin{tenkz}[rows={ket}, cols=4, bonds=none]
        \tn[at=(1,1), name=P, skin=box, wide=4, ports={90@1:physical:$1$, 90@2:physical:$j$, 90@3:physical:$j{+}1$, 90@4:physical:$N$}]{\Psi_{\mathrm{gauged}}}
    \end{tenkz}
\end{center}

\begin{proof}\leanok
    \uses{thm:mpug_completion_transport,
        thm:mpug_chain_label_action_mul,
        thm:mpug_placed_gauss_projector_average,
        def:mpug_gauss_operator}
    Fix an output chain configuration and let $\alpha$ be its gauge-label part,
    with $(a,b)=(a_j,a_{j+1})$ the two labels adjacent to the window. Among the
    columns contributing to this fixed output row of
    $\widetilde{\mathcal G}^{j,j+1}_g(R)$, only the gauge labels
    $T_{j,g^{-1}}\alpha$ survive. Thus, with
    $\beta=T_{j,g^{-1}}\alpha$, the component of
    $\widetilde{\mathcal G}^{j,j+1}_g(R)\ket{\Psi_{\mathrm{gauged}}}$ equals
    the component of
    \begin{align}
        \left(U_{a,b}^\dagger U_{ag,g^{-1}b}
        \otimes\id_{\mathrm{rest}}\right)\ket{\Psi_\beta},
    \end{align}
    since the labels of $\beta$ adjacent to the window are $(ag,g^{-1}b)$ and
    $T_g(ag,g^{-1}b)=(a,b)$, whence $T_{j,g}\beta=\alpha$ by the composition
    law~(\ref{eq:mpug_chain_label_action_mul}). Local defect support puts the
    matter slices of $\ket{\Psi_\beta}$ and $\ket{\Psi_\alpha}$ over that
    window in $\mathcal K_{ag,g^{-1}b}$ and $\mathcal K_{a,b}$, and exact
    covariance~(\ref{eq:mpug_exact_defect_covariance}) relates them, so the
    comparison identity~(\ref{eq:mpug_completion_transport}) turns that vector
    into $\ket{\Psi_\alpha}$. Identifying the label configuration $\beta$ with
    $\alpha$ through the bijection $T_{j,g}$ is the reindexing of the sum
    over gauge-label configurations in~(\ref{eq:mpug_gauged_vector}). Averaging~(\ref{eq:mpug_universal_invariant_vector})
    over $g$ via~(\ref{eq:mpug_placed_gauss_projector_average}) yields~(\ref{eq:mpug_universal_invariant_proj}).
\end{proof}

\begin{definition}[Gauge-invariant subspace of a completion]
    \label{def:mpug_gauge_invariant_subspace}
    \lean{MPOTensor.gaugeInvariantSubspace}
    \leanok
    \discussion{7569}
    \uses{def:mpug_common_fixed_subspace,
        def:mpug_completion_class,
        def:mpug_placed_gauss_projector}
    For $R\in\mathfrak C(D)$ and a chain of length $N\geq2$, the
    gauge-invariant subspace is the common fixed subspace
    \begin{align}
        \mathcal V_N(R)
         & =\left\{\ket{\phi}:
        \widetilde P_j(R)\ket{\phi}=\ket{\phi}
        \text{ for every }j\right\}.
        \label{eq:mpug_gauge_invariant_subspace}
    \end{align}
    As noted in~\cite[lines~4325--4335]{FrancoRubio2025MPUGauging}, the
    modified projectors on neighbouring windows need not commute.
\end{definition}

\begin{lemma}[The gauged vector spans at least one dimension]
    \label{thm:mpug_gauge_invariant_subspace_finrank}
    \lean{MPOTensor.gaugedVector_mem_gaugeInvariantSubspace,
        MPOTensor.one_le_finrank_gaugeInvariantSubspace}
    \leanok
    \discussion{7569}
    \uses{def:mpug_gauge_invariant_subspace,
        def:mpug_completion_class,
        def:mpug_gauged_vector,
        def:mpug_defect_state_hypotheses}
    Under local defect support and exact defect
    covariance~(\ref{eq:mpug_exact_defect_covariance}),
    \begin{align}
        \ket{\Psi_{\mathrm{gauged}}} & \in\mathcal V_N(R)
    \end{align}
    for every $R\in\mathfrak C(D)$. If $\ket{\Psi_{\mathrm{gauged}}}\neq0$, then
    \begin{align}
        \dim_{\C}\mathcal V_N(R) & \geq1.
    \end{align}
    Whereas \cite[line~5198]{FrancoRubio2025MPUGauging} asks how
    $\dim_{\C}\mathcal V_N(R)$ grows with $N$, neither the minimum nor the
    maximum of $\dim_{\C}\mathcal V_N(R)$ over $\mathfrak C(D)$ is asserted
    here.
\end{lemma}

\begin{proof}\leanok
    \uses{thm:mpug_universal_invariant_vector,
        thm:mpug_mem_common_fixed_subspace,
        thm:mpug_common_fixed_subspace_finrank}
    Membership is the projector identity in
    Lemma~\ref{thm:mpug_universal_invariant_vector} at every window. A
    nonzero member of a subspace forces its dimension to be at least one.
\end{proof}

\begin{definition}[Scalar cochains and normalization]
    \label{def:mpug_scalar_cochains}
    \lean{TNLean.Algebra.ScalarThreeCochain,
        TNLean.Algebra.ScalarCocycle.IsNormalized,
        TNLean.Algebra.ScalarThreeCochain.IsNormalized}
    \leanok
    \uses{def:scalar_cocycle}
    By Definition~\ref{def:scalar_cocycle}, a scalar $2$-cochain
    is a function $\beta\colon G^2\to\C^\times$. A scalar $3$-cochain is a
    function $\omega\colon G^3\to\C^\times$. Following the paper's later
    standing convention, it is normalized when, for all $g,h,k\in G$,
    \begin{align}
        \omega(e,h,k) & = 1, \\
        \omega(g,e,k) & = 1, \\
        \omega(g,h,e) & = 1.
        \label{eq:mpug_three_normalized}
    \end{align}
    The standing convention of~\cite{FrancoRubio2025MPUGauging} makes every
    fusion tensor involving the identity trivial, which is what forces all
    three identity axes. As recorded in
    \cite{gap:fbc25_three_cocycle_normalization_typo}, the three-axis form
    corrects the printed display, which repeats the third axis in place of
    the first two.
    % Source anchor: arXiv:2502.20257, standing convention at lines 1936--1937
    %   and its display eq:triv_omegas, whose printed form repeats
    %   omega(g,h,e) three times.
    % Local fix: docs/paper-gaps/fbc25_three_cocycle_normalization_typo.tex.
    The scalar $2$-cochain $\beta$ is normalized when
    $\beta(e,g)=\beta(g,e)=1$ for every $g\in G$.
\end{definition}

\begin{definition}[Scalar three-cocycle]
    \label{def:mpug_three_cocycle}
    \lean{TNLean.Algebra.ScalarThreeCochain.IsCocycle}
    \leanok
    \uses{def:mpug_scalar_cochains}
    As in~\cite{FrancoRubio2025MPUGauging}, a scalar $3$-cochain is a
    $3$-cocycle when, for all $g,h,k,l\in G$,
    \begin{align}
        \omega(gh,k,l)\omega(g,h,kl)
         & =\omega(g,h,k)\omega(g,hk,l)\omega(h,k,l).
        \label{eq:mpug_three_cocycle}
    \end{align}
    % Source anchor: arXiv:2502.20257, eq:3-cocycle.
\end{definition}

\begin{definition}[Three-coboundary and fusion-tensor gauge action]
    \label{def:mpug_fusion_gauge}
    \lean{TNLean.Algebra.ScalarThreeCochain.coboundary,
        TNLean.Algebra.ScalarThreeCochain.fusionGauge}
    \leanok
    \uses{def:mpug_scalar_cochains}
    For a scalar $2$-cochain $\beta$, define
    \begin{align}
        (d\beta)(g,h,k)
         & =\frac{\beta(g,hk)\beta(h,k)}
        {\beta(g,h)\beta(gh,k)},
        \label{eq:mpug_coboundary}        \\
        (\beta\mathbin{\triangleright}\omega)(g,h,k)
         & =(d\beta)(g,h,k)\omega(g,h,k).
        \label{eq:mpug_gauge_action}
    \end{align}
    % Source anchor: arXiv:2502.20257, eq:omegagauge.
    As in~\cite{FrancoRubio2025MPUGauging}, formula~(\ref{eq:mpug_gauge_action})
    is the change of $\omega$ induced by (\ref{eq:mpug_fusion_gauge}).
\end{definition}

\begin{lemma}[Identity and composition of fusion gauges]
    \label{thm:mpug_gauge_action}
    \lean{TNLean.Algebra.ScalarThreeCochain.fusionGauge_one,
        TNLean.Algebra.ScalarThreeCochain.fusionGauge_comp}
    \leanok
    \uses{def:mpug_fusion_gauge}
    Let $1$ denote the constant scalar $2$-cochain and let
    $(\beta\gamma)(g,h)=\beta(g,h)\gamma(g,h)$. Then
    \begin{align}
        1\mathbin{\triangleright}\omega & =\omega,
        \label{eq:mpug_gauge_identity}                                                  \\
        \gamma\mathbin{\triangleright}
        (\beta\mathbin{\triangleright}\omega)
                                        & =(\beta\gamma)\mathbin{\triangleright}\omega.
        \label{eq:mpug_gauge_composition}
    \end{align}
\end{lemma}

\begin{proof}\leanok
    \uses{def:mpug_fusion_gauge}
    The constant cochain has $d1=1$. Commutativity of $\C^\times$ gives
    $d(\beta\gamma)=(d\beta)(d\gamma)$, and substitution into
    (\ref{eq:mpug_gauge_action}) gives
    (\ref{eq:mpug_gauge_identity}) and
    (\ref{eq:mpug_gauge_composition}).
\end{proof}

\begin{lemma}[A three-coboundary is a three-cocycle]
    \label{thm:mpug_coboundary_cocycle}
    \lean{TNLean.Algebra.ScalarThreeCochain.isCocycle_coboundary}
    \leanok
    \uses{def:mpug_three_cocycle, def:mpug_fusion_gauge}
    For every scalar $2$-cochain $\beta$, the cochain $d\beta$ satisfies
    (\ref{eq:mpug_three_cocycle}).
\end{lemma}

\begin{proof}\leanok
    \uses{def:mpug_three_cocycle, def:mpug_fusion_gauge}
    Expanding the two sides gives the same product after cancellation:
    \begin{align}
         & (d\beta)(gh,k,l)(d\beta)(g,h,kl) \\
         & \quad =
        \frac{\beta(k,l)\beta(h,kl)\beta(g,hkl)}
        {\beta(ghk,l)\beta(gh,k)\beta(g,h)} \\
         & \quad =
        (d\beta)(g,h,k)(d\beta)(g,hk,l)(d\beta)(h,k,l).
    \end{align}
\end{proof}

\begin{lemma}[Products and fusion gauges preserve three-cocycles]
    \label{thm:mpug_gauge_preserves_cocycle}
    \lean{TNLean.Algebra.ScalarThreeCochain.IsCocycle.mul,
        TNLean.Algebra.ScalarThreeCochain.IsCocycle.fusionGauge}
    \leanok
    \uses{def:mpug_three_cocycle, def:mpug_fusion_gauge}
    The pointwise product of two scalar $3$-cocycles is a scalar
    $3$-cocycle. Consequently, if $\omega$ satisfies
    (\ref{eq:mpug_three_cocycle}), then so does
    $\beta\mathbin{\triangleright}\omega$.
\end{lemma}

\begin{proof}\leanok
    \uses{def:mpug_three_cocycle, thm:mpug_coboundary_cocycle}
    Multiplying the two instances of (\ref{eq:mpug_three_cocycle}) and
    reordering their factors in $\C^\times$ shows that the pointwise product of
    two $3$-cocycles is a $3$-cocycle. Applying this product identity to
    $(d\beta)\omega$, together with Lemma~\ref{thm:mpug_coboundary_cocycle}
    for $d\beta$, yields the second claim.
\end{proof}

\begin{lemma}[Normalization under a fusion gauge]
    \label{thm:mpug_gauge_preserves_normalization}
    \lean{TNLean.Algebra.ScalarThreeCochain.isNormalized_coboundary_iff,
        TNLean.Algebra.ScalarThreeCochain.isNormalized_coboundary,
        TNLean.Algebra.ScalarThreeCochain.IsNormalized.mul,
        TNLean.Algebra.ScalarThreeCochain.IsNormalized.fusionGauge}
    \leanok
    \uses{def:mpug_scalar_cochains, def:mpug_fusion_gauge}
    The $3$-coboundary $d\beta$ is normalized if and only if there is a
    constant $c\in\C^\times$ such that, for every $g\in G$,
    \begin{align}
        \beta(e,g) & =c=\beta(g,e).
        \label{eq:mpug_coboundary_normalized_iff}
    \end{align}
    In particular, a normalized scalar $2$-cochain has normalized
    $3$-coboundary. Hence $\beta\mathbin{\triangleright}\omega$ is normalized
    whenever both $\beta$ and $\omega$ are normalized.
\end{lemma}

\begin{proof}\leanok
    \uses{def:mpug_scalar_cochains, def:mpug_fusion_gauge}
    Setting one argument at a time equal to $e$ in
    (\ref{eq:mpug_coboundary}), if $d\beta$ is normalized, then
    (\ref{eq:mpug_three_normalized}) for $d\beta$ at $(e,e,g)$ and at
    $(g,e,e)$ gives $\beta(e,g)=\beta(e,e)=\beta(g,e)$. Conversely,
    (\ref{eq:mpug_coboundary_normalized_iff}) gives
    \begin{align}
        (d\beta)(e,h,k) & =1, \\
        (d\beta)(g,e,k) & =1, \\
        (d\beta)(g,h,e) & =1.
        \label{eq:mpug_coboundary_normalized_converse}
    \end{align}
    Taking the common constant to be $1$ proves the normalized-$\beta$
    corollary. Since pointwise products preserve
    (\ref{eq:mpug_coboundary_normalized_converse}), the gauged cochain
    $\beta\mathbin{\triangleright}\omega$ of (\ref{eq:mpug_gauge_action}) is
    normalized whenever both $\beta$ and $\omega$ are normalized.
\end{proof}

\begin{definition}[Cohomologous scalar three-cochains]
    \label{def:mpug_cohomologous_to}
    \lean{TNLean.Algebra.ScalarThreeCochain.CohomologousTo}
    \leanok
    \uses{def:mpug_fusion_gauge}
    Two scalar $3$-cochains $\omega$ and $\eta$ are cohomologous when
    $\omega=\beta\mathbin{\triangleright}\eta$ for some scalar
    $2$-cochain $\beta$.
\end{definition}

\begin{lemma}[Cohomologous-to is an equivalence relation]
    \label{thm:mpug_cohomologous_to_equivalence}
    \lean{TNLean.Algebra.ScalarThreeCochain.CohomologousTo.refl,
        TNLean.Algebra.ScalarThreeCochain.CohomologousTo.symm,
        TNLean.Algebra.ScalarThreeCochain.CohomologousTo.trans,
        TNLean.Algebra.ScalarThreeCochain.CohomologousTo.equivalence}
    \leanok
    \uses{def:mpug_cohomologous_to, thm:mpug_gauge_action}
    Cohomologous-to is reflexive, symmetric, and transitive.
\end{lemma}

\begin{proof}\leanok
    \uses{thm:mpug_gauge_action}
    Reflexivity uses the constant cochain $1$. If $\omega=
        \beta\mathbin{\triangleright}\eta$, then
    \begin{align}
        \beta^{-1}\mathbin{\triangleright}\omega
         & =(\beta\beta^{-1})\mathbin{\triangleright}\eta
        =\eta,
    \end{align}
    which proves symmetry. Composition formula
    (\ref{eq:mpug_gauge_composition}) proves transitivity.
\end{proof}

\begin{definition}[Trivial scalar gauge class]
    \label{def:mpug_trivial_gauge_class}
    \lean{TNLean.Algebra.ScalarThreeCochain.IsTrivialGaugeClass}
    \leanok
    \uses{def:mpug_cohomologous_to}
    A scalar $3$-cochain has trivial gauge class when it is cohomologous to
    the constant cochain $1$.
\end{definition}

\begin{definition}[Scalar three-cocycles in degree-three group cohomology]
    \label{def:mpug_h3_group_cohomology}
    \lean{TNLean.Algebra.scalarRepresentationEquiv,
        TNLean.Algebra.ScalarCocycle.toInhomogeneousCochain,
        TNLean.Algebra.ScalarCocycle.ofInhomogeneousCochain,
        TNLean.Algebra.ScalarThreeCochain.toInhomogeneousCochain,
        TNLean.Algebra.ScalarThreeCochain.ofInhomogeneousCochain,
        TNLean.Algebra.ScalarThreeCochain.toCocycles,
        TNLean.Algebra.ScalarThreeCochain.ofCocycles,
        TNLean.Algebra.ScalarThreeCochain.cocyclesEquiv,
        TNLean.Algebra.ScalarThreeCochain.anomalyClass}
    \leanok
    \uses{def:mpug_three_cocycle, thm:mpug_h2_group_cohomology_comparison}
    Giving $\C^\times$ the trivial $G$-action and writing it additively as in
    Lemma~\ref{thm:mpug_h2_group_cohomology_comparison}, a scalar
    $n$-cochain for $n=2,3$ is an inhomogeneous $n$-cochain
    $G^n\to\C^\times$ of this representation, and conversely. Under this
    identification a scalar $3$-cocycle is an inhomogeneous $3$-cocycle, and
    every inhomogeneous $3$-cocycle arises this way. The class of a scalar
    $3$-cocycle $\omega$ is its image
    \begin{align}
        [\omega] & \in H^3\bigl(G,\C^\times_{\mathrm{triv}}\bigr)
        \label{eq:mpug_h3_class}
    \end{align}
    in the degree-three cohomology of the inhomogeneous cochain complex.
\end{definition}

\begin{theorem}[Fusion-gauge classes are degree-three cohomology classes]
    \label{thm:mpug_h3_group_cohomology}
    \lean{TNLean.Algebra.groupCohomology_π_eq_zero_iff,
        TNLean.Algebra.scalarRepresentationEquiv_ofMul_eq_zero_iff,
        TNLean.Algebra.ScalarCohomology.ofMulDistribMulAction_ρ_apply,
        TNLean.Algebra.ScalarThreeCochain.d_toInhomogeneousCochain,
        TNLean.Algebra.ScalarThreeCochain.isCocycle_iff_d_eq_zero,
        TNLean.Algebra.ScalarCocycle.d_toInhomogeneousCochain_apply,
        TNLean.Algebra.ScalarCocycle.d_toInhomogeneousCochain,
        TNLean.Algebra.ScalarThreeCochain.fusionGauge_eq_iff,
        TNLean.Algebra.ScalarThreeCochain.cohomologousTo_iff_exists_d_eq,
        TNLean.Algebra.ScalarThreeCochain.cohomologousTo_iff_anomalyClass_eq,
        TNLean.Algebra.ScalarThreeCochain.isTrivialGaugeClass_iff_anomalyClass_eq_zero,
        TNLean.Algebra.ScalarThreeCochain.anomalyClass_surjective}
    \leanok
    \uses{def:mpug_h3_group_cohomology, def:mpug_fusion_gauge,
        def:mpug_cohomologous_to, def:mpug_trivial_gauge_class}
    With the trivial action, the differentials of the inhomogeneous cochain
    complex in degrees two and three are, multiplicatively,
    \begin{align}
        (d\beta)(g,h,k)
         & =\frac{\beta(h,k)\beta(g,hk)}{\beta(gh,k)\beta(g,h)},
        \label{eq:mpug_h3_d2}                                    \\
        (d\omega)(g,h,k,l)
         & =\frac{\omega(h,k,l)\omega(g,hk,l)\omega(g,h,k)}
        {\omega(gh,k,l)\omega(g,h,kl)}.
        \label{eq:mpug_h3_d3}
    \end{align}
    Hence (\ref{eq:mpug_h3_d2}) is the three-coboundary
    (\ref{eq:mpug_coboundary}), and $d\omega=1$ is the three-cocycle equation
    (\ref{eq:mpug_three_cocycle}). For scalar $3$-cocycles $\omega$ and
    $\eta$,
    \begin{align}
        \omega\text{ and }\eta\text{ are cohomologous}
         & \iff [\omega]=[\eta],
        \label{eq:mpug_h3_cohomologous} \\
        \omega\text{ has trivial gauge class}
         & \iff [\omega]=0,
        \label{eq:mpug_h3_trivial}
    \end{align}
    and every class in $H^3(G,\C^\times_{\mathrm{triv}})$ is the class of a
    scalar $3$-cocycle. No normalization of the cochains is assumed.
\end{theorem}

\begin{proof}\leanok
    \uses{def:mpug_h3_group_cohomology, def:mpug_fusion_gauge,
        def:mpug_cohomologous_to, def:mpug_trivial_gauge_class}
    In degree $n$ the inhomogeneous differential of a cochain $f$ with
    trivial action is the alternating sum of $f$ at the $n+2$ faces
    obtained by dropping the first argument, multiplying two adjacent
    arguments, or dropping the last argument. For $n=2,3$ this is
    (\ref{eq:mpug_h3_d2}) and (\ref{eq:mpug_h3_d3}) written multiplicatively.
    A cocycle has zero class exactly when it lies in the image of the
    preceding differential, since the cohomology is the cokernel of that
    differential into the cocycles. By (\ref{eq:mpug_gauge_action}),
    $\beta\mathbin{\triangleright}\eta=\omega$ holds exactly when
    $d\beta=\omega\eta^{-1}$, whose additive form is the difference of the
    two cochains. This gives (\ref{eq:mpug_h3_cohomologous}), and taking
    $\eta=1$ gives (\ref{eq:mpug_h3_trivial}). Surjectivity holds because
    every inhomogeneous $3$-cocycle is a scalar $3$-cocycle.
\end{proof}

\begin{definition}[Cyclic invariant of a scalar three-cochain]
    \label{def:mpug_cyclic_invariant}
    \lean{TNLean.Algebra.ScalarThreeCochain.cyclicInvariant}
    \leanok
    \uses{def:mpug_scalar_cochains}
    For a scalar $3$-cochain $\omega$, an element $g\in G$, and a natural
    number $n$, set
    \begin{align}
        I_{g,n}(\omega)
         & =\prod_{k=0}^{n-1}\omega(g,g^k,g).
        \label{eq:mpug_cyclic_invariant}
    \end{align}
\end{definition}

\begin{lemma}[Gauge invariance of the cyclic invariant]
    \label{thm:mpug_cyclic_invariant_gauge}
    \lean{TNLean.Algebra.ScalarThreeCochain.cyclicInvariant_coboundary,
        TNLean.Algebra.ScalarThreeCochain.cyclicInvariant_fusionGauge,
        TNLean.Algebra.ScalarThreeCochain.CohomologousTo.cyclicInvariant_eq,
        TNLean.Algebra.ScalarThreeCochain.cyclicInvariant_eq_one_of_isTrivialGaugeClass,
        TNLean.Algebra.ScalarThreeCochain.not_isTrivialGaugeClass_of_cyclicInvariant_ne_one}
    \leanok
    \uses{def:mpug_cyclic_invariant, def:mpug_fusion_gauge,
        def:mpug_cohomologous_to, def:mpug_trivial_gauge_class}
    Let $g^n=e$. For every scalar $2$-cochain $\beta$,
    \begin{align}
        I_{g,n}(d\beta)                              & =1,               \\
        I_{g,n}(\beta\mathbin{\triangleright}\omega) & =I_{g,n}(\omega).
    \end{align}
    Consequently cohomologous $3$-cochains have the same cyclic invariant,
    every $3$-cochain with trivial gauge class has $I_{g,n}(\omega)=1$, and a
    cyclic invariant different from $1$ shows that the gauge class of
    $\omega$ is nontrivial. No normalization of $\beta$ or $\omega$ is
    assumed.
\end{lemma}

\begin{proof}\leanok
    \uses{def:mpug_cyclic_invariant, def:mpug_fusion_gauge,
        def:mpug_cohomologous_to, def:mpug_trivial_gauge_class}
    By (\ref{eq:mpug_coboundary}),
    \begin{align}
        (d\beta)(g,g^k,g)
         & =\frac{\beta(g,g^{k+1})}{\beta(g,g^k)}\cdot
        \frac{\beta(g^k,g)}{\beta(g^{k+1},g)}.
    \end{align}
    Both quotients telescope over $k=0,\dots,n-1$, giving
    $\beta(g,g^n)/\beta(g,e)$ and $\beta(e,g)/\beta(g^n,g)$, which equal $1$
    because $g^n=e$. Multiplying by $\omega$ pointwise proves the second
    identity, and the remaining statements follow from
    Definitions~\ref{def:mpug_cohomologous_to}
    and~\ref{def:mpug_trivial_gauge_class}.
\end{proof}

\begin{definition}[Restriction of scalar three-cochains]
    \label{def:mpug_cochain_restriction}
    \lean{TNLean.Algebra.ScalarThreeCochain.comap,
        TNLean.Algebra.ScalarThreeCochain.comap_apply}
    \leanok
    \uses{def:mpug_scalar_cochains}
    For a group homomorphism $f\colon H\to G$, the restriction of a scalar
    $3$-cochain $\omega$ on $G$ is
    $(f^*\omega)(a,b,c)=\omega(f(a),f(b),f(c))$.
\end{definition}

\begin{lemma}[Restriction preserves cocycles and gauge classes]
    \label{thm:mpug_cochain_restriction}
    \lean{TNLean.Algebra.ScalarThreeCochain.IsCocycle.comap,
        TNLean.Algebra.ScalarThreeCochain.comap_fusionGauge,
        TNLean.Algebra.ScalarThreeCochain.CohomologousTo.comap,
        TNLean.Algebra.ScalarThreeCochain.IsTrivialGaugeClass.comap}
    \leanok
    \uses{def:mpug_cochain_restriction, def:mpug_three_cocycle,
        def:mpug_fusion_gauge, def:mpug_cohomologous_to,
        def:mpug_trivial_gauge_class}
    Restriction along $f$ sends $3$-cocycles to $3$-cocycles and satisfies
    $f^*(\beta\mathbin{\triangleright}\omega)
        =(f^*\beta)\mathbin{\triangleright}f^*\omega$, where
    $(f^*\beta)(a,b)=\beta(f(a),f(b))$. Hence it preserves the relation of
    being cohomologous and sends cochains with trivial gauge class to
    cochains with trivial gauge class.
\end{lemma}

\begin{proof}\leanok
    \uses{def:mpug_cochain_restriction, def:mpug_three_cocycle,
        def:mpug_fusion_gauge, def:mpug_cohomologous_to,
        def:mpug_trivial_gauge_class}
    Since $f(ab)=f(a)f(b)$, for all $a,b,c,d\in H$
    \begin{align}
        (f^*\omega)(ab,c,d)\,(f^*\omega)(a,b,cd)
         & =\omega\bigl(f(a)f(b),f(c),f(d)\bigr)\,
        \omega\bigl(f(a),f(b),f(c)f(d)\bigr),
        \notag                                                              \\
        (d f^*\beta)(a,b,c)
         & =\frac{\beta\bigl(f(a),f(b)f(c)\bigr)\beta\bigl(f(b),f(c)\bigr)}
        {\beta\bigl(f(a),f(b)\bigr)\beta\bigl(f(a)f(b),f(c)\bigr)}
        =(d\beta)\bigl(f(a),f(b),f(c)\bigr).
        \notag
    \end{align}
    The first line and the corresponding expansion of the right-hand side
    of (\ref{eq:mpug_three_cocycle}) turn the cocycle relation for $f^*\omega$
    into (\ref{eq:mpug_three_cocycle}) at $f(a),f(b),f(c),f(d)$. The second
    line, which is (\ref{eq:mpug_coboundary}) evaluated at $f(a),f(b),f(c)$,
    gives $d(f^*\beta)=f^*(d\beta)$, hence
    $f^*(\beta\mathbin{\triangleright}\omega)
        =(f^*\beta)\mathbin{\triangleright}f^*\omega$. The last two claims
    follow from Definitions~\ref{def:mpug_cohomologous_to}
    and~\ref{def:mpug_trivial_gauge_class}.
\end{proof}

\begin{definition}[Restriction of a group of matrix product operators]
    \label{def:mpug_family_restriction}
    \lean{MPOTensor.GroupFamily.comap,
        MPOTensor.GroupFamily.FusionData.comap}
    \leanok
    \uses{def:mpug_anomaly_three_cochain}
    For a group homomorphism $f\colon H\to G$ and a family
    $\{O_g\}_{g\in G}$ as in Definition~\ref{def:mpug_anomaly_three_cochain},
    the restricted family is $\{O_{f(a)}\}_{a\in H}$. A choice of fusion
    tensors $(F^<_{g,h},F^>_{g,h})$ of the family restricts to the fusion
    tensors $(F^<_{f(a),f(b)},F^>_{f(a),f(b)})$ of the restricted family,
    where the bond space of $O_{f(a)f(b)}$ is identified with that of
    $O_{f(ab)}$ through the equality $f(a)f(b)=f(ab)$.
\end{definition}

\begin{lemma}[Anomaly three-cochains of restrictions]
    \label{thm:mpug_family_restriction_omega}
    \lean{MPOTensor.GroupFamily.IsNormalRepresentation.comap,
        MPOTensor.GroupFamily.castMat_comap,
        MPOTensor.GroupFamily.ext_of_heq,
        MPOTensor.GroupFamily.FusionData.comap_leftV,
        MPOTensor.GroupFamily.FusionData.comap_rightV,
        MPOTensor.GroupFamily.FusionData.omega_comap,
        MPOTensor.GroupFamily.FusionData.cyclicInvariant_omega_map,
        MPOTensor.GroupFamily.FusionData.cyclicInvariant_omega_map_of_comap_eq}
    \leanok
    \uses{def:mpug_family_restriction, def:mpug_cochain_restriction,
        def:mpug_cyclic_invariant}
    The restriction of a family satisfying the hypotheses of
    Definition~\ref{def:mpug_anomaly_three_cochain} along $f\colon H\to G$
    satisfies them for $H$, and the anomaly three-cochain of the restricted
    fusion tensors is $f^*\omega$. Consequently, if $a\in H$ and $a^n=e$,
    then for every choice of fusion tensors of the restricted family, with
    anomaly three-cochain $\omega_H$,
    \begin{align}
        I_{f(a),n}(\omega) & =I_{a,n}(\omega_H).
    \end{align}
\end{lemma}

\begin{proof}\leanok
    \uses{def:mpug_family_restriction, thm:mpug_anomaly_three_cochain_spec,
        thm:mpug_anomaly_three_cocycle_gauge, thm:mpug_cyclic_invariant_gauge}
    The operator law of the restricted family is that of the family at the
    elements $f(a)$, $f(b)$ and $f(ab)=f(a)f(b)$. The two fusion trees of the
    restricted fusion tensors at $(a,b,c)$ are those of the family at
    $(f(a),f(b),f(c))$, followed by one identification of the bond space of
    $O_{f(a)f(b)f(c)}$ with that of $O_{f(abc)}$. Left multiplication by this
    permutation matrix preserves the dressed identity, so
    Lemma~\ref{thm:mpug_anomaly_three_cochain_spec} gives
    $\omega_H(a,b,c)=\omega(f(a),f(b),f(c))$ for the restricted fusion
    tensors. For any other choice of fusion tensors of the restricted family,
    Lemma~\ref{thm:mpug_anomaly_three_cocycle_gauge} makes the two
    three-cochains cohomologous, and
    Lemma~\ref{thm:mpug_cyclic_invariant_gauge} makes their cyclic
    invariants at $a$ equal. Finally $I_{a,n}(f^*\omega)=I_{f(a),n}(\omega)$
    because $f(a^k)=f(a)^k$.
\end{proof}

\begin{lemma}[Values of the anomaly three-cochain from finitely many identities]%
    \label{thm:mpug_omega_integer_certificate}
    \lean{MPSTensor.isDressedProportional_complexOfInt,
        MPSTensor.isDressedProportional_complexOfInt_of_mem,
        MPOTensor.GroupFamily.FusionData.omega_eq_one_of_leftV_eq_rightV,
        MPOTensor.GroupFamily.ofBondOne,
        MPOTensor.GroupFamily.FusionData.ofBondOne,
        MPOTensor.GroupFamily.FusionData.omega_ofBondOne}
    \leanok
    \uses{def:mpug_anomaly_three_cochain}
    Let $T$ be a tensor with integer letters, $A$ a family of integer
    matrices, $L$ and $R$ integer matrices, and $z\in\{1,-1\}$. Suppose
    that $LT^aT^b=A^aLT^b$ for all letters $a,b$ and that $LT^a=zRT^a$ for
    every letter $a$. Then $LT^{\mathbf w}=zRT^{\mathbf w}$ for every nonempty
    word $\mathbf w$. It suffices to check these identities on a list of pairs
    $(T^a,A^a)$ that contains the pair of every letter. In particular, for the
    anomaly three-cochain of Definition~\ref{def:mpug_anomaly_three_cochain}:
    if the two fusion trees of a triple coincide, then $\omega$ is one at
    that triple; and a family of bond-one tensors whose stacked products
    $O_gO_h$ equal $O_{gh}$ has, with all fusion tensors equal to one, the
    constant three-cochain one.
\end{lemma}

\begin{proof}\leanok
    \uses{def:mpug_anomaly_three_cochain, thm:mpug_anomaly_three_cochain_spec}
    Since $z^2=1$, the identity $LT^a=zRT^a$ gives $RT^a=zLT^a$, and hence
    $RT^aT^b=zA^aLT^b=A^aRT^b$. Induction on the word gives the dressed
    identity, and the entrywise image in $\C$ preserves all products. When
    the two fusion trees coincide, the scalar $1$ satisfies the dressed
    identity, so Lemma~\ref{thm:mpug_anomaly_three_cochain_spec} gives
    $\omega=1$. For bond dimension one, both fusion trees are products of
    one-by-one identity matrices and identifications of one-dimensional
    spaces, so they coincide.
\end{proof}

\begin{definition}[L-symbols and compatibility]
    \label{def:mpug_l_symbols}
    \lean{TNLean.Algebra.LSymbol,
        TNLean.Algebra.LSymbol.IsCompatible}
    \leanok
    \uses{def:mpug_scalar_cochains}
    Let $G$ act on a set $\mathsf X$. An $L$-symbol is a function
    $L\colon \mathsf X\times G^2\to\C^\times$, written $L^x_{g,h}$.
    It is compatible with a scalar $3$-cochain $\omega$ when, for every
    $x\in\mathsf X$ and $g,h,k\in G$,
    \begin{align}
        L^x_{g,hk}L^x_{h,k}
         & =\omega(g,h,k)L^{k x}_{g,h}L^x_{gh,k}.
        \label{eq:mpug_l_compatibility}
    \end{align}
    The $L$-symbols were introduced by
    \cite{GarreRubio2023MPOSym}; this is the scalar compatibility relation of
    \cite{FrancoRubio2025MPUGauging}. No transitivity or finiteness assumption
    is imposed on the action.
    % Source anchor: arXiv:2502.20257, eq:omega_and_Ls.
\end{definition}

\begin{definition}[Joint scalar gauge of L-symbols]
    \label{def:mpug_l_gauge}
    \lean{TNLean.Algebra.ActionTensorGauge,
        TNLean.Algebra.LSymbol.gauge}
    \leanok
    \uses{def:mpug_l_symbols, def:mpug_fusion_gauge}
    An action-tensor scalar gauge is a function
    $\gamma\colon G\times\mathsf X\to\C^\times$, written $\gamma^x_g$.
    Together with a fusion-tensor scalar gauge $\beta\colon G^2\to\C^\times$,
    it changes an $L$-symbol by
    \begin{align}
        (\beta,\gamma)\mathbin{\triangleright}L^x_{g,h}
         & =\frac{\gamma^x_{gh}\beta_{g,h}}
        {\gamma^{h x}_g\gamma^x_h}L^x_{g,h}.
        \label{eq:mpug_l_gauge}
    \end{align}
    This is the joint fusion-tensor and action-tensor gauge freedom in
    \cite{FrancoRubio2025MPUGauging}, written with the block label as a
    superscript throughout, as that source does for its block-dependent
    cochains.
    % Source anchor: arXiv:2502.20257, eq:Lsymbgauge (which spells the block
    %   label as a second subscript); the superscript spelling gamma^x_g is
    %   the paper's own at lines 5908 and 5988--5990.
\end{definition}

\begin{lemma}[Identity and composition of joint scalar gauges]
    \label{thm:mpug_l_gauge_composition}
    \lean{TNLean.Algebra.LSymbol.gauge_one,
        TNLean.Algebra.LSymbol.gauge_comp}
    \leanok
    \uses{def:mpug_l_gauge}
    The constant fusion and action gauges leave every $L$-symbol unchanged.
    Successive joint gauges compose by pointwise multiplication:
    \begin{align}
        (\delta,\varepsilon)\mathbin{\triangleright}
        \bigl((\beta,\gamma)\mathbin{\triangleright}L\bigr)
         & = (\beta\delta,\gamma\varepsilon)
        \mathbin{\triangleright}L.
        \label{eq:mpug_l_gauge_composition}
    \end{align}
\end{lemma}

\begin{proof}\leanok
    \uses{def:mpug_l_gauge}
    Substituting (\ref{eq:mpug_l_gauge}) twice and cancelling the intermediate
    factors yields the composition law, and the identity case
    follows by taking all four cochains equal to $1$.
\end{proof}

\begin{definition}[Normalized L-symbols and action-tensor gauges]
    \label{def:mpug_l_normalized}
    \lean{TNLean.Algebra.LSymbol.IsNormalized,
        TNLean.Algebra.ActionTensorGauge.IsNormalized}
    \leanok
    \uses{def:mpug_l_symbols, def:mpug_l_gauge}
    As in \cite{FrancoRubio2025MPUGauging}, the $L$-symbol is normalized when,
    for every $x\in\mathsf X$ and $g\in G$,
    \begin{align}
        L^x_{g,e} & =1.
        \label{eq:mpug_l_normalized_right}
    \end{align}
    \begin{align}
        L^x_{e,g} & =1.
        \label{eq:mpug_l_normalized_left}
    \end{align}
    Alongside (\ref{eq:mpug_l_normalized_right}) and
    (\ref{eq:mpug_l_normalized_left}), the action-tensor gauge is normalized
    when
    \begin{align}
        \gamma^x_e & =1
        \label{eq:mpug_action_gauge_normalized}
    \end{align}
    for every $x\in\mathsf X$. This last condition is not a displayed
    equation of the source: it is the scalar shadow of the standing
    convention that every action tensor involving the identity is trivial.
    % Source anchor: arXiv:2502.20257, eq:triv_Ls for the L-symbol axes;
    %   gamma^x_e = 1 is not a labelled source equation but the standing
    %   convention at lines 1936--1937 that any fusion or action tensor
    %   involving e is trivial, applied to the action-tensor gauge scalar of
    %   eq:scalar_act_ten.
\end{definition}

\begin{definition}[Restricted scalar gauges]
    \label{def:mpug_restricted_scalar_gauge}
    \lean{TNLean.Algebra.ScalarCocycle.IsInverseNormalized,
        TNLean.Algebra.RestrictedFusionGauge,
        TNLean.Algebra.RestrictedScalarGauge}
    \leanok
    \discussion{7700}
    \uses{def:mpug_fusion_gauge, def:mpug_l_gauge, def:mpug_l_normalized}
    A fusion scalar $\beta\colon G^2\to\C^\times$ is
    \emph{inverse-normalized} when, for every $g\in G$,
    \begin{align}
        \beta_{g,e}=\beta_{e,g}=\beta_{g,g^{-1}}
         & =1.
        \label{eq:mpug_restricted_fusion_gauge}
    \end{align}
    Following the standing restriction imposed at
    \cite[lines~2050--2054]{FrancoRubio2025MPUGauging} after fixing the
    fusion and action gauges of Definition~\ref{def:mpug_l_gauge}, a
    restricted scalar gauge is a pair $(\beta,\gamma)$ in which $\beta$ is
    inverse-normalized and, for every $x\in\mathsf X$,
    \begin{align}
        \gamma^x_e & =1.
        \label{eq:mpug_restricted_action_gauge}
    \end{align}
    % Source anchor: arXiv:2502.20257, unlabelled display at lines 2050--2054,
    % immediately after eq:zetas_in_gauge.
\end{definition}

\begin{lemma}[Joint scalar gauges preserve L-symbol compatibility]
    \label{thm:mpug_l_gauge_preservation}
    \lean{TNLean.Algebra.LSymbol.IsCompatible.gauge,
        TNLean.Algebra.LSymbol.gauge_apply_right_one,
        TNLean.Algebra.LSymbol.gauge_apply_left_one,
        TNLean.Algebra.LSymbol.isNormalized_gauge_iff,
        TNLean.Algebra.LSymbol.IsNormalized.gauge}
    \leanok
    \uses{def:mpug_fusion_gauge, def:mpug_l_symbols,
        def:mpug_l_gauge, def:mpug_l_normalized}
    If $L$ is compatible with $\omega$, then
    $(\beta,\gamma)\mathbin{\triangleright}L$ is compatible with
    $\beta\mathbin{\triangleright}\omega$. On the identity axes, the gauge
    formula is
    \begin{align}
        ((\beta,\gamma)\mathbin{\triangleright}L)^x_{g,e}
         & =\frac{\beta_{g,e}}{\gamma^x_e}L^x_{g,e}.
        \label{eq:mpug_l_gauge_right}
    \end{align}
    \begin{align}
        ((\beta,\gamma)\mathbin{\triangleright}L)^x_{e,g}
         & =\frac{\beta_{e,g}}{\gamma^{g x}_e}L^x_{e,g}.
        \label{eq:mpug_l_gauge_left}
    \end{align}
    Thus the gauged $L$-symbol is normalized exactly when
    $\beta_{g,e}L^x_{g,e}=\gamma^x_e$ and
    $\beta_{e,g}L^x_{e,g}=\gamma^{g x}_e$ for all $x$ and $g$. In particular,
    normalized $\beta$, $\gamma$, and $L$ give a normalized gauged $L$-symbol.
\end{lemma}

\begin{proof}\leanok
    \uses{def:mpug_fusion_gauge, def:mpug_l_symbols, def:mpug_l_gauge,
        def:mpug_l_normalized}
    With $L'=(\beta,\gamma)\mathbin{\triangleright}L$, substituting
    (\ref{eq:mpug_l_gauge}) into the left side of
    (\ref{eq:mpug_l_compatibility}) for $L'$ cancels the factor
    $\gamma^x_{hk}$ and leaves
    \begin{align}
        {L'}^x_{g,hk}{L'}^x_{h,k}
         & =\frac{\gamma^x_{ghk}}
        {\gamma^{(hk)x}_g\gamma^{k x}_h\gamma^x_k}\,
        \beta_{g,hk}\beta_{h,k}\,L^x_{g,hk}L^x_{h,k}.
        \label{eq:mpug_l_gauge_left_side}
    \end{align}
    The same substitution on the right side cancels
    $\gamma^{k x}_{gh}$ and leaves
    \begin{align}
        {L'}^{k x}_{g,h}{L'}^x_{gh,k}
         & =\frac{\gamma^x_{ghk}}
        {\gamma^{h(k x)}_g\gamma^{k x}_h\gamma^x_k}\,
        \beta_{g,h}\beta_{gh,k}\,L^{k x}_{g,h}L^x_{gh,k}.
        \label{eq:mpug_l_gauge_right_side}
    \end{align}
    The action law $h(k x)=(hk)x$ makes the two action-tensor prefactors of
    (\ref{eq:mpug_l_gauge_left_side}) and
    (\ref{eq:mpug_l_gauge_right_side}) equal. Applying
    (\ref{eq:mpug_l_compatibility}) for $L$ and $\omega$ to the right of
    (\ref{eq:mpug_l_gauge_left_side}) therefore gives
    \begin{align}
        {L'}^x_{g,hk}{L'}^x_{h,k}
         & =\frac{\beta_{g,hk}\beta_{h,k}}
        {\beta_{g,h}\beta_{gh,k}}\,\omega(g,h,k)\,
        {L'}^{k x}_{g,h}{L'}^x_{gh,k},
        \label{eq:mpug_l_gauge_compatibility}
    \end{align}
    whose scalar factor is
    $(\beta\mathbin{\triangleright}\omega)(g,h,k)$ by
    (\ref{eq:mpug_coboundary}) and
    (\ref{eq:mpug_gauge_action}). Setting either group argument equal
    to $e$ in (\ref{eq:mpug_l_gauge}) gives
    (\ref{eq:mpug_l_gauge_right}) and
    (\ref{eq:mpug_l_gauge_left}), from which the normalization statements
    follow.
\end{proof}

\begin{definition}[Unimodular square root and modulus of a phase]
    \label{def:mpug_phase_root}
    \lean{TNLean.Algebra.phaseRoot,
        TNLean.Algebra.modulusUnit}
    \leanok
    For $u\in\C^\times$, let $s_u=e^{i\arg(u)/2}$ be the unimodular square root
    of the phase of $u$, with $|u|$ viewed as an element of $\C^\times$.
\end{definition}

\begin{lemma}[Unimodular square roots of phases]
    \label{thm:mpug_phase_root}
    \lean{TNLean.Algebra.star_phaseRoot,
        TNLean.Algebra.phaseRoot_mul_phaseRoot_mul_modulusUnit,
        TNLean.Algebra.star_eq_self_of_star_eq_ofReal_mul}
    \leanok
    \uses{def:mpug_phase_root}
    For $u\in\C^\times$, $\overline{s_u}=s_u^{-1}$ and $u=s_u^2\,|u|$. If
    $w\in\C^\times$ and $\overline w=p\,w$ for a positive real number $p$,
    then $p=1$ and $w$ is real.
\end{lemma}

\begin{proof}\leanok
    The first two claims are $\overline{e^{i\theta}}=e^{-i\theta}$ and the
    polar form $u=|u|e^{i\arg u}$. For the third, taking absolute values in
    $\overline w=p\,w$ gives $p=1$.
\end{proof}

\begin{lemma}[Time reversal makes the anomaly real]
    \label{thm:mpug_time_reversal_real_anomaly}
    \lean{TNLean.Algebra.ScalarThreeCochain.coboundary_mul,
        TNLean.Algebra.ScalarThreeCochain.star_fusionGauge,
        TNLean.Algebra.ScalarThreeCochain.coboundary_modulusUnit_pos,
        TNLean.Algebra.ScalarThreeCochain.star_fusionGauge_phaseRoot,
        TNLean.Algebra.ScalarThreeCochain.exists_cohomologousTo_star_eq_self}
    \leanok
    \uses{def:mpug_fusion_gauge, def:mpug_cohomologous_to,
        def:mpug_phase_root}
    Let $\omega$ be a scalar $3$-cochain and $\beta$ a scalar $2$-cochain with
    $\beta\mathbin{\triangleright}\omega=\overline\omega$. Then
    $s\mathbin{\triangleright}\omega$ is real-valued, where
    $s(g,h)=s_{\beta(g,h)}$ is the unimodular square root of the phase of
    $\beta(g,h)$ from Definition~\ref{def:mpug_phase_root}. Consequently, a
    scalar $3$-cochain cohomologous to its complex conjugate is cohomologous
    to a real-valued one. This is the conclusion
    of~\cite[Section~4.3]{GarreRubio2023MPOSym} that, when time reversal
    commutes with a group matrix product operator symmetry, the three-cocycle
    can be gauge transformed to a real number; no finiteness of $G$ is
    needed.
    % Source anchor: arXiv:2203.12563, lines 873--891.
\end{lemma}

\begin{proof}\leanok
    \uses{def:mpug_fusion_gauge, thm:mpug_phase_root}
    The coboundary is multiplicative in $\beta$ and commutes with complex
    conjugation. Writing $\beta=s^2N$ with $N(g,h)=|\beta(g,h)|$ and using
    $\overline s=s^{-1}$ gives
    \begin{align}
        \overline{s\mathbin{\triangleright}\omega}
         & =s^{-1}\mathbin{\triangleright}\overline\omega
        =s^{-1}\mathbin{\triangleright}(s^2N\mathbin{\triangleright}\omega)
        =N\mathbin{\triangleright}(s\mathbin{\triangleright}\omega).
    \end{align}
    The values of $dN$ are positive real numbers, so
    Lemma~\ref{thm:mpug_phase_root} gives $dN=1$ and
    $\overline{s\mathbin{\triangleright}\omega}
        =s\mathbin{\triangleright}\omega$. If
    $\omega=\beta\mathbin{\triangleright}\overline\omega$, conjugating gives
    $\overline\beta\mathbin{\triangleright}\omega=\overline\omega$, and the
    first part applies.
\end{proof}

\begin{lemma}[Time reversal makes the anomaly and the L-symbols jointly real]
    \label{thm:mpug_time_reversal_real_l_symbols}
    \lean{TNLean.Algebra.LSymbol.star_gauge,
        TNLean.Algebra.LSymbol.gauge_modulusUnit_pos,
        TNLean.Algebra.LSymbol.star_gauge_phaseRoot,
        TNLean.Algebra.LSymbol.exists_gauge_star_eq_self}
    \leanok
    \uses{def:mpug_fusion_gauge, def:mpug_l_gauge, def:mpug_phase_root}
    Let $\omega$ be a scalar $3$-cochain, $L$ an $L$-symbol, $\beta$ a
    scalar $2$-cochain and $\gamma$ an action-tensor scalar gauge with
    $\beta\mathbin{\triangleright}\omega=\overline\omega$ and
    $(\beta,\gamma)\mathbin{\triangleright}L=\overline L$. Let $s$ and $t$ be
    the unimodular square roots of the phases of $\beta$ and $\gamma$. Then
    $s\mathbin{\triangleright}\omega$ and
    $(s,t)\mathbin{\triangleright}L$ are both real-valued. This result is not stated
    in~\cite[Section~4.3]{GarreRubio2023MPOSym}.
    The source derives $\beta\mathbin{\triangleright}\omega=\overline\omega$
    and, with the fusion-tensor phase omitted,
    $(1,\gamma)\mathbin{\triangleright}L=\overline L$; the hypothesis here,
    with the same $\beta$ in both relations, is the gauge-consistent pair
    and is assumed rather than derived from time reversal.
    % Source anchor: arXiv:2203.12563, lines 873--891 and 935--957.
\end{lemma}

\begin{proof}\leanok
    \uses{def:mpug_l_gauge, thm:mpug_l_gauge_composition,
        thm:mpug_time_reversal_real_anomaly, thm:mpug_phase_root}
    The claim for $\omega$ is Lemma~\ref{thm:mpug_time_reversal_real_anomaly}.
    For $L$, writing $\beta=s^2N$ and $\gamma=t^2M$ with
    $N_{g,h}=|\beta_{g,h}|$ and $M^x_g=|\gamma^x_g|$ gives, as in that proof,
    $\overline{(s,t)\mathbin{\triangleright}L}
        =(N,M)\mathbin{\triangleright}((s,t)\mathbin{\triangleright}L)$, whose
    prefactor
    $M^x_{gh}N_{g,h}/(M^{h x}_gM^x_h)$ is a positive real number, hence
    equal to one by Lemma~\ref{thm:mpug_phase_root}.
\end{proof}

\begin{lemma}[The L-symbol relation imposed by time reversal]
    \label{thm:mpug_time_reversal_l_symbol_relation}
    \lean{TNLean.Algebra.LSymbol.star_apply_eq_gauge_one_apply,
        TNLean.Algebra.LSymbol.phase_sq_eq_of_gauge_eq_star}
    \leanok
    \uses{def:mpug_l_gauge}
    Let $T\neq0$ and $P$ be vectors in a complex vector space and let
    $L^x_{g,h}$ and $\gamma$ be as in Definition~\ref{def:mpug_l_gauge}.
    If $P=\overline{L^x_{g,h}}\,T$ and
    $\gamma^x_h\gamma^{h x}_gP=L^x_{g,h}\gamma^x_{gh}T$, then
    \begin{align}
        \frac{\overline{L^x_{g,h}}}{L^x_{g,h}}
         & =\frac{\gamma^x_{gh}}{\gamma^x_h\gamma^{h x}_g}.
        \label{eq:mpug_time_reversal_l_symbol}
    \end{align}
    More generally, if $(\beta,\gamma)\mathbin{\triangleright}L=\overline L$,
    then for all $x$, $g$ and $h$,
    \begin{align}
        \left(\frac{L^x_{g,h}}{|L^x_{g,h}|}\right)^2
         & =\frac{\gamma^{h x}_g\gamma^x_h}{\gamma^x_{gh}\beta_{g,h}}.
    \end{align}
    In~\cite[Section~4.3]{GarreRubio2023MPOSym}, $T$ and $P$ are the
    complex conjugates of the two composite reductions compared by the
    $L$-symbol, the first hypothesis is the conjugate of the defining
    relation of $L$, and the second is the display derived from time
    reversal. The source prints the consequence as
    $|L^x_{g,h}|^2=\gamma^x_{gh}/(\gamma^x_h\gamma^{h x}_g)$; that form fails
    for an on-site $\mathbb Z_2\times\mathbb Z_2$ symmetry, and
    (\ref{eq:mpug_time_reversal_l_symbol}) is the local correction recorded
    in~\cite{gap:glm23_time_reversal_l_symbol}. The second display is the
    corrected sense in which the $L$-symbols squared are gauge
    transformations: the square of the phase of $L$ is, while its modulus is
    unconstrained.
    % Source anchor: arXiv:2203.12563, lines 935--957 (TRSF1group at
    %   line 954).
    % Local fix: docs/paper-gaps/glm23_time_reversal_l_symbol.tex.
\end{lemma}

\begin{proof}\leanok
    \uses{def:mpug_l_gauge}
    Substituting the first hypothesis into the second gives
    $\gamma^x_h\gamma^{h x}_g\overline{L^x_{g,h}}\,T
        =L^x_{g,h}\gamma^x_{gh}T$, and $T\neq0$ gives
    (\ref{eq:mpug_time_reversal_l_symbol}). For the second claim,
    $\overline L=cL$ with
    $c=\gamma^x_{gh}\beta_{g,h}/(\gamma^{h x}_g\gamma^x_h)$, so
    $|L|^2=\overline LL=cL^2$ and $(L/|L|)^2=c^{-1}$.
\end{proof}

\begin{lemma}[A block fixed by a subgroup trivializes the restricted three-cocycle]
    \label{thm:mpug_fixed_block_trivializes}
    \lean{TNLean.Algebra.LSymbol.isTrivialGaugeClass_comap_of_isCompatible,
        TNLean.Algebra.LSymbol.isTrivialGaugeClass_of_isCompatible_of_forall_smul_eq}
    \leanok
    \uses{def:mpug_l_symbols, def:mpug_cochain_restriction,
        def:mpug_fusion_gauge, def:mpug_trivial_gauge_class}
    Let $L$ be compatible with $\omega$, let $f\colon H\to G$ be a group
    homomorphism, and let $x\in\mathsf X$ satisfy $f(a)x=x$ for every
    $a\in H$. Then $f^*\omega$ has trivial gauge class; explicitly,
    $f^*\omega=d\lambda$ with $\lambda(a,b)=L^x_{f(a),f(b)}$. In particular,
    if some $x$ is fixed by all of $G$, then $\omega$ has trivial gauge class.
    This provides the scalar step both for the statement of
    \cite[lines~740--745]{GarreRubio2023MPOSym} that a subgroup leaving a block
    invariant trivializes the three-cocycle and for the statement of
    \cite[line~1925]{FrancoRubio2025MPUGauging} that a single invariant block
    reduces (\ref{eq:mpug_l_compatibility}) to the three-coboundary condition.
\end{lemma}

\begin{proof}\leanok
    \uses{def:mpug_l_symbols, def:mpug_fusion_gauge,
        def:mpug_cochain_restriction, def:mpug_trivial_gauge_class}
    For $a,b,c\in H$, with $g=f(a)$, $h=f(b)$, and $k=f(c)$, the invariance
    $kx=x$ turns (\ref{eq:mpug_l_compatibility}) into
    \begin{align}
        \omega(g,h,k)
         & =\frac{L^x_{g,hk}L^x_{h,k}}{L^x_{g,h}L^x_{gh,k}},
    \end{align}
    which is $(d\lambda)(a,b,c)$ by (\ref{eq:mpug_coboundary}) and the
    multiplicativity of $f$. Thus
    $f^*\omega=\lambda\mathbin{\triangleright}1$. For a block fixed by all
    of $G$, taking $f$ to be the identity completes the proof.
\end{proof}

\begin{lemma}[Action of the restricted scalar gauge group]
    \label{thm:mpug_restricted_scalar_gauge_action}
    \lean{TNLean.Algebra.ScalarCocycle.isInverseNormalized_one,
        TNLean.Algebra.ScalarCocycle.IsInverseNormalized.mul,
        TNLean.Algebra.ScalarCocycle.IsInverseNormalized.inv,
        TNLean.Algebra.ActionTensorGauge.isNormalized_one,
        TNLean.Algebra.ActionTensorGauge.IsNormalized.mul,
        TNLean.Algebra.ActionTensorGauge.IsNormalized.inv,
        TNLean.Algebra.RestrictedFusionGauge.one,
        TNLean.Algebra.RestrictedFusionGauge.mul,
        TNLean.Algebra.RestrictedFusionGauge.inv,
        TNLean.Algebra.RestrictedScalarGauge.one,
        TNLean.Algebra.RestrictedScalarGauge.mul,
        TNLean.Algebra.RestrictedScalarGauge.inv,
        TNLean.Algebra.RestrictedFusionGauge.fusionGauge_one,
        TNLean.Algebra.RestrictedFusionGauge.isNormalized_fusionGauge,
        TNLean.Algebra.RestrictedFusionGauge.isCocycle_fusionGauge,
        TNLean.Algebra.RestrictedFusionGauge.fusionGauge_mul,
        TNLean.Algebra.RestrictedFusionGauge.fusionGauge_inv,
        TNLean.Algebra.RestrictedScalarGauge.gauge_one,
        TNLean.Algebra.RestrictedScalarGauge.isNormalized_gauge,
        TNLean.Algebra.RestrictedScalarGauge.isCompatible_gauge,
        TNLean.Algebra.RestrictedScalarGauge.gauge_mul,
        TNLean.Algebra.RestrictedScalarGauge.gauge_inv}
    \leanok
    \discussion{7700}
    \uses{def:mpug_restricted_scalar_gauge}
    Let $\kappa=(\beta,\gamma)$ be a restricted scalar gauge. If $\omega$ is
    a normalized scalar $3$-cocycle and $L$ is a normalized $L$-symbol
    compatible with $\omega$, then
    \begin{align}
        \omega' & =\beta\mathbin{\triangleright}\omega,     \\
        L'      & =(\beta,\gamma)\mathbin{\triangleright}L.
        \label{eq:mpug_restricted_gauge_action}
    \end{align}
    are again normalized, $\omega'$ is a scalar $3$-cocycle, and $L'$ is
    compatible with $\omega'$. The identity gauge acts trivially, successive
    restricted gauges compose by pointwise multiplication, and applying
    $\kappa^{-1}$ after $\kappa$ recovers $(\omega,L)$.
    These are scalar statements only.
    % The theorem does not choose fusion tensors, action tensors, or
    % tensor-level gauges satisfying the same normalization. Tensor-level
    % restricted gauge choices remain issue #7325.
\end{lemma}

\begin{proof}\leanok
    \uses{def:mpug_restricted_scalar_gauge,
        thm:mpug_gauge_preserves_cocycle,
        thm:mpug_gauge_preserves_normalization,
        thm:mpug_l_gauge_composition,
        thm:mpug_l_gauge_preservation}
    The three specified values of $\beta$ and the identity value of $\gamma$
    remain equal to $1$ under pointwise products and inverses, so the general
    preservation laws apply to~(\ref{eq:mpug_restricted_gauge_action}).
    The group laws and the composition formulas for the two gauge actions then
    give the identity, composition, and inverse assertions.
\end{proof}

\begin{lemma}[Cocycles and classes for trivial coefficients]
    \label{thm:mpug_h2_trivial_coefficients}
    \lean{TNLean.Algebra.isMulCocycle₂_iff,
        TNLean.Algebra.H2π_eq_iff_isMulCoboundary₂_div}
    \leanok
    \uses{def:is_coboundary, def:cohomologous_to, def:is_cocycle, def:h2_cx}
    Let $G$ act trivially on $\C^\times$, and let
    $f(g,h)=\omega(g,h)$. The multiplicative cocycle equation
    \begin{align}
        f(gh,k)f(g,h)
         & =f(h,k)f(g,hk)
        \label{eq:mpug_h2_cocycle_comparison}
    \end{align}
    is equivalent to Definition~\ref{def:is_cocycle}.

    If $\pi$ denotes the canonical map from multiplicative $2$-cocycles to
    degree-two group cohomology, then for any two cocycles $\omega_1$ and
    $\omega_2$,
    \begin{align}
        \omega_1\sim\omega_2
         & \quad\Longleftrightarrow\quad
        \pi(\omega_1)=\pi(\omega_2).
        \label{eq:mpug_h2_class_comparison}
    \end{align}
    % Source anchor: arXiv:2502.20257, lines 1927--1931.
\end{lemma}

\begin{proof}\leanok
    \uses{def:is_coboundary, def:cohomologous_to, def:is_cocycle,
        thm:coboundary_iff_cohomologous_one}
    Since $\C^\times$ is commutative,
    (\ref{eq:mpug_h2_cocycle_comparison}) is obtained from the cocycle
    equation by interchanging the two factors on each side.
    By Lemma~\ref{thm:coboundary_iff_cohomologous_one},
    $\omega_1\sim\omega_2$ holds when $\omega_1\omega_2^{-1}$ is a
    coboundary. In additive notation this ratio is the difference of the
    corresponding cochains, so the quotient criterion for $\pi$ gives
    (\ref{eq:mpug_h2_class_comparison}).
\end{proof}

\begin{lemma}[Every class has a scalar cocycle representative]
    \label{thm:mpug_h2_group_cohomology_comparison}
    \lean{TNLean.Algebra.ScalarCocycle.exists_H2π_eq}
    \leanok
    \uses{def:h2_cx, thm:mpug_h2_trivial_coefficients}
    Endowing $\C^\times$ with the trivial $G$-action, every element of
    $H^2\bigl(G,\C^\times_{\mathrm{triv}}\bigr)$ is the class $\pi(\omega)$ of
    a scalar $2$-cocycle $\omega$. Together with
    (\ref{eq:mpug_h2_class_comparison}), this identifies the
    $H^2(H,\C^\times)$ appearing in the pair $(H,\psi)$ of
    \cite[lines~1927--1931]{FrancoRubio2025MPUGauging} with the scalar
    $2$-cocycles modulo $\sim$:
    \begin{align}
        Z^2(G,\C^\times)/{\sim}
         & \simeq H^2\bigl(G,\C^\times_{\mathrm{triv}}\bigr).
        \label{eq:mpug_h2_equivalence}
    \end{align}
\end{lemma}

\begin{proof}\leanok
    \uses{thm:mpug_h2_trivial_coefficients}
    Currying and uncurrying identify a multiplicative representative
    $f:G\times G\to\C^\times$ in the inhomogeneous cochain complex with a
    scalar $2$-cochain, and by (\ref{eq:mpug_h2_cocycle_comparison}) the
    cocycle conditions agree.
\end{proof}

\begin{definition}[Pullback of scalar cohomology]
    \label{def:mpug_h2_pullback}
    \lean{TNLean.Algebra.ScalarCocycle.comap,
        TNLean.Algebra.ScalarCocycle.comap_apply,
        TNLean.Algebra.ScalarCocycle.mapCocycles₂_cocyclesOfIsMulCocycle₂,
        TNLean.Algebra.ScalarCocycle.map_H2π}
    \leanok
    \uses{def:h2_cx}
    Let $f:K\to G$ be a group homomorphism. For a scalar $2$-cochain
    $\omega:G\times G\to\C^\times$, define
    $(f^*\omega)(a,b)=\omega(f(a),f(b))$.
    The restriction map
    $f^*:H^2(G,\C^\times)\to H^2(K,\C^\times)$ induced by $f$, with the
    identity of $\C^\times$ as coefficient map, satisfies
    $f^*[\omega]=[f^*\omega]$ for every $2$-cocycle $\omega$.
\end{definition}

\begin{lemma}[Pullback preserves cocycles and cohomology]
    \label{thm:mpug_h2_pullback_preserves}
    \lean{TNLean.Algebra.ScalarCocycle.isMulCocycle₂_comap,
        TNLean.Algebra.isMulCoboundary₂_div_comap}
    \leanok
    \uses{def:mpug_h2_pullback, def:cohomologous_to, def:is_cocycle}
    Let $f:K\to G$ be a group homomorphism. If $\omega$ is a scalar
    $2$-cocycle on $G$, then $f^*\omega$ is a scalar $2$-cocycle on $K$.
    If $\omega_1$ and $\omega_2$ are cohomologous scalar $2$-cochains on $G$,
    then $f^*\omega_1$ and $f^*\omega_2$ are cohomologous on $K$.
\end{lemma}

\begin{proof}\leanok
    \uses{def:mpug_h2_pullback, def:cohomologous_to, def:is_cocycle}
    For $a,b,c\in K$, substituting $f(a),f(b),f(c)$ into the cocycle equation
    for $\omega$ and using $f(ab)=f(a)f(b)$ and $f(bc)=f(b)f(c)$ shows that
    $f^*\omega$ is a $2$-cocycle.
    If $\omega_1(g,h)=\varphi(g)\varphi(h)\varphi(gh)^{-1}\omega_2(g,h)$,
    then with $\eta=\varphi\circ f$, substituting $g=f(a)$ and $h=f(b)$ gives
    $(f^*\omega_1)(a,b)=\eta(a)\eta(b)\eta(ab)^{-1}(f^*\omega_2)(a,b)$,
    so $\eta$ gives the required cohomology relation.
\end{proof}

\begin{lemma}[Restriction is the inhomogeneous restriction map]
    \label{lem:mpug_restriction_inhomogeneous}
    \lean{TNLean.Algebra.scalarRepresentationRes,
        TNLean.Algebra.ScalarCocycle.toInhomogeneousCochain_comap,
        TNLean.Algebra.ScalarThreeCochain.toInhomogeneousCochain_comap}
    \leanok
    \uses{def:mpug_h2_pullback, def:mpug_cochain_restriction,
        def:mpug_h3_group_cohomology}
    Let $f:K\to G$ be a group homomorphism. Under the identification of scalar
    $2$- and $3$-cochains with inhomogeneous cochains valued in the trivial
    representation $\C^\times_{\mathrm{triv}}$, the pullback $f^*$ in degrees
    two and three is the restriction map of inhomogeneous cochains along $f$,
    $x\mapsto x\circ(f\times\cdots\times f)$, whose coefficient map is the
    identity of $\C^\times$.
\end{lemma}

\begin{proof}\leanok
    \uses{def:mpug_h2_pullback, def:mpug_cochain_restriction,
        def:mpug_h3_group_cohomology}
    Both sides send a cochain $x$ to $(g_1,\dots,g_n)\mapsto x(f(g_1),\dots,f(g_n))$.
\end{proof}

% Partial support for #7361 (child #7754), not the full prop:BI_psi criterion.
% No U(1) restriction naturality or tensor-level block-independence claim.
\begin{lemma}[Exact extension of a subgroup cocycle]
    \label{thm:mpug_h2_exact_extension}
    \lean{TNLean.Algebra.ScalarCocycle.H2π_mem_range_map_subtype_iff,
        TNLean.Algebra.ScalarCocycle.isMulCocycle₂_coboundary_mul}
    \leanok
    \uses{def:mpug_h2_pullback, def:is_cocycle}
    Let $H\le G$ be any subgroup, let $\iota:H\hookrightarrow G$ be its
    inclusion, and let $\psi:H\times H\to\C^\times$ be a $2$-cocycle.
    Then $[\psi]$ belongs to the image of $\iota^*$ if and only if there is
    a $2$-cocycle $\Psi:G\times G\to\C^\times$ with
    $\Psi(a,b)=\psi(a,b)$ for every $a,b\in H$.
    No normalization of either cocycle is required.
    This is the representative adjustment implicit in the extension
    condition of \cite[lines~7042--7064]{FrancoRubio2025MPUGauging},
    here stated with $\C^\times$ coefficients.
    % Source: Papers/2502.20257/main.tex:7042--7064.
\end{lemma}

\begin{proof}\leanok
    \uses{def:mpug_h2_pullback, def:cohomologous_to, def:is_cocycle}
    If $[\psi]=\iota^*[\omega]$, there is a cochain
    $\varphi:H\to\C^\times$ such that
    $\psi(a,b)=\varphi(a)\varphi(b)\varphi(ab)^{-1}\omega(a,b)$.
    Extending $\varphi$ to a function $\Phi:G\to\C^\times$ that takes value $1$
    outside $H$, define
    $\Psi(g,h)=\Phi(g)\Phi(h)\Phi(gh)^{-1}\omega(g,h)$.
    In the cocycle equation for $\Psi$, cancellation leaves the common
    factor $\Phi(g)\Phi(h)\Phi(k)\Phi(ghk)^{-1}$ multiplied by the two sides
    of the cocycle equation for $\omega$. Hence $\Psi$ is a cocycle, and
    $\Phi|_H=\varphi$ gives $\Psi|_{H\times H}=\psi$ exactly.
    Conversely, any exact extension $\Psi$ satisfies
    $\iota^*[\Psi]=[\psi]$.
\end{proof}

\begin{lemma}[Circle and nonzero-complex coefficients in degree two]
    \label{thm:mpug_circle_complex_units_h2}
    \lean{TNLean.Algebra.Complex.div_toUnits_unitsPhase,
        TNLean.Algebra.ScalarCocycle.circlePhase_isMulCocycle₂,
        TNLean.Algebra.ScalarCocycle.circlePhaseInclusion_isMulCocycle₂,
        TNLean.Algebra.ScalarCocycle.H2π_circlePhaseInclusion,
        TNLean.Algebra.ScalarCocycle.circlePhaseMathlibCocycle,
        TNLean.Algebra.ScalarCocycle.H2π_map_circlePhaseMathlibCocycle,
        TNLean.Algebra.circleH2ToComplexUnitsH2_injective,
        TNLean.Algebra.circleH2ToComplexUnitsH2_surjective,
        TNLean.Algebra.circleH2EquivComplexUnitsH2}
    \leanok
    \uses{thm:mpug_h2_group_cohomology_comparison,
        thm:mpug_positive_generalized_cocycle}
    Let $G$ be finite, with $U(1)$ and $\C^\times$ carrying the trivial
    $G$-actions.  Inclusion of the unit circle induces an equivalence
    \begin{align}
        H^2(G,U(1))
         & \simeq H^2(G,\C^\times).
        \label{eq:mpug_circle_complex_units_h2}
    \end{align}
    More concretely, for a $\C^\times$-valued cocycle $\omega$, define its
    pointwise phase and radial residual by
    \begin{align}
        u(g,h)
         & =\frac{\omega(g,h)}{|\omega(g,h)|}\in U(1),
         &
        r(g,h)
         & =\frac{\omega(g,h)}{u(g,h)}=|\omega(g,h)|.
        \label{eq:mpug_circle_phase_residual}
    \end{align}
    Then $u$ is a cocycle, $r$ is a positive coboundary, and hence $\omega$
    is cohomologous to the inclusion of $u$.  The inclusion map is injective
    for every group because pointwise phase is its left inverse; finiteness is
    used only to prove surjectivity.

    As recorded in \cite{gap:fbc25_circle_complex_units_cohomology}, this
    algebraic group-cohomology statement for a finite discrete group replaces
    the comparison of sheaf cohomology with topological coefficient groups
    invoked in the footnote at
    \cite[line~1931]{FrancoRubio2025MPUGauging}.
\end{lemma}

\begin{proof}\leanok
    \uses{thm:mpug_positive_generalized_cocycle}
    The phase map $q(z)=z/|z|$ is multiplicative and satisfies
    $q(u)=u$ for $u\in U(1)$, so it preserves the cocycle equation and
    retracts inclusion.  The polar identity in
    (\ref{eq:mpug_circle_phase_residual}) identifies the residual $r$ with a
    positive cocycle.  Applying
    Lemma~\ref{thm:mpug_positive_generalized_cocycle} to the one-point block
    space writes $r=d\chi$.  Therefore $\omega\sim u$, which proves
    surjectivity of coefficient inclusion on degree-two classes.  Its left
    inverse induced by $q$ proves injectivity.
\end{proof}

\begin{definition}[Cocycle-twisted regular matrices]
    \label{def:mpug_twisted_regular_matrices}
    \lean{TNLean.Algebra.ScalarCocycle.phase,
        TNLean.Algebra.ScalarCocycle.twistedLeftRegular,
        TNLean.Algebra.ScalarCocycle.twistedRightRegular}
    \leanok
    \uses{def:scalar_cocycle, thm:mpug_circle_complex_units_h2}
    Let $G$ be a group and let $\omega$ be a $\C^\times$-valued scalar
    cochain.  Write
    \begin{align}
        \phi(g,h) & =\frac{\omega(g,h)}{|\omega(g,h)|}\in U(1).
    \end{align}
    On the standard basis of $\C^G$, define
    \begin{align}
        L_g^\phi
         & =\sum_{h\in G}\phi(g,h)^{-1}|gh\rangle\langle h|,
         &
        R_g^\phi
         & =\sum_{h\in G}\phi(h,g^{-1})^{-1}|hg^{-1}\rangle\langle h|.
        \label{eq:mpug_twisted_regular_matrices}
    \end{align}
    Whereas \cite[lines~412--418]{FrancoRubio2025MPUGauging} starts from a
    $U(1)$-valued cocycle $\Omega$, taking $\Omega=\phi$ reproduces its leg and
    inverse conventions.
\end{definition}

\begin{lemma}[Cocycle equation for the circle phase]
    \label{lem:mpug_circle_phase_cocycle}
    \lean{TNLean.Algebra.ScalarCocycle.phase_isCocycle}
    \leanok
    \uses{def:is_cocycle, thm:mpug_circle_complex_units_h2}
    If $\omega$ is a cocycle, then its circle phase satisfies
    \begin{align}
        \phi(g,h)\phi(gh,k)
         & =\phi(g,hk)\phi(h,k)
    \end{align}
    for all $g,h,k\in G$.
\end{lemma}

\begin{proof}\leanok
    \uses{thm:mpug_circle_complex_units_h2}
    The pointwise phase map $z\mapsto z/|z|$ is multiplicative, so it preserves
    the cocycle equation.
\end{proof}

\begin{lemma}[Action of the twisted regular matrices]
    \label{thm:mpug_twisted_regular_action}
    \lean{TNLean.Algebra.ScalarCocycle.twistedLeftRegular_mulVec_apply,
        TNLean.Algebra.ScalarCocycle.twistedRightRegular_mulVec_apply}
    \leanok
    \uses{def:mpug_twisted_regular_matrices}
    If $G$ is finite, then for $v\in\C^G$ and $g,h\in G$,
    \begin{align}
        (L_g^\phi v)(gh)
         & =\phi(g,h)^{-1}v(h),
         &
        (R_g^\phi v)(hg^{-1})
         & =\phi(h,g^{-1})^{-1}v(h).
    \end{align}
    % Source anchor: arXiv:2502.20257, lines 412--418.
\end{lemma}

\begin{proof}\leanok
    \uses{def:mpug_twisted_regular_matrices}
    In each matrix-vector product, exactly one summand is nonzero: the column
    indexed by $h$.
\end{proof}

\begin{lemma}[Multiplication laws for twisted regular matrices]
    \label{thm:mpug_twisted_regular_multiplication}
    \lean{TNLean.Algebra.ScalarCocycle.twistedLeftRegular_mul,
        TNLean.Algebra.ScalarCocycle.twistedRightRegular_mul}
    \leanok
    \uses{def:mpug_twisted_regular_matrices,
        lem:mpug_circle_phase_cocycle}
    If $G$ is finite and $\omega$ is a cocycle, then
    \begin{align}
        L_g^\phi L_h^\phi
         & =\phi(g,h)^{-1}L_{gh}^\phi,           \\
        R_g^\phi R_h^\phi
         & =\phi(h^{-1},g^{-1})^{-1}R_{gh}^\phi.
        \label{eq:mpug_twisted_regular_multiplication}
    \end{align}
    % Source anchor: arXiv:2502.20257, lines 406--418; the laws follow from
    % the displayed twisted regular matrices and the cocycle equation.
\end{lemma}

\begin{proof}\leanok
    \uses{lem:mpug_circle_phase_cocycle}
    Expanding the matrix products entrywise reduces the two scalar identities
    to the cocycle equation at $(g,h,k)$ and $(k,h^{-1},g^{-1})$, respectively.
\end{proof}

\begin{lemma}[Commutation of twisted left and right actions]
    \label{thm:mpug_twisted_regular_commutation}
    \lean{TNLean.Algebra.ScalarCocycle.twistedLeftRegular_commute}
    \leanok
    \uses{def:mpug_twisted_regular_matrices}
    If $G$ is finite and $\omega$ is a cocycle, then
    \begin{align}
        L_g^\phi R_h^\phi
         & =R_h^\phi L_g^\phi.
        \label{eq:mpug_twisted_regular_commutation}
    \end{align}
    % Source anchor: arXiv:2502.20257, lines 406--418 and 457; the latter
    % explicitly attributes commuting local symmetries to commutativity of the
    % left and right group actions.
\end{lemma}

\begin{proof}\leanok
    \uses{lem:mpug_circle_phase_cocycle}
    Expanding both matrix products entrywise and applying the cocycle equation
    at $(g,k,h^{-1})$ yields the equality.
\end{proof}

\begin{lemma}[Unitarity of twisted regular matrices]
    \label{thm:mpug_twisted_regular_unitary}
    \lean{TNLean.Algebra.ScalarCocycle.twistedLeftRegular_mem_unitaryGroup,
        TNLean.Algebra.ScalarCocycle.twistedRightRegular_mem_unitaryGroup}
    \leanok
    \uses{def:mpug_twisted_regular_matrices}
    If $G$ is finite, the matrices $L_g^\phi$ and $R_g^\phi$ are unitary for
    every $g\in G$.
    % Source anchor: arXiv:2502.20257, lines 406--418, where these matrices
    % form the two gauge legs of the local unitary representation.
\end{lemma}

\begin{proof}\leanok
    \uses{def:mpug_twisted_regular_matrices}
    Each matrix is a permutation matrix whose nonzero entries have modulus
    one.  Distinct columns therefore have disjoint support, while every column
    has squared norm one.
\end{proof}

\begin{lemma}[Onsite Gauss coefficient identity]
    \label{thm:mpug_onsite_gauss_phase}
    \lean{TNLean.Algebra.ScalarCocycle.onsiteGauss_phase_div}
    \leanok
    \uses{def:mpug_twisted_regular_matrices,
        lem:mpug_circle_phase_cocycle}
    Suppose that $\omega$ is a cocycle and impose the additional normalization
    \begin{align}
        \phi(t^{-1},t) & =1
    \end{align}
    for every $t\in G$.  Then
    \begin{align}
        \frac{\phi(ag^{-1},gb)}{\phi(a,b)}
         & =\frac{1}{\phi(a,g^{-1})\phi(g,b)}.
        \label{eq:mpug_onsite_gauss_phase}
    \end{align}
    The scalar convention $S_{t^{-1}}=S_t^{-1}$ appears in
    \cite[line~238]{FrancoRubio2025MPUGauging}, and the coefficient rewrite is
    at lines 3763--3766. The displayed normalization remains an explicit
    assumption here; it is not a consequence of cocyclicity.
\end{lemma}

\begin{proof}\leanok
    \uses{lem:mpug_circle_phase_cocycle}
    Applying the cocycle equation first at $(a,g^{-1},gb)$ and then at
    $(g^{-1},g,b)$, together with $\phi(e,b)=1$ from the normalization at
    $t=e$, yields the identity.
\end{proof}

\begin{definition}[Block scalar and block independence]
    \label{def:mpug_block_independence_factor}
    \lean{TNLean.Algebra.LSymbol.ell,
        TNLean.Algebra.LSymbol.IsBlockIndependentEll}
    \leanok
    \uses{def:mpug_l_symbols}
    For an $L$-symbol, define
    \begin{align}
        \ell^x_{a,b;r}
         & =\frac{L^x_{ar,r^{-1}b}}{L^x_{a,b}}.
        \label{eq:mpug_block_scalar}
    \end{align}
    The scalar $\ell$ is block-independent when
    \begin{align}
        \ell^x_{a,b;r}
         & =\ell^y_{a,b;r}
        \label{eq:mpug_block_independence}
    \end{align}
    for all $a,b,r\in G$ and $x,y\in\mathsf X$.
    % Source anchor: arXiv:2502.20257, app:block_indep, lines 6965--6974.
\end{definition}

\begin{definition}[Four scalar gauge formulations]
    \label{def:mpug_block_independence_gauge_forms}
    \lean{TNLean.Algebra.LSymbol.IsBlockConstant,
        TNLean.Algebra.LSymbol.IsTrivialEll,
        TNLean.Algebra.LSymbol.IsTrivial,
        TNLean.Algebra.LSymbol.HasTrivialGauge,
        TNLean.Algebra.LSymbol.HasBlockConstantGauge,
        TNLean.Algebra.LSymbol.HasTrivialEllGauge,
        TNLean.Algebra.LSymbol.IsBlockIndependent}
    \leanok
    \uses{def:mpug_l_gauge, def:mpug_block_independence_factor}
    Writing $L'=(\beta,\gamma)\mathbin{\triangleright}L$ for a joint scalar
    gauge of an $L$-symbol, the four scalar $L$-symbol properties used in the
    proof of the source proposition are, respectively, the existence of $\beta$
    and $\gamma$ such that
    \begin{align}
        (L')^x_{g,h}       & =1
                           &                     & \text{for all $x,g,h$},
        \label{eq:mpug_block_form_trivial}                                     \\
        (L')^x_{a,b}       & =(L')^y_{a,b}
                           &                     & \text{for all $x,y,a,b$},
        \label{eq:mpug_block_form_constant}                                    \\
        \ell[L']^x_{a,b;r} & =1
                           &                     & \text{for all $x,a,b,r$},
        \label{eq:mpug_block_form_trivial_ell}                                 \\
        \ell[L']^x_{a,b;r} & =\ell[L']^y_{a,b;r}
                           &                     & \text{for all $x,y,a,b,r$}.
        \label{eq:mpug_block_form_independent_ell}
    \end{align}
    % Source anchor: arXiv:2502.20257, prop:technical01, lines 6975--6989.
\end{definition}

\begin{lemma}[Corrected block factorization]
    \label{thm:mpug_block_factorization}
    \lean{TNLean.Algebra.LSymbol.IsCompatible.apply_right_one,
        TNLean.Algebra.LSymbol.blockFactor,
        TNLean.Algebra.LSymbol.factorization_of_isBlockIndependentEll}
    \leanok
    \uses{def:mpug_l_symbols, def:mpug_block_independence_factor}
    Suppose that $L$ is compatible with $\omega$ and that $\ell$ is
    block-independent. Fix $x_0\in\mathsf X$ and set
    \begin{align}
        A(g,h)
         & =\bigl(\omega(gh,e,e)\ell^{x_0}_{g,h;h}\bigr)^{-1}.
        \label{eq:mpug_block_factor}
    \end{align}
    Then, with the reciprocal factor recorded in
    \cite{gap:fbc25_block_independence_reciprocal_errors}, for every
    $x\in\mathsf X$ and $g,h\in G$,
    \begin{align}
        L^x_{g,h}
         & =A(g,h)L^x_{e,e}.
        \label{eq:mpug_corrected_block_factorization}
    \end{align}
    % Local correction:
    % docs/paper-gaps/fbc25_block_independence_reciprocal_errors.tex.
    % Source anchor: arXiv:2502.20257, eq:Lfact, lines 6994--7002.
\end{lemma}

\begin{proof}\leanok
    \uses{def:mpug_block_independence_factor}
    Definition~\ref{def:mpug_block_independence_factor} gives
    \begin{align}
        L^x_{g,h}
         & =\frac{L^x_{gh,e}}{\ell^x_{g,h;h}}.
        \label{eq:mpug_block_factorization_first}
    \end{align}
    Block independence replaces $x$ by $x_0$ in the denominator.
    Compatibility at $(gh,e,e)$ gives
    $L^x_{gh,e}=L^x_{e,e}/\omega(gh,e,e)$, which proves
    (\ref{eq:mpug_corrected_block_factorization}).
\end{proof}

\begin{lemma}[Relative block character]
    \label{thm:mpug_relative_block_character}
    \lean{TNLean.Algebra.LSymbol.IsCompatible.apply_left_one,
        TNLean.Algebra.LSymbol.relativeCharacter,
        TNLean.Algebra.LSymbol.relativeCharacter_mul_base,
        TNLean.Algebra.LSymbol.relativeCharacter_mul}
    \leanok
    \uses{thm:mpug_block_factorization}
    Under the hypotheses of Lemma~\ref{thm:mpug_block_factorization}, define
    \begin{align}
        \rho(g)
         & =\bigl(\ell^{x_0}_{e,g;g}\omega(g,e,e)
        \omega(e,e,g)\bigr)^{-1}.
        \label{eq:mpug_relative_block_character}
    \end{align}
    As recorded in \cite{gap:fbc25_block_independence_reciprocal_errors}, the
    reciprocal factor $(\ell^{x_0}_{e,g;g})^{-1}$ corrects the corresponding
    printed factor. Then
    \begin{align}
        L^{g x}_{e,e}
         & =\rho(g)L^x_{e,e},
        \label{eq:mpug_relative_block_transport} \\
        \rho(gh)
         & =\rho(g)\rho(h).
        \label{eq:mpug_relative_block_multiplicative}
    \end{align}
    Thus $\rho\colon G\to\C^\times$ is a one-dimensional character. No
    transitivity or finiteness assumption on the action is required.
    % Local correction:
    % docs/paper-gaps/fbc25_block_independence_reciprocal_errors.tex.
    % Source anchor: arXiv:2502.20257, eq:defrho, lines 7003--7018.
\end{lemma}

\begin{proof}\leanok
    \uses{thm:mpug_block_factorization}
    Compatibility at $(e,e,g)$ gives
    $L^x_{e,g}=\omega(e,e,g)L^{g x}_{e,e}$, and comparing this with
    (\ref{eq:mpug_corrected_block_factorization}) at $(e,g)$ yields
    (\ref{eq:mpug_relative_block_transport}). Applying that identity first
    to $h$ and then to $g$, and comparing with its application to $gh$, proves
    (\ref{eq:mpug_relative_block_multiplicative}) by cancellation.
\end{proof}

\begin{lemma}[Explicit inverse gauge]
    \label{thm:mpug_block_independence_inverse_gauge}
    \lean{TNLean.Algebra.LSymbol.factorFusionGauge,
        TNLean.Algebra.LSymbol.factorActionGauge,
        TNLean.Algebra.LSymbol.gauge_one_eq_of_isBlockIndependentEll,
        TNLean.Algebra.LSymbol.gauge_inv_eq_one_of_isBlockIndependentEll,
        TNLean.Algebra.LSymbol.hasTrivialGauge_of_isBlockIndependentEll}
    \leanok
    \uses{def:mpug_l_gauge, thm:mpug_l_gauge_composition,
        thm:mpug_block_factorization, thm:mpug_relative_block_character}
    With $A$ and $\rho$ as above, define the forward gauges by
    \begin{align}
        \beta^+_{g,h}
         & =\frac{A(g,h)}{\rho(h)},
        \label{eq:mpug_forward_fusion_gauge} \\
        \gamma^{+,x}_g
         & =\bigl(L^{g x}_{e,e}\bigr)^{-1}.
        \label{eq:mpug_forward_action_gauge}
    \end{align}
    Their joint gauge multiplier equals $L$:
    \begin{align}
        (\beta^+,\gamma^+)\mathbin{\triangleright}1
         & =L.
        \label{eq:mpug_forward_gauge_multiplier}
    \end{align}
    Consequently the pointwise inverse gauges trivialize $L$:
    \begin{align}
        ((\beta^+)^{-1},(\gamma^+)^{-1})
        \mathbin{\triangleright}L
         & =1.
        \label{eq:mpug_inverse_gauge_trivialization}
    \end{align}
    As recorded in \cite{gap:fbc25_block_independence_reciprocal_errors}, the
    distinction between the forward multiplier and its pointwise inverse
    corrects the orientation of the displayed source calculation.
    If $\mathsf X$ is nonempty, a choice of $x_0$ therefore gives a
    trivializing joint scalar gauge for every compatible block-independent
    $L$-symbol.
    % Local correction:
    % docs/paper-gaps/fbc25_block_independence_reciprocal_errors.tex.
    % Source anchor: arXiv:2502.20257, eq:Lgauge and lines 7019--7027.
\end{lemma}

\begin{proof}\leanok
    \uses{thm:mpug_l_gauge_composition, thm:mpug_block_factorization,
        thm:mpug_relative_block_character}
    Substituting (\ref{eq:mpug_forward_fusion_gauge}) and
    (\ref{eq:mpug_forward_action_gauge}) into
    (\ref{eq:mpug_l_gauge}), equation
    (\ref{eq:mpug_relative_block_transport}) reduces the multiplier to
    $A(g,h)L^x_{e,e}$, which equals $L^x_{g,h}$ by
    (\ref{eq:mpug_corrected_block_factorization}). Joint gauges compose by
    pointwise multiplication, so the inverse gauges give
    (\ref{eq:mpug_inverse_gauge_trivialization}).
\end{proof}

\begin{lemma}[Trivial and block-constant gauges]
    \label{thm:mpug_trivial_iff_block_constant}
    \lean{TNLean.Algebra.LSymbol.IsTrivial.isBlockConstant,
        TNLean.Algebra.LSymbol.IsBlockConstant.isBlockIndependentEll,
        TNLean.Algebra.LSymbol.HasTrivialGauge.hasBlockConstantGauge,
        TNLean.Algebra.LSymbol.HasBlockConstantGauge.isBlockIndependent,
        TNLean.Algebra.LSymbol.hasTrivialGauge_iff_hasBlockConstantGauge}
    \leanok
    \uses{def:mpug_block_independence_gauge_forms,
        thm:mpug_l_gauge_composition, thm:mpug_l_gauge_preservation,
        thm:mpug_block_independence_inverse_gauge}
    Let $L$ be an $L$-symbol compatible with a scalar $3$-cochain,
    and suppose that $\mathsf X$ is nonempty. A joint scalar gauge trivializes
    $L$ if and only if a joint scalar gauge makes $L^x_{g,h}$ independent of
    $x$.
    % Source anchor: arXiv:2502.20257, prop:technical01, items 1 and 2.
\end{lemma}

\begin{proof}\leanok
    \uses{thm:mpug_l_gauge_composition, thm:mpug_l_gauge_preservation,
        thm:mpug_block_independence_inverse_gauge}
    A trivial $L$-symbol is block-constant. Conversely, a block-constant
    representative has block-independent $\ell$, so applying
    Lemma~\ref{thm:mpug_block_independence_inverse_gauge} to that
    representative and composing the two joint gauges trivializes $L$.
\end{proof}

\begin{lemma}[Trivial gauges and trivial block scalars]
    \label{thm:mpug_trivial_iff_trivial_ell}
    \lean{TNLean.Algebra.LSymbol.IsTrivial.isTrivialEll,
        TNLean.Algebra.LSymbol.IsTrivialEll.isBlockIndependentEll,
        TNLean.Algebra.LSymbol.HasTrivialGauge.hasTrivialEllGauge,
        TNLean.Algebra.LSymbol.HasTrivialEllGauge.isBlockIndependent,
        TNLean.Algebra.LSymbol.hasTrivialGauge_iff_hasTrivialEllGauge}
    \leanok
    \uses{def:mpug_block_independence_gauge_forms,
        thm:mpug_l_gauge_composition, thm:mpug_l_gauge_preservation,
        thm:mpug_block_independence_inverse_gauge}
    Let $L$ be an $L$-symbol compatible with a scalar $3$-cochain,
    and suppose that $\mathsf X$ is nonempty. A joint scalar gauge trivializes
    $L$ if and only if a joint scalar gauge makes every
    $\ell^x_{a,b;r}$ equal to $1$.
    % Source anchor: arXiv:2502.20257, prop:technical01, items 1 and 3.
\end{lemma}

\begin{proof}\leanok
    \uses{thm:mpug_l_gauge_composition, thm:mpug_l_gauge_preservation,
        thm:mpug_block_independence_inverse_gauge}
    A trivial $L$-symbol has $\ell=1$. Conversely, the identity $\ell=1$
    implies block independence, so applying
    Lemma~\ref{thm:mpug_block_independence_inverse_gauge} and composing the
    gauges trivializes $L$.
\end{proof}

\begin{lemma}[Scalar gauge triviality and block independence]
    \label{thm:mpug_trivial_iff_block_independent}
    \lean{TNLean.Algebra.LSymbol.IsBlockIndependent.hasTrivialGauge,
        TNLean.Algebra.LSymbol.hasTrivialGauge_iff_isBlockIndependent}
    \leanok
    \uses{def:mpug_block_independence_gauge_forms,
        thm:mpug_l_gauge_composition, thm:mpug_l_gauge_preservation,
        thm:mpug_block_independence_inverse_gauge}
    Let $L$ be an $L$-symbol compatible with a scalar $3$-cochain,
    and suppose that $\mathsf X$ is nonempty. A joint scalar gauge trivializes
    $L$ if and only if a joint scalar gauge makes $\ell^x_{a,b;r}$ independent
    of $x$.
    % Source anchor: arXiv:2502.20257, prop:technical01, items 1 and 4.
\end{lemma}

\begin{proof}\leanok
    \uses{thm:mpug_l_gauge_composition, thm:mpug_l_gauge_preservation,
        thm:mpug_block_independence_inverse_gauge}
    Gauge triviality implies block constancy and hence block independence.
    For the converse, apply
    Lemma~\ref{thm:mpug_block_independence_inverse_gauge} to the
    block-independent representative and compose its trivializing gauge with
    the original gauge.
\end{proof}

\begin{lemma}[Four scalar gauge formulations of block independence]
    \label{thm:mpug_block_independence_four_forms}
    \lean{TNLean.Algebra.LSymbol.hasTrivialGauge_tfae}
    \leanok
    \uses{def:mpug_block_independence_gauge_forms,
        thm:mpug_trivial_iff_block_constant,
        thm:mpug_trivial_iff_trivial_ell,
        thm:mpug_trivial_iff_block_independent}
    Let $L$ be an $L$-symbol compatible with a scalar $3$-cochain,
    and suppose that $\mathsf X$ is nonempty. The four scalar gauge
    properties (\ref{eq:mpug_block_form_trivial}),
    (\ref{eq:mpug_block_form_constant}),
    (\ref{eq:mpug_block_form_trivial_ell}), and
    (\ref{eq:mpug_block_form_independent_ell}) of $L$ are equivalent. This is
    the scalar core of the block-independence proposition of
    \cite{FrancoRubio2025MPUGauging}, whose fourth condition is stated for an
    MPS--MPU pair.
    % Source anchor: arXiv:2502.20257, prop:technical01, lines 6975--6989.
\end{lemma}

\begin{proof}\leanok
    \uses{thm:mpug_trivial_iff_block_constant,
        thm:mpug_trivial_iff_trivial_ell,
        thm:mpug_trivial_iff_block_independent}
    Lemmas~\ref{thm:mpug_trivial_iff_block_constant},
    \ref{thm:mpug_trivial_iff_trivial_ell}, and
    \ref{thm:mpug_trivial_iff_block_independent} identify each of the last
    three properties with the first.
\end{proof}

For a compatible scalar $L$-symbol,
Lemma~\ref{thm:mpug_block_independence_four_forms} establishes the
algebraic equivalence among its four scalar gauge properties; the corresponding
tensor-level bridge to an underlying MPS--MPU pair is not asserted here.
% Scope boundary: the tensor-facing bridge for arXiv:2502.20257,
% prop:technical01 is not contained in the scalar L-symbol declarations.

\begin{lemma}[Stabilizer transition elements]
    \label{thm:mpug_stabilizer_transition}
    \lean{TNLean.Algebra.StabilizerRepresentatives.transitionElement_mul,
        TNLean.Algebra.StabilizerRepresentatives,
        TNLean.Algebra.StabilizerRepresentatives.ofPretransitive,
        TNLean.Algebra.StabilizerRepresentatives.transitionElement,
        TNLean.Algebra.StabilizerRepresentatives.coe_transitionElement,
        TNLean.Algebra.StabilizerRepresentatives.transitionElement_stabilizer}
    \leanok
    Suppose that $G$ acts transitively on $\mathsf X$, and fix
    $x_0\in\mathsf X$. There are representatives $k_x\in G$ satisfying
    $k_xx_0=x$ and $k_{x_0}=e$. With
    $H=\operatorname{Stab}_G(x_0)$, define
    \begin{align}
        t_k(g,x)
         & =k_{g x}^{-1}gk_x.
        \label{eq:mpug_stabilizer_transition}
    \end{align}
    Then $t_k(g,x)\in H$ and
    \begin{align}
        t_k(g_1g_2,x)
         & =t_k(g_1,g_2x)t_k(g_2,x).
        \label{eq:mpug_stabilizer_transition_product}
    \end{align}
    Consequently, the source factors are
    \begin{align}
        h_2
         & =t_k(g_2,x)=k_{g_2x}^{-1}g_2k_x,
        \label{eq:mpug_stabilizer_h2}                  \\
        h_1
         & =t_k(g_1,g_2x)=k_{g_1g_2x}^{-1}g_1k_{g_2x}.
        \label{eq:mpug_stabilizer_h1}
    \end{align}
    As recorded in \cite{gap:fbc25_stabilizer_transition_index_typo}, the last
    representative in $h_1$ corrects the printed $k_{g_1x}$. No finiteness or
    normality assumption is required.
    % Local correction:
    % docs/paper-gaps/fbc25_stabilizer_transition_index_typo.tex.
    % Source anchor: arXiv:2502.20257, lemma:h1h2 and eq:h1h2,
    % lines 7028--7042.
\end{lemma}

\begin{proof}\leanok
    Transitivity gives one element carrying $x_0$ to each $x$, with the
    identity chosen at $x=x_0$. Acting with (\ref{eq:mpug_stabilizer_transition})
    on $x_0$ gives
    $k_{gx}^{-1}g(k_xx_0)=k_{gx}^{-1}(gx)=x_0$.
    For (\ref{eq:mpug_stabilizer_transition_product}), substituting
    (\ref{eq:mpug_stabilizer_transition}) and cancelling
    $k_{g_2x}k_{g_2x}^{-1}$ yields the product formula. The two formulas for
    $h_1$ and $h_2$ are the corresponding specializations.
\end{proof}

\begin{lemma}[L-symbol induced by a stabilizer cocycle]
    \label{thm:mpug_stabilizer_cocycle_l_symbol}
    \lean{TNLean.Algebra.StabilizerRepresentatives.inducedLSymbol_isCompatible,
        TNLean.Algebra.StabilizerRepresentatives.inducedLSymbol,
        TNLean.Algebra.StabilizerRepresentatives.inducedLSymbol_base,
        TNLean.Algebra.StabilizerRepresentatives.inducedActionGauge,
        TNLean.Algebra.StabilizerRepresentatives.inducedLSymbol_eq_gauge,
        TNLean.Algebra.StabilizerRepresentatives.isMulCoboundary₂_div_iff_exists_actionGauge_inducedLSymbol_eq}
    \leanok
    \uses{def:mpug_l_symbols, def:mpug_l_gauge, def:is_cocycle,
        def:cohomologous_to, thm:mpug_stabilizer_transition}
    Let $H=\operatorname{Stab}_G(x_0)$ and let
    $\psi\colon H\times H\to\C^\times$ be a scalar $2$-cocycle. Define the
    $\omega=1$ scalar specialization of \cite[Eq.~(20)]{GarreRubio2023MPOSym}
    by
    \begin{align}
        L^x_{g,h}
         & =\psi\bigl(t_k(g,hx),t_k(h,x)\bigr).
        \label{eq:mpug_induced_stabilizer_l_symbol}
    \end{align}
    The $L$-symbol (\ref{eq:mpug_induced_stabilizer_l_symbol}) satisfies
    (\ref{eq:mpug_l_compatibility}) with $\omega=1$, and its base-point
    restriction recovers $\psi$:
    \begin{align}
        L^{x_0}_{a,b} & =\psi(a,b)
        \label{eq:mpug_induced_stabilizer_base}
    \end{align}
    for all $a,b\in H$.

    More precisely, two stabilizer cocycles are cohomologous exactly when
    their induced $L$-symbols differ by an action-tensor gauge and the trivial
    fusion-tensor gauge:
    \begin{align}
        \psi_1\sim\psi_2
         & \quad\Longleftrightarrow\quad
        \exists\gamma,\quad
        L[\psi_1]=(1,\gamma)\mathbin{\triangleright}L[\psi_2].
        \label{eq:mpug_induced_cohomology_iff}
    \end{align}
    If the cohomology relation is given by
    $\varphi\colon H\to\C^\times$, with
    \begin{align}
        \psi_1(a,b)
         & =\varphi(a)\varphi(b)\varphi(ab)^{-1}\psi_2(a,b),
        \label{eq:mpug_stabilizer_cocycle_gauge}
    \end{align}
    then the forward action-tensor gauge is
    \begin{align}
        \gamma_{g,x} & =\varphi\bigl(t_k(g,x)\bigr)^{-1}.
        \label{eq:mpug_induced_action_gauge}
    \end{align}
    Thus the induced action-gauge class depends only on the cohomology class
    of $\psi$.

    This establishes the scalar part of the $L$-symbol reduction in
    \cite[lines~7028--7042]{FrancoRubio2025MPUGauging}; the
    representation-theoretic claim for a tensor-derived $L$-symbol is not
    asserted here.
    % Source anchor: arXiv:2203.12563, Eq. (20).
\end{lemma}

\begin{proof}\leanok
    \uses{def:mpug_l_gauge, thm:mpug_stabilizer_transition}
    With
    $a=t_k(g,hkx)$, $b=t_k(h,kx)$, and $c=t_k(k,x)$,
    the product law (\ref{eq:mpug_stabilizer_transition_product}) gives
    $t_k(gh,kx)=ab$ and $t_k(hk,x)=bc$. Compatibility is therefore
    \begin{align}
        \psi(a,b)\psi(ab,c)
         & =\psi(a,bc)\psi(b,c),
    \end{align}
    which is the cocycle equation. If $p,q\in H$, then
    $t_k(p,x_0)=p$ and $t_k(q,x_0)=q$ because $k_{x_0}=e$, proving
    (\ref{eq:mpug_induced_stabilizer_base}). For the forward gauge
    calculation, setting $a=t_k(g,hx)$ and $b=t_k(h,x)$ gives
    $t_k(gh,x)=ab$. Substituting (\ref{eq:mpug_induced_action_gauge}) into
    (\ref{eq:mpug_l_gauge}) gives the multiplier
    \begin{align}
        \frac{\gamma_{gh,x}}{\gamma_{g,hx}\gamma_{h,x}}
         & =\varphi(ab)^{-1}
        \bigl(\varphi(a)^{-1}\varphi(b)^{-1}\bigr)^{-1}
        =\varphi(a)\varphi(b)\varphi(ab)^{-1}.
    \end{align}
    Formula (\ref{eq:mpug_stabilizer_cocycle_gauge}) now identifies the two
    induced $L$-symbols. Conversely, if
    $L[\psi_1]=(1,\gamma)\mathbin{\triangleright}L[\psi_2]$ and
    $\varphi(a)=\gamma_{a,x_0}^{-1}$ for $a\in H$, evaluating the gauge
    equality at $x_0$ and $a,b\in H$ gives
    \begin{align}
        \psi_1(a,b)
         & =\frac{\gamma_{ab,x_0}}
        {\gamma_{a,x_0}\gamma_{b,x_0}}\psi_2(a,b)
        =\varphi(a)\varphi(b)\varphi(ab)^{-1}\psi_2(a,b).
    \end{align}
    Hence $\psi_1\sim\psi_2$, which proves the reverse implication in
    (\ref{eq:mpug_induced_cohomology_iff}).
\end{proof}

\begin{lemma}[Explicit gauge from a stabilizer cocycle extension]
    \label{thm:mpug_stabilizer_extension_gauge}
    \lean{TNLean.Algebra.StabilizerRepresentatives.extensionActionGauge,
        TNLean.Algebra.StabilizerRepresentatives.inducedLSymbol_eq_extensionActionGauge,
        TNLean.Algebra.StabilizerRepresentatives.inducedLSymbol_hasTrivialGauge_of_extension}
    \leanok
    \uses{thm:mpug_stabilizer_cocycle_l_symbol, def:mpug_l_gauge,
        def:mpug_block_independence_gauge_forms}
    Let $G$ act on $\mathsf X$, with basepoint $x_0$, stabilizer
    $H=\operatorname{Stab}_G(x_0)$, and representatives $k_x$ satisfying
    $k_xx_0=x$ and $k_{x_0}=e$.
    Let $\psi\colon H\times H\to\C^\times$ and let
    $\Psi\colon G\times G\to\C^\times$ be a two-cocycle whose restriction
    to $H\times H$ is $\psi$. Define
    \begin{align}
        \gamma_g^x
         & =\frac{\Psi(k_{gx}^{-1},gk_x)\Psi(g,k_x)}
        {\Psi(e,e)\Psi(k_x,k_x^{-1})}.
        \label{eq:mpug_extension_gamma}
    \end{align}
    With $h_1=t_k(g_1,g_2x)$ and $h_2=t_k(g_2,x)$, the induced symbols satisfy
    \begin{align}
        L^x_{g_1,g_2}=\psi(h_1,h_2)
         & =\Psi(g_1,g_2)\frac{\gamma_{g_1g_2}^x}
        {\gamma_{g_1}^{g_2x}\gamma_{g_2}^x}.
        \label{eq:mpug_extension_identity}
    \end{align}
    Consequently the inverse fusion and action gauges
    $(\Psi^{-1},\gamma^{-1})$ make $L$ identically one.
    This is the scalar representative calculation in the forward proof of
    the block-independence extension proposition of
    \cite{FrancoRubio2025MPUGauging};
    here extension means the equality $\Psi|_{H\times H}=\psi$ of chosen
    cocycle representatives, rather than only equality of their classes.
    See \cite{gap:fbc25_mpu_action_stabilizer_source_audit} for the
    relationship between cocycle representatives, cohomology classes, and
    tensor data.
\end{lemma}

\begin{proof}\leanok
    With $a=k_{g_1g_2x}$, $b=k_{g_2x}$, $c=k_x$, $u=g_1$, and $v=g_2$,
    the cocycle equation gives $\Psi(z,e)=\Psi(e,z)=\Psi(e,e)$ and
    $\Psi(b^{-1},b)=\Psi(b,b^{-1})$. Applying it successively to the triples
    \begin{align}
         & (a^{-1}ub,b^{-1},vc),\quad (a^{-1}ub,b^{-1},b),\quad
        (a^{-1},u,vc),\quad (a^{-1},u,b),\quad (u,v,c)
    \end{align}
    and solving these equalities for the corresponding factors of
    $\Psi(a^{-1}ub,b^{-1}vc)$ yields
    (\ref{eq:mpug_extension_identity}) with
    (\ref{eq:mpug_extension_gamma}) after cancellation. The restriction
    hypothesis identifies its left side with $\psi(h_1,h_2)$. Substitution into
    (\ref{eq:mpug_l_gauge}) cancels the multiplier under the inverse gauges.
\end{proof}

\begin{lemma}[Reconstruction of stabilizer L-symbols]
    \label{thm:mpug_stabilizer_reconstruction}
    \lean{TNLean.Algebra.StabilizerRepresentatives.reconstructedLSymbol,
        TNLean.Algebra.StabilizerRepresentatives.reconstructedLSymbol_isCompatible,
        TNLean.Algebra.StabilizerRepresentatives.reconstructedLSymbol_one_isCompatible,
        TNLean.Algebra.StabilizerRepresentatives.reconstructedLSymbol_eq_mul,
        TNLean.Algebra.StabilizerRepresentatives.reconstructedLSymbol_base,
        TNLean.Algebra.StabilizerRepresentatives.gauge_reconstructedLSymbol_isCompatible,
        TNLean.Algebra.ScalarThreeCochain.exists_fusionGauge_eq_one_of_isTrivialGaugeClass_comap,
        TNLean.Algebra.LSymbol.IsCompatible.mul}
    \leanok
    \uses{def:mpug_l_symbols, def:mpug_l_gauge, def:mpug_three_cocycle,
        def:mpug_fusion_gauge, def:mpug_trivial_gauge_class, def:is_cocycle,
        thm:mpug_stabilizer_transition}
    Let $G$ act on $\mathsf X$ with basepoint $x_0$, stabilizer
    $H=\operatorname{Stab}_G(x_0)$, and representatives $k_x$ as in
    Lemma~\ref{thm:mpug_stabilizer_transition}. Let $\omega$ be a scalar
    $3$-cocycle on $G$ with $\omega(a,b,c)=1$ for all $a,b,c\in H$, and let
    $\psi$ be a scalar $2$-cocycle on $H$. With $h_2=t_k(g_2,x)$ and
    $h_1=t_k(g_1,g_2x)$, define
    \begin{align}
        L^x_{g_1,g_2}
         & =\frac{\omega(g_1g_2k_x,h_2^{-1},h_1^{-1})\,\psi(h_1,h_2)}
        {\omega(g_1,g_2k_x,h_2^{-1})\,\omega(g_1,g_2,k_x)}.
        \label{eq:mpug_stabilizer_reconstruction}
    \end{align}
    Then the $L$-symbol (\ref{eq:mpug_stabilizer_reconstruction}) of
    \cite[Eq.~(20)]{GarreRubio2023MPOSym} satisfies
    (\ref{eq:mpug_l_compatibility}), with
    $L^{x_0}_{a,b}=\psi(a,b)$ for $a,b\in H$.

    If only $\omega|_H$ is a coboundary, some fusion gauge $\beta$ on $G$
    makes $\beta\mathbin{\triangleright}\omega$ equal to one on
    $H\times H\times H$, and applying the inverse gauge
    $(\beta^{-1},1)$ to (\ref{eq:mpug_stabilizer_reconstruction}) built from
    $\beta\mathbin{\triangleright}\omega$ gives an $L$-symbol compatible with
    $\omega$. The hypothesis that $\omega$ is one on the stabilizer fixes the
    fusion gauge in which the source's formula holds; as recorded in
    \cite{gap:glm23_eq20_fusion_gauge}, for a general
    $\omega$ with $\omega|_H=d\beta$ the printed formula can fail.
    No normalization of $\omega$ is required.
    % Local fix: docs/paper-gaps/glm23_eq20_fusion_gauge.tex.
    % Source anchor: arXiv:2203.12563, REsubmission.tex lines 740--761.
\end{lemma}

\begin{proof}\leanok
    \uses{thm:mpug_stabilizer_transition, thm:mpug_stabilizer_cocycle_l_symbol}
    The formula factors as its value at $\psi=1$ times the induced
    $L$-symbol (\ref{eq:mpug_induced_stabilizer_l_symbol}), and pointwise
    products of compatible $L$-symbols are compatible with the product of
    the $3$-cochains. By Lemma~\ref{thm:mpug_stabilizer_cocycle_l_symbol}
    it suffices to treat $\psi=1$. For fixed $x$ and $g,h,k$, with $p=k_x$,
    $a=t_k(k,x)$, $b=t_k(h,kx)$, and $c=t_k(g,hkx)$, we have $k_{kx}=kpa^{-1}$
    and, by (\ref{eq:mpug_stabilizer_transition_product}),
    $t_k(hk,x)=ba$ and $t_k(gh,kx)=cb$. The compatibility equation then
    follows from the $3$-cocycle equation at the four quadruples
    \begin{align}
         & (g,h,k,p),\quad (g,h,kp,a^{-1}),\quad (g,hkp,a^{-1},b^{-1}),\quad
        (ghkp,a^{-1},b^{-1},c^{-1}),
    \end{align}
    together with $\omega(a^{-1},b^{-1},c^{-1})=1$. At $x_0$ and $a,b\in H$
    every argument of $\omega$ in (\ref{eq:mpug_stabilizer_reconstruction})
    lies in $H$, which gives $L^{x_0}_{a,b}=\psi(a,b)$. If
    $\omega|_H=d\beta_0$, extending $\beta_0^{-1}$ by one outside
    $H\times H$ produces a coboundary that cancels $\omega$ on
    $H\times H\times H$. Finally, $(\beta^{-1},1)$ changes
    $\beta\mathbin{\triangleright}\omega$ back to $\omega$.
\end{proof}

\begin{lemma}[Classification and torsor of stabilizer L-symbols]
    \label{thm:mpug_stabilizer_torsor}
    \lean{TNLean.Algebra.StabilizerRepresentatives.exists_actionGauge_eq_inducedLSymbol,
        TNLean.Algebra.StabilizerRepresentatives.exists_actionGauge_eq_reconstructedLSymbol,
        TNLean.Algebra.StabilizerRepresentatives.exists_actionGauge_gauge_eq_reconstructedLSymbol,
        TNLean.Algebra.StabilizerRepresentatives.isCompatible_iff_exists_actionGauge_mul_inducedLSymbol,
        TNLean.Algebra.StabilizerRepresentatives.actionGaugeEquiv_mul_inducedLSymbol_iff,
        TNLean.Algebra.StabilizerRepresentatives.solutionSetoid,
        TNLean.Algebra.StabilizerRepresentatives.h2EquivSolutionClasses,
        TNLean.Algebra.StabilizerRepresentatives.h2EquivSolutionClasses_H2π,
        TNLean.Algebra.StabilizerRepresentatives.h2EquivSolutionClasses_apply,
        TNLean.Algebra.StabilizerRepresentatives.solutionAddAction,
        TNLean.Algebra.StabilizerRepresentatives.existsUnique_vadd_eq,
        TNLean.Algebra.LSymbol.ActionGaugeEquiv,
        TNLean.Algebra.LSymbol.actionGaugeEquiv_equivalence,
        TNLean.Algebra.LSymbol.IsCompatible.inv,
        TNLean.Algebra.LSymbol.IsCompatible.mul_inv_one,
        TNLean.Algebra.LSymbol.IsCompatible.gauge_one,
        TNLean.Algebra.LSymbol.IsCompatible.isCocycle_restrict,
        TNLean.Algebra.LSymbol.gauge_one_mul}
    \leanok
    \uses{thm:mpug_stabilizer_reconstruction, def:mpug_l_gauge,
        thm:mpug_h2_group_cohomology_comparison, def:cohomologous_to}
    Keep the setting of Lemma~\ref{thm:mpug_stabilizer_reconstruction},
    and call two $L$-symbols action-gauge equivalent when they differ by
    $(1,\gamma)\mathbin{\triangleright}$ for some action-tensor gauge
    $\gamma$.
    \begin{enumerate}
        \item If $R$ is compatible with the trivial $3$-cochain, then
              $R=(1,\gamma)\mathbin{\triangleright}L[\psi]$ for the induced
              $L$-symbol (\ref{eq:mpug_induced_stabilizer_l_symbol}) of the
              $2$-cocycle $\psi(a,b)=R^{x_0}_{a,b}$, with
              $\gamma^x_g=R^{x_0}_{k_{gx},t_k(g,x)}/R^{x_0}_{g,k_x}$.
        \item If $\omega$ is a $3$-cocycle equal to one on
              $H\times H\times H$, every $L$-symbol compatible with $\omega$
              is action-gauge equivalent to
              (\ref{eq:mpug_stabilizer_reconstruction}) for some $2$-cocycle
              $\psi$ on $H$. For a general fusion gauge $\beta$ making
              $\beta\mathbin{\triangleright}\omega$ trivial on $H$, the
              same holds for $(\beta,1)\mathbin{\triangleright}L$ and the
              formula built from $\beta\mathbin{\triangleright}\omega$.
        \item Let $L_0$ be compatible with an arbitrary scalar $3$-cochain
              $\omega$. An $L$-symbol $L$ is compatible with $\omega$ if and
              only if $L$ is action-gauge equivalent to $L_0L[\psi]$ for
              some $2$-cocycle $\psi$ on $H$, and $L_0L[\psi_1]$ and
              $L_0L[\psi_2]$ are action-gauge equivalent if and only if
              $\psi_1\sim\psi_2$. Hence sending the class of $\psi$ to the
              class of $L_0L[\psi]$ is a bijection from $H^2(H,\C^\times)$
              onto the action-gauge classes of $L$-symbols compatible with
              $\omega$.
        \item Sending the class of $\psi$ and the class of $L$ to the
              class of $L\,L[\psi]$ is a well-defined action of the group
              cohomology $H^2(H,\C^\times)$, for the trivial action of $H$
              on $\C^\times$, on these classes, and for any two classes
              exactly one cohomology class carries the first to the second.
    \end{enumerate}
    The last two items are the torsor statement of \cite{GarreRubio2023MPOSym}
    for fixed $(G,\omega,H)$, with the fusion gauge held fixed as in the
    source. In all four items $G$, and hence $H$, is a small group: its
    underlying set lies in a fixed set-theoretic universe containing $\Z$,
    in which the group cohomology $H^2(H,\C^\times)$ is formed. In the last
    two items the set $\mathsf X$ lies in the same universe.
    % Source anchor: arXiv:2203.12563, REsubmission.tex lines 749--761.
\end{lemma}

\begin{proof}\leanok
    \uses{thm:mpug_stabilizer_reconstruction,
        thm:mpug_stabilizer_cocycle_l_symbol}
    For the first item, compatibility at $x_0$ with arguments
    $(g,h,k_x)$ expresses $R^x_{g,h}$ through basepoint values. Rewriting
    $hk_x=k_{hx}t_k(h,x)$ and $gk_{hx}=k_{ghx}t_k(g,hx)$ and using
    compatibility at $x_0$ with arguments $(g,k_{hx},t_k(h,x))$ and
    $(k_{ghx},t_k(g,hx),t_k(h,x))$ combines the three equations into
    $R^x_{g,h}=\gamma^x_{gh}(\gamma^{hx}_g\gamma^x_h)^{-1}
        \psi(t_k(g,hx),t_k(h,x))$. Compatibility at $x_0$ with arguments in
    $H$ is the cocycle equation for $\psi$.

    For the third item, if $L$ and $L_0$ are compatible with $\omega$,
    then $LL_0^{-1}$ is compatible with the trivial $3$-cochain, and the first
    item applies; conversely $L_0L[\psi]$ is compatible with $\omega$, and an
    action gauge with trivial fusion part preserves compatibility and
    commutes with multiplication by $L_0$. Cancelling $L_0$ reduces the
    freeness statement to (\ref{eq:mpug_induced_cohomology_iff}). The second
    item is the third with $L_0$ the formula at $\psi=1$, which is
    compatible by Lemma~\ref{thm:mpug_stabilizer_reconstruction}.
    For the last item, cocycles modulo cohomology form the second group
    cohomology of $H$ with coefficients in $\C^\times$, and
    \begin{align}
        [\psi_1]+[\psi_2] & =[\psi_1\psi_2],
        \notag                                     \\
        L[\psi_1\psi_2]   & =L[\psi_1]\,L[\psi_2].
        \notag
    \end{align}
    Cohomologous cocycles give induced $L$-symbols related by an action
    gauge, so the rule is well defined, and by the two identities acting by
    $[\psi_1]+[\psi_2]$ is acting by $[\psi_2]$ and then by $[\psi_1]$. With
    base point $L$ the action is the bijection of the third item.
\end{proof}

\begin{definition}[Generalized coboundary and cocycle]
    \label{def:mpug_generalized_cocycle}
    \lean{TNLean.Algebra.GeneralizedThreeCochain,
        TNLean.Algebra.GeneralizedFourCochain,
        TNLean.Algebra.ActionTensorGauge.coboundary,
        TNLean.Algebra.LSymbol.coboundary,
        TNLean.Algebra.GeneralizedThreeCochain.coboundary,
        TNLean.Algebra.LSymbol.IsGeneralizedCocycle}
    \leanok
    \uses{def:mpug_l_symbols, def:mpug_l_gauge}
    Let $\chi\colon G\times\mathsf X\to\C^\times$ be block-dependent. Its
    generalized coboundary is
    \begin{align}
        (d\chi)^x_{g,h}
         & =\frac{\chi^{h x}_g\chi^x_h}{\chi^x_{gh}}.
        \label{eq:mpug_generalized_coboundary}
    \end{align}
    For a block-dependent scalar $2$-cochain
    $f\colon\mathsf X\times G^2\to\C^\times$, set
    \begin{align}
        (df)^x_{g,h,k}
         & =\frac{f^x_{g,hk}f^x_{h,k}}
        {f^{k x}_{g,h}f^x_{gh,k}}.
        \label{eq:mpug_generalized_two_coboundary}
    \end{align}
    As recorded in \cite{gap:fbc25_two_cochain_coboundary_index_typo}, the
    final group index in the denominator is $k$, correcting the printed
    formula. For a block-dependent scalar $3$-cochain
    $\omega\colon\mathsf X\times G^3\to\C^\times$, the corresponding next
    coboundary is
    \begin{align}
        (d\omega)^x_{g,h,k,l}
         & =\frac{\omega^{l x}_{g,h,k}\omega^x_{g,hk,l}
            \omega^x_{h,k,l}}
        {\omega^x_{gh,k,l}\omega^x_{g,h,kl}}.
        \label{eq:mpug_generalized_three_coboundary}
    \end{align}
    A function $\Omega\colon\mathsf X\times G^2\to\C^\times$ is a
    generalized $2$-cocycle when
    \begin{align}
        \Omega^x_{g,hk}\Omega^x_{h,k}
         & =\Omega^{k x}_{g,h}\Omega^x_{gh,k}.
        \label{eq:mpug_generalized_cocycle}
    \end{align}
    The differential identities $d\Omega=1$ and $d(df)=1$ follow the group
    cohomology conventions of
    \cite[lines~5914--5922]{FrancoRubio2025MPUGauging}, with the
    three-to-four-cochain coboundary given explicitly by
    (\ref{eq:mpug_generalized_three_coboundary}).
    No transitivity or finiteness assumption is imposed on $\mathsf X$.

    % Local fix: docs/paper-gaps/fbc25_two_cochain_coboundary_index_typo.tex.
    % Source anchor: arXiv:2502.20257, lines 5914--5922 and lemma:G_cocycle.
\end{definition}

\begin{lemma}[Generalized coboundaries square to one]
    \label{thm:mpug_generalized_coboundary_squared}
    \lean{TNLean.Algebra.ActionTensorGauge.isGeneralizedCocycle_coboundary,
        TNLean.Algebra.GeneralizedThreeCochain.coboundary_coboundary}
    \leanok
    \uses{def:mpug_generalized_cocycle}
    The generalized coboundary of every block-dependent scalar $1$-cochain
    is a generalized $2$-cocycle. With the index correction of
    \cite{gap:fbc25_two_cochain_coboundary_index_typo}, every block-dependent
    scalar $2$-cochain $f$ satisfies the identity of
    \cite[line~5922]{FrancoRubio2025MPUGauging},
    \begin{align}
        d(df) & =1.
        \label{eq:mpug_generalized_coboundary_squared}
    \end{align}
\end{lemma}

\begin{proof}\leanok
    \uses{def:mpug_generalized_cocycle}
    Substituting the coboundary formulas, the action law
    $k(lx)=(kl)x$ and associativity identify the shifted block labels and
    group products. Every value of the original cochain then occurs once in
    the numerator and once in the denominator, so all factors cancel.
\end{proof}

\begin{definition}[Positive block-dependent cochains]
    \label{def:mpug_positive_generalized_cochains}
    \lean{TNLean.Algebra.PositiveUnits.embedding,
        TNLean.Algebra.PositiveUnits.IsPositive,
        TNLean.Algebra.ActionTensorGauge.IsPositiveValued,
        TNLean.Algebra.LSymbol.IsPositiveValued}
    \leanok
    \uses{def:mpug_l_symbols, def:mpug_l_gauge}
    Let $\iota\colon\mathbb R_{>0}\to\C^\times$ be the natural
    multiplicative embedding. A block-dependent scalar $1$-cochain $\chi$ is
    positive when
    \begin{align}
        \chi^x_g & \in\iota(\mathbb R_{>0})
        \label{eq:mpug_positive_one_cochain}
    \end{align}
    for every $x\in\mathsf X$ and $g\in G$. A block-dependent scalar
    $2$-cochain $\Omega$ is positive when
    \begin{align}
        \Omega^x_{g,h} & \in\iota(\mathbb R_{>0})
        \label{eq:mpug_positive_two_cochain}
    \end{align}
    for every $x\in\mathsf X$ and $g,h\in G$. Thus positivity means
    membership in the embedded subgroup of strictly positive real numbers,
    not positivity of the real part.
\end{definition}

\begin{definition}[Hat operation on block-dependent cochains]
    \label{def:mpug_generalized_hat}
    \lean{TNLean.Algebra.ActionTensorGauge.hat,
        TNLean.Algebra.LSymbol.hat}
    \leanok
    \uses{def:mpug_l_symbols, def:mpug_l_gauge}
    Following the one-index and two-index cases of the hat operation in
    \cite[line~5912]{FrancoRubio2025MPUGauging}, for block-dependent scalar
    cochains define
    \begin{align}
        \widehat\chi^{\,x}_g
         & =\chi^{\,g x}_{g^{-1}},
        \label{eq:mpug_generalized_hat_one}
    \end{align}
    and
    \begin{align}
        \widehat\Omega^{\,x}_{g,h}
         & =\Omega^{\,(gh)x}_{h^{-1},g^{-1}}.
        \label{eq:mpug_generalized_hat_two}
    \end{align}
\end{definition}

\begin{lemma}[Hat involution and generalized coboundaries]
    \label{thm:mpug_generalized_hat_coboundary}
    \lean{TNLean.Algebra.ActionTensorGauge.hat_hat,
        TNLean.Algebra.LSymbol.hat_hat,
        TNLean.Algebra.LSymbol.hat_coboundary}
    \leanok
    \uses{def:mpug_generalized_cocycle, def:mpug_generalized_hat}
    The block-dependent hat operation is an involution:
    \begin{align}
        \widehat{\widehat\chi} & =\chi,
        \label{eq:mpug_generalized_hat_hat_one}
    \end{align}
    and
    \begin{align}
        \widehat{\widehat\Omega} & =\Omega.
        \label{eq:mpug_generalized_hat_hat_two}
    \end{align}
    It commutes with the generalized coboundary:
    \begin{align}
        \widehat{d\chi} & =d\widehat\chi.
        \label{eq:mpug_generalized_hat_coboundary}
    \end{align}
\end{lemma}

\begin{proof}\leanok
    \uses{def:mpug_generalized_cocycle, def:mpug_generalized_hat}
    Applying (\ref{eq:mpug_generalized_hat_one}) twice and cancelling each
    group element with its inverse gives
    (\ref{eq:mpug_generalized_hat_hat_one}), and applying
    (\ref{eq:mpug_generalized_hat_two}) twice gives
    (\ref{eq:mpug_generalized_hat_hat_two}). For
    (\ref{eq:mpug_generalized_hat_coboundary}), the action law
    $g(h x)=(gh)x$ and $g^{-1}\bigl((gh)x\bigr)=h x$ give
    \begin{align}
        (d\widehat\chi)^x_{g,h}
         & =\frac{\widehat\chi^{\,h x}_g\widehat\chi^{\,x}_h}
        {\widehat\chi^{\,x}_{gh}}
        =\frac{\chi^{(gh)x}_{g^{-1}}\chi^{h x}_{h^{-1}}}
        {\chi^{(gh)x}_{(gh)^{-1}}}
        =(d\chi)^{(gh)x}_{h^{-1},g^{-1}}
        =\widehat{d\chi}^{\,x}_{g,h}.
    \end{align}
    The two sides therefore carry the same three values of $\chi$, reordered
    in $\C^\times$.
\end{proof}

\begin{definition}[Regular-action matrices]
    \label{def:mpug_generalized_regular_matrices}
    \lean{TNLean.Algebra.LSymbol.regularMatrix}
    \leanok
    \uses{def:mpug_l_symbols}
    Let $\Omega\colon\mathsf X\times G^2\to\C^\times$ be an $L$-symbol.
    As in \cite[Lemma~\texttt{lemma:G\_cocycle}]{FrancoRubio2025MPUGauging},
    define the regular-action matrix $X^x_g$ on the group algebra $\C G$ by
    \begin{align}
        X^x_g\ket{k}
         & =\Omega^{k^{-1}x}_{g,k}\ket{gk}.
        \label{eq:mpug_generalized_regular_matrix}
    \end{align}
    % Source anchor: arXiv:2502.20257, lines 5938--5942.
\end{definition}

\begin{lemma}[Generalized projective multiplication law]
    \label{thm:mpug_generalized_regular_multiplication}
    \lean{TNLean.Algebra.LSymbol.regularMatrix_mul}
    \leanok
    \uses{def:mpug_generalized_cocycle,
        def:mpug_generalized_regular_matrices}
    Let $G$ be finite. If $\Omega$
    satisfies~(\ref{eq:mpug_generalized_cocycle}), then
    \begin{align}
        X^{h x}_gX^x_h
         & =\Omega^x_{g,h}X^x_{gh}.
        \label{eq:mpug_generalized_regular_multiplication}
    \end{align}
    % Source anchor: arXiv:2502.20257, lemma:G_cocycle, lines 5931--5943,
    %   which states this multiplication law "for G being a finite group".
    %   The blueprint statement previously omitted that hypothesis, which the
    %   formal statement carries.
\end{lemma}

\begin{proof}\leanok
    \uses{def:mpug_generalized_cocycle,
        def:mpug_generalized_regular_matrices}
    Applying the left side to $\ket{k}$ yields the coefficient
    \begin{align}
        \Omega^{(hk)^{-1}h x}_{g,hk}\Omega^{k^{-1}x}_{h,k}
         & =\Omega^{k^{-1}x}_{g,hk}\Omega^{k^{-1}x}_{h,k} \\
         & =\Omega^x_{g,h}\Omega^{k^{-1}x}_{gh,k},
    \end{align}
    where the last equality is
    (\ref{eq:mpug_generalized_cocycle}) at the block $k^{-1}x$.
    Formula~(\ref{eq:mpug_generalized_regular_matrix}) now gives
    (\ref{eq:mpug_generalized_regular_multiplication}).
\end{proof}

\begin{lemma}[Finite-group power trivialization]
    \label{thm:mpug_generalized_cocycle_power}
    \lean{TNLean.Algebra.LSymbol.regularMatrix_eq_permutation_mul_diagonal,
        TNLean.Algebra.LSymbol.regularDeterminant_ne_zero,
        TNLean.Algebra.LSymbol.determinantCochain,
        TNLean.Algebra.LSymbol.pow_card_eq_coboundary_apply,
        TNLean.Algebra.LSymbol.pow_card_eq_coboundary}
    \leanok
    \uses{def:mpug_generalized_cocycle,
        def:mpug_generalized_regular_matrices}
    Let $G$ be finite and let $\Omega$ be a generalized $2$-cocycle. Then
    $\Omega^{|G|}$ is a generalized coboundary. More precisely, with
    \begin{align}
        \chi^x_g
         & =\det X^x_g,
        \label{eq:mpug_generalized_determinant_cochain}
    \end{align}
    one has
    \begin{align}
        \left(\Omega^x_{g,h}\right)^{|G|}
         & =\frac{\det X^{h x}_g\det X^x_h}{\det X^x_{gh}}
        =(d\chi)^x_{g,h}.
        \label{eq:mpug_generalized_power_coboundary}
    \end{align}
    This establishes
    \cite[Lemma~\texttt{lemma:G\_cocycle}]{FrancoRubio2025MPUGauging}; as
    recorded in \cite{gap:fbc25_generalized_cocycle_determinant_typo}, the
    second numerator determinant in
    (\ref{eq:mpug_generalized_power_coboundary}) carries the index $h$,
    correcting the printed formula.

    % Local fix: docs/paper-gaps/fbc25_generalized_cocycle_determinant_typo.tex.
    % Source anchor: arXiv:2502.20257, lemma:G_cocycle, lines 5931--5948.
\end{lemma}

\begin{proof}\leanok
    \uses{def:mpug_generalized_regular_matrices,
        thm:mpug_generalized_regular_multiplication}
    Let $P_g$ be the regular-action permutation matrix defined by
    $P_g\ket{k}=\ket{gk}$, and set
    $D^x_g=\operatorname{diag}(\Omega^{k^{-1}x}_{g,k})_{k\in G}$.
    Then
    \begin{align}
        X^x_g
         & =P_gD^x_g.
    \end{align}
    Since $\det P_g\in\{1,-1\}$ and every diagonal entry of $D^x_g$
    is nonzero, $\det X^x_g\ne0$. Thus
    (\ref{eq:mpug_generalized_determinant_cochain}) takes values in
    $\C^\times$. Taking determinants in
    (\ref{eq:mpug_generalized_regular_multiplication}), using
    multiplicativity of the determinant on the left and
    $\det(cX)=c^{|G|}\det X$ for the $|G|\times|G|$ matrix $X^x_{gh}$ on the
    right, gives
    \begin{align}
        \det X^{h x}_g\det X^x_h
         & =\left(\Omega^x_{g,h}\right)^{|G|}\det X^x_{gh}.
        \label{eq:mpug_generalized_determinant_product}
    \end{align}
    Dividing (\ref{eq:mpug_generalized_determinant_product}) by
    $\det X^x_{gh}$ gives
    (\ref{eq:mpug_generalized_power_coboundary}).
\end{proof}

\begin{lemma}[Positive generalized cocycles are coboundaries]
    \label{thm:mpug_positive_generalized_cocycle}
    \lean{TNLean.Algebra.PositiveUnits.positiveRpow,
        TNLean.Algebra.PositiveUnits.positiveRpow_inv_natCast_pow,
        TNLean.Algebra.ActionTensorGauge.positiveRpow,
        TNLean.Algebra.LSymbol.positivePrimitive,
        TNLean.Algebra.LSymbol.isPositiveValued_positivePrimitive,
        TNLean.Algebra.LSymbol.exists_positive_coboundary}
    \leanok
    \uses{def:mpug_generalized_cocycle,
        def:mpug_positive_generalized_cochains,
        thm:mpug_generalized_cocycle_power}
    Let $G$ be finite and let $G$ act on a set $\mathsf X$. If $\Omega$ is a
    positive generalized $2$-cocycle, then there is a positive
    block-dependent scalar $1$-cochain $\chi_0$ such that
    \begin{align}
        \Omega^x_{g,h}
         & =(d\chi_0)^x_{g,h}
        =\frac{\chi^{\,h x}_{0,g}\chi^x_{0,h}}
        {\chi^x_{0,gh}}
        \label{eq:mpug_positive_coboundary}
    \end{align}
    for all $x\in\mathsf X$ and $g,h\in G$.
    This establishes the first assertion of
    \cite[Corollary at lines~5950--5960]{FrancoRubio2025MPUGauging}; as
    recorded in \cite{gap:fbc25_positive_generalized_cocycle_block_index},
    whereas the source writes $\chi_g$,
    Lemma~\ref{thm:mpug_generalized_cocycle_power} constructs $\chi^x_g$
    through (\ref{eq:mpug_generalized_determinant_cochain}) and the block label
    $x$ is therefore retained.
    % Local correction:
    % docs/paper-gaps/fbc25_positive_generalized_cocycle_block_index.tex.
    % Source anchor: arXiv:2502.20257, cor:positive_trivial, lines 5950--5956.
\end{lemma}

\begin{proof}\leanok
    \uses{def:mpug_positive_generalized_cochains,
        thm:mpug_generalized_cocycle_power}
    With $\delta^x_g=\det X^x_g$ denoting the determinant cochain from
    Lemma~\ref{thm:mpug_generalized_cocycle_power}, set
    \begin{align}
        \chi^x_{0,g}
         & =\lvert\delta^x_g\rvert^{1/\lvert G\rvert}.
        \label{eq:mpug_positive_determinant_root}
    \end{align}
    As recorded in
    \cite{gap:fbc25_positive_generalized_cocycle_determinant_sign}, because
    the regular-action permutation can contribute a sign to $\delta^x_g$, its
    norm is taken before the positive $\lvert G\rvert$-th root.
    Since $\Omega$ is positive and the complex norm is multiplicative, taking
    norms in (\ref{eq:mpug_generalized_power_coboundary}) gives
    \begin{align}
        \Omega^{\lvert G\rvert} & =d\lvert\delta\rvert.
        \label{eq:mpug_positive_norm_power}
    \end{align}
    Positive real powers are multiplicative, so
    (\ref{eq:mpug_positive_determinant_root}) and
    (\ref{eq:mpug_positive_norm_power}) give
    \begin{align}
        (d\chi_0)^{\lvert G\rvert}
         & =d\lvert\delta\rvert=\Omega^{\lvert G\rvert},
        \label{eq:mpug_positive_root_power}
    \end{align}
    and the $\lvert G\rvert$-th power is injective on $\mathbb R_{>0}$. Hence
    $d\chi_0=\Omega$.
    % Consolidated clarification:
    % docs/paper-gaps/fbc25_positive_generalized_cocycle_determinant_sign.tex.
\end{proof}

\begin{lemma}[Hat-compatible positive primitive]
    \label{thm:mpug_positive_generalized_cocycle_hat}
    \lean{TNLean.Algebra.ActionTensorGauge.coboundary_div,
        TNLean.Algebra.ActionTensorGauge.coboundary_positiveRpow,
        TNLean.Algebra.LSymbol.exists_positive_coboundary_hat_mul_eq_one}
    \leanok
    \uses{def:mpug_generalized_hat,
        def:mpug_positive_generalized_cochains,
        thm:mpug_generalized_hat_coboundary,
        thm:mpug_positive_generalized_cocycle}
    Under the hypotheses of
    Lemma~\ref{thm:mpug_positive_generalized_cocycle}, suppose additionally
    that
    \begin{align}
        \Omega^x_{g,h}\widehat\Omega^{\,x}_{g,h}
         & =1
        \label{eq:mpug_positive_hat_hypothesis}
    \end{align}
    for every $x\in\mathsf X$ and $g,h\in G$. Then, as in the second assertion
    of \cite[Corollary at lines~5950--5960]{FrancoRubio2025MPUGauging}, the
    positive cochain can be chosen so that
    \begin{align}
        \Omega & =d\chi,
        \label{eq:mpug_positive_hat_coboundary}
    \end{align}
    and
    \begin{align}
        \chi^x_g\widehat\chi^{\,x}_g
         & =1.
        \label{eq:mpug_positive_hat_witness_identity}
    \end{align}
    % Source anchor: arXiv:2502.20257, cor:positive_trivial, lines 5955--5960.
\end{lemma}

\begin{proof}\leanok
    \uses{thm:mpug_generalized_hat_coboundary,
        thm:mpug_positive_generalized_cocycle}
    Let $\chi_0$ satisfy $d\chi_0=\Omega$ as in
    Lemma~\ref{thm:mpug_positive_generalized_cocycle}. Define
    \begin{align}
        \chi^x_g
         & =\sqrt{\frac{\chi^x_{0,g}}
            {\widehat\chi^{\,x}_{0,g}}}
        =\sqrt{\frac{\chi^x_{0,g}}
            {\chi^{\,g x}_{0,g^{-1}}}}.
        \label{eq:mpug_positive_hat_witness}
    \end{align}
    All values are positive, so the positive square roots are unambiguous.
    Using (\ref{eq:mpug_generalized_hat_coboundary}) and the hypothesis
    (\ref{eq:mpug_positive_hat_hypothesis}) gives
    \begin{align}
        d\chi
         & =\sqrt{d\chi_0/d\widehat\chi_0}
        =\sqrt{\Omega/\widehat\Omega}
        =\sqrt{\Omega^2}
        =\Omega.
    \end{align}
    The same substitution in
    (\ref{eq:mpug_positive_hat_witness}) gives
    \begin{align}
        \widehat\chi
         & =\sqrt{\widehat\chi_0/\chi_0}
        =\chi^{-1},
    \end{align}
    which proves
    (\ref{eq:mpug_positive_hat_witness_identity}).
\end{proof}

\begin{definition}[Inversion cochains]
    \label{def:mpug_inversion_cochains}
    \lean{TNLean.Algebra.ScalarOneCochain,
        TNLean.Algebra.ScalarOneCochain.coboundary,
        TNLean.Algebra.ScalarCocycle.hat,
        TNLean.Algebra.ScalarThreeCochain.hat,
        TNLean.Algebra.ScalarThreeCochain.Xi,
        TNLean.Algebra.ScalarThreeCochain.sigma}
    \leanok
    \uses{def:mpug_scalar_cochains}
    These conventions and inversion cochains follow
    \cite[lines~5910--5916 and 5967--5973]{FrancoRubio2025MPUGauging}.
    For scalar cochains with one, two, and three arguments, set
    \begin{align}
        (d\sigma)(g,h) & =\frac{\sigma(g)\sigma(h)}{\sigma(gh)}.
        \label{eq:mpug_one_coboundary}
    \end{align}
    The hat reverses the arguments and replaces each group element by its
    inverse:
    \begin{align}
        \widehat\beta(g,h) & =\beta(h^{-1},g^{-1}).
        \label{eq:mpug_two_hat}
    \end{align}
    \begin{align}
        \widehat\omega(g,h,k) & =\omega(k^{-1},h^{-1},g^{-1}).
        \label{eq:mpug_three_hat}
    \end{align}
    Define the inversion cochains by
    \begin{align}
        \Xi(g,h) & =
        \frac{\omega(h^{-1},g^{-1},g)}
        {\omega(h^{-1}g^{-1},g,h)}.
        \label{eq:mpug_intro_Xi}
    \end{align}
    \begin{align}
        \sigma(g) & =\omega(g,g^{-1},g).
        \label{eq:mpug_intro_sigma}
    \end{align}
    % Source anchors: arXiv:2502.20257, lines 5910--5916 and eq:introXi.
\end{definition}

\begin{lemma}[The hat operation is an involution]
    \label{thm:mpug_scalar_hat_involution}
    \lean{TNLean.Algebra.ScalarCocycle.hat_hat,
        TNLean.Algebra.ScalarThreeCochain.hat_hat}
    \leanok
    \uses{def:mpug_inversion_cochains}
    Let $G$ be an arbitrary group. Every scalar $2$-cochain $\beta$ and every
    scalar $3$-cochain $\omega$ satisfy
    \begin{align}
        \widehat{\widehat\beta} & =\beta,
        \label{eq:mpug_two_hat_hat}
    \end{align}
    \begin{align}
        \widehat{\widehat\omega} & =\omega.
        \label{eq:mpug_three_hat_hat}
    \end{align}
    % Source anchor: arXiv:2502.20257, line 5912, where the hat operation is
    %   introduced as an involution for cochains of any arity.
\end{lemma}

\begin{proof}\leanok
    \uses{def:mpug_inversion_cochains}
    Reversing the arguments of (\ref{eq:mpug_two_hat}) and inverting each of
    them twice restores the original arguments, and likewise for
    (\ref{eq:mpug_three_hat}).
\end{proof}

\begin{lemma}[Two inversion-cochain formulas]
    \label{thm:mpug_inversion_cochain_formulas}
    \lean{TNLean.Algebra.ScalarThreeCochain.Xi_eq_second_formula,
        TNLean.Algebra.ScalarThreeCochain.sigma_inv}
    \leanok
    \uses{def:mpug_scalar_cochains, def:mpug_three_cocycle,
        def:mpug_inversion_cochains}
    These formulas are part of the cocycle calculation in
    \cite[lines~5967--5979]{FrancoRubio2025MPUGauging}.
    Let $G$ be an arbitrary group and let $\omega$ be a normalized scalar
    $3$-cocycle. Then
    \begin{align}
        \Xi(g,h) & =
        \frac{\omega(h^{-1},g^{-1},g)}
        {\omega(h^{-1}g^{-1},g,h)}
        =
        \frac{\omega(h^{-1},g^{-1},gh)}
        {\omega(g^{-1},g,h)}.
        \label{eq:mpug_Xi_two_formulas}
    \end{align}
    The inversion one-cochain satisfies
    \begin{align}
        \sigma(g^{-1}) & =\sigma(g)^{-1}.
        \label{eq:mpug_sigma_inverse}
    \end{align}
    % Source anchor: arXiv:2502.20257, eq:introXi and lemma:omegaXi.
\end{lemma}

\begin{proof}\leanok
    \uses{def:mpug_scalar_cochains, def:mpug_three_cocycle,
        def:mpug_inversion_cochains}
    Applying (\ref{eq:mpug_three_cocycle}) to $(h^{-1},g^{-1},g,h)$ and using
    (\ref{eq:mpug_three_normalized}) gives the second quotient in
    (\ref{eq:mpug_Xi_two_formulas}). Applying
    (\ref{eq:mpug_three_cocycle}) to
    $(g,g^{-1},g,g^{-1})$ gives
    $1=\sigma(g)\sigma(g^{-1})$, hence
    (\ref{eq:mpug_sigma_inverse}).
\end{proof}

\begin{lemma}[Inversion identities for a normalized three-cocycle]
    \label{thm:mpug_omega_Xi}
    \lean{TNLean.Algebra.ScalarThreeCochain.hat_eq_mul_coboundary_Xi,
        TNLean.Algebra.ScalarThreeCochain.Xi_eq_hat_mul_coboundary_sigma}
    \leanok
    \uses{def:mpug_scalar_cochains, def:mpug_three_cocycle,
        def:mpug_fusion_gauge, def:mpug_inversion_cochains}
    These inversion identities and their cocycle substitutions are given in
    \cite[lines~5967--5979]{FrancoRubio2025MPUGauging}.
    Let $\omega$ be a normalized scalar $3$-cocycle. Then
    \begin{align}
        \widehat\omega & =\omega\,d\Xi.
        \label{eq:mpug_hat_omega}
    \end{align}
    The inversion two-cochain also satisfies
    \begin{align}
        \Xi & =\widehat\Xi\,d\sigma.
        \label{eq:mpug_Xi_hat_sigma}
    \end{align}
    % Source anchor: arXiv:2502.20257, lemma:omegaXi, lines 5967--5979.
\end{lemma}

\begin{proof}\leanok
    \uses{def:mpug_three_cocycle, def:mpug_fusion_gauge,
        def:mpug_inversion_cochains, thm:mpug_inversion_cochain_formulas}
    For (\ref{eq:mpug_hat_omega}), multiplying the cocycle identities at
    \begin{align}
         & (k^{-1},h^{-1},g^{-1},g),\quad
        (k^{-1},h^{-1}g^{-1},g,h),\quad
        (k^{-1}h^{-1}g^{-1},g,h,k),
    \end{align}
    and cancelling the normalized factors gives the equality. For
    (\ref{eq:mpug_Xi_hat_sigma}), substituting the second expression in
    (\ref{eq:mpug_Xi_two_formulas}) for $\widehat\Xi$ and combining the
    identities at
    \begin{align}
         & (h^{-1}g^{-1},g,h,h^{-1}g^{-1}),\quad
        (h^{-1},h,h^{-1},g^{-1}),\quad
        (h^{-1},g^{-1},g,g^{-1}),\quad
        (g,g^{-1},g,g^{-1})
    \end{align}
    gives (\ref{eq:mpug_sigma_inverse}) in the three required places at $g$,
    $h$, and $gh$.
\end{proof}

Let $\Z_2=\{e,s\}$ with $s^2=e$.

\begin{lemma}[Normalized $L$-symbol ratio]
    \label{thm:mpug_sigma_l_ratio}
    \lean{TNLean.Algebra.LSymbol.sigma_eq_div}
    \leanok
    \uses{def:mpug_l_symbols, def:mpug_l_normalized,
        def:mpug_inversion_cochains}
    Let $G$ act on a set $\mathsf X$. Suppose that a normalized $L$-symbol
    $L$ is compatible with a scalar $3$-cochain $\omega$. Then, as in
    \cite[lines~5327--5329]{FrancoRubio2025MPUGauging}, for every
    $x\in\mathsf X$ and $g\in G$,
    \begin{align}
        \sigma(g)
         & =\frac{L^x_{g^{-1},g}}{L^{g x}_{g,g^{-1}}}.
        \label{eq:mpug_sigma_l_ratio}
    \end{align}
    % Source anchor: arXiv:2502.20257, lines 5327--5329.
\end{lemma}

\begin{proof}\leanok
    \uses{def:mpug_l_symbols, def:mpug_l_normalized,
        def:mpug_inversion_cochains}
    Applying the compatibility relation
    (\ref{eq:mpug_l_compatibility}) to $(g,g^{-1},g)$, the two factors with an
    identity argument are $1$ by normalization, leaving
    (\ref{eq:mpug_sigma_l_ratio}).
\end{proof}

\begin{lemma}[Trivial scalar gauge class for $\Z_2$]
    \label{thm:mpug_z2_trivial_scalar_class}
    \lean{TNLean.Algebra.ScalarThreeCochain.isTrivialGaugeClass_iff_sigma_generator_eq_one}
    \leanok
    \uses{def:mpug_scalar_cochains, def:mpug_fusion_gauge,
        def:mpug_trivial_gauge_class, def:mpug_inversion_cochains}
    Let $\omega$ be a normalized scalar $3$-cochain on $\Z_2$. Its gauge
    class is trivial if and only if
    \begin{align}
        \sigma(s) & =1.
        \label{eq:mpug_z2_trivial_sigma}
    \end{align}
    This strengthens the scalar step in
    \cite[lines~5339--5341]{FrancoRubio2025MPUGauging}: only normalization,
    rather than the cocycle identity, is needed for this equivalence.
\end{lemma}

\begin{proof}\leanok
    \uses{def:mpug_fusion_gauge,
        thm:mpug_gauge_preserves_normalization,
        thm:mpug_cohomologous_to_equivalence}
    Suppose first that $\omega=d\beta$ for an arbitrary scalar $2$-cochain
    $\beta$. Although $\beta$ need not be normalized, the normalization of
    $d\beta$ implies that there is a constant $c\in\C^\times$ such that
    $\beta(e,h)=\beta(h,e)=c$ for every $h\in\Z_2$. Since $s^2=e$,
    \begin{align}
        (d\beta)(s,s,s)
         & =\frac{\beta(s,e)\beta(s,s)}
        {\beta(s,s)\beta(e,s)}=\frac{c}{c}=1.
        \label{eq:mpug_z2_coboundary_generator}
    \end{align}
    Thus every scalar $2$-cochain whose coboundary is normalized contributes
    the factor $1$ at $(s,s,s)$. In particular, an arbitrary gauge between
    normalized scalar $3$-cochains cannot change this value, and
    $\sigma(s)=1$. Conversely, normalization determines every
    value of $\omega$ except $\omega(s,s,s)=\sigma(s)$. Hence
    (\ref{eq:mpug_z2_trivial_sigma}) makes $\omega$ the constant cochain $1$.
\end{proof}

\begin{lemma}[Sign dichotomy for the $\Z_2$ inversion scalar]
    \label{thm:mpug_z2_sigma_sign}
    \lean{TNLean.Algebra.ScalarThreeCochain.sigma_generator_eq_one_or_eq_neg_one}
    \leanok
    \uses{def:mpug_three_cocycle, def:mpug_inversion_cochains}
    Let $\omega$ be a normalized scalar $3$-cocycle on $\Z_2$. Then
    \begin{align}
        \sigma(s) & \in\{1,-1\}.
        \label{eq:mpug_z2_sigma_sign}
    \end{align}
\end{lemma}

\begin{proof}\leanok
    \uses{thm:mpug_inversion_cochain_formulas}
    Since $s^{-1}=s$, formula (\ref{eq:mpug_sigma_inverse}) gives
    $\sigma(s)=\sigma(s)^{-1}$. Hence $\sigma(s)^2=1$, whose two solutions in
    $\C^\times$ are $1$ and $-1$.
\end{proof}

For an element $g$ whose square is the identity, the source condition
$\xi_g(x)\xi_{g^{-1}}(g\mathbin{\cdot}x)=-1$ reduces to
$\xi_g(x)\xi_g(g\mathbin{\cdot}x)=-1$ because $g^{-1}=g$.

\begin{definition}[Anomalous sign for an involution]
    \label{def:mpug_anomalous_involution_sign}
    \lean{TNLean.Algebra.IsComplexUnitSign,
        TNLean.Algebra.HasAnomalousInvolutionSign}
    \leanok
    Let $f\colon\mathsf X\to\mathsf X$. An anomalous sign for $f$ is a
    function $\xi\colon\mathsf X\to\C^\times$ such that
    \begin{align}
        \xi(x) & \in\{1,-1\},
        \label{eq:mpug_anomalous_sign_range}
    \end{align}
    and
    \begin{align}
        \xi(x)\xi(f(x)) & =-1
        \label{eq:mpug_anomalous_sign_pairing}
    \end{align}
    for every $x\in\mathsf X$.
    % Source anchor: arXiv:2502.20257, prop:def_tens_anomalous,
    % eq:defects_Z2b, lines 2979--2981.
\end{definition}

\begin{lemma}[Anomalous signs and fixed points]
    \label{thm:mpug_anomalous_involution_sign_iff}
    \lean{TNLean.Algebra.hasAnomalousInvolutionSign_of_fixedPointFree,
        TNLean.Algebra.not_hasAnomalousInvolutionSign_of_fixedPoint,
        TNLean.Algebra.hasAnomalousInvolutionSign_iff_fixedPointFree,
        TNLean.Algebra.involutive_smul_of_mul_self_eq_one,
        TNLean.Algebra.hasAnomalousInvolutionSign_smul_iff}
    \leanok
    \uses{def:mpug_anomalous_involution_sign}
    Let $f\colon\mathsf X\to\mathsf X$ be an involution. An anomalous sign
    exists if and only if $f$ has no fixed point. In particular, if a group
    acts on $\mathsf X$ from the left and $g^2=e$, the sign condition for
    anomalous defect tensors in
    \cite[lines~2979--2981]{FrancoRubio2025MPUGauging} holds for the action
    $x\mapsto g\mathbin{\cdot}x$ if and only if
    $g\mathbin{\cdot}x\neq x$ for every $x\in\mathsf X$.
    % Source anchor: arXiv:2502.20257, prop:def_tens_anomalous,
    % eq:defects_Z2b, lines 2979--2981.
\end{lemma}

\begin{proof}\leanok
    \uses{def:mpug_anomalous_involution_sign}
    Choosing a well-order on $\mathsf X$ and assigning $1$ to the smaller
    element and $-1$ to the larger in each two-element orbit $\{x,f(x)\}$,
    involutivity exchanges the two elements so that their signs multiply to
    $-1$. Conversely, at a fixed point the product is the square of either
    $1$ or $-1$, hence equals $1$ rather than $-1$. For the action statement,
    $g\mathbin{\cdot}(g\mathbin{\cdot}x)=(g^2)\mathbin{\cdot}x=x$.
\end{proof}

\begin{lemma}[Nontrivial scalar gauge class for $\Z_2$]
    \label{thm:mpug_z2_nontrivial_scalar_class}
    \lean{TNLean.Algebra.ScalarThreeCochain.not_isTrivialGaugeClass_iff_sigma_generator_eq_neg_one}
    \leanok
    \uses{def:mpug_three_cocycle, def:mpug_trivial_gauge_class,
        def:mpug_inversion_cochains}
    For a normalized scalar $3$-cocycle $\omega$ on $\Z_2$, the scalar
    cohomology classification used in the proof at
    \cite[lines~5339--5341]{FrancoRubio2025MPUGauging} states that its gauge
    class is nontrivial if and only if
    \begin{align}
        \sigma(s) & =-1.
        \label{eq:mpug_z2_nontrivial_sigma}
    \end{align}
    % Source anchor: arXiv:2502.20257, lines 5339--5341.
\end{lemma}

\begin{proof}\leanok
    \uses{thm:mpug_z2_trivial_scalar_class, thm:mpug_z2_sigma_sign}
    By Lemma~\ref{thm:mpug_z2_trivial_scalar_class}, a nontrivial class has
    $\sigma(s)\neq1$. The dichotomy
    (\ref{eq:mpug_z2_sigma_sign}) then forces $\sigma(s)=-1$, and the converse
    follows from the same two statements.
\end{proof}

This completes the classification of the normalized scalar $3$-cocycle by
$\sigma(s)$ in the setting of \cite[lines~5207--5338]{FrancoRubio2025MPUGauging}.

\section{The GHZ--blocked-cluster tensor}
\label{sec:mpug_ghz_blocked_cluster}

\begin{definition}[GHZ--blocked-cluster tensor and its two blocks]
    \label{def:mpug_ghz_blocked_cluster_tensor}
    \lean{MPSTensor.clusterBlocked,
        MPSTensor.ghzSectorTensor,
        MPSTensor.ghzSectorTensor_apply_same,
        MPSTensor.z4z2GHZClusterTensor,
        MPSTensor.z4z2GHZClusterBlock,
        MPSTensor.z4z2GHZClusterBlock_apply,
        MPSTensor.z4z2GHZClusterDirectSum}
    \leanok
    \discussion{7699}
    \uses{def:ghz_tensor, def:cluster_tensor,
        def:mps_independent_tensor_product}
    Let $C=\{C^q\}_{q=0}^3$ be the once-blocked cluster-state tensor. The tensor
    obtained by taking its independent tensor product with the GHZ tensor has
    physical dimension $8$ and bond dimension $4$:
    \begin{align}
        A & =A^{\mathrm{GHZ}}\boxtimes C.
        \label{eq:mpug_ghz_blocked_cluster_tensor}
    \end{align}
    For $x\in\{0,1\}$, let $E_x$ be the one-dimensional GHZ sector supported
    at the physical value $x$. The two bond-dimension-$2$ tensors
    \begin{align}
        A_x & =E_x\boxtimes C
    \end{align}
    satisfy, for $p,x\in\{0,1\}$ and $q\in\{0,1,2,3\}$,
    \begin{align}
        A_x^{(p,q)} & =\delta_{p,x}C^q.
        \label{eq:mpug_ghz_blocked_cluster_blocks}
    \end{align}
    As in \cite[lines~4427--4479]{FrancoRubio2025MPUGauging}, their canonical
    direct sum is denoted $A_0\oplus A_1$.
    % Source anchor: arXiv:2502.20257, example:z4z2, eq:z4z2MPS and
    % lines 4427--4479.
\end{definition}

\begin{lemma}[Injective and locally orthogonal GHZ--cluster blocks]
    \label{thm:mpug_ghz_blocked_cluster_blocks}
    \lean{MPSTensor.ghzSectorTensor_isInjective,
        MPSTensor.z4z2GHZClusterBlock_isInjective,
        MPSTensor.z4z2GHZClusterBlocks_isBNTLocallyOrthogonal,
        MPSTensor.z4z2GHZClusterDirectSum_eq_tensor}
    \leanok
    \discussion{7699}
    \uses{def:mpug_ghz_blocked_cluster_tensor, def:injective,
        def:is_bnt_locally_orthogonal}
    Each tensor $A_x$ is injective. For $x\neq y$ and every
    $X\in\operatorname{Mat}_2(\C)$, their one-site mixed transfer map satisfies
    \begin{align}
        \mathcal E_{x,y}(X)
         & =\sum_{p=0}^1\sum_{q=0}^3
        A_x^{(p,q)}X\bigl(A_y^{(p,q)}\bigr)^\dagger
        =0.
        \label{eq:mpug_ghz_blocked_cluster_local_orthogonality}
    \end{align}
    Thus the two blocks are locally orthogonal as observed at
    \cite[line~4481]{FrancoRubio2025MPUGauging}, and
    \begin{align}
        A_0\oplus A_1 & =A
        \label{eq:mpug_ghz_blocked_cluster_direct_sum}
    \end{align}
    after identifying the direct-sum and product bond
    coordinates.
    % The theorem does not construct the Z4 x Z2 action, its virtual action
    % tensors, its L-symbols, or its gauging data. The group action and gauging
    % of this example remain issue #7352.
\end{lemma}

\begin{proof}\leanok
    \uses{def:mpug_ghz_blocked_cluster_tensor, def:injective,
        def:is_bnt_locally_orthogonal, thm:mps_word_span_tensor_product,
        thm:cluster_normal}
    The one-dimensional sector $E_x$ is injective, and injectivity is preserved
    by independent tensor products, so $A_x=E_x\boxtimes C$ is injective.
    If $x\neq y$, formula~(\ref{eq:mpug_ghz_blocked_cluster_blocks}) makes
    every summand in
    (\ref{eq:mpug_ghz_blocked_cluster_local_orthogonality}) vanish: a physical
    value $p$ cannot equal both $x$ and $y$. For $p,x,y\in\{0,1\}$,
    $q\in\{0,1,2,3\}$, and $a,b\in\{0,1\}$, the direct-sum and product tensors
    have the same entry in their canonical bond coordinates:
    \begin{align}
        (A_0\oplus A_1)^{(p,q)}_{(x,a),(y,b)}
         & =A^{(p,q)}_{(x,a),(y,b)}
        =\delta_{x,y}\delta_{p,x}C^q_{a,b}.
    \end{align}
    Equality of all entries proves
    (\ref{eq:mpug_ghz_blocked_cluster_direct_sum}).
\end{proof}

\begin{lemma}[GHZ--cluster physical representation and virtual action table]
    \label{thm:mpug_ghz_cluster_action_table}
    % Principal declaration: MPSTensor.z4z2GHZClusterBlock_covariance.
    \lean{MPSTensor.z4z2GHZClusterBlock_covariance,
        MPSTensor.z4z2GHZClusterAction,
        MPSTensor.z4z2GHZClusterAction_unitary,
        MPSTensor.z4z2GHZClusterAction_generators,
        MPSTensor.z4z2GHZClusterIndex,
        MPSTensor.z4z2GHZClusterBlockAction,
        MPSTensor.z4z2GHZClusterVirtual,
        MPSTensor.z4z2GHZClusterVirtual_table,
        MPSTensor.z4z2GHZClusterVirtual_adjoint,
        MPSTensor.z4z2GHZClusterVirtual_sq,
        MPSTensor.z4z2GHZClusterVirtual_stabilizer}
    \leanok
    \discussion{7747}
    \uses{def:mpug_ghz_blocked_cluster_tensor}
    As in \cite[lines~4481--4489]{FrancoRubio2025MPUGauging}, write the
    physical basis as $|p,s,r\rangle$, with integer index $4p+2s+r$ for
    $p,s,r\in\{0,1\}$. Thus the blocked cluster index is $q=2s+r$. The
    operators
    $U_{(1,0)}=(X_1\otimes X_2)\mathrm{CNOT}_{1\to2}$ and
    $U_{(0,1)}=X_3$ generate a unitary representation of
    $G=\Z_4\times\Z_2$, with $U_{(2,0)}=X_2$.
    Put $g\cdot x=x+a\pmod 2$ for $g=(a,b)$ and $x\in\{0,1\}$.
    Let $V_{g,x}$ be given by the following table, where
    $Y=\left(\begin{smallmatrix}0&-i\\i&0\end{smallmatrix}\right)$:
    \begin{center}
        \begin{tabular}{ccccccccc}
            $g$   & $(0,0)$ & $(2,0)$ & $(0,1)$ & $(2,1)$
                  & $(1,0)$ & $(3,0)$ & $(1,1)$ & $(3,1)$                             \\
            $x=0$ & $\Id$   & $X$     & $Z$     & $Y$     & $X$   & $\Id$ & $Y$ & $Z$ \\
            $x=1$ & $\Id$   & $X$     & $Z$     & $Y$     & $\Id$ & $X$   & $Z$ & $Y$ \\
        \end{tabular}
    \end{center}
    Every $V_{g,x}$ is Hermitian and squares to the identity. For every
    physical index $i$, the blocks satisfy
    \begin{align}
        \sum_j (U_g)_{i,j} A_x^j
         & =V_{g,x}A_{g\cdot x}^iV_{g,x}^{\dagger}.
        \label{eq:mpug_ghz_cluster_block_covariance}
    \end{align}
    On the subgroup $H=\{(2a,b):a,b\in\{0,1\}\}$, the comparison with the
    cluster virtual representatives $W_{(b,a)}=Z^bX^a$ is
    $V_{(2a,b),x}=(-i)^{ab}W_{(b,a)}$.
    No claim is made about defect normalization, $L$-symbols, cocycle
    nonextension, block independence, or gauging.
    % Scoped calculation only: no defect normalization, L-symbols,
    % cocycle nonextension, block independence, or gauging is asserted.
    % Those remaining parts of example:z4z2 stay open in #7352.
\end{lemma}

\begin{proof}\leanok
    \uses{def:mpug_ghz_blocked_cluster_tensor, thm:mpug_monomial_matrix_laws}
    On basis vectors the generators send $(p,s,r)$ to $(p+1,s+p+1,r)$
    and $(p,s,r+1)$, with coordinates added modulo two. Their permutations
    have orders four and two and commute. The corresponding permutation
    matrices are unitary and preserve multiplication.
    Substituting $A_x^{(p,q)}=\delta_{p,x}C^q$ into
    (\ref{eq:mpug_ghz_cluster_block_covariance}) and multiplying the matrices
    for each of the eight group elements and two block labels gives the
    displayed table. Hermiticity and $V_{g,x}^2=\Id$ turn the forward action
    on the physical index into the stated conjugation identity.
    Finally, $ZX=iY$, so the only phase needed to compare the restriction
    with $Z^bX^a$ is $-i$ when $a=b=1$.
\end{proof}

\section{The gauge-invariant subspace of the CZX circuit tuple}
\label{sec:mpug_czx_gauge_invariant}

The CZX model of \cite[lines~4503--5183]{FrancoRubio2025MPUGauging} is the
anomalous $\Z_2=\{e,g\}$ matrix product unitary acting on the GHZ state.
After blocking, one site carries two matter qubits, and the fusion operators on
two neighboring blocked sites are $\lambda^R_{e,e}=\lambda^R_{e,g}=\Id$,
$\lambda^R_{g,e}=w_R$, together with the modified fusion operator
$\tilde\lambda^R_{g,g}=\lambda^R_{g,g}Z_1$ introduced at the end of that
passage.
For the displayed circuit tuple $(\Id,\Id,w,\tilde\lambda)$, the common $+1$
eigenspace of the resulting Gauss projectors has dimension $2^N$ on
$N\geq3$ blocked sites, answering the scaling question of
\cite[lines~5198--5204]{FrancoRubio2025MPUGauging} for this tuple only. The
displayed tuple is not identified with a member of the full physical
completion class of the CZX defect data; that identification requires the
three defect domains and maps, which the displayed matrices do not specify.

\subsection{Monomial operators and their common fixed space}

\begin{definition}[Monomial matrix]
    \label{def:mpug_monomial_matrix}
    \lean{Matrix.monomial,
        Matrix.monomial_apply,
        Matrix.monomial_apply_perm,
        Matrix.monomial_apply_of_ne}
    \leanok
    Let $\sigma$ be a permutation of a finite index set $S$ and let
    $\varphi\colon S\to\C$. The monomial matrix $M(\sigma,\varphi)$ is the
    operator on $\C^S$ with
    \begin{align}
        M(\sigma,\varphi)\ket{s} & =\varphi(s)\ket{\sigma(s)}.
        \label{eq:mpug_monomial_matrix}
    \end{align}
    Its entry in row $\sigma(s)$ and column $s$ is $\varphi(s)$, and every
    other entry vanishes.
\end{definition}

\begin{lemma}[Calculus of monomial matrices]
    \label{thm:mpug_monomial_matrix_laws}
    \lean{Matrix.monomial_mulVec,
        Matrix.monomial_mulVec_eq_self_iff,
        Matrix.monomial_mul_monomial,
        Matrix.conjTranspose_monomial,
        Matrix.smul_monomial,
        Matrix.neg_monomial,
        Matrix.monomial_one,
        Matrix.reindex_monomial,
        Matrix.monomial_mem_unitaryGroup,
        Matrix.trace_monomial,
        Matrix.trace_monomial_eq_zero}
    \leanok
    \uses{def:mpug_monomial_matrix}
    A vector $v\in\C^S$ is fixed by $M(\sigma,\varphi)$ if and only if
    $v_{\sigma(s)}=\varphi(s)v_s$ for every $s\in S$. Products, adjoints,
    scalar multiples, negatives, and simultaneous relabellings of monomial
    matrices are monomial:
    \begin{align}
        M(\sigma,\varphi)M(\tau,\psi)
         & =M\bigl(\sigma\tau,\ s\mapsto\varphi(\tau(s))\psi(s)\bigr),
        \label{eq:mpug_monomial_mul}                                               \\
        M(\sigma,\varphi)^\dagger
         & =M\bigl(\sigma^{-1},\ t\mapsto\overline{\varphi(\sigma^{-1}(t))}\bigr).
        \label{eq:mpug_monomial_adjoint}
    \end{align}
    The identity is $M(\id,1)$, and $M(\sigma,\varphi)$ is unitary when
    $|\varphi(s)|=1$ for every $s$. The trace is
    $\tr M(\sigma,\varphi)=\sum_{s:\,\sigma(s)=s}\varphi(s)$, which vanishes
    when $\sigma$ has no fixed point.
\end{lemma}

\begin{proof}\leanok
    \uses{def:mpug_monomial_matrix}
    The action of $M(\sigma,\varphi)$ on a vector is
    $(M(\sigma,\varphi)v)_t=\varphi(\sigma^{-1}(t))v_{\sigma^{-1}(t)}$, which
    gives the fixed-vector criterion after the substitution $t=\sigma(s)$.
    The remaining identities are entrywise: in a product only the column
    $\tau(s)$ of the first factor meets the row $\tau(s)$ of the second, the
    adjoint transposes the single nonzero entry of each column, and a
    relabelling conjugates the permutation. For unitarity,
    (\ref{eq:mpug_monomial_mul}) and (\ref{eq:mpug_monomial_adjoint}) give
    $M(\sigma,\varphi)^\dagger M(\sigma,\varphi)=M(\id,|\varphi|^2)=\Id$.
    The diagonal entry at $s$ is $\varphi(s)$ if $\sigma(s)=s$ and zero
    otherwise, which gives the trace.
\end{proof}

\begin{lemma}[Flat phases on the hypercube have a potential]
    \label{thm:mpug_hypercube_potential}
    \lean{TNLean.Algebra.hypercube_induction,
        TNLean.Algebra.exists_potential_of_flat}
    \leanok
    Let $M$ be a commutative group and let
    $\varphi_j\colon\mathbb F_2^r\to M$ for $0\leq j<r$ be phases attached
    to the edges $(\gamma,\gamma+e_j)$ of the hypercube. Assume that
    \begin{align}
        \varphi_j(\gamma)\varphi_j(\gamma+e_j)
         & =1,
        \label{eq:mpug_hypercube_involutive}       \\
        \varphi_j(\gamma)\varphi_k(\gamma+e_j)
         & =\varphi_k(\gamma)\varphi_j(\gamma+e_k)
        \label{eq:mpug_hypercube_flat}
    \end{align}
    for all $j,k$ and $\gamma$. Then there is $\Phi\colon\mathbb F_2^r\to M$
    with $\Phi(0)=1$ and $\Phi(\gamma+e_j)=\varphi_j(\gamma)\Phi(\gamma)$ for
    every $j$ and $\gamma$. In addition, a property of vectors of
    $\mathbb F_2^r$ that holds at $0$ and is preserved by every coordinate
    flip holds everywhere.
\end{lemma}

\begin{proof}\leanok
    For the connectivity claim, every $\gamma$ is the sum of the $e_j$ over
    its support, so induction on the support carries the property from $0$ to
    $\gamma$.

    The potential is constructed by induction on $r$. Writing
    $\gamma=(a,\gamma')$ with $a\in\mathbb F_2$ and
    $\gamma'\in\mathbb F_2^{r-1}$, the phases
    $\gamma'\mapsto\varphi_{j+1}(0,\gamma')$ for $0\leq j<r-1$ satisfy
    (\ref{eq:mpug_hypercube_involutive}) and (\ref{eq:mpug_hypercube_flat})
    on $\mathbb F_2^{r-1}$, so the induction hypothesis gives a potential
    $\Phi'$ for them. Setting
    \begin{align}
        \Phi(0,\gamma') & =\Phi'(\gamma'),
        \qquad
        \Phi(1,\gamma')=\varphi_0(0,\gamma')\Phi'(\gamma'),
        \label{eq:mpug_hypercube_potential_step}
    \end{align}
    the edge relation in direction $0$ at $(0,\gamma')$ holds by definition,
    and at $(1,\gamma')$ it follows from
    (\ref{eq:mpug_hypercube_involutive}). The edge relation in a direction
    $j+1$ at $(0,\gamma')$ is the relation for $\Phi'$, and at $(1,\gamma')$
    it is the relation for $\Phi'$ combined with
    (\ref{eq:mpug_hypercube_flat}) for the square spanned by the directions
    $0$ and $j+1$ at $(0,\gamma')$.
\end{proof}

\begin{definition}[Trivial holonomy on a fiber]
    \label{def:mpug_trivial_holonomy}
    \lean{TNLean.Algebra.fiberFlip,
        TNLean.Algebra.fiberFlip_apply,
        TNLean.Algebra.IsTrivialHolonomy}
    \leanok
    Let $X$ be a finite set and consider the basis $X\times\mathbb F_2^r$.
    For $0\leq j<r$ let $\sigma_j$ be the coordinate flip
    $(x,\gamma)\mapsto(x,\gamma+e_j)$ and let
    $\varphi_j\colon X\times\mathbb F_2^r\to\C$. The phases have trivial
    holonomy on the fiber over $x\in X$ when
    \begin{align}
        \varphi_j(x,\gamma)\varphi_k(x,\gamma+e_j)
         & =\varphi_k(x,\gamma)\varphi_j(x,\gamma+e_k)
        \label{eq:mpug_trivial_holonomy}
    \end{align}
    for all $j,k$ and $\gamma\in\mathbb F_2^r$.
\end{definition}

\begin{lemma}[Monomial fixed-space criterion]
    \label{thm:mpug_monomial_fixed_space_criterion}
    \lean{TNLean.Algebra.reindex_mulVec,
        TNLean.Algebra.commonFixedSubmodule_reindex,
        TNLean.Algebra.finrank_commonFixedSubmodule_reindex,
        TNLean.Algebra.add_self_eq_zero_pi_zmod_two,
        TNLean.Algebra.eq_zero_of_mem_commonFixedSubmodule_monomial_fiberFlip,
        TNLean.Algebra.finrank_commonFixedSubmodule_monomial_fiberFlip}
    \leanok
    \uses{def:mpug_monomial_matrix, def:mpug_trivial_holonomy,
        def:mpug_common_fixed_subspace}
    In the setting of Definition~\ref{def:mpug_trivial_holonomy}, let
    $T_j=M(\sigma_j,\varphi_j)$ and assume that every $T_j$ is an involution,
    that is, $\varphi_j(s)\varphi_j(\sigma_j(s))=1$ for every basis element
    $s$. Then
    \begin{align}
        \dim\mathcal V(T)
         & =\#\{x\in X:\text{the phases have trivial holonomy on the fiber over }x\}.
        \label{eq:mpug_monomial_fixed_space_count}
    \end{align}
    The dimension of a common fixed subspace is also unchanged by a
    simultaneous relabelling of the basis.
\end{lemma}

\begin{proof}\leanok
    \uses{thm:mpug_monomial_matrix_laws, thm:mpug_hypercube_potential,
        thm:mpug_mem_common_fixed_subspace}
    By Lemma~\ref{thm:mpug_monomial_matrix_laws}, a vector $v$ lies in
    $\mathcal V(T)$ exactly when
    $v_{(x,\gamma+e_j)}=\varphi_j(x,\gamma)v_{(x,\gamma)}$ for all $j$, $x$,
    and $\gamma$. Let $X_0$ be the set of fibers with trivial holonomy and
    consider the linear map $E\colon\mathcal V(T)\to\C^{X_0}$,
    $v\mapsto(v_{(x,0)})_{x\in X_0}$.

    The map $E$ is injective: if $v\in\mathcal V(T)$ vanishes at one point
    of a fiber, the edge relation and the connectivity statement of
    Lemma~\ref{thm:mpug_hypercube_potential} make it vanish on the whole
    fiber. On a fiber with trivial holonomy, $v$ therefore vanishes as soon
    as $v_{(x,0)}=0$. On a fiber where (\ref{eq:mpug_trivial_holonomy})
    fails at some $(j,k,\gamma_0)$, the two paths from $(x,\gamma_0)$ to
    $(x,\gamma_0+e_j+e_k)$ give
    \begin{align}
        v_{(x,\gamma_0+e_j+e_k)} & =\varphi_j(x,\gamma_0)\varphi_k(x,\gamma_0+e_j)\,v_{(x,\gamma_0)}, \\
        v_{(x,\gamma_0+e_j+e_k)} & =\varphi_k(x,\gamma_0)\varphi_j(x,\gamma_0+e_k)\,v_{(x,\gamma_0)}.
    \end{align}
    The two phase products differ, so $v_{(x,\gamma_0)}=0$ and again $v$
    vanishes on the fiber.

    The map $E$ is also surjective: given $w\in\C^{X_0}$,
    Lemma~\ref{thm:mpug_hypercube_potential} provides a potential $\Phi_x$
    on each fiber with trivial holonomy for the invertible phases
    $\gamma\mapsto\varphi_j(x,\gamma)$, and setting
    $v_{(x,\gamma)}=w_x\Phi_x(\gamma)$ there and $v=0$ on the other fibers
    satisfies the potential edge relation, so $v\in\mathcal V(T)$,
    and $\Phi_x(0)=1$ gives $Ev=w$.

    Hence $\dim\mathcal V(T)=\dim\C^{X_0}=\#X_0$. For the relabelling
    statement, a relabelled matrix acts by relabelling the vector, acting,
    and relabelling back, so the common fixed subspace is the image of the
    original one under the coordinate relabelling.
\end{proof}

\begin{lemma}[Local Gauss operator of monomial matter unitaries]
    \label{thm:mpug_gauss_operator_monomial}
    \lean{TNLean.Algebra.gaussMonomialPerm,
        TNLean.Algebra.gaussMonomialPerm_apply,
        TNLean.Algebra.gaussMonomialPhase,
        TNLean.Algebra.gaussOperator_eq_monomial}
    \leanok
    \uses{def:mpug_gauss_operator, def:mpug_monomial_matrix}
    Suppose that every matter unitary of the tuple $R$ is monomial,
    $R_{a,b}=M(\sigma_{a,b},\varphi_{a,b})$. Then the local Gauss operator
    of Definition~\ref{def:mpug_gauss_operator} is the monomial operator
    \begin{align}
        \mathcal G_g\ket{i,(a,b)}
         & =\overline{\varphi_{T_g(a,b)}\bigl(\sigma_{T_g(a,b)}^{-1}\sigma_{a,b}(i)\bigr)}\,
        \varphi_{a,b}(i)\,
        \ket{\sigma_{T_g(a,b)}^{-1}\sigma_{a,b}(i),\ T_g(a,b)}.
        \label{eq:mpug_gauss_operator_monomial}
    \end{align}
\end{lemma}

\begin{proof}\leanok
    \uses{thm:mpug_monomial_matrix_laws, def:mpug_gauss_operator}
    The block of $\mathcal G_g$ from $(a,b)$ to $T_g(a,b)$ is
    $R_{T_g(a,b)}^\dagger R_{a,b}$, which is monomial by
    (\ref{eq:mpug_monomial_mul}) and (\ref{eq:mpug_monomial_adjoint}); all
    other blocks vanish.
\end{proof}

\begin{lemma}[Periodic placement of monomial operators]
    \label{thm:mpug_embed_monomial}
    \lean{MPOTensor.windowPerm,
        MPOTensor.windowPerm_apply,
        MPOTensor.embedLocalOperator_monomial,
        MPOTensor.extractWindow_two,
        MPOTensor.replaceWindow_two_apply}
    \leanok
    \uses{def:mpug_monomial_matrix}
    Let $M(\sigma,\varphi)$ be a monomial operator on $L$ consecutive sites.
    Its placement on the periodic window beginning at site $i$ of an $N$-site
    chain is the monomial operator whose permutation replaces the window of a
    configuration $s$ by $\sigma$ applied to that window and whose phase is
    $\varphi$ evaluated on the window of $s$. For $L=2$ and $N\geq2$ the
    window at $j$ is $(s_j,s_{j+1})$, and replacing it changes exactly the
    sites $j$ and $j+1$.
\end{lemma}

\begin{proof}\leanok
    \uses{def:mpug_monomial_matrix}
    The placed operator has entry $(t,s)$ equal to the window entry
    $M(\sigma,\varphi)_{t_W,s_W}$ when $t$ and $s$ agree outside the window
    and $0$ otherwise. Both conditions together say that $t$ is $s$ with its
    window replaced by $\sigma(s_W)$, and the entry is then
    $\varphi(s_W)$. The two-site statements are the cyclic offset
    computations $(j+N-j)\bmod N=0$ and $(j+1+N-j)\bmod N=1$.
\end{proof}

\subsection{The exact CZX matrix product unitary}

\begin{definition}[The once-blocked CZX tensor]
    \label{def:mpug_czx_exact_tensor}
    % Principal declaration: MPOTensor.CZX.tensor.
    \lean{MPOTensor.CZX.decoratedSiteTensor, MPOTensor.CZX.tensor,
        MPOTensor.CZX.complementSite, MPOTensor.CZX.complement,
        MPOTensor.CZX.edgeExponent, MPOTensor.CZX.cyclicExponent}
    \leanok
    \uses{def:mpdo_two_site_blocking}
    Write a physical index as $2a+b$, with $a,b\in\{0,1\}$, and let
    $\overline{(a,b)}=(1-a,1-b)$. The tensor in
    \cite[lines~4503--4547]{FrancoRubio2025MPUGauging}
    is the two-site block $B^{ab,cd}=C^{a,c}C^{b,d}$ of the bond-two tensor
    \begin{align}
        C^{u,d}_{\ell r}
         & =\delta_{u,1-d}\delta_{\ell,u}(-1)^{u+\ell r}.
    \end{align}
    Here the constituent indices $u,d,\ell,r$ are binary, the weight is
    $W_{\ell r}=(-1)^{\ell r}=\sqrt{2}H_{\ell r}$, where $H$ is the
    normalized Hadamard matrix, and the factor $(-1)^u$ is the output
    Pauli $Z$. Thus, if $P$ denotes the normalized tensor printed in
    \cite[lines~1906--1943]{Cirac2017MPU}, then
    $B=(Z\otimes Z)\operatorname{block}_2(\sqrt{2}P)$, not
    $\operatorname{block}_2(P)$.
    Define the local integer exponent $e((a,b),r)=a+b+ab+br$.
    For a periodic configuration $s_j=(a_j,b_j)$ of positive length $N$,
    let $\kappa_N(s)_j=\overline{s_j}$ and set
    \begin{align}
        E(s) & =\sum_{j=0}^{N-1}e(s_j,a_{j+1})
        =\sum_{j=0}^{N-1}(a_j+b_j+a_jb_j+b_ja_{j+1}),
    \end{align}
    where $j+1$ is interpreted modulo $N$ and the bit values in the sum
    are integers.
\end{definition}

\begin{lemma}[Coordinates and unitarity of the exact CZX tensor]
    \label{thm:mpug_czx_exact_tensor_unitary}
    % Principal declaration: MPOTensor.CZX.tensor_isMPU.
    \lean{MPOTensor.CZX.tensor_apply, MPOTensor.CZX.tensor_apply_bits,
        MPOTensor.CZX.mpo_tensor_apply, MPOTensor.CZX.mpo_tensor,
        MPOTensor.CZX.mpo_tensor_mem_unitaryGroup, MPOTensor.CZX.tensor_isMPU}
    \leanok
    \uses{def:mpug_czx_exact_tensor, def:mpug_monomial_matrix, def:mpu}
    The tensor $B$ has bond dimension two and coordinates
    \begin{align}
        B^{ab,cd}_{\ell r}
         & =\delta_{a,1-c}\delta_{b,1-d}\delta_{\ell,a}
        (-1)^{a+b+ab+br}.
    \end{align}
    Its periodic operator is unitary for every $N>0$, with entries and
    monomial form
    \begin{align}
        O_N(B)_{s,t} & =\delta_{s,\kappa_N(t)}(-1)^{E(s)},                     \\
        O_N(B)       & =M\bigl(\kappa_N,\ t\mapsto(-1)^{E(\kappa_N(t))}\bigr).
    \end{align}
    In particular, $B$ is a matrix product unitary. This assertion does
    not include its phase classification, GHZ action, or defect data.
\end{lemma}
\begin{proof}
    \leanok
    \uses{thm:trace_eval_word_eq_sum_cyclic, thm:mpug_monomial_matrix_laws}
    The two-site block is the matrix product contracted over a binary
    intermediate bond $m$. The second factor vanishes unless $m=b$, so
    \begin{align}
        B^{ab,cd}_{\ell r}
         & =\sum_{m\in\{0,1\}}C^{a,c}_{\ell m}C^{b,d}_{mr} \\
         & =C^{a,c}_{\ell b}C^{b,d}_{br}                   \\
         & =\delta_{a,1-c}\delta_{b,1-d}\delta_{\ell,a}
        (-1)^{a+\ell b+b+br}                               \\
         & =\delta_{a,1-c}\delta_{b,1-d}\delta_{\ell,a}
        (-1)^{a+b+ab+br}.
    \end{align}
    For the periodic contraction with $s_j=(a_j,b_j)$ and
    $t_j=(c_j,d_j)$, the cyclic trace expansion is
    \begin{align}
        O_N(B)_{s,t}
         & =\sum_{g:\Z/N\Z\to\{0,1\}}
        \prod_{j=0}^{N-1}B^{s_j,t_j}_{g_j,g_{j+1}}.
    \end{align}
    The factor at $j$ vanishes unless $g_j=a_j$, so only the path
    $g_j=a_j$ can contribute. Substituting it gives
    \begin{align}
        O_N(B)_{s,t}
         & =\prod_{j=0}^{N-1}B^{s_j,t_j}_{a_j,a_{j+1}} \\
         & =\prod_{j=0}^{N-1}
        \delta_{a_j,1-c_j}\delta_{b_j,1-d_j}
        (-1)^{a_j+b_j+a_jb_j+b_ja_{j+1}}               \\
         & =\delta_{s,\kappa_N(t)}(-1)^{E(s)}.
    \end{align}
    Thus the operator is the monomial matrix for complementation, with
    column phase $(-1)^{E(\kappa_N(t))}$. Each phase has absolute value one,
    so Lemma~\ref{thm:mpug_monomial_matrix_laws} proves unitarity.
\end{proof}

\begin{lemma}[Injectivity of the once-blocked CZX tensor]
    \label{thm:mpug_czx_exact_tensor_injective}
    \lean{MPOTensor.CZX.tensor_isInjective}
    \leanok
    \discussion{7354}
    \uses{def:mpug_czx_exact_tensor, def:injective}
    As in \cite[lines~4505--4547]{FrancoRubio2025MPUGauging}, the
    once-blocked CZX tensor $B$ of Definition~\ref{def:mpug_czx_exact_tensor},
    with the unnormalized Hadamard weights, both output $Z$ gates, and
    physical indexing $B^{2a+b,2c+d}=B^{ab,cd}$ for binary $a,b,c,d$, is
    injective:
    \begin{align}
        \operatorname{span}_{\C}\{B^{i,j}:0\leq i,j<4\}
         & =\operatorname{Mat}_2(\C).
    \end{align}
    % Source: eq:MPU_CZX. This proves only injectivity of the shared tensor;
    % it does not complete the GHZ, fusion, or action-data scope of #7354.
\end{lemma}

\begin{proof}\leanok
    \uses{thm:mpug_czx_exact_tensor_unitary}
    Let $E_{ab}$ be the bond matrix units, with $a,b\in\{0,1\}$.
    The coordinate formula in
    Lemma~\ref{thm:mpug_czx_exact_tensor_unitary} gives
    \begin{align}
        E_{00}
         & =\tfrac12(B^{0,3}-B^{1,2}),  \\
        E_{01}
         & =\tfrac12(B^{0,3}+B^{1,2}),  \\
        E_{10}
         & =-\tfrac12(B^{2,1}+B^{3,0}), \\
        E_{11}
         & =\tfrac12(B^{3,0}-B^{2,1}).
    \end{align}
    Thus every bond matrix unit lies in the span of the physical letters.
    Since every $2\times2$ matrix is a linear combination of these units,
    the letters span the full bond matrix algebra.
\end{proof}

\subsection{The explicit CZX adjoint gauge}

\begin{definition}[The Pauli adjoint gauge]
    \label{def:mpug_czx_dagger_gauge}
    \lean{MPOTensor.CZX.daggerGauge}
    \leanok
    \discussion{7354}
    On the bond space $\C^2$, with basis indexed by $0,1$, the matrix on
    both virtual legs in the adjoint diagram of
    \cite[lines~4547--4559]{FrancoRubio2025MPUGauging} is the Pauli matrix
    \begin{align}
        Y & =\begin{pmatrix}0&-\mathrm{i}\\ \mathrm{i}&0\end{pmatrix}.
    \end{align}
\end{definition}

\begin{lemma}[Elementary identities for the Pauli adjoint gauge]
    \label{thm:mpug_czx_dagger_gauge_laws}
    \lean{MPOTensor.CZX.daggerGauge_conjTranspose,
        MPOTensor.CZX.daggerGauge_mul_self,
        MPOTensor.CZX.daggerGauge_mem_unitaryGroup}
    \leanok
    \uses{def:mpug_czx_dagger_gauge}
    As used for the unitary gauge in
    \cite[lines~4547--4559]{FrancoRubio2025MPUGauging}, the matrix $Y$ of
    Definition~\ref{def:mpug_czx_dagger_gauge} satisfies
    \begin{align}
        Y^\dagger & =Y,   \\
        Y^2       & =I_2,
    \end{align}
    and in particular $Y\in\mathrm{U}(2)$.
\end{lemma}

\begin{proof}\leanok
    \uses{def:mpug_czx_dagger_gauge}
    Conjugate-transposing the displayed matrix gives $Y^\dagger=Y$, and
    the Pauli multiplication rule $Y^2=I_2$ then yields $YY^\dagger=I_2$.
\end{proof}

\begin{lemma}[The physical adjoint of the exact CZX tensor]
    \label{thm:mpug_czx_physical_adjoint}
    \lean{MPOTensor.CZX.physicalAdjointTensor_tensor}
    \leanok
    \discussion{7354}
    \uses{def:mpug_czx_exact_tensor, def:mpug_czx_dagger_gauge,
        def:mpdo_physical_adjoint}
    Let $B$ be the exact once-blocked CZX tensor of
    Definition~\ref{def:mpug_czx_exact_tensor}, with physical adjoint
    $(B^\dagger)^{u,v}_{\ell r}=\overline{B^{v,u}_{\ell r}}$ in which the
    virtual indices are not transposed. As in the adjoint diagram of
    \cite[lines~4547--4559]{FrancoRubio2025MPUGauging} with no scalar
    prefactor, for $0\leq u,v<4$,
    \begin{align}
        (B^\dagger)^{u,v} & =YB^{u,v}Y.
    \end{align}
\end{lemma}

\begin{proof}\leanok
    \uses{thm:mpug_czx_exact_tensor_unitary,
        def:mpug_czx_dagger_gauge, def:mpdo_physical_adjoint}
    Inserting the coordinate formula from
    Lemma~\ref{thm:mpug_czx_exact_tensor_unitary} on both sides, expanding
    the two matrix products over the binary bond indices, and substituting
    the four physical values and the two bond values verifies each entry.
\end{proof}

\begin{lemma}[An explicit unitary adjoint gauge for CZX]
    \label{thm:mpug_czx_unitary_dagger_gauge}
    \lean{MPOTensor.CZX.daggerGaugeUnitary,
        MPOTensor.CZX.physicalAdjointTensor_tensor_eq_unitaryGauge}
    \leanok
    \uses{def:mpug_czx_exact_tensor, def:mpug_czx_dagger_gauge,
        def:mpdo_physical_adjoint}
    In the unitary gauge convention of
    \cite[lines~1552--1557 and 4547--4559]{FrancoRubio2025MPUGauging}, for
    the exact once-blocked CZX tensor $B$, the matrix $T=Y$, regarded as an
    element of $\mathrm{U}(2)$, satisfies, for $0\leq u,v<4$,
    \begin{align}
        (B^\dagger)^{u,v} & =T^\dagger B^{u,v}T.
    \end{align}
    % This exhibits the explicit witness, not an identification with the
    % choice-selected generic dagger-inverse gauge.
\end{lemma}

\begin{proof}\leanok
    \uses{thm:mpug_czx_dagger_gauge_laws, thm:mpug_czx_physical_adjoint}
    Lemma~\ref{thm:mpug_czx_dagger_gauge_laws} gives $Y\in\mathrm{U}(2)$
    and $Y^\dagger=Y$, and substituting these identities into
    Lemma~\ref{thm:mpug_czx_physical_adjoint} yields the claim.
\end{proof}

\begin{lemma}[The CZX adjoint-gauge sign]
    \label{thm:mpug_czx_dagger_sign}
    \lean{MPOTensor.CZX.daggerGauge_map_star,
        MPOTensor.CZX.daggerGauge_mul_map_star}
    \leanok
    \discussion{7354}
    \uses{def:mpug_czx_dagger_gauge}
    For the Pauli adjoint gauge $Y$, entrywise complex conjugation gives
    \begin{align}
        \overline{Y}  & =-Y,   \\
        Y\overline{Y} & =-I_2.
    \end{align}
    Hence, as stated in \cite[lines~4547--4559]{FrancoRubio2025MPUGauging},
    the sign in $Y\overline{Y}=\sigma_g I_2$ is $\sigma_g=-1$.
    % These entries do not determine fusion tensors, local action tensors,
    % L-symbols, or a gauge selected by the generic existence theorem.
\end{lemma}

\begin{proof}\leanok
    \uses{def:mpug_czx_dagger_gauge, thm:mpug_czx_dagger_gauge_laws}
    Conjugating each entry of the displayed Pauli matrix gives
    $\overline{Y}=-Y$. Then $Y\overline{Y}=-Y^2=-I_2$ by
    Lemma~\ref{thm:mpug_czx_dagger_gauge_laws}.
\end{proof}

\subsection{The displayed CZX source factors}

\begin{definition}[Four explicit CZX source factors]
    \label{def:mpug_czx_displayed_factors}
    \lean{MPOTensor.CZX.displayedX₂, MPOTensor.CZX.displayedY₂,
        MPOTensor.CZX.displayedX₁, MPOTensor.CZX.displayedY₁}
    \leanok
    \discussion{7354}
    \uses{def:mpug_czx_dagger_gauge}
    For $0\leq i<4$, write $i=2a(i)+b(i)$ with $a(i),b(i)\in\{0,1\}$,
    and put $\overline i=3-i$. Let $Y$ be the Pauli matrix of
    Definition~\ref{def:mpug_czx_dagger_gauge}. Define matrices
    $X_2\in\operatorname{Mat}_{(2\times4)\times4}(\C)$ and
    $Y_2\in\operatorname{Mat}_{4\times(4\times2)}(\C)$ by
    \begin{align}
        (X_2)_{(\alpha,i),k}
         & =\delta_{i,k}\delta_{\alpha,a(i)}(-1)^{a(i)+b(i)}, \\
        (Y_2)_{k,(j,\beta)}
         & =\delta_{k,\overline j}(-1)^{a(k)b(k)+b(k)\beta},
    \end{align}
    where $0\leq i,j,k<4$ and $\alpha,\beta\in\{0,1\}$. Define
    $X_1\in\operatorname{Mat}_{(4\times2)\times4}(\C)$ and
    $Y_1\in\operatorname{Mat}_{4\times(2\times4)}(\C)$ by
    \begin{align}
        (X_1)_{(i,\beta),k}
         & =\sum_{t=0}^1Y_{\beta t}\overline{(Y_2)_{k,(i,t)}}, \\
        (Y_1)_{k,(\alpha,j)}
         & =\sum_{t=0}^1\overline{(X_2)_{(t,j),k}}Y_{t\alpha}.
    \end{align}
    Product indices in these matrix spaces retain their displayed order.
    As in the filled and empty factors of
    \cite[lines~4559--4659]{FrancoRubio2025MPUGauging}, the weights
    $(-1)^{ab}$ are unnormalized Hadamard weights, $(-1)^{a(i)+b(i)}$
    includes both output $Z$ gates, and the Pauli entry in the formula for
    $X_1$ is $Y_{\beta t}$ rather than $Y_{t\beta}$.
\end{definition}

\begin{lemma}[The displayed factors recover both CZX cuts]
    \label{thm:mpug_czx_displayed_cuts}
    \lean{MPOTensor.CZX.displayed_cuts}
    \leanok
    \uses{def:mpug_czx_displayed_factors, def:mpug_czx_exact_tensor,
        def:mpu_source_matrix_one, def:mpu_source_matrix_two}
    As in the two decompositions of
    \cite[lines~4559--4659]{FrancoRubio2025MPUGauging}, for the exact
    once-blocked CZX tensor $B$ with source matrices
    $(M_1(B))_{(i,\beta),(\alpha,j)}=B^{i,j}_{\alpha\beta}$ and
    $(M_2(B))_{(\alpha,i),(j,\beta)}=B^{i,j}_{\alpha\beta}$, the factors of
    Definition~\ref{def:mpug_czx_displayed_factors} satisfy
    \begin{align}
        X_1Y_1 & =M_1(B), \\
        X_2Y_2 & =M_2(B).
    \end{align}
\end{lemma}

\begin{proof}\leanok
    \uses{def:mpug_czx_displayed_factors, def:mpug_czx_dagger_gauge,
        thm:mpug_czx_exact_tensor_unitary,
        def:mpu_source_matrix_one, def:mpu_source_matrix_two}
    Expanding each matrix product over its four-dimensional factor index,
    substituting the displayed formulas and the coordinate formula from
    Lemma~\ref{thm:mpug_czx_exact_tensor_unitary}, and evaluating the
    binary bond indices and the four physical values proves both entry
    identities.
\end{proof}

\begin{lemma}[Four Gram matrices of the displayed CZX factors]
    \label{thm:mpug_czx_displayed_grams}
    \lean{MPOTensor.CZX.displayed_grams}
    \leanok
    \uses{def:mpug_czx_displayed_factors}
    With the factors of two retaining the unnormalized Hadamard weights of
    \cite[lines~4559--4659]{FrancoRubio2025MPUGauging}, the four explicit
    factors of Definition~\ref{def:mpug_czx_displayed_factors} satisfy
    \begin{align}
        X_1^\dagger X_1 & =2I_4,
        \label{eq:mpug_czx_displayed_gram_x1} \\
        X_2^\dagger X_2 & =I_4,
        \label{eq:mpug_czx_displayed_gram_x2} \\
        Y_1Y_1^\dagger  & =I_4,               \\
        Y_2Y_2^\dagger  & =2I_4.
    \end{align}
\end{lemma}

\begin{proof}\leanok
    \uses{def:mpug_czx_displayed_factors, thm:mpug_czx_dagger_gauge_laws}
    The empty factors are Pauli-gauge contractions of the filled factors,
    $X_1=(I_4\otimes Y)Y_2^\dagger$ and $Y_1=X_2^\dagger(Y\otimes I_4)$.
    By Lemma~\ref{thm:mpug_czx_dagger_gauge_laws}, $Y$ is unitary, hence so
    are $I_4\otimes Y$ and $Y\otimes I_4$; composing a unitary matrix with
    $Y_2^\dagger$, respectively with $X_2^\dagger$, leaves the Gram matrix
    unchanged, so $X_1^\dagger X_1=Y_2Y_2^\dagger$ and
    $Y_1Y_1^\dagger=X_2^\dagger X_2$. Expanding the two filled
    Gram matrices $X_2^\dagger X_2$ and $Y_2Y_2^\dagger$ over their
    product-index spaces and substituting the factor coordinates at each
    pair of factor indices gives the two displayed scalar identities.
\end{proof}

\begin{definition}[Right inverses of the displayed CZX factors]
    \label{def:mpug_czx_displayed_right_inverses}
    \lean{MPOTensor.CZX.displayedZ₁, MPOTensor.CZX.displayedZ₂}
    \leanok
    \uses{def:mpug_czx_displayed_factors}
    For the factors of Definition~\ref{def:mpug_czx_displayed_factors}, the
    right inverses associated with the explicit decompositions in
    \cite[lines~4559--4659]{FrancoRubio2025MPUGauging} are
    \begin{align}
        Z_1 & =Y_1^\dagger,         \\
        Z_2 & =\tfrac12Y_2^\dagger.
    \end{align}
\end{definition}

\begin{lemma}[Right inverses and weighted normalization]
    \label{thm:mpug_czx_displayed_normalizations}
    \lean{MPOTensor.CZX.displayed_normalizations}
    \leanok
    \uses{def:mpug_czx_displayed_factors,
        def:mpug_czx_displayed_right_inverses, def:mpu_source_weight}
    Set $\rho=I_2/2$. The displayed factors and right inverses satisfy
    \begin{align}
        Y_1Z_1                         & =I_4,
        \label{eq:mpug_czx_displayed_right_inverse_y1} \\
        Y_2Z_2                         & =I_4,
        \label{eq:mpug_czx_displayed_right_inverse_y2} \\
        X_1^\dagger(I_4\otimes\rho)X_1 & =I_4.
    \end{align}
\end{lemma}

\begin{proof}\leanok
    \uses{thm:mpug_czx_displayed_grams,
        def:mpug_czx_displayed_right_inverses, def:mpu_source_weight}
    The two inverse identities follow directly from
    Lemma~\ref{thm:mpug_czx_displayed_grams}. Since
    $I_4\otimes\rho=I_8/2$, the weighted identity is
    $\tfrac12X_1^\dagger X_1=I_4$.
\end{proof}

\begin{lemma}[Both CZX source-cut ranks are four]
    \label{thm:mpug_czx_displayed_ranks}
    \lean{MPOTensor.CZX.displayed_ranks}
    \leanok
    \uses{def:mpug_czx_exact_tensor, def:mpu_source_matrix_one,
        def:mpu_source_matrix_two}
    For the exact once-blocked CZX tensor $B$,
    \begin{align}
        d_r=\operatorname{rank}M_1(B) & =4, \\
        d_l=\operatorname{rank}M_2(B) & =4.
    \end{align}
    Thus the four-dimensional legs in
    \cite[lines~4559--4659]{FrancoRubio2025MPUGauging} have the
    source-cut dimensions.
\end{lemma}

\begin{proof}\leanok
    \uses{thm:mpug_czx_displayed_cuts, thm:mpug_czx_displayed_grams,
        thm:mpug_czx_displayed_normalizations}
    Each cut factors through a four-dimensional space, so its rank is at
    most four. Setting $L_1=X_1^\dagger/2$ for the first cut and
    $L_2=X_2^\dagger$ for the second, (\ref{eq:mpug_czx_displayed_gram_x1})
    and (\ref{eq:mpug_czx_displayed_gram_x2}) give $L_iX_i=I_4$, while
    (\ref{eq:mpug_czx_displayed_right_inverse_y1}) and
    (\ref{eq:mpug_czx_displayed_right_inverse_y2}) give $Y_iZ_i=I_4$, hence
    $L_i(X_iY_i)Z_i=I_4$. Multiplication cannot increase rank, so each cut
    has rank at least four.
\end{proof}

\begin{definition}[The displayed factors in source-rank coordinates]
    \label{def:mpug_czx_displayed_source_factors}
    \lean{MPOTensor.CZX.displayedRightEquiv, MPOTensor.CZX.displayedLeftEquiv,
        MPOTensor.CZX.displayedSourceFactors,
        MPOTensor.CZX.displayedSourceFactors_coordinates,
        MPOTensor.CZX.displayedSourceFactors_inverse_coordinates}
    \leanok
    \uses{def:mpug_czx_displayed_factors,
        def:mpug_czx_displayed_right_inverses, thm:mpug_czx_displayed_ranks,
        def:mpug_czx_exact_tensor, def:mpu_weighted_source_factors,
        def:mpug_czx_dagger_gauge}
    Let $d_r,d_l$ be the two source-cut ranks of the exact CZX tensor $B$.
    By Lemma~\ref{thm:mpug_czx_displayed_ranks}, both are four. Let
    \begin{align}
        e_r & :\{0,\ldots,d_r-1\}\longrightarrow\{0,1,2,3\}, \\
        e_l & :\{0,\ldots,d_l-1\}\longrightarrow\{0,1,2,3\}
    \end{align}
    be the bijections that preserve each integer index. Pulling back the
    four-dimensional factor index of $X_1,Y_1,Z_1$ along $e_r$ and that
    of $X_2,Y_2,Z_2$ along $e_l$ yields matrices
    $\widehat X_i,\widehat Y_i,\widehat Z_i$. Together with $\rho=I_2/2$,
    these give source factors for $B$ in the sense of
    Definition~\ref{def:mpu_weighted_source_factors}, satisfying both cut
    factorizations, the two left-factor normalizations, and both
    right-inverse identities.

    Explicitly, for $0\leq k,i<4$ and $\alpha,\beta\in\{0,1\}$, evaluating
    the pullbacks gives
    \begin{align}
        (\widehat X_1)_{(i,\alpha),e_r^{-1}(k)} & =(X_1)_{(i,\alpha),k}, \\
        (\widehat Y_1)_{e_r^{-1}(k),(\alpha,i)} & =(Y_1)_{k,(\alpha,i)}, \\
        (\widehat Z_1)_{(\alpha,i),e_r^{-1}(k)} & =(Z_1)_{(\alpha,i),k}, \\
        (\widehat X_2)_{(\alpha,i),e_l^{-1}(k)} & =(X_2)_{(\alpha,i),k}, \\
        (\widehat Y_2)_{e_l^{-1}(k),(i,\alpha)} & =(Y_2)_{k,(i,\alpha)}, \\
        (\widehat Z_2)_{(i,\alpha),e_l^{-1}(k)} & =(Z_2)_{(i,\alpha),k},
    \end{align}
    and substituting the definitions of $X_1,Y_1,Z_1,Z_2$ gives the
    Pauli and inverse coordinates
    \begin{align}
        (\widehat X_1)_{(i,\beta),e_r^{-1}(k)}
         & =\sum_{t=0}^1Y_{\beta t}
        \overline{(\widehat Y_2)_{e_l^{-1}(k),(i,t)}},                           \\
        (\widehat Y_1)_{e_r^{-1}(k),(\beta,i)}
         & =\sum_{t=0}^1\overline{(\widehat X_2)_{(t,i),e_l^{-1}(k)}}Y_{t\beta}, \\
        (\widehat Z_1)_{(\beta,i),e_r^{-1}(k)}
         & =\overline{(\widehat Y_1)_{e_r^{-1}(k),(\beta,i)}},                   \\
        (\widehat Z_2)_{(i,\beta),e_l^{-1}(k)}
         & =\tfrac12\overline{(\widehat Y_2)_{e_l^{-1}(k),(i,\beta)}}.
    \end{align}
    In particular, as in the first factor of
    \cite[lines~4611--4659]{FrancoRubio2025MPUGauging}, pulling back
    preserves the Pauli orientation $Y_{\beta t}$.
    % This is an explicit SourceFactors witness, not an equality with the
    % choice-selected compact SVD factors. It supplies no canonical-form-II
    % presentation, fusion tensors, local action tensors, or anomaly data.
\end{definition}

\begin{proof}\leanok
    \uses{thm:mpug_czx_displayed_cuts, thm:mpug_czx_displayed_grams,
        thm:mpug_czx_displayed_normalizations, def:mpug_czx_displayed_factors,
        def:mpug_czx_displayed_right_inverses}
    Pulling back the cut factorizations and inverse identities along the
    two index bijections preserves products, conjugate transposes, and
    identities, so the weighted first normalization and the ordinary
    second normalization hold in the source-rank coordinates.
    Evaluating the pullbacks, using $e_r(e_r^{-1}(k))=k$ and
    $e_l(e_l^{-1}(k))=k$, gives the coordinate identities; substituting
    the definitions of $X_1,Y_1,Z_1,Z_2$ and cancelling the inverse
    bijections gives the Pauli and inverse coordinates.
\end{proof}

\subsection{Normalization of the displayed CZX comparison}

\begin{definition}[Reflected candidates and the two displayed comparisons]
    \label{def:mpug_czx_displayed_comparison}
    \lean{MPOTensor.CZX.displayedReflectedX₁, MPOTensor.CZX.displayedReflectedY₁,
        MPOTensor.CZX.displayedRawComparison,
        MPOTensor.CZX.displayedWeightedComparison}
    \leanok
    \uses{def:mpug_czx_displayed_factors, def:mpug_czx_dagger_gauge,
        def:mpu_source_weight}
    Let $Y$ be the Pauli adjoint gauge of
    Definition~\ref{def:mpug_czx_dagger_gauge} and let
    $X_1,Y_1,X_2,Y_2$ be the displayed factors of
    Definition~\ref{def:mpug_czx_displayed_factors}, with $d=4$ and $D=2$.
    Put $A=I_4\otimes\overline Y$ and $B=Y^{\mathsf T}\otimes I_4$, and
    define the reflected candidates
    \begin{align}
        \widehat X_1 & =AY_2^\dagger,  \\
        \widehat Y_1 & =X_2^\dagger B.
    \end{align}
    Evaluating the candidates and comparison of
    \cite[lines~5432--5441]{FrancoRubio2025MPUGauging} on the explicit
    factors of lines~4559--4659 with $\rho=I_2/2$ defines the unweighted
    and the density-weighted comparisons
    \begin{align}
        K_{\mathrm{unw}} & =X_1^\dagger\widehat X_1,                 \\
        K_\rho           & =X_1^\dagger(I_4\otimes\rho)\widehat X_1.
    \end{align}
\end{definition}

\begin{lemma}[The reflected candidates are the negatives of the displayed factors]
    \label{thm:mpug_czx_reflected_neg}
    \lean{MPOTensor.CZX.displayedReflected_eq_neg}
    \leanok
    \uses{def:mpug_czx_displayed_comparison, def:mpug_czx_displayed_factors}
    For the maintained Pauli gauge,
    \begin{align}
        \widehat X_1 & =-X_1, \\
        \widehat Y_1 & =-Y_1.
    \end{align}
\end{lemma}
\begin{proof}\leanok
    \uses{thm:mpug_czx_dagger_gauge_laws, thm:mpug_czx_dagger_sign,
        def:mpug_czx_displayed_factors, def:mpug_czx_displayed_comparison}
    Lemma~\ref{thm:mpug_czx_dagger_sign} gives $\overline Y=-Y$, and
    $Y^{\mathsf T}=\overline{Y^\dagger}=\overline Y=-Y$ by
    Lemma~\ref{thm:mpug_czx_dagger_gauge_laws}. Expanding both entries
    of the displayed factors against these two gauges gives the sign.
\end{proof}

\begin{lemma}[Coordinates of the two comparisons and the right overlap]
    \label{thm:mpug_czx_comparison_coordinates}
    \lean{MPOTensor.CZX.displayedComparison_coordinates,
        MPOTensor.CZX.displayedRawComparison_gram,
        MPOTensor.CZX.displayedRawComparison_not_isUnitaryBetween,
        MPOTensor.CZX.displayedRawComparison_ne_rightOverlap}
    \leanok
    \uses{def:mpug_czx_displayed_comparison}
    On the displayed factors,
    \begin{align}
        K_{\mathrm{unw}}        & =-2I_4,
        \label{eq:mpug_czx_raw_comparison}      \\
        K_\rho                  & =-I_4,
        \label{eq:mpug_czx_weighted_comparison} \\
        Y_1\widehat Y_1^\dagger & =-I_4.
        \label{eq:mpug_czx_right_overlap}
    \end{align}
    Consequently $K_{\mathrm{unw}}^\dagger K_{\mathrm{unw}}=4I_4$, so
    $K_{\mathrm{unw}}$ is not unitary, and
    $K_{\mathrm{unw}}\neq Y_1\widehat Y_1^\dagger$: the unweighted left
    overlap and the right overlap of
    \cite[lines~5438--5441]{FrancoRubio2025MPUGauging} disagree on these
    factors.
\end{lemma}
\begin{proof}\leanok
    \uses{thm:mpug_czx_reflected_neg, thm:mpug_czx_displayed_grams,
        thm:mpug_czx_displayed_normalizations,
        def:mpug_czx_displayed_comparison}
    Substituting Lemma~\ref{thm:mpug_czx_reflected_neg} turns each
    overlap into the negative of a Gram matrix of the displayed factors.
    The unweighted Gram $X_1^\dagger X_1=2I_4$ of
    Lemma~\ref{thm:mpug_czx_displayed_grams} gives
    (\ref{eq:mpug_czx_raw_comparison}); the weighted normalization
    $X_1^\dagger(I_4\otimes\rho)X_1=I_4$ of
    Lemma~\ref{thm:mpug_czx_displayed_normalizations} gives
    (\ref{eq:mpug_czx_weighted_comparison}); and $Y_1Y_1^\dagger=I_4$ gives
    (\ref{eq:mpug_czx_right_overlap}). Squaring
    (\ref{eq:mpug_czx_raw_comparison}) gives the Gram matrix $4I_4$, whose
    diagonal entry $4$ differs from $1$, and comparing the diagonal entries
    of (\ref{eq:mpug_czx_raw_comparison}) and
    (\ref{eq:mpug_czx_right_overlap}) separates the two overlaps.
\end{proof}

\begin{lemma}[The weighted comparison is unitary and recovers both factors]
    \label{thm:mpug_czx_weighted_comparison_spec}
    \lean{MPOTensor.CZX.displayedWeightedComparison_spec}
    \leanok
    \uses{def:mpug_czx_displayed_comparison}
    The weighted comparison $K_\rho$ is unitary, and
    \begin{align}
        X_1 & =\widehat X_1K_\rho^\dagger, \\
        Y_1 & =K_\rho\widehat Y_1.
    \end{align}
    The weight is the one carried by the normalization of the first left
    factor. This calculation concerns the explicit displayed factors of
    Definition~\ref{def:mpug_czx_displayed_factors}, not choice-selected
    factors, so the unequal overlaps of
    Lemma~\ref{thm:mpug_czx_comparison_coordinates} diagnose the
    intermediate formula of
    \cite[lines~5438--5441]{FrancoRubio2025MPUGauging} on these factors
    rather than contradicting the proposition being proved there. The
    weighted comparison is the one retained in
    \cite{gap:fbc25_inverse_compatible_tilde_omission}; no gauge phase is
    changed.
\end{lemma}
\begin{proof}\leanok
    \uses{thm:mpug_czx_comparison_coordinates, thm:mpug_czx_reflected_neg}
    By (\ref{eq:mpug_czx_weighted_comparison}), $K_\rho=-I_4$, which is
    unitary. Lemma~\ref{thm:mpug_czx_reflected_neg} then gives
    $\widehat X_1K_\rho^\dagger=(-X_1)(-I_4)=X_1$ and
    $K_\rho\widehat Y_1=(-I_4)(-Y_1)=Y_1$.
\end{proof}

\subsection{The blocked GHZ state and its CZX invariance}

\begin{definition}[The GHZ tensor in the blocked CZX alphabet]
    \label{def:mpug_czx_blocked_ghz}
    \lean{MPOTensor.CZX.blockedGHZTensor}
    \leanok
    \discussion{7354}
    \uses{def:ghz_tensor, def:blocked_tensor}
    Let $A^0=\ketbra{0}{0}$ and $A^1=\ketbra{1}{1}$ be the GHZ tensor,
    and identify the blocked physical alphabet with $\{0,1,2,3\}$ by
    $\varphi(2a+b)=(a,b)$ for $a,b\in\{0,1\}$. As in
    \cite[lines~4660--4670]{FrancoRubio2025MPUGauging}, with the same
    physical-index order as the CZX tensor $B$, the once-blocked GHZ
    tensor is
    \begin{align}
        \widehat A^{2a+b} & =A^aA^b=(A^{[2]})^{\varphi(2a+b)}.
    \end{align}
\end{definition}

\begin{lemma}[The two blocked GHZ sectors]
    \label{thm:mpug_czx_blocked_ghz_sectors}
    \lean{MPOTensor.CZX.blockTensor_ghzSectorTensor_apply}
    \leanok
    \uses{def:mpug_ghz_blocked_cluster_tensor, def:blocked_tensor}
    For $x\in\{0,1\}$, let $A_x^a=\delta_{a,x}$ be the
    bond-one GHZ sector. Under $\varphi(2a+b)=(a,b)$ its two-site block is
    \begin{align}
        (A_x^{[2]})^{\varphi(i)} & =\delta_{i,3x},
        \qquad i\in\{0,1,2,3\}.
    \end{align}
    Thus, as in \cite[lines~4660--4670]{FrancoRubio2025MPUGauging}, the
    blocked sector vectors are $\ket{00}$ and $\ket{11}$.
\end{lemma}

\begin{proof}\leanok
    For $i=2a+b$, the block coefficient is
    $A_x^aA_x^b=\delta_{a,x}\delta_{b,x}$, which equals one exactly when
    $a=b=x$, or equivalently $i=3x$.
\end{proof}

\begin{lemma}[The matrix product vector of the blocked GHZ tensor]
    \label{thm:mpug_czx_blocked_ghz_mpv}
    \lean{MPOTensor.CZX.mpv_blockedGHZTensor}
    \leanok
    \uses{def:mpug_czx_blocked_ghz, def:mpv}
    For every $N\geq0$, the periodic matrix product vector of
    $\widehat A$ is
    \begin{align}
        \ket{\psi_N}
         & =\ket{0}^{\otimes N}+\ket{3}^{\otimes N}.
        \label{eq:mpug_czx_blocked_ghz_mpv}
    \end{align}
    In the original two-qubit alphabet the two summands are
    $\ket{00}^{\otimes N}$ and $\ket{11}^{\otimes N}$.
    At $N=0$ they are the same empty configuration, with total coefficient
    two.
\end{lemma}

\begin{proof}\leanok
    The blocked letters are diagonal, with
    $\widehat A^i=\operatorname{diag}(\delta_{i,0},\delta_{i,3})$.
    Hence, for $s\in\{0,1,2,3\}^N$,
    \begin{align}
        \operatorname{tr}\bigl(\widehat A^{s_1}\cdots
        \widehat A^{s_N}\bigr)
         & =\prod_{j=1}^N\delta_{s_j,0}
        +\prod_{j=1}^N\delta_{s_j,3}.
    \end{align}
    These are the coefficients of the two constant configurations.
\end{proof}

\begin{lemma}[CZX invariance of the blocked GHZ state]
    \label{thm:mpug_czx_blocked_ghz_invariant}
    \lean{MPOTensor.CZX.mpo_tensor_mulVec_blockedGHZ}
    \leanok
    \discussion{7354}
    \uses{def:mpug_czx_exact_tensor, def:mpug_czx_blocked_ghz, def:mpv}
    Let $B$ be the once-blocked CZX tensor, with both output $Z$ gates and
    the unnormalized Hadamard weights, and let $N>0$.
    Reindex the matrix product vector of the two-site-blocked GHZ tensor
    $A^{[2]}$ by the alphabet bijection $\varphi(2a+b)=(a,b)$:
    \begin{align}
        \psi_N(s) & =V^{(N)}(A^{[2]})_{(\varphi(s_1),\ldots,\varphi(s_N))},
        \qquad s\in\{0,1,2,3\}^N.
    \end{align}
    Then, as in the GHZ-invariance assertion of
    \cite[lines~4660--4670]{FrancoRubio2025MPUGauging}, the periodic CZX
    operator satisfies
    \begin{align}
        O_N(B)\ket{\psi_N} & =\ket{\psi_N}.
    \end{align}
    % Only the actual blocked GHZ MPV and its positive-length invariance
    % are asserted. Fusion, action tensors, and classification remain open.
\end{lemma}

\begin{proof}\leanok
    \uses{thm:mpug_czx_blocked_ghz_mpv, thm:mpug_czx_exact_tensor_unitary}
    Reindexing the alphabet gives the vector in
    (\ref{eq:mpug_czx_blocked_ghz_mpv}). The CZX operator complements
    every matter bit. Its phase on either constant configuration is one:
    the cyclic exponent is zero on $0^N$ and $4N$ on $3^N$.
    Consequently
    \begin{align}
        O_N(B)\ket{0}^{\otimes N} & =\ket{3}^{\otimes N}, &
        O_N(B)\ket{3}^{\otimes N} & =\ket{0}^{\otimes N}.
    \end{align}
    Their sum is invariant.
\end{proof}

\subsection{The printed CZX action pair and its empty-word boundary}

The action tensors displayed in
\cite[lines~4690--4694]{FrancoRubio2025MPUGauging} satisfy the interior
identity for nonempty words and the exterior identity with nonempty buffers.
Their empty contraction is one half, whereas the general interior identity
in \cite[lines~1841--1870]{FrancoRubio2025MPUGauging} includes length zero.
The distinction and a local replacement are recorded in
\cite{gap:fbc25_czx_action_empty_word}. The replacement below is a different
pair, not a scalar gauge of the printed pair or a substitution in the
unbuffered defect formulas.

\begin{definition}[The acted GHZ sectors]
    \label{def:mpug_czx_acted_sectors}
    \lean{MPOTensor.CZX.blockedSector, MPOTensor.CZX.actedSector,
        MPOTensor.CZX.actionTarget, MPOTensor.CZX.actionLetter,
        MPOTensor.CZX.actionMatrix}
    \leanok
    \uses{def:mpug_czx_exact_tensor, def:mpug_ghz_blocked_cluster_tensor,
        def:blocked_tensor}
    For $x\in\{0,1\}$, let $A_x$ be the two-site block of the
    bond-one GHZ sector, with physical indices identified by
    $\varphi(2a+b)=(a,b)$. Thus $A_x^i=\delta_{i,3x}$.
    Write $B^{i,j}$ for the CZX tensor and define
    \begin{align}
        B_x^i & =\sum_{j=0}^3 B^{i,j}A_x^j, \\
        k_x   & =3(1-x),                    \\
        C_x   & =B^{k_x,3x}.
        \label{eq:mpug_czx_acted_sectors}
    \end{align}
    As in the physical contraction and exchanged GHZ blocks of
    \cite[lines~4660--4694]{FrancoRubio2025MPUGauging}, the target sector
    is $A_{1-x}$.
\end{definition}

\begin{lemma}[Coordinates of the acted sectors]
    \label{thm:mpug_czx_acted_sector_coordinates}
    \lean{MPOTensor.CZX.actedSector_apply, MPOTensor.CZX.actionTarget_apply,
        MPOTensor.CZX.actionMatrix_coordinates}
    \leanok
    \uses{def:mpug_czx_acted_sectors}
    For $x\in\{0,1\}$ and $i\in\{0,1,2,3\}$,
    \begin{align}
        B_x^i     & =\delta_{i,k_x}C_x,                     \\
        A_{1-x}^i & =\delta_{i,k_x},                        \\
        C_0       & =\begin{pmatrix}0&0\\-1&1\end{pmatrix}, \\
        C_1       & =\begin{pmatrix}1&1\\0&0\end{pmatrix}.
        \label{eq:mpug_czx_action_matrices}
    \end{align}
\end{lemma}
\begin{proof}\leanok
    \uses{def:mpug_czx_acted_sectors, thm:mpug_czx_blocked_ghz_sectors,
        def:mpug_czx_exact_tensor}
    Since $A_x^j=\delta_{j,3x}$, the defining sum of
    Definition~\ref{def:mpug_czx_acted_sectors} collapses to
    $B_x^i=\sum_{j=0}^3 B^{i,j}A_x^j=B^{i,3x}$. The exact CZX tensor of
    Definition~\ref{def:mpug_czx_exact_tensor} vanishes off its support,
    giving
    \begin{align}
        B^{i,3x} & =\delta_{i,3(1-x)}B^{3(1-x),3x}=\delta_{i,k_x}C_x,
    \end{align}
    so $B_x^i=\delta_{i,k_x}C_x$; likewise $A_x^j=\delta_{j,3x}$ at
    $x\mapsto1-x$ gives $A_{1-x}^i=\delta_{i,k_x}$. Expanding the defining
    CZX matrices at $B^{k_x,3x}$ evaluates the two displayed matrices
    $C_0,C_1$.
\end{proof}

\begin{definition}[The printed action pair]
    \label{def:mpug_czx_printed_action_pair}
    \lean{MPOTensor.CZX.printedActionBra, MPOTensor.CZX.printedActionKet}
    \leanok
    For $x\in\{0,1\}$, denote the printed left and right action tensors
    by $p_x$ and $q_x$. With $|\pm\rangle=(|0\rangle\pm|1\rangle)/\sqrt2$,
    the formulas of \cite[lines~4690--4694]{FrancoRubio2025MPUGauging} read
    \begin{align}
        p_0 & =(0,-i),                                    \\
        p_1 & =(1,0),                                     \\
        q_0 & =\tfrac12\begin{pmatrix}-i\\i\end{pmatrix}, \\
        q_1 & =\tfrac12\begin{pmatrix}1\\1\end{pmatrix}.
        \label{eq:mpug_czx_printed_action_pair}
    \end{align}
    The normalization of the kets is not specified explicitly in that
    passage. These coordinates agree with the later defect evaluations;
    see \cite{gap:fbc25_czx_action_empty_word}.
\end{definition}

\begin{lemma}[Finite contractions of the printed action pair]
    \label{thm:mpug_czx_printed_action_contractions}
    \lean{MPOTensor.CZX.printedAction_finite_identities}
    \leanok
    \uses{def:mpug_czx_acted_sectors, def:mpug_czx_printed_action_pair}
    For each $x\in\{0,1\}$,
    \begin{align}
        C_x^2        & =C_x,      \\
        p_xq_x       & =\tfrac12, \\
        p_xC_xq_x    & =1,        \\
        C_xq_xp_xC_x & =C_x.
        \label{eq:mpug_czx_printed_action_contractions}
    \end{align}
\end{lemma}
\begin{proof}\leanok
    \uses{thm:mpug_czx_acted_sector_coordinates, def:mpug_czx_printed_action_pair}
    Multiplying the two-by-two matrices and the printed rows and columns
    yields the four identities.
\end{proof}

\begin{lemma}[Nonempty acted and target words]
    \label{thm:mpug_czx_action_words}
    \lean{MPOTensor.CZX.evalWord_actedSector, MPOTensor.CZX.evalWord_actionTarget}
    \leanok
    \uses{def:mpug_czx_acted_sectors}
    For a word $\mathbf a$ in $\{0,1,2,3\}$, write $B_x^{\mathbf a}$
    and $A_{1-x}^{\mathbf a}$ for ordered products of letters. If
    $|\mathbf a|>0$, then
    \begin{align}
        B_x^{\mathbf a}     & =
        \begin{cases}C_x&\text{if every letter is }k_x,\\0&\text{otherwise},\end{cases} \\
        A_{1-x}^{\mathbf a} & =
        \begin{cases}1&\text{if every letter is }k_x,\\0&\text{otherwise}.\end{cases}
        \label{eq:mpug_czx_action_words}
    \end{align}
\end{lemma}
\begin{proof}\leanok
    \uses{thm:mpug_czx_acted_sector_coordinates, thm:mpug_czx_printed_action_contractions}
    An unsupported letter makes the product zero. Otherwise every acted
    letter is $C_x$, whose positive powers equal $C_x$ by idempotency.
\end{proof}

\begin{lemma}[Printed interior and buffered exterior identities]
    \label{thm:mpug_czx_printed_action_words}
    \lean{MPOTensor.CZX.printedAction_interior, MPOTensor.CZX.printedAction_exterior}
    \leanok
    \uses{def:mpug_czx_acted_sectors, def:mpug_czx_printed_action_pair}
    Let $x\in\{0,1\}$. For every nonempty word $\mathbf c$,
    \begin{align}
        p_xB_x^{\mathbf c}q_x & =A_{1-x}^{\mathbf c}.
        \label{eq:mpug_czx_printed_interior}
    \end{align}
    If $\mathbf a,\mathbf b$ are also nonempty, then
    \begin{align}
        B_x^{\mathbf a}q_xA_{1-x}^{\mathbf c}p_xB_x^{\mathbf b}
         & =B_x^{\mathbf a\mathbf c\mathbf b}.
        \label{eq:mpug_czx_printed_exterior}
    \end{align}
    No equality of buffer lengths is needed: while the equal-length case
    matches the exterior diagram of
    \cite[lines~1788--1838]{FrancoRubio2025MPUGauging}, the interior
    equation covers only the nonempty part of that source assertion; see
    \cite{gap:fbc25_czx_action_empty_word}.
\end{lemma}
\begin{proof}\leanok
    \uses{thm:mpug_czx_action_words, thm:mpug_czx_printed_action_contractions}
    If an unsupported letter occurs, both sides vanish. Otherwise the
    interior is $p_xC_xq_x=1$ and the exterior is $C_xq_xp_xC_x=C_x$.
\end{proof}

\begin{lemma}[The printed pair is not an all-word reduction]
    \label{thm:mpug_czx_printed_action_empty}
    \lean{MPOTensor.CZX.printedAction_empty_ne, MPOTensor.CZX.printedAction_not_isReduction}
    \leanok
    \uses{def:mpug_czx_acted_sectors, def:mpug_czx_printed_action_pair,
        def:mps_rectangular_reduction}
    As recorded in the empty-word obstruction of
    \cite{gap:fbc25_czx_action_empty_word}, for each $x\in\{0,1\}$, the
    printed pair $(p_x,q_x)$ is not a reduction from $B_x$ to $A_{1-x}$.
    Indeed, if $\varepsilon$ is the empty word, then
    \begin{align}
        p_xB_x^{\varepsilon}q_x & =\tfrac12\ne1=A_{1-x}^{\varepsilon}.
        \label{eq:mpug_czx_printed_action_empty}
    \end{align}
\end{lemma}
\begin{proof}\leanok
    \uses{thm:mpug_czx_printed_action_contractions}
    Empty products are identities, whereas $p_xq_x=1/2$.
\end{proof}

\begin{definition}[A separately dressed left action tensor]
    \label{def:mpug_czx_dressed_action}
    \lean{MPOTensor.CZX.dressedActionBra}
    \leanok
    \uses{def:mpug_czx_acted_sectors, def:mpug_czx_printed_action_pair}
    As the local replacement of \cite{gap:fbc25_czx_action_empty_word}
    rather than the printed pair, for $x\in\{0,1\}$, define $V_x=p_xC_x$
    and retain $W_x=q_x$. In coordinates $V_0=(i,-i)$ and
    $V_1=(1,1)$, so neither is a scalar multiple of the corresponding
    printed row. No replacement in an unbuffered defect formula is asserted.
\end{definition}

\begin{lemma}[Dressing is absorbed by a nonempty buffer]
    \label{thm:mpug_czx_dressed_action_absorption}
    \lean{MPOTensor.CZX.dressedActionBra_absorption}
    \leanok
    \uses{def:mpug_czx_dressed_action}
    For $x\in\{0,1\}$ and every nonempty word $\mathbf b$,
    \begin{align}
        V_xB_x^{\mathbf b} & =p_xB_x^{\mathbf b}.
        \label{eq:mpug_czx_dressed_action_absorption}
    \end{align}
\end{lemma}
\begin{proof}\leanok
    \uses{thm:mpug_czx_action_words, thm:mpug_czx_printed_action_contractions,
        def:mpug_czx_dressed_action}
    The word evaluates to either zero or $C_x$, and in the latter case
    $C_x^2=C_x$ gives the equality.
\end{proof}

\begin{lemma}[An all-word reduction with exterior buffer length one]
    \label{thm:mpug_czx_dressed_action_reduction}
    \lean{MPOTensor.CZX.dressedAction_isReduction,
        MPOTensor.CZX.dressedAction_isReductionExteriorBufferLength}
    \leanok
    \uses{def:mpug_czx_dressed_action, def:mps_rectangular_reduction,
        def:mps_reduction_exterior_buffer_length}
    For each $x\in\{0,1\}$, $(V_x,W_x)$ is a reduction from $B_x$
    to $A_{1-x}$ with exterior buffer length one. The reduction identity
    holds for every word $\mathbf c$, including the empty word; the exterior
    identity holds for nonempty words $\mathbf a$, $\mathbf b$, and
    $\mathbf c$. Thus
    \begin{align}
        V_xW_x                & =1,                                 \\
        V_xB_x^{\mathbf c}W_x & =A_{1-x}^{\mathbf c},               \\
        B_x^{\mathbf a}W_xA_{1-x}^{\mathbf c}V_xB_x^{\mathbf b}
                              & =B_x^{\mathbf a\mathbf c\mathbf b}.
        \label{eq:mpug_czx_dressed_action_reduction}
    \end{align}
    This proves the all-word relation for the separately dressed pair,
    not the empty-word claim for the printed pair.
\end{lemma}
\begin{proof}\leanok
    \uses{thm:mpug_czx_printed_action_contractions,
        thm:mpug_czx_printed_action_words, thm:mpug_czx_dressed_action_absorption,
        def:mpug_czx_dressed_action}
    The empty contraction is $V_xW_x=p_xC_xq_x=1$. For a nonempty
    interior word, absorbing $C_x$ and applying the printed interior
    identity yields the second equation, while in the exterior identity,
    absorbing $C_x$ into the right buffer and applying
    (\ref{eq:mpug_czx_printed_exterior}) completes the proof.
\end{proof}

\subsection{The CZX \texorpdfstring{$g,g$}{g,g} fusion caps}

\begin{definition}[The CZX $g,g$ fusion caps]
    \label{def:mpug_czx_fusion_caps}
    \lean{MPOTensor.CZX.fusionV, MPOTensor.CZX.fusionW}
    \leanok
    As in the concrete caps of
    \cite[lines~4671--4689]{FrancoRubio2025MPUGauging} (the $g=h=g^{-1}$
    specialization of the general fusion-tensor display in
    \cite[lines~2000--2037]{FrancoRubio2025MPUGauging}), in the ordered bond
    coordinates $0,1,2,3$, define the row vector
    $F^<_{g,g}\colon\C^4\to\C$ and the column vector
    $F^>_{g,g}\colon\C\to\C^4$ by
    \begin{align}
        F^<_{g,g} & =(0,-i,i,0),                     \\
        F^>_{g,g} & =\tfrac12(0,i,-i,0)^{\mathsf T}.
    \end{align}
    % Only the concrete g,g slice of #7354 is covered here.
\end{definition}

\begin{lemma}[Product coordinates of the CZX fusion caps]
    \label{thm:mpug_czx_fusion_caps_apply}
    \lean{MPOTensor.CZX.fusion_caps_apply}
    \leanok
    \uses{def:mpug_czx_fusion_caps, def:mpug_czx_dagger_gauge}
    Let $F^<_{g,g},F^>_{g,g}$ be the caps of
    Definition~\ref{def:mpug_czx_fusion_caps}, and let $Y$ be the Pauli
    matrix of Definition~\ref{def:mpug_czx_dagger_gauge}. Identify the
    product-bond coordinate $(a,b)$ with $2a+b$. For all bond bits
    $a,b\in\{0,1\}$,
    \begin{align}
        F^<_{g,g}(a,b) & =Y_{a,b},
        \label{eq:mpug_czx_fusion_caps_apply} \\
        F^>_{g,g}(a,b) & =\tfrac12Y_{b,a}.
        \label{eq:mpug_czx_fusion_caps_apply_right}
    \end{align}
    Thus, as in the printed diagram of
    \cite[lines~4671--4689]{FrancoRubio2025MPUGauging}, both caps vanish at
    the diagonal bond coordinates and carry $Y$ in opposite orientations.
\end{lemma}
\begin{proof}\leanok
    \uses{def:mpug_czx_fusion_caps, def:mpug_czx_dagger_gauge}
    Evaluating both caps at $2a+b$ for the four pairs $(a,b)$ and comparing
    with the entries of $Y$ yields
    (\ref{eq:mpug_czx_fusion_caps_apply}) and
    (\ref{eq:mpug_czx_fusion_caps_apply_right}).
\end{proof}

\begin{lemma}[The CZX fusion reduction]
    \label{thm:mpug_czx_fusion_reduction}
    \lean{MPOTensor.CZX.fusion_isReduction,
        MPOTensor.CZX.fusion_isReductionExteriorBufferLength_one}
    \leanok
    \uses{def:mpug_czx_fusion_caps, def:mpdo_operator_product,
        def:mpug_czx_exact_tensor, threeMPU, def:mpo_to_mps,
        def:mps_rectangular_reduction,
        def:mps_reduction_exterior_buffer_length}
    Let $B$ be the doubled-index tensor of the product
    $\mathcal U_g\cdot\mathcal U_g$
    (Definition~\ref{def:mpdo_operator_product} applied to
    Definition~\ref{def:mpug_czx_exact_tensor} with itself, viewed
    through Definition~\ref{def:mpo_to_mps}), and let $I$ be the
    doubled-index bond-one identity tensor of
    Definition~\ref{threeMPU}. Then $(F^<_{g,g},F^>_{g,g})$ is a
    reduction from $B$ to $I$ with exterior buffer length one. For every
    word $\mathbf c$, including the empty word,
    \begin{align}
        F^<_{g,g}F^>_{g,g}              & =1,             \\
        F^<_{g,g}B^{\mathbf c}F^>_{g,g} & =I^{\mathbf c}.
    \end{align}
    For nonempty words $\mathbf p,\mathbf q,\mathbf c$, this gives the
    concrete $g,g$ instance of the fusion reduction
    \cite[lines~1405--1498]{FrancoRubio2025MPUGauging} with the fusion
    tensors of \cite[lines~2000--2037]{FrancoRubio2025MPUGauging}
    specialized to the CZX display
    \cite[lines~4671--4689]{FrancoRubio2025MPUGauging}:
    \begin{align}
        B^{\mathbf p}F^>_{g,g}I^{\mathbf c}F^<_{g,g}B^{\mathbf q}
         & =B^{\mathbf p\mathbf c\mathbf q}.
        \label{eq:mpug_czx_fusion_reduction}
    \end{align}
\end{lemma}
\begin{proof}\leanok
    \uses{def:mpug_czx_fusion_caps, def:mpug_czx_exact_tensor,
        def:mpdo_operator_product, def:mpo_to_mps, threeMPU}
    With $V=F^<_{g,g}$ and $W=F^>_{g,g}$, substituting the cap and
    tensor coordinates gives the three local identities
    \begin{align}
        VW       & =1,
        \label{eq:mpug_czx_fusion_cap_product} \\
        VB^iW    & =I^i,
        \label{eq:mpug_czx_fusion_letter}      \\
        B^iWVB^j & =B^iB^j.
        \label{eq:mpug_czx_fusion_letter_insertion}
    \end{align}
    The all-word identity $VB^{\mathbf c}W=I^{\mathbf c}$ follows by
    induction on $\mathbf c$, removing its first letter. The empty word
    is (\ref{eq:mpug_czx_fusion_cap_product}) and a single letter is
    (\ref{eq:mpug_czx_fusion_letter}). For $\mathbf c=ij\mathbf c'$,
    inserting $WV$ after the first letter by
    (\ref{eq:mpug_czx_fusion_letter_insertion}), splitting the caps off both
    factors, and applying (\ref{eq:mpug_czx_fusion_letter}) together with
    the induction hypothesis for $j\mathbf c'$ gives
    \begin{align}
        VB^{ij\mathbf c'}W
         & =(VB^iW)(VB^{j\mathbf c'}W)
        =I^iI^{j\mathbf c'}=I^{ij\mathbf c'}.
    \end{align}
    For nonempty words $\mathbf p,\mathbf q$, a second induction on the left
    word $\mathbf p$, again removing its first letter, gives
    \begin{align}
        B^{\mathbf p}WVB^{\mathbf q} & =B^{\mathbf p\mathbf q}.
        \label{eq:mpug_czx_fusion_word_insertion}
    \end{align}
    For the singleton case $\mathbf p=i$ and
    $\mathbf q=j\mathbf q'$,
    (\ref{eq:mpug_czx_fusion_letter_insertion}) gives
    \begin{align}
        B^iWVB^{j\mathbf q'}
         & =(B^iWVB^j)B^{\mathbf q'}
        =B^iB^jB^{\mathbf q'}=B^{i j\mathbf q'}.
    \end{align}
    If $\mathbf p=i\mathbf p'$ with $\mathbf p'$ nonempty, the induction
    hypothesis applied to the tail gives
    \begin{align}
        B^{i\mathbf p'}WVB^{\mathbf q}
         & =B^i(B^{\mathbf p'}WVB^{\mathbf q})
        =B^iB^{\mathbf p'\mathbf q}=B^{i\mathbf p'\mathbf q}.
    \end{align}
    For nonempty buffers and middle word, substituting
    $I^{\mathbf c}=VB^{\mathbf c}W$ and applying
    (\ref{eq:mpug_czx_fusion_word_insertion}) twice yields
    \begin{align}
        B^{\mathbf p}WI^{\mathbf c}VB^{\mathbf q}
         & =B^{\mathbf p}WVB^{\mathbf c}WVB^{\mathbf q}
        =B^{\mathbf p\mathbf c\mathbf q}.
    \end{align}
\end{proof}

\subsection{Printed fusion--action coordinate contractions}

The contractions below evaluate the printed action pair and right fusion
cap of \cite[lines~4671--4694]{FrancoRubio2025MPUGauging} in fixed product-bond
order. They are compared with the fusion--action relation of
\cite[lines~1875--1913]{FrancoRubio2025MPUGauging}, but do not identify its
$L$-symbol or its retained physical scalar. Compatibility of the oriented
legs, density insertion, and fusion gauge remains open
\cite{gap:fbc25_czx_fusion_coordinate_compatibility}. No cap, action tensor,
block label, or physical gate is changed.

\begin{definition}[Sequential action covector]
    \label{def:mpug_czx_sequential_action_bra}
    \lean{MPOTensor.CZX.sequentialActionBra}
    \leanok
    \uses{def:mpug_czx_printed_action_pair}
    For $x\in\{0,1\}$, let $p_x$ be the printed action row of
    Definition~\ref{def:mpug_czx_printed_action_pair}. Define the row
    $P_x\colon\C^4\to\C$ by the following coordinates, for
    $\alpha,\gamma\in\{0,1\}$:
    \begin{align}
        P_x                    & =p_{1-x}\otimes p_x,         &
        (P_x)_{2\alpha+\gamma} & =p_{1-x}(\alpha)p_x(\gamma).
    \end{align}
    The first action factor occupies the first product-bond coordinate.
\end{definition}

\begin{definition}[Sequential action vector]
    \label{def:mpug_czx_sequential_action_ket}
    \lean{MPOTensor.CZX.sequentialActionKet}
    \leanok
    \uses{def:mpug_czx_printed_action_pair}
    For $x\in\{0,1\}$, let $q_x$ be the printed action column of
    Definition~\ref{def:mpug_czx_printed_action_pair}. Define the column
    $Q_x\colon\C\to\C^4$ by the following coordinates, for
    $\alpha,\gamma\in\{0,1\}$:
    \begin{align}
        Q_x                    & =q_{1-x}\otimes q_x,         &
        (Q_x)_{2\alpha+\gamma} & =q_{1-x}(\alpha)q_x(\gamma).
    \end{align}
    Rows contract with columns by ordinary matrix multiplication, with no
    additional complex conjugation.
\end{definition}

\begin{definition}[Supported product letter]
    \label{def:mpug_czx_fusion_action_letter}
    \lean{MPOTensor.CZX.fusionActionLetter}
    \leanok
    \uses{def:mpug_czx_exact_tensor, def:mpdo_operator_product}
    Let $T$ be the exact blocked CZX tensor of
    Definition~\ref{def:mpug_czx_exact_tensor}. For $x\in\{0,1\}$, define
    $B_x\in\operatorname{Mat}_4(\C)$ to be the physical $(3x,3x)$ letter
    of $T\mathbin{\cdot}T$. Explicitly, for
    $\alpha,\gamma,\beta,\delta\in\{0,1\}$,
    \begin{align}
        (B_x)_{2\alpha+\gamma,2\beta+\delta}
         & =\sum_{j=0}^3 T^{3x,j}_{\alpha\beta}T^{j,3x}_{\gamma\delta}.
    \end{align}
    The physical index $3x$ encodes the two bits $(x,x)$.
\end{definition}

\begin{lemma}[Coordinates of the supported product letters]
    \label{thm:mpug_czx_supported_letter_coordinates}
    \lean{MPOTensor.CZX.mulTensor_diagonal_letter_coordinates}
    \leanok
    \uses{def:mpug_czx_fusion_action_letter}
    The two supported product letters have coordinates
    \begin{align}
        B_0 & =\begin{pmatrix}
                   0 & 0 & 0 & 0 \\-1&1&-1&1\\0&0&0&0\\0&0&0&0
               \end{pmatrix}, \\
        B_1 & =\begin{pmatrix}
                   0 & 0 & 0 & 0 \\0&0&0&0\\-1&-1&1&1\\0&0&0&0
               \end{pmatrix}.
        \label{eq:mpug_czx_supported_letters}
    \end{align}
\end{lemma}
\begin{proof}\leanok
    \uses{def:mpug_czx_exact_tensor, def:mpdo_operator_product}
    The letter of $T\mathbin{\cdot}T$ at a physical pair $(i,i')$ vanishes
    unless $i=i'$, because the blocked tensor complements the physical index
    it carries. At $i=i'=3x$ the summation index of
    Definition~\ref{def:mpug_czx_fusion_action_letter} retains the single
    term $j=3(1-x)$, whose two factors force the left virtual pair
    $(\alpha,\gamma)=(x,1-x)$, the row of index $x+1$. On that row the
    surviving sign is $(-1)^{1+(1-x)\delta+x\beta}$, which is the row
    displayed above.
\end{proof}

\begin{definition}[Printed fusion--action coefficient]
    \label{def:mpug_czx_fusion_action_coefficient}
    \lean{MPOTensor.CZX.fusionActionCoefficient}
    \leanok
    \uses{def:mpug_czx_sequential_action_bra,
        def:mpug_czx_fusion_action_letter, def:mpug_czx_fusion_caps}
    For $x\in\{0,1\}$, let $P_x$ and $B_x$ be the sequential action row
    and supported product letter, and put
    $W=F^>_{g,g}=\tfrac12(0,i,-i,0)^{\mathsf T}$, in the product-bond
    order of Definition~\ref{def:mpug_czx_fusion_caps}. Define
    \begin{align}
        c_x & =P_xB_xW\in\C.
    \end{align}
    The one-by-one matrix on the right is identified with its sole entry.
\end{definition}

\begin{lemma}[Values of the printed coefficient]
    \label{thm:mpug_czx_fusion_action_coefficient_coordinates}
    \lean{MPOTensor.CZX.fusionActionCoefficient_coordinates}
    \leanok
    \uses{def:mpug_czx_fusion_action_coefficient}
    The printed coefficients satisfy $c_0=1$ and $c_1=-1$.
\end{lemma}
\begin{proof}\leanok
    \uses{thm:mpug_czx_supported_letter_coordinates,
        def:mpug_czx_sequential_action_bra, def:mpug_czx_printed_action_pair,
        def:mpug_czx_fusion_caps}
    The printed rows are $P_0=(0,-i,0,0)$ and $P_1=(0,0,-i,0)$, and the
    column is $W=\tfrac12(0,i,-i,0)^{\mathsf T}$. With $e_0,\dots,e_3$ the
    standard basis of the product bond, the single nonzero row of each
    supported letter of (\ref{eq:mpug_czx_supported_letters}) contracts with
    $W$ as
    \begin{align}
        B_0W & =\tfrac12(i+i)e_1=ie_1,   \\
        B_1W & =\tfrac12(-i-i)e_2=-ie_2,
    \end{align}
    so the two finite matrix products are
    \begin{align}
        c_0 & =P_0B_0W=(0,-i,0,0)\,ie_1=1,     \\
        c_1 & =P_1B_1W=(0,0,-i,0)\,(-ie_2)=-1.
    \end{align}
\end{proof}

\begin{lemma}[Normalization, proportionality, and idempotence]
    \label{thm:mpug_czx_fusion_action_finite_identities}
    \lean{MPOTensor.CZX.fusionAction_finite_identities}
    \leanok
    \uses{def:mpug_czx_sequential_action_bra,
        def:mpug_czx_sequential_action_ket, def:mpug_czx_fusion_action_letter,
        def:mpug_czx_fusion_action_coefficient, def:mpug_czx_fusion_caps}
    For each $x\in\{0,1\}$, the printed rows, columns, and supported
    product letter satisfy
    \begin{align}
        P_xB_xQ_x & =1, & B_xW & =c_xB_xQ_x, & B_x^2 & =B_x.
    \end{align}
\end{lemma}
\begin{proof}\leanok
    \uses{thm:mpug_czx_supported_letter_coordinates,
        def:mpug_czx_fusion_action_letter,
        def:mpug_czx_sequential_action_bra, def:mpug_czx_sequential_action_ket,
        def:mpug_czx_printed_action_pair, def:mpug_czx_fusion_caps,
        def:mpug_czx_fusion_action_coefficient}
    Substituting the printed action columns gives
    $Q_0=\tfrac14(-i,i,-i,i)^{\mathsf T}$ and
    $Q_1=\tfrac14(-i,-i,i,i)^{\mathsf T}$. With $e_0,\dots,e_3$ the standard
    basis of the product bond, the two supported letters of
    (\ref{eq:mpug_czx_supported_letters}) contract with these columns and
    with the cap $W$ as
    \begin{align}
        B_0Q_0 & =ie_1,  \\
        B_0W   & =ie_1,  \\
        B_1Q_1 & =ie_2,  \\
        B_1W   & =-ie_2.
        \label{eq:mpug_czx_fusion_action_contractions}
    \end{align}
    The printed rows are $P_0=(0,-i,0,0)$ and $P_1=(0,0,-i,0)$, so applying
    $P_x$ to (\ref{eq:mpug_czx_fusion_action_contractions}) gives
    $P_xB_xQ_x=1$ and $P_xB_xW=c_x$, with $c_0=1$ and $c_1=-1$. Comparing
    the two right-hand sides for each $x$ gives $B_xW=c_xB_xQ_x$. Each
    matrix of (\ref{eq:mpug_czx_supported_letters}) has a single nonzero row,
    whose entry in the supporting column is $1$, so $B_x^2=B_x$.
\end{proof}

\begin{lemma}[Positive repetitions of the supported letter]
    \label{thm:mpug_czx_fusion_action_positive_power}
    \lean{MPOTensor.CZX.fusionAction_positive_power}
    \leanok
    \uses{def:mpug_czx_sequential_action_bra,
        def:mpug_czx_sequential_action_ket, def:mpug_czx_fusion_action_letter,
        def:mpug_czx_fusion_action_coefficient, def:mpug_czx_fusion_caps}
    For every $x\in\{0,1\}$ and every integer $n\geq0$,
    \begin{align}
        B_x^{n+1}       & =B_x, & P_xB_x^{n+1}W & =c_x,             \\
        P_xB_x^{n+1}Q_x & =1,   & B_x^{n+1}W    & =c_xB_x^{n+1}Q_x.
    \end{align}
    These statements concern repetitions of this single supported letter,
    not arbitrary words or an identification with the source $L$-symbol.
\end{lemma}
\begin{proof}\leanok
    \uses{thm:mpug_czx_fusion_action_finite_identities,
        def:mpug_czx_fusion_action_coefficient}
    By induction on $n$, the first power is $B_x$, and multiplication by
    $B_x$ together with idempotence proves the successor step.
    Substituting $B_x^{n+1}=B_x$ into the remaining three identities and
    applying the definition of $c_x$ along with the normalization and
    proportionality of
    Lemma~\ref{thm:mpug_czx_fusion_action_finite_identities} completes
    the proof.
\end{proof}
\subsection{The displayed CZX circuit tuple}

\begin{definition}[The displayed CZX circuit tuple]
    \label{def:mpug_czx_circuit_tuple}
    \lean{MPOTensor.CZX.bitFlip,
        MPOTensor.CZX.pauliX,
        MPOTensor.CZX.pauliZ,
        MPOTensor.CZX.controlledZ,
        MPOTensor.CZX.w,
        MPOTensor.CZX.lambda,
        MPOTensor.CZX.tildeLambda,
        MPOTensor.CZX.siteBits,
        MPOTensor.CZX.localBits,
        MPOTensor.CZX.localBits_apply_zero,
        MPOTensor.CZX.localBits_apply_one,
        MPOTensor.CZX.localBits_apply_two,
        MPOTensor.CZX.localBits_apply_three,
        MPOTensor.CZX.matterMatrix,
        MPOTensor.CZX.gen,
        MPOTensor.CZX.circuitTuple,
        MPOTensor.CZX.circuitTuple_one,
        MPOTensor.CZX.circuitTuple_gen_one,
        MPOTensor.CZX.circuitTuple_gen_gen}
    \leanok
    \uses{def:mpug_monomial_matrix}
    Number the four matter qubits of two neighboring blocked sites by
    $0,1,2,3$, so that $(x_0,x_1)$ are the qubits of the first site and
    $(x_2,x_3)$ those of the second, and write
    $\ket{x}=\ket{x_0}\otimes\ket{x_1}\otimes\ket{x_2}\otimes\ket{x_3}$ for
    $x\in\mathbb F_2^4$. The elementary gates act by
    \begin{align}
        X_r\ket{x}              & =\ket{x+e_r},          \\
        Z_r\ket{x}              & =(-1)^{x_r}\ket{x},    \\
        \mathrm{CZ}_{rs}\ket{x} & =(-1)^{x_rx_s}\ket{x},
    \end{align}
    and operator products act from right to left. In this qubit
    numbering, the fusion operators $\lambda^R_{g,e}=w_R$ and
    $\lambda^R_{g,g}$ and the modified fusion operator
    $\tilde\lambda^R_{g,g}=\lambda^R_{g,g}Z_1$ of
    \cite[lines~4503--5183]{FrancoRubio2025MPUGauging} are
    \begin{align}
        w             & =(X_0X_1)\,\mathrm{CZ}_{01}\mathrm{CZ}_{12}\,(Z_2Z_3),
        \label{eq:mpug_czx_w}                                                       \\
        \lambda       & =-i\,(X_0X_1X_2Z_3)\,\mathrm{CZ}_{01}\mathrm{CZ}_{12}\,X_2,
        \label{eq:mpug_czx_lambda}                                                  \\
        \tilde\lambda & =\lambda Z_0.
        \label{eq:mpug_czx_tilde_lambda}
    \end{align}
    The displayed circuit tuple, indexed by the ordered pair of
    gauge labels in $\Z_2$ with $e\leftrightarrow0$ and $g\leftrightarrow1$,
    is
    \begin{align}
        R^{\mathrm{CZX}}
         & =\bigl(U_{0,0},U_{0,1},U_{1,0},U_{1,1}\bigr)
        =\bigl(\Id,\Id,w,\tilde\lambda\bigr).
        \label{eq:mpug_czx_tuple}
    \end{align}
    Each site index in $\{0,1,2,3\}$ is read as the pair of matter bits
    $(m^{(1)},m^{(2)})$ with $m^{(1)}$ the leading binary digit, and the
    four-qubit operators are reindexed to the two-site matter space by
    this identification.
    % Source anchor: arXiv:2502.20257, lines 4503--5183 (CZX model and modified
    % fusion operator).
\end{definition}

\begin{lemma}[Phase tables of the CZX fusion operators]
    \label{thm:mpug_czx_phase_tables}
    \lean{MPOTensor.CZX.barFlip,
        MPOTensor.CZX.barFlip_apply,
        MPOTensor.CZX.barFlip_mul_barFlip,
        MPOTensor.CZX.barFlip_symm,
        MPOTensor.CZX.eExponent,
        MPOTensor.CZX.hExponent,
        MPOTensor.CZX.fExponent,
        ZMod.neg_one_pow_val_add,
        MPOTensor.CZX.w_eq,
        MPOTensor.CZX.lambda_eq,
        MPOTensor.CZX.tildeLambda_eq,
        MPOTensor.CZX.fExponent_barFlip,
        MPOTensor.CZX.eExponent_barFlip,
        MPOTensor.CZX.matterBarFlip,
        MPOTensor.CZX.matterBarFlip_symm,
        MPOTensor.CZX.localBits_matterBarFlip,
        MPOTensor.CZX.matterMatrix_monomial,
        MPOTensor.CZX.w_mem_unitaryGroup,
        MPOTensor.CZX.tildeLambda_mem_unitaryGroup}
    \leanok
    \uses{def:mpug_czx_circuit_tuple}
    Write $\overline x=(x_0+1,x_1+1,x_2,x_3)$ and
    \begin{align}
        e(x) & =x_0x_1+x_1x_2+x_2+x_3,     \\
        h(x) & =x_0x_1+x_1x_2+x_1+x_3,     \\
        f(x) & =x_0x_1+x_1x_2+x_0+x_1+x_3.
    \end{align}
    Then
    \begin{align}
        w\ket{x}             & =(-1)^{e(x)}\ket{\overline x},
        \label{eq:mpug_czx_w_action}                            \\
        \lambda\ket{x}       & =-i(-1)^{h(x)}\ket{\overline x},
        \label{eq:mpug_czx_lambda_action}                       \\
        \tilde\lambda\ket{x} & =-i(-1)^{f(x)}\ket{\overline x},
        \label{eq:mpug_czx_tilde_lambda_action}
    \end{align}
    the exponents satisfy $f(\overline x)=e(x)+1$ and
    $e(\overline x)=f(x)+1$. After reindexing to the two-site matter coordinates,
    $w$ and $\tilde\lambda$ are unitary.
\end{lemma}

\begin{proof}\leanok
    \uses{thm:mpug_monomial_matrix_laws, def:mpug_czx_circuit_tuple}
    Each gate is monomial, so (\ref{eq:mpug_monomial_mul}) expresses each
    product as a monomial operator. Reading (\ref{eq:mpug_czx_w}) from right
    to left, the diagonal gates contribute the phase
    $(-1)^{x_0x_1+x_1x_2+x_2+x_3}$ and the two bit flips send $x$ to
    $\overline x$. In (\ref{eq:mpug_czx_lambda}) the first $X_2$ replaces
    $x_2$ by $x_2+1$ in the phase of $\mathrm{CZ}_{12}$, giving
    $x_1(x_2+1)=x_1x_2+x_1$, and the second $X_2$ restores $x_2$, so the
    permutation is again $\overline{\,\cdot\,}$. The factor $Z_0$ in
    (\ref{eq:mpug_czx_tilde_lambda}) adds $x_0$ to the exponent. The
    identities between $e$ and $f$ are checked on the sixteen bit strings.
    After relabelling the computational basis as two four-level matter sites,
    unitarity follows from Lemma~\ref{thm:mpug_monomial_matrix_laws} because
    all phases are unimodular.
\end{proof}

\begin{lemma}[Right movement coordinates of the displayed CZX factors]
    \label{thm:mpug_czx_displayed_right_movement_coordinates}
    \lean{MPOTensor.CZX.displayedSourceFactors_sourceWR_coordinates}
    \leanok
    \discussion{7355}
    \uses{def:mpug_czx_displayed_source_factors, def:mpug_staircase_gates,
        def:mpug_czx_circuit_tuple}
    Let $w_R$ be the crossed contraction of the displayed source factors
    of the exact CZX tensor $B$, and let $e_r$ be their right-rank
    identification with $\{0,1,2,3\}$. Let $W$ denote the matrix of
    $w=(X_0X_1)\mathrm{CZ}_{01}\mathrm{CZ}_{12}(Z_2Z_3)$ on two
    four-level matter sites, with each site read in leading-bit-first
    binary order. Identifying the right movement contraction in
    \cite[lines~811--869, 4559--4659, 4860--4888]{FrancoRubio2025MPUGauging}
    without additional hypotheses, for every $i,r,a,j\in\{0,1,2,3\}$,
    \begin{align}
        (w_R)_{(i,e_r^{-1}(r)),(e_r^{-1}(a),j)}
         & =W_{(i,r),(a,j)}.
    \end{align}
\end{lemma}

\begin{proof}\leanok
    \uses{def:mpug_czx_displayed_source_factors, def:mpug_staircase_gates,
        def:mpug_czx_displayed_factors, def:mpug_czx_circuit_tuple,
        thm:mpug_czx_phase_tables}
    Pulling back the rank indices expands the crossed contraction as
    $\sum_{\beta=0}^1(X_1)_{(i,\beta),a}(Y_1)_{r,(\beta,j)}$.
    The sparse action (\ref{eq:mpug_czx_w_action}) gives
    \begin{align}
        W_{(i,r),(a,j)}
         & =\begin{cases}
                (-1)^{e(x)}, & i=3-a\text{ and }r=j, \\
                0,           & \text{otherwise},
            \end{cases}
    \end{align}
    where $x=(x_0,x_1,x_2,x_3)\in\{0,1\}^4$ satisfies
    $a=2x_0+x_1$ and $j=2x_2+x_3$.
    Substituting the displayed factors and first discarding the zero
    entries of the empty factor $X_1$ by comparing $i$ and $a$ leaves
    only the comparison of $r$ and $j$, and on the surviving support the
    phases agree.
\end{proof}

\begin{lemma}[Left movement coordinates of the displayed CZX factors]
    \label{thm:mpug_czx_displayed_left_movement_coordinates}
    \lean{MPOTensor.CZX.displayedSourceFactors_sourceWL_coordinates}
    \leanok
    \discussion{7355}
    \uses{def:mpug_czx_displayed_source_factors, def:mpug_staircase_gates,
        def:mpug_czx_circuit_tuple}
    Let $w_L$ be the crossed contraction of the displayed source factors
    of the exact CZX tensor $B$, and let $e_l$ be their left-rank
    identification with $\{0,1,2,3\}$. Let $W$ denote the matrix of
    $w=(X_0X_1)\mathrm{CZ}_{01}\mathrm{CZ}_{12}(Z_2Z_3)$ on two
    four-level matter sites, with each site read in leading-bit-first
    binary order. In the same matter coordinates and without additional
    hypotheses, the left movement identification in
    \cite[lines~811--869, 4559--4659, 4860--4888]{FrancoRubio2025MPUGauging}
    holds for every $l,i,j,k\in\{0,1,2,3\}$:
    \begin{align}
        (w_L)_{(e_l^{-1}(l),i),(j,e_l^{-1}(k))}
         & =(W^\dagger)_{(l,i),(j,k)}.
    \end{align}
\end{lemma}

\begin{proof}\leanok
    \uses{def:mpug_czx_displayed_source_factors, def:mpug_staircase_gates,
        def:mpug_czx_displayed_factors, def:mpug_czx_circuit_tuple,
        thm:mpug_czx_phase_tables}
    Pulling back the rank indices expands the crossed contraction as
    $\sum_{\beta=0}^1(Y_2)_{l,(j,\beta)}(X_2)_{(\beta,i),k}$, while the
    conjugated reversed entry is
    $(W^\dagger)_{(l,i),(j,k)}=\overline{W_{(j,k),(l,i)}}$.
    By (\ref{eq:mpug_czx_w_action}), it vanishes unless $j=3-l$ and
    $k=i$; on this support it is $\overline{(-1)^{e(x)}}$, where
    $x=(x_0,x_1,x_2,x_3)\in\{0,1\}^4$ satisfies
    $l=2x_0+x_1$ and $i=2x_2+x_3$.
    Substituting the displayed factors and first discarding the zero
    entries of the filled factor $Y_2$ by comparing $l$ and $j$ leaves
    the comparison of $i$ and $k$, on which the remaining phases agree.
\end{proof}

% Scope of #7355 here: only the two displayed movement-coordinate identities.
% No full truncated-symmetry formula, associator, sign, or parent issue
% completion is asserted.

\begin{lemma}[Local monomial action of the CZX Gauss operator]
    \label{thm:mpug_czx_local_monomial_action}
    \lean{MPOTensor.CZX.iPow,
        MPOTensor.CZX.iPow_add,
        MPOTensor.CZX.iPow_natCast,
        MPOTensor.CZX.iPow_zero,
        MPOTensor.CZX.iPow_two,
        MPOTensor.CZX.iPow_three,
        MPOTensor.CZX.iPow_two_mul_val,
        MPOTensor.CZX.iPow_two_mul_cast,
        MPOTensor.CZX.uExponent,
        MPOTensor.CZX.phaseExponent,
        MPOTensor.CZX.phaseExponent_zero_zero,
        MPOTensor.CZX.phaseExponent_zero_one,
        MPOTensor.CZX.phaseExponent_one_zero,
        MPOTensor.CZX.phaseExponent_one_one,
        MPOTensor.CZX.localPerm,
        MPOTensor.CZX.localPhase,
        MPOTensor.CZX.gen_inv,
        MPOTensor.CZX.gen_mul_gen,
        MPOTensor.CZX.toAdd_gen_ne_zero,
        MPOTensor.CZX.gaussOperator_circuitTuple_gen}
    \leanok
    \uses{def:mpug_czx_circuit_tuple, def:mpug_gauss_operator}
    Let $a,a'\in\mathbb F_2$ be the two gauge bits and write
    $s=(x_0,x_1,a,x_2,x_3,a')$ and
    $\sigma(s)=(x_0+1,x_1+1,a+1,x_2,x_3,a'+1)$. The nontrivial local Gauss
    operator obtained from (\ref{eq:mpug_czx_tuple}) and
    (\ref{eq:mpug_gauss_operator}) is the monomial operator
    \begin{align}
        G\ket{s} & =i^{q(s)}\ket{\sigma(s)}
        \label{eq:mpug_czx_monomial_action}
    \end{align}
    with, for $u(x)=2(x_0x_1+x_1x_2+x_3)\bmod4$, the phase table
    \begin{align}
        q(x_0,x_1,0,x_2,x_3,0) & =3+u(x)+2x_2,      \\
        q(x_0,x_1,0,x_2,x_3,1) & =2+u(x)+2x_0+2x_1, \\
        q(x_0,x_1,1,x_2,x_3,0) & =u(x)+2x_2,        \\
        q(x_0,x_1,1,x_2,x_3,1) & =3+u(x)+2x_0+2x_1
        \label{eq:mpug_czx_phase_table}
    \end{align}
    modulo $4$.
\end{lemma}

\begin{proof}\leanok
    \uses{thm:mpug_gauss_operator_monomial, thm:mpug_czx_phase_tables}
    By Lemma~\ref{thm:mpug_gauss_operator_monomial} applied to the monomial
    unitaries of Lemma~\ref{thm:mpug_czx_phase_tables}, the operator
    $\mathcal G_g$ is monomial. The permutation part is
    $\overline{\,\cdot\,}$ on the matter bits in all four gauge sectors,
    since $\overline{\,\cdot\,}$ is an involution, and $T_g$ flips both
    gauge bits. The phases are read off
    (\ref{eq:mpug_czx_w_action})--(\ref{eq:mpug_czx_tilde_lambda_action})
    and the adjoint formula (\ref{eq:mpug_monomial_adjoint}) together with
    $f(\overline x)=e(x)+1$ and $e(\overline x)=f(x)+1$: for example the
    first row is $-i(-1)^{e(x)}=i^{3+2e(x)}=i^{3+u(x)+2x_2}$, and the other
    three rows follow in the same way. The four identities between the phase
    table and $e$, $f$ are checked on the sixteen bit strings.
\end{proof}

\subsection{Orbits, holonomy, and the exact dimension}

\begin{definition}[Chain configurations, bond flips, and labels]
    \label{def:mpug_czx_chain_labels}
    \lean{MPOTensor.CZX.Site,
        MPOTensor.CZX.siteDecode,
        MPOTensor.CZX.siteDecode_equivFin,
        MPOTensor.CZX.chainDecode,
        MPOTensor.CZX.chainDecode_apply,
        MPOTensor.CZX.flipPattern,
        MPOTensor.CZX.chainFlip,
        MPOTensor.CZX.chainFlip_apply,
        MPOTensor.CZX.succ_ne_self,
        MPOTensor.CZX.succ_succ_ne_self,
        MPOTensor.CZX.flipPattern_apply_self,
        MPOTensor.CZX.flipPattern_apply_succ,
        MPOTensor.CZX.flipPattern_apply_of_ne,
        MPOTensor.CZX.bondExponent,
        MPOTensor.CZX.bLabel,
        MPOTensor.CZX.pLabel}
    \leanok
    \uses{thm:mpug_czx_local_monomial_action}
    A computational-basis configuration of the periodic chain of $N$ blocked
    sites is
    \begin{align}
        s & =\bigl(m^{(1)}_k,m^{(2)}_k,g_k\bigr)_{k\in\Z/N\Z}\in S_N:=(\mathbb F_2^3)^N,
        \label{eq:mpug_czx_configuration}
    \end{align}
    read from the enumerated site indices of the placed Gauss projector. The
    bond-$j$ flip $\sigma_j$ adds $1$ to $m^{(1)}_j$, $m^{(2)}_j$, $g_j$, and
    $g_{j+1}$. The bond-$j$ phase exponent $q_j(s)$ is the table
    (\ref{eq:mpug_czx_phase_table}) evaluated at
    $x=(m^{(1)}_j,m^{(2)}_j,m^{(1)}_{j+1},m^{(2)}_{j+1})$ and
    $(a,a')=(g_j,g_{j+1})$. The affine labels are
    \begin{align}
        p_k(s) & =m^{(1)}_k+m^{(2)}_k,
        \label{eq:mpug_czx_p_label}            \\
        b_k(s) & =g_k+m^{(1)}_k+m^{(1)}_{k-1}.
        \label{eq:mpug_czx_b_label}
    \end{align}
\end{definition}

\begin{lemma}[Permutation-orbit classification]
    \label{thm:mpug_czx_orbit_classification}
    \lean{MPOTensor.CZX.fiberEquiv,
        MPOTensor.CZX.fiberEquiv_apply,
        MPOTensor.CZX.pLabel_fiberEquiv,
        MPOTensor.CZX.bLabel_fiberEquiv,
        MPOTensor.CZX.chainFlip_fiberEquiv}
    \leanok
    \uses{def:mpug_czx_chain_labels}
    For $N\geq2$, the map
    \begin{align}
        \bigl((p,b),\gamma\bigr)
         & \longmapsto
        \bigl(m^{(1)}_k,m^{(2)}_k,g_k\bigr)_k
        =\bigl(\gamma_k,\ p_k+\gamma_k,\ b_k+\gamma_k+\gamma_{k-1}\bigr)_k
        \label{eq:mpug_czx_fiber_coordinates}
    \end{align}
    is a bijection $\mathbb F_2^N\times\mathbb F_2^N\times\mathbb F_2^N\to S_N$
    whose inverse reads off $p(s)$, $b(s)$, and $\gamma=m^{(1)}$. In these
    coordinates $\sigma_j$ is $\gamma\mapsto\gamma+e_j$ and leaves $(p,b)$
    unchanged.
\end{lemma}

\begin{proof}\leanok
    \uses{def:mpug_czx_chain_labels}
    Both compositions of (\ref{eq:mpug_czx_fiber_coordinates}) with the
    label map are the identity, by the identity $a+a=0$ in $\mathbb F_2$.
    The flip $\sigma_j$ adds $((1,1),1)$ at site $j$ and $((0,0),1)$ at
    site $j+1$; in fiber coordinates this is $\gamma_j\mapsto\gamma_j+1$
    because $k-1=j$ exactly when $k=j+1$.
\end{proof}

\begin{lemma}[Involution and neighboring holonomy]
    \label{thm:mpug_czx_holonomy}
    \lean{MPOTensor.CZX.bondExponent_add_bondExponent_chainFlip,
        MPOTensor.CZX.bondExponent_holonomy,
        MPOTensor.CZX.bondExponent_chainFlip_of_disjoint,
        MPOTensor.CZX.fiberPhase,
        MPOTensor.CZX.fiberPhase_mul_fiberPhase_flip,
        MPOTensor.CZX.fiberPhase_flip_of_disjoint,
        MPOTensor.CZX.fiberPhase_holonomy,
        MPOTensor.CZX.fiberPhase_ne_zero}
    \leanok
    \uses{def:mpug_czx_chain_labels}
    For every $N\geq3$, every bond $j$, and every configuration $s$,
    \begin{align}
        q_j(s)+q_j(\sigma_js)
         & =0,
        \label{eq:mpug_czx_involution}                \\
        q_j(s)+q_{j+1}(\sigma_js)
         & =q_{j+1}(s)+q_j(\sigma_{j+1}s)+2b_{j+1}(s)
        \label{eq:mpug_czx_holonomy}
    \end{align}
    modulo $4$, and $q_j(\sigma_ks)=q_j(s)$ whenever
    $k\notin\{j-1,j,j+1\}$. For the phases
    $\varphi_j(s)=i^{q_j(s)}$, (\ref{eq:mpug_czx_involution}) gives
    $\varphi_j(s)\varphi_j(\sigma_js)=1$, while
    (\ref{eq:mpug_czx_holonomy}) gives the neighboring square relation with
    factor $(-1)^{b_{j+1}(s)}$.
\end{lemma}

\begin{proof}\leanok
    \uses{def:mpug_czx_chain_labels}
    The exponent $q_j$ depends only on the six bits at the sites $j$ and
    $j+1$, and $\sigma_k$ changes only the sites $k$ and $k+1$; this gives
    the disjoint-bond statement. For (\ref{eq:mpug_czx_involution}) the
    flip $\sigma_j$ changes $x_0$, $x_1$, $a$, and $a'$, and the identity is
    verified on the $64$ values of the six bits. For
    (\ref{eq:mpug_czx_holonomy}) the sites $j$, $j+1$, $j+2$ are distinct
    because $N\geq3$; $\sigma_j$ changes the left gauge bit seen by the bond
    $(j+1,j+2)$, and $\sigma_{j+1}$ changes $x_2$, $x_3$, and $a'$ seen by
    the bond $(j,j+1)$, so the identity is verified on the $512$ values of
    the nine bits at the three sites.
\end{proof}

\begin{lemma}[Fixed spaces on the orbit fibers]
    \label{thm:mpug_czx_good_fibers}
    \lean{MPOTensor.CZX.isTrivialHolonomy_fiberPhase_iff}
    \leanok
    \uses{def:mpug_trivial_holonomy, thm:mpug_czx_orbit_classification}
    Let $N\geq3$. In the fiber coordinates
    (\ref{eq:mpug_czx_fiber_coordinates}), the phases $i^{q_j}$ have trivial
    holonomy on the fiber over $(p,b)$ if and only if $b=0$.
\end{lemma}

\begin{proof}\leanok
    \uses{thm:mpug_czx_holonomy, thm:mpug_czx_orbit_classification}
    If $b=0$, then for $k=j+1$ or $j=k+1$ the square identity is
    (\ref{eq:mpug_czx_holonomy}) with $b_{j+1}=0$, for $k=j$ it is trivial,
    and otherwise the bonds are disjoint and both sides agree. Conversely,
    trivial holonomy at the square spanned by the bonds $k-1$ and $k$ at
    $\gamma=0$ gives $(-1)^{b_k}=1$ by (\ref{eq:mpug_czx_holonomy}), since the
    phases are nonzero, hence $b_k=0$ for every $k$.
\end{proof}

\begin{lemma}[The placed CZX Gauss operators]
    \label{thm:mpug_czx_placed_operators}
    \lean{MPOTensor.CZX.windowCoordinates,
        MPOTensor.CZX.placedGaussOperator,
        MPOTensor.CZX.sum_multiplicative_zmod_two,
        MPOTensor.CZX.gaussProjector_circuitTuple,
        MPOTensor.CZX.placedGaussProjector_circuitTuple,
        MPOTensor.CZX.inv_two_smul_add_mulVec_eq_iff,
        MPOTensor.CZX.commonFixedSubmodule_placedGaussProjector_circuitTuple,
        MPOTensor.CZX.placedGaussOperator_eq_monomial,
        MPOTensor.CZX.siteBits_eq_localBits_zero,
        MPOTensor.CZX.siteBits_eq_localBits_one,
        MPOTensor.CZX.barFlip_apply_zero,
        MPOTensor.CZX.barFlip_apply_one,
        MPOTensor.CZX.barFlip_apply_two,
        MPOTensor.CZX.barFlip_apply_three,
        MPOTensor.CZX.toAdd_mul_gen_inv,
        MPOTensor.CZX.toAdd_gen_mul,
        MPOTensor.CZX.gaussLegAction_fst,
        MPOTensor.CZX.gaussLegAction_snd,
        MPOTensor.CZX.chainDecode_window,
        MPOTensor.CZX.chainDecode_windowPerm,
        MPOTensor.CZX.localPhase_window,
        MPOTensor.CZX.fiberCoordinates,
        MPOTensor.CZX.placedGaussOperator_eq_reindex}
    \leanok
    \uses{def:mpug_placed_gauss_projector, def:mpug_czx_circuit_tuple,
        def:mpug_czx_chain_labels, def:mpug_common_fixed_subspace}
    Let $N\geq2$ and let $G_j$ be the local Gauss operator
    (\ref{eq:mpug_czx_monomial_action}) placed on the bond $(j,j+1)$. For
    $\Z_2$ the placed Gauss projector of
    Definition~\ref{def:mpug_placed_gauss_projector} is
    \begin{align}
        \widetilde P_j & =\tfrac12\bigl(\Id+G_j\bigr),
        \label{eq:mpug_czx_placed_projector}
    \end{align}
    so the common fixed subspace $\mathcal V_N(R^{\mathrm{CZX}})$ of the
    $\widetilde P_j$ is the common fixed subspace of the $G_j$. In the
    coordinates (\ref{eq:mpug_czx_configuration}), $G_j$ is the monomial
    operator $G_j\ket{s}=i^{q_j(s)}\ket{\sigma_js}$, and in the fiber
    coordinates (\ref{eq:mpug_czx_fiber_coordinates}) it is the monomial
    operator with permutation $\gamma\mapsto\gamma+e_j$ and phase
    $i^{q_j}$.
\end{lemma}

\begin{proof}\leanok
    \uses{def:mpug_gauss_projector, thm:mpug_gauss_representation,
        thm:mpug_czx_local_monomial_action, thm:mpug_embed_monomial,
        thm:mpug_czx_orbit_classification, thm:mpug_mem_common_fixed_subspace}
    The average over $\Z_2=\{e,g\}$ is $\frac12(\mathcal G_e+\mathcal G_g)$
    with $\mathcal G_e=\Id$, and relabelling and placement are linear and
    unital, which gives (\ref{eq:mpug_czx_placed_projector}). A vector $v$
    satisfies $\frac12(v+G_jv)=v$ exactly when $G_jv=v$.
    Lemma~\ref{thm:mpug_czx_local_monomial_action} and
    Lemma~\ref{thm:mpug_embed_monomial} express $G_j$ as a monomial
    operator; decoding the two-site window at $j$ into the bits of the sites
    $j$ and $j+1$ identifies its permutation with $\sigma_j$ and its phase
    with $i^{q_j(s)}$. Finally
    Lemma~\ref{thm:mpug_czx_orbit_classification} reindexes both to the
    fiber coordinates.
\end{proof}

\begin{lemma}[Exact dimension for the displayed CZX circuit tuple]
    \label{thm:mpug_czx_dimension}
    \lean{MPOTensor.CZX.finrank_commonFixedSubmodule_placedGaussProjector_circuitTuple}
    \leanok
    \uses{def:mpug_placed_gauss_projector, def:mpug_czx_circuit_tuple,
        def:mpug_common_fixed_subspace}
    Answering the growth question of
    \cite[lines~5198--5204]{FrancoRubio2025MPUGauging} for the displayed
    circuit tuple, on a periodic chain of $N\geq3$ blocked sites,
    \begin{align}
        \dim\mathcal V_N(R^{\mathrm{CZX}}) & =2^N.
        \label{eq:mpug_czx_dimension}
    \end{align}
\end{lemma}

\begin{proof}\leanok
    \uses{thm:mpug_czx_placed_operators,
        thm:mpug_monomial_fixed_space_criterion,
        thm:mpug_czx_holonomy, thm:mpug_czx_good_fibers}
    By Lemma~\ref{thm:mpug_czx_placed_operators},
    $\mathcal V_N(R^{\mathrm{CZX}})$ is the common fixed subspace of monomial
    operators which, in fiber coordinates, have permutations
    $\gamma\mapsto\gamma+e_j$ on the fibers over
    $(p,b)\in\mathbb F_2^N\times\mathbb F_2^N$, and by
    (\ref{eq:mpug_czx_involution}) these operators are involutions. The
    monomial fixed-space criterion
    (\ref{eq:mpug_monomial_fixed_space_count}) counts the fibers with trivial
    holonomy, which by Lemma~\ref{thm:mpug_czx_good_fibers} are exactly
    the fibers with $b=0$. There is one such $b$ and there are $2^N$ values
    of $p$, so
    \begin{align}
        \dim\mathcal V_N(R^{\mathrm{CZX}})
         & =\sum_{p\in\mathbb F_2^N}1=2^N.
    \end{align}
\end{proof}

\begin{lemma}[The CZX dimension in the gauge-invariant subspace of the displayed tuple]
    \label{thm:mpug_czx_dimension_of_completion}
    \lean{MPOTensor.CZX.finrank_gaugeInvariantSubspace_circuitTuple}
    \leanok
    \uses{def:mpug_gauge_invariant_subspace, def:mpug_czx_circuit_tuple}
    For the gauge-invariant subspace~(\ref{eq:mpug_gauge_invariant_subspace}) at
    $R=R^{\mathrm{CZX}}$ and $N\geq3$,
    $\dim_{\C}\mathcal V_N(R^{\mathrm{CZX}})\allowbreak=2^N$. The value is thereby
    recorded for the subspace that
    Definition~\ref{def:mpug_gauge_invariant_subspace} attaches to an arbitrary
    tuple, so exhibiting $R^{\mathrm{CZX}}$ as a completion of prescribed
    defect maps with matter defect states satisfying local defect support,
    exact defect covariance, and nonvanishing of the gauged vector would put
    the exact value beside the lower bound $\dim_{\C}\mathcal V_N(R)\geq1$ of
    Lemma~\ref{thm:mpug_gauge_invariant_subspace_finrank}. None of those
    hypotheses is established here. The statement is about the displayed
    circuit tuple and asserts no membership in the full physical completion
    class of the CZX defect data.
\end{lemma}

\begin{proof}\leanok
    \uses{thm:mpug_czx_dimension, def:mpug_placed_gauss_projector}
    Both subspaces are the common fixed subspace of the placed Gauss
    projectors $\widetilde P_j(R^{\mathrm{CZX}})$ of
    Definition~\ref{def:mpug_placed_gauss_projector} acting on the
    configuration space of the chain, so
    Lemma~\ref{thm:mpug_czx_dimension} gives the value.
\end{proof}

\subsection{A second displayed circuit tuple at three and four sites}

Constraining the block $U_{1,1}$ of (\ref{eq:mpug_czx_tuple}) on the defect
subspace $\spn_{\C}\{\ket{0011},\ket{1100}\}$ leaves gauge freedom in
$\tilde\lambda$. Replacing $\tilde\lambda$ by another unitary with identical
action on this subspace yields an alternate circuit tuple whose common $+1$
eigenspace dimension differs at three and four blocked sites.

\begin{definition}[The second displayed circuit tuple]
    \label{def:mpug_czx_star_tuple}
    \lean{MPOTensor.CZX.rStarExponent,
        MPOTensor.CZX.rStar,
        MPOTensor.CZX.rStar_eq,
        MPOTensor.CZX.rStarExponent_zero_zero_one_one,
        MPOTensor.CZX.rStarExponent_one_one_zero_zero,
        MPOTensor.CZX.rStar_mulVec_eq_self,
        MPOTensor.CZX.tildeLambdaStar,
        MPOTensor.CZX.tildeLambdaStar_eq,
        MPOTensor.CZX.tildeLambdaStar_mem_unitaryGroup,
        MPOTensor.CZX.circuitTupleStar,
        MPOTensor.CZX.circuitTupleStar_one,
        MPOTensor.CZX.circuitTupleStar_gen_one,
        MPOTensor.CZX.circuitTupleStar_gen_gen}
    \leanok
    \uses{def:mpug_czx_circuit_tuple, thm:mpug_czx_phase_tables}
    In the qubit numbering of Definition~\ref{def:mpug_czx_circuit_tuple}, set
    \begin{align}
        r_\star(x) & =x_1+x_3+x_0x_1+x_0x_3+x_1x_2+x_2x_3,
        \label{eq:mpug_czx_star_phase}                                         \\
        R_\star    & =Z_1Z_3\,\mathrm{CZ}_{01}\mathrm{CZ}_{03}\mathrm{CZ}_{12}
        \mathrm{CZ}_{23}.
        \label{eq:mpug_czx_star_gate}
    \end{align}
    Then $R_\star\ket{x}=(-1)^{r_\star(x)}\ket{x}$, and
    $r_\star(0011)=r_\star(1100)=0$, so $R_\star$ fixes the prescribed defect
    subspace $\spn_{\C}\{\ket{0011},\ket{1100}\}$ pointwise. The second
    displayed circuit tuple is
    \begin{align}
        R^\star
         & =\bigl(U^\star_{0,0},U^\star_{0,1},U^\star_{1,0},U^\star_{1,1}\bigr)
        =\bigl(\Id,\Id,w,\tilde\lambda R_\star\bigr),
        \label{eq:mpug_czx_star_tuple}
    \end{align}
    with
    $\tilde\lambda R_\star\ket{x}=-i(-1)^{f(x)+r_\star(x)}\ket{\overline x}$;
    after reindexing to the two-site matter coordinates it is unitary. It has the
    same action as $R^{\mathrm{CZX}}$ on that defect subspace. The four defect
    domains and prescribed maps of
    Definition~\ref{def:mpug_czx_defect_maps} determine a four-sector
    completion class, and both tuples extend all four prescribed maps: the
    first by Lemma~\ref{thm:mpug_czx_tuple_completion} and the second by
    Lemma~\ref{thm:mpug_czx_star_tuple_completion}, because $R_\star$ fixes
    the defect domain of the pair $(g,g)$ pointwise and the other three
    operators are those of $R^{\mathrm{CZX}}$.
    % Source anchor: arXiv:2502.20257, lines 4503--5183 (CZX model, defect and
    % fusion data, modified Gauss law).
\end{definition}

\begin{lemma}[Local monomial action of the second tuple's Gauss operator]
    \label{thm:mpug_czx_star_local_monomial_action}
    \lean{MPOTensor.CZX.rStarExponent_barFlip,
        MPOTensor.CZX.starDefect,
        MPOTensor.CZX.phaseExponentStar,
        MPOTensor.CZX.phaseExponentStar_zero_zero,
        MPOTensor.CZX.phaseExponentStar_zero_one,
        MPOTensor.CZX.phaseExponentStar_one_zero,
        MPOTensor.CZX.phaseExponentStar_one_one,
        MPOTensor.CZX.localPhaseStar,
        MPOTensor.CZX.gaussOperator_circuitTupleStar_gen}
    \leanok
    \uses{def:mpug_czx_star_tuple, def:mpug_gauss_operator}
    The nontrivial local Gauss operator obtained from
    (\ref{eq:mpug_czx_star_tuple}) and (\ref{eq:mpug_gauss_operator}) is the
    monomial operator $G^\star\ket{s}=i^{q^\star(s)}\ket{\sigma(s)}$ with the
    same permutation part $\sigma$ as in
    (\ref{eq:mpug_czx_monomial_action}) and with
    \begin{align}
        q^\star(x_0,x_1,a,x_2,x_3,a') & =q(x_0,x_1,a,x_2,x_3,a')
        +2\delta_\star(x,a,a')
        \label{eq:mpug_czx_star_phase_table}
    \end{align}
    modulo $4$, where $\delta_\star(x,0,0)=r_\star(\overline x)$,
    $\delta_\star(x,1,1)=r_\star(x)$, and $\delta_\star$ vanishes in the two
    remaining gauge sectors. The phase polynomial satisfies
    \begin{align}
        r_\star(\overline x) & =r_\star(x)+x_0+x_1+x_2+x_3.
        \label{eq:mpug_czx_star_bar}
    \end{align}
\end{lemma}

\begin{proof}\leanok
    \uses{thm:mpug_gauss_operator_monomial, thm:mpug_czx_phase_tables,
        def:mpug_czx_star_tuple}
    By Lemma~\ref{thm:mpug_gauss_operator_monomial} the operator is monomial,
    and the two blocks $U^\star_{0,1}=\Id$ and $U^\star_{1,0}=w$ agree with
    those of (\ref{eq:mpug_czx_tuple}), so the $(0,1)$ and $(1,0)$ rows of the
    phase table are unchanged. In the $(1,1)$ sector the phase acquires the
    factor $(-1)^{r_\star(x)}$ from
    $\tilde\lambda R_\star\ket{x}=-i(-1)^{f(x)+r_\star(x)}\ket{\overline x}$; in
    the $(0,0)$ sector the adjoint formula
    (\ref{eq:mpug_monomial_adjoint}) evaluates that phase at $\overline x$ and
    conjugates it, which contributes $(-1)^{r_\star(\overline x)}$. Since
    $R_\star$ is diagonal, the permutation part is that of
    (\ref{eq:mpug_czx_monomial_action}). Identity
    (\ref{eq:mpug_czx_star_bar}) and the four rows of
    (\ref{eq:mpug_czx_star_phase_table}) are checked on the sixteen bit
    strings.
\end{proof}

\begin{lemma}[Involution and neighboring holonomy for the second tuple]
    \label{thm:mpug_czx_star_holonomy}
    \lean{MPOTensor.CZX.bondExponentStar,
        MPOTensor.CZX.starHolonomyLabel,
        MPOTensor.CZX.bondExponentStar_add_bondExponentStar_chainFlip,
        MPOTensor.CZX.bondExponentStar_holonomy,
        MPOTensor.CZX.bondExponentStar_chainFlip_of_disjoint,
        MPOTensor.CZX.starFiberLabel,
        MPOTensor.CZX.starHolonomyLabel_fiberEquiv,
        MPOTensor.CZX.fiberPhaseStar,
        MPOTensor.CZX.fiberPhaseStar_mul_fiberPhaseStar_flip,
        MPOTensor.CZX.fiberPhaseStar_flip_of_disjoint,
        MPOTensor.CZX.fiberPhaseStar_holonomy,
        MPOTensor.CZX.fiberPhaseStar_ne_zero}
    \leanok
    \uses{def:mpug_czx_chain_labels, def:mpug_czx_star_tuple}
    Let $q^\star_j(s)$ be the table (\ref{eq:mpug_czx_star_phase_table})
    evaluated on the bond $(j,j+1)$ as in
    Definition~\ref{def:mpug_czx_chain_labels}, and set
    \begin{align}
        c_j(s) & =b_{j+1}(s)\bigl(1+p_j(s)+p_{j+1}(s)\bigr)
        +b_{j+2}(s)\bigl(p_{j+1}(s)+p_{j+2}(s)\bigr).
        \label{eq:mpug_czx_star_holonomy_label}
    \end{align}
    For every $N\geq3$, every bond $j$, and every configuration $s$,
    \begin{align}
        q^\star_j(s)+q^\star_j(\sigma_js)
         & =0,
        \label{eq:mpug_czx_star_involution}                   \\
        q^\star_j(s)+q^\star_{j+1}(\sigma_js)
         & =q^\star_{j+1}(s)+q^\star_j(\sigma_{j+1}s)+2c_j(s)
        \label{eq:mpug_czx_star_holonomy}
    \end{align}
    modulo $4$, and $q^\star_j(\sigma_ks)=q^\star_j(s)$ whenever
    $k\notin\{j-1,j,j+1\}$. In the fiber coordinates
    (\ref{eq:mpug_czx_fiber_coordinates}) the label $c_j$ depends only on the
    orbit labels, $c_j(p,b)=b_{j+1}(1+p_j+p_{j+1})+b_{j+2}(p_{j+1}+p_{j+2})$,
    and the phases $\varphi^\star_j=i^{q^\star_j}$ are involutive and satisfy
    the neighboring square relation with factor $(-1)^{c_j(p,b)}$.
\end{lemma}

\begin{proof}\leanok
    \uses{thm:mpug_czx_star_local_monomial_action,
        thm:mpug_czx_orbit_classification}
    The exponent $q^\star_j$ depends only on the six bits at the sites $j$ and
    $j+1$, which gives the disjoint-bond statement. For
    (\ref{eq:mpug_czx_star_involution}) the flip $\sigma_j$ changes $x_0$,
    $x_1$, $a$, and $a'$, and the identity is verified on the $64$ values of
    those six bits. For (\ref{eq:mpug_czx_star_holonomy}) the sites $j$,
    $j+1$, $j+2$ are distinct because $N\geq3$, and the identity is verified on
    the $512$ values of the nine bits at those three sites. In fiber
    coordinates $p$ and $b$ are read off by
    Lemma~\ref{thm:mpug_czx_orbit_classification}, which turns $c_j$ into the
    displayed polynomial in the orbit labels alone.
\end{proof}

\begin{lemma}[Fixed spaces on the orbit fibers for the second tuple]
    \label{thm:mpug_czx_star_good_fibers}
    \lean{MPOTensor.CZX.isTrivialHolonomy_fiberPhaseStar_iff}
    \leanok
    \uses{def:mpug_trivial_holonomy, thm:mpug_czx_orbit_classification,
        thm:mpug_czx_star_holonomy}
    Let $N\geq3$. In the fiber coordinates
    (\ref{eq:mpug_czx_fiber_coordinates}), the phases $i^{q^\star_j}$ have
    trivial holonomy on the fiber over $(p,b)$ if and only if $c_j(p,b)=0$ for
    every $j\in\Z/N\Z$.
\end{lemma}

\begin{proof}\leanok
    \uses{thm:mpug_czx_star_holonomy, thm:mpug_czx_orbit_classification}
    If all $c_j(p,b)$ vanish, then for $k=j+1$ or $j=k+1$ the square identity is
    (\ref{eq:mpug_czx_star_holonomy}) with vanishing correction, for $k=j$ it is
    trivial, and otherwise the bonds are disjoint and both sides agree.
    Conversely, trivial holonomy at the square spanned by the bonds $j$ and
    $j+1$ at $\gamma=0$ gives $(-1)^{c_j(p,b)}=1$ by
    (\ref{eq:mpug_czx_star_holonomy}), since the phases are nonzero, hence
    $c_j(p,b)=0$.
\end{proof}

\begin{lemma}[The placed Gauss operators of the second tuple]
    \label{thm:mpug_czx_star_placed_operators}
    \lean{MPOTensor.CZX.placedGaussOperatorStar,
        MPOTensor.CZX.gaussProjector_circuitTupleStar,
        MPOTensor.CZX.placedGaussProjector_circuitTupleStar,
        MPOTensor.CZX.commonFixedSubmodule_placedGaussProjector_circuitTupleStar,
        MPOTensor.CZX.placedGaussOperatorStar_eq_monomial,
        MPOTensor.CZX.localPhaseStar_window,
        MPOTensor.CZX.placedGaussOperatorStar_eq_reindex}
    \leanok
    \uses{def:mpug_placed_gauss_projector, def:mpug_czx_star_tuple,
        def:mpug_czx_chain_labels, def:mpug_common_fixed_subspace}
    Let $N\geq2$ and let $G^\star_j$ be the local Gauss operator of
    Lemma~\ref{thm:mpug_czx_star_local_monomial_action} placed on the bond
    $(j,j+1)$. The placed Gauss projector of
    Definition~\ref{def:mpug_placed_gauss_projector} is
    $\widetilde P^\star_j=\tfrac12(\Id+G^\star_j)$, so the common fixed subspace
    $\mathcal V_N(R^\star)$ of the $\widetilde P^\star_j$ is the common fixed
    subspace of the $G^\star_j$. In the coordinates
    (\ref{eq:mpug_czx_configuration}), $G^\star_j$ is the monomial operator
    $G^\star_j\ket{s}=i^{q^\star_j(s)}\ket{\sigma_js}$, and in the fiber
    coordinates (\ref{eq:mpug_czx_fiber_coordinates}) it is the monomial
    operator with permutation $\gamma\mapsto\gamma+e_j$ and phase
    $i^{q^\star_j}$.
\end{lemma}

\begin{proof}\leanok
    \uses{def:mpug_gauss_projector, thm:mpug_czx_star_local_monomial_action,
        thm:mpug_embed_monomial, thm:mpug_czx_orbit_classification,
        thm:mpug_czx_placed_operators}
    The average over $\Z_2=\{e,g\}$ is $\frac12(\mathcal G_e+\mathcal G_g)$ with
    $\mathcal G_e=\Id$, and relabelling and placement are linear and unital. A
    vector $v$ satisfies $\frac12(v+G^\star_jv)=v$ exactly when
    $G^\star_jv=v$. Lemma~\ref{thm:mpug_czx_star_local_monomial_action} and
    Lemma~\ref{thm:mpug_embed_monomial} express $G^\star_j$ as a monomial
    operator; the permutation part is that of
    Lemma~\ref{thm:mpug_czx_placed_operators}, and decoding the two-site
    window at $j$ identifies the phase with $i^{q^\star_j(s)}$. Finally
    Lemma~\ref{thm:mpug_czx_orbit_classification} reindexes both to the
    fiber coordinates.
\end{proof}

\begin{lemma}[Reduction of the dimension to a finite count]
    \label{thm:mpug_czx_star_reduction}
    \lean{MPOTensor.CZX.finrank_commonFixedSubmodule_placedGaussProjector_circuitTupleStar}
    \leanok
    \uses{def:mpug_placed_gauss_projector, def:mpug_czx_star_tuple,
        def:mpug_common_fixed_subspace}
    On a periodic chain of $N\geq3$ blocked sites,
    \begin{align}
        \dim\mathcal V_N(R^\star)
         & =\#\bigl\{(p,b)\in\mathbb F_2^N\times\mathbb F_2^N:
        c_j(p,b)=0\text{ for all }j\in\Z/N\Z\bigr\}.
        \label{eq:mpug_czx_star_count}
    \end{align}
\end{lemma}

\begin{proof}\leanok
    \uses{thm:mpug_czx_star_placed_operators,
        thm:mpug_monomial_fixed_space_criterion,
        thm:mpug_czx_star_holonomy, thm:mpug_czx_star_good_fibers}
    By Lemma~\ref{thm:mpug_czx_star_placed_operators},
    $\mathcal V_N(R^\star)$ is the common fixed subspace of monomial operators
    which, in fiber coordinates, have permutations $\gamma\mapsto\gamma+e_j$ on
    the fibers over $(p,b)$, and by (\ref{eq:mpug_czx_star_involution}) these
    operators are involutions. The monomial fixed-space criterion
    (\ref{eq:mpug_monomial_fixed_space_count}) counts the fibers with trivial
    holonomy, which by Lemma~\ref{thm:mpug_czx_star_good_fibers} are exactly
    the fibers whose labels satisfy the $N$ equations $c_j(p,b)=0$.
\end{proof}

\begin{lemma}[Dimensions for the second displayed tuple at three and four sites]
    \label{thm:mpug_czx_star_dimension}
    \lean{MPOTensor.CZX.finrank_commonFixedSubmodule_placedGaussProjector_circuitTupleStar_three,
        MPOTensor.CZX.finrank_commonFixedSubmodule_placedGaussProjector_circuitTupleStar_four}
    \leanok
    \uses{def:mpug_placed_gauss_projector, def:mpug_czx_star_tuple,
        def:mpug_common_fixed_subspace}
    On periodic chains of three and four blocked sites,
    \begin{align}
        \dim\mathcal V_3(R^\star) & =14,
        \label{eq:mpug_czx_star_dim_three} \\
        \dim\mathcal V_4(R^\star) & =36.
        \label{eq:mpug_czx_star_dim_four}
    \end{align}
    These two values are certificates at two system sizes. They carry no
    asymptotic content, and no growth law for $\dim\mathcal V_N(R^\star)$
    follows from them. Since $R^\star$ extends the four prescribed maps of
    Definition~\ref{def:mpug_czx_defect_maps} by
    Lemma~\ref{thm:mpug_czx_star_tuple_completion}, these are two values of
    the dimension on that completion class, recorded as lower bounds on its
    supremum in Lemma~\ref{thm:mpug_czx_star_class_values}; they give no
    value for the minimum or the maximum of the dimension over completions.
    Compared with
    (\ref{eq:mpug_czx_dimension}), which gives $8$ and $16$ at the same two
    sizes, they exhibit dependence of the dimension on the displayed circuit
    data at those two sizes.
\end{lemma}

\begin{proof}\leanok
    \uses{thm:mpug_czx_star_reduction}
    Evaluating the right-hand side of (\ref{eq:mpug_czx_star_count}) by
    exhaustive enumeration gives the two dimensions: at $N=3$ the $64$
    pairs $(p,b)$ yield $14$ solutions of the three equations
    $c_j(p,b)=0$, and at $N=4$ the $256$ pairs yield $36$ solutions of
    the four equations.
\end{proof}

\section{The CZX defect domains and prescribed defect maps}
\label{sec:mpug_czx_defect_maps}

The modified-fusion lemma of
\cite[lines~4215--4254]{FrancoRubio2025MPUGauging} attaches to each pair
$(a,b)$ of gauge labels and each injective block $x$ a two-site defect vector
$v(a,b,x)$ mapped unitarily to $v(e,ab,x)$. For the CZX model, the defect
subspaces and prescribed maps are evaluated by finite contraction in
\cite{gap:fbc25_czx_defect_domains}.

\begin{definition}[Single-site CZX defects]
    \label{def:mpug_czx_site_defects}
    \lean{MPOTensor.CZX.identityDefect,
        MPOTensor.CZX.flipDefect,
        MPOTensor.CZX.siteDefect,
        MPOTensor.CZX.identityDefect_zero,
        MPOTensor.CZX.identityDefect_one,
        MPOTensor.CZX.flipDefect_zero,
        MPOTensor.CZX.flipDefect_one,
        MPOTensor.CZX.siteDefect_one,
        MPOTensor.CZX.siteDefect_gen}
    \leanok
    \discussion{7739}
    One blocked site of the CZX model carries two matter qubits, and the two
    injective blocks of the once-blocked GHZ state are the product vectors
    $\ket{\circ}=\ket{00}$ and $\ket{\bullet}=\ket{11}$
    \cite[lines~4660--4670]{FrancoRubio2025MPUGauging}, both of bond dimension
    one. Writing $e\leftrightarrow0$ and $g\leftrightarrow1$ for the labels of
    $\Z_2=\{e,g\}$, the identity defect and the domain-wall defect of
    \cite[lines~2831--2845]{FrancoRubio2025MPUGauging} at the block
    $x\in\mathbb F_2$ are
    \begin{align}
        e[x] & =\ket{xx},   \label{eq:mpug_czx_identity_defect}  \\
        g[0] & =\ket{00},   \label{eq:mpug_czx_flip_defect_zero} \\
        g[1] & =-i\ket{11}. \label{eq:mpug_czx_flip_defect_one}
    \end{align}
    As carried out in~\cite{gap:fbc25_czx_defect_domains} and displayed
    in \cite[lines~4800--4859]{FrancoRubio2025MPUGauging}, contracting the
    action tensors $A^{\Gamma}_{g,0}=-i\bra{1}$ and
    $A^{\Gamma}_{g,1}=\bra{0}$ of
    \cite[lines~4690--4694]{FrancoRubio2025MPUGauging} against the defect
    diagram gives (\ref{eq:mpug_czx_flip_defect_zero}) and
    (\ref{eq:mpug_czx_flip_defect_one}).
\end{definition}

\begin{definition}[Two-site CZX defect vectors]
    \label{def:mpug_czx_defect_vectors}
    \lean{MPOTensor.CZX.defectVector,
        MPOTensor.CZX.defectVector_one_one,
        MPOTensor.CZX.defectVector_one_gen,
        MPOTensor.CZX.defectVector_gen_one,
        MPOTensor.CZX.defectVector_gen_gen}
    \leanok
    \uses{def:mpug_czx_site_defects}
    The two-site defect vector of the label pair $(a,b)$ and the block $x$ is
    the left defect $a[bx]$ tensored with the right defect $b[x]$, in that
    order \cite[lines~4234--4245]{FrancoRubio2025MPUGauging},
    \begin{align}
        v(a,b,x) & =a[bx]\otimes b[x].
        \label{eq:mpug_czx_defect_vector}
    \end{align}
    The four qubits of two neighboring blocked sites are numbered $0,1$ on the
    left site and $2,3$ on the right, and the computational basis of the
    sixteen-dimensional two-site matter space is written $\ket{x_0x_1x_2x_3}$.
\end{definition}

\begin{lemma}[The eight CZX defect vectors]
    \label{thm:mpug_czx_defect_vector_values}
    \lean{MPOTensor.CZX.defectVector_one_one_zero,
        MPOTensor.CZX.defectVector_one_one_one,
        MPOTensor.CZX.defectVector_one_gen_zero,
        MPOTensor.CZX.defectVector_one_gen_one,
        MPOTensor.CZX.defectVector_gen_one_zero,
        MPOTensor.CZX.defectVector_gen_one_one,
        MPOTensor.CZX.defectVector_gen_gen_zero,
        MPOTensor.CZX.defectVector_gen_gen_one,
        MPOTensor.CZX.star_defectVector_dotProduct_defectVector}
    \leanok
    \uses{def:mpug_czx_defect_vectors}
    The eight two-site defect vectors of the CZX model are
    \begin{align}
        v(e,e,0) & =\ket{0000},   & v(e,e,1) & =\ket{1111},
        \label{eq:mpug_czx_v_ee}                              \\
        v(e,g,0) & =\ket{1100},   & v(e,g,1) & =-i\ket{0011},
        \label{eq:mpug_czx_v_eg}                              \\
        v(g,e,0) & =\ket{0000},   & v(g,e,1) & =-i\ket{1111},
        \label{eq:mpug_czx_v_ge}                              \\
        v(g,g,0) & =-i\ket{1100}, & v(g,g,1) & =-i\ket{0011}.
        \label{eq:mpug_czx_v_gg}
    \end{align}
    For every label pair the two vectors $v(a,b,0)$ and $v(a,b,1)$ are
    orthonormal.
\end{lemma}

\begin{proof}\leanok
    \uses{def:mpug_czx_site_defects, def:mpug_czx_defect_vectors}
    Substituting (\ref{eq:mpug_czx_identity_defect}),
    (\ref{eq:mpug_czx_flip_defect_zero}), and
    (\ref{eq:mpug_czx_flip_defect_one}) into
    (\ref{eq:mpug_czx_defect_vector}), with $bx=x+b$ in $\mathbb F_2$, gives
    the eight values. The two vectors of a label pair differ in the last two
    bits, hence are orthogonal, and each carries a scalar of modulus one,
    hence is a unit vector.
\end{proof}

\begin{lemma}[Reindexed monomial operators on the computational basis]
    \label{lem:mpug_czx_matter_monomial_action}
    \lean{MPOTensor.CZX.matterKet,
        MPOTensor.CZX.matterMatrix_monomial_mulVec_matterKet}
    \leanok
    \uses{def:mpug_monomial_matrix, def:mpug_czx_circuit_tuple,
        def:mpug_czx_defect_vectors}
    Write $\ket{x}$ for the computational basis vector of the two-site matter
    space attached to $x\in\mathbb F_2^4$ by the qubit numbering of
    Definition~\ref{def:mpug_czx_defect_vectors}. After reindexing to that
    space a monomial four-qubit operator still acts on the computational basis
    by its phase and its permutation, so that for every $x\in\mathbb F_2^4$
    \begin{align}
        M(\sigma,\varphi)\ket{x} & =\varphi(x)\ket{\sigma(x)}.
        \label{eq:mpug_czx_matter_monomial_action}
    \end{align}
\end{lemma}

\begin{proof}\leanok
    \uses{thm:mpug_monomial_matrix_laws, def:mpug_czx_circuit_tuple}
    Reindexing is the simultaneous relabelling of rows and columns along the
    qubit numbering. It carries $M(\sigma,\varphi)$ to the monomial matrix of
    the relabelled permutation and phase, and carries the basis vector of $x$
    to $\ket{x}$, so reading (\ref{eq:mpug_monomial_matrix}) in the new labels
    gives (\ref{eq:mpug_czx_matter_monomial_action}).
\end{proof}

\begin{lemma}[Unmodified CZX fusion contractions]
    \label{thm:mpug_czx_unmodified_fusion}
    \lean{MPOTensor.CZX.matterMatrix_lambda_mulVec_defectVector_zero,
        MPOTensor.CZX.matterMatrix_lambda_mulVec_defectVector_one}
    \leanok
    \discussion{7355}
    \uses{def:mpug_czx_circuit_tuple, def:mpug_czx_defect_vectors}
    Let $\Z_2=\{e,g\}$, and let $v(a,b,x)$ be the two-site defect vectors
    of Definition~\ref{def:mpug_czx_defect_vectors}, where $x\in\{0,1\}$.
    As in the two contractions displayed in
    \cite[lines~4989--5067]{FrancoRubio2025MPUGauging}, on the four matter
    qubits numbered $0,1,2,3$, the unmodified operator (as distinct from
    $\tilde\lambda=\lambda Z_0$ used in the displayed circuit tuple)
    \begin{align}
        \lambda & =\lambda^R_{g,g}
        =-i\,(X_0X_1X_2Z_3) \mathrm{CZ}_{01}\mathrm{CZ}_{12}\,X_2
    \end{align}
    satisfies
    \begin{align}
        \lambda\,v(g,g,0) & =-v(e,e,0), \\
        \lambda\,v(g,g,1) & =v(e,e,1).
    \end{align}
\end{lemma}

\begin{proof}\leanok
    \uses{thm:mpug_czx_phase_tables, thm:mpug_czx_defect_vector_values,
        lem:mpug_czx_matter_monomial_action}
    The phase formula~(\ref{eq:mpug_czx_lambda_action}), read on the two-site
    matter space through (\ref{eq:mpug_czx_matter_monomial_action}), gives
    \begin{align}
        \lambda\ket{1100} & =-i\ket{0000},
        \label{eq:mpug_czx_lambda_ket_even} \\
        \lambda\ket{0011} & =i\ket{1111},
        \label{eq:mpug_czx_lambda_ket_odd}
    \end{align}
    since $h(1100)=0$ and $h(0011)=1$.
    By Lemma~\ref{thm:mpug_czx_defect_vector_values}, the input defect
    vectors are $-i\ket{1100}$ and $-i\ket{0011}$, while the target
    vectors are $\ket{0000}$ and $\ket{1111}$. Multiplying
    (\ref{eq:mpug_czx_lambda_ket_even}) and
    (\ref{eq:mpug_czx_lambda_ket_odd}) by $-i$ gives the claimed signs.
\end{proof}

\begin{definition}[CZX defect domains]
    \label{def:mpug_czx_defect_domains}
    \lean{MPOTensor.CZX.defectDomain,
        MPOTensor.CZX.evenTarget,
        MPOTensor.CZX.oddTarget}
    \leanok
    \uses{def:mpug_czx_defect_vectors}
    The defect domain of the label pair $(a,b)$ and the two target planes are
    \begin{align}
        \mathcal K_{a,b} & =\spn_{\C}\{v(a,b,x):x\in\mathbb F_2\},
        \label{eq:mpug_czx_domain}                                 \\
        \mathcal L_0     & =\spn_{\C}\{\ket{0000},\ket{1111}\},
        \label{eq:mpug_czx_even_target}                            \\
        \mathcal L_1     & =\spn_{\C}\{\ket{1100},\ket{0011}\}.
        \label{eq:mpug_czx_odd_target}
    \end{align}
\end{definition}

\begin{lemma}[The four CZX defect domains]
    \label{thm:mpug_czx_defect_domains}
    \lean{MPOTensor.CZX.defectDomain_one_one,
        MPOTensor.CZX.defectDomain_gen_one,
        MPOTensor.CZX.defectDomain_one_gen,
        MPOTensor.CZX.defectDomain_gen_gen,
        MPOTensor.CZX.defectDomain_eq,
        MPOTensor.CZX.star_dotProduct_eq_zero_of_mem_evenTarget_of_mem_oddTarget,
        MPOTensor.CZX.finrank_evenTarget,
        MPOTensor.CZX.finrank_oddTarget,
        MPOTensor.CZX.finrank_defectDomain}
    \leanok
    \uses{def:mpug_czx_defect_domains}
    The four CZX defect domains are
    \begin{align}
        \mathcal K_{e,e}=\mathcal K_{g,e} & =\mathcal L_0,
        \label{eq:mpug_czx_even_domains}                   \\
        \mathcal K_{e,g}=\mathcal K_{g,g} & =\mathcal L_1,
        \label{eq:mpug_czx_odd_domains}
    \end{align}
    so the domain depends only on the right label. The two planes are
    orthogonal, $\mathcal L_0\perp\mathcal L_1$, and each has dimension two.
\end{lemma}

\begin{proof}\leanok
    \uses{thm:mpug_czx_defect_vector_values, def:mpug_czx_defect_domains}
    By (\ref{eq:mpug_czx_v_ee})--(\ref{eq:mpug_czx_v_gg}) the two spanning
    vectors of $\mathcal K_{a,b}$ are the two spanning vectors of
    $\mathcal L_b$ up to nonzero scalars, and rescaling a spanning vector by a
    nonzero scalar leaves the span unchanged, which gives
    (\ref{eq:mpug_czx_even_domains}) and (\ref{eq:mpug_czx_odd_domains}). The
    four basis vectors occurring in (\ref{eq:mpug_czx_even_target}) and
    (\ref{eq:mpug_czx_odd_target}) are pairwise distinct members of the
    computational basis, so expanding a vector of $\mathcal L_0$ and a vector
    of $\mathcal L_1$ in that basis makes every one of the four cross terms of
    their inner product vanish, and the two spanning vectors of each plane are
    linearly independent.
\end{proof}

\begin{definition}[The prescribed CZX defect maps]
    \label{def:mpug_czx_prescribed_maps}
    \lean{MPOTensor.CZX.prescribedMap}
    \leanok
    \uses{def:mpug_czx_defect_vectors, def:mpug_czx_defect_domains}
    The prescribed defect map of the label pair $(a,b)$ is the partial
    isometry
    \begin{align}
        D_{a,b} & =\sum_{x\in\mathbb F_2}\ketbra{v(e,ab,x)}{v(a,b,x)}
        \label{eq:mpug_czx_prescribed_map}
    \end{align}
    of the two-site matter space. It is formed from the defect vectors
    (\ref{eq:mpug_czx_v_ee})--(\ref{eq:mpug_czx_v_gg}) and not from the
    circuits $w$ and $\tilde\lambda$ of
    Definition~\ref{def:mpug_czx_circuit_tuple}.
\end{definition}

\begin{lemma}[Action of the prescribed CZX defect maps]
    \label{thm:mpug_czx_prescribed_action}
    \lean{MPOTensor.CZX.prescribedMap_mulVec_defectVector,
        MPOTensor.CZX.prescribedMap_mem_defectDomain,
        MPOTensor.CZX.isometry_on_defectDomain,
        MPOTensor.CZX.prescribedMap_one_one_mulVec,
        MPOTensor.CZX.prescribedMap_one_gen_mulVec,
        MPOTensor.CZX.prescribedMap_gen_one_mulVec_zero,
        MPOTensor.CZX.prescribedMap_gen_one_mulVec_one,
        MPOTensor.CZX.prescribedMap_gen_gen_mulVec_zero,
        MPOTensor.CZX.prescribedMap_gen_gen_mulVec_one}
    \leanok
    \uses{def:mpug_czx_prescribed_maps}
    For every label pair and every block,
    \begin{align}
        D_{a,b}v(a,b,x) & =v(e,ab,x),
        \label{eq:mpug_czx_prescription}
    \end{align}
    the map $D_{a,b}$ carries $\mathcal K_{a,b}$ into $\mathcal K_{e,ab}$, and
    it is an isometry on $\mathcal K_{a,b}$, that is,
    $D_{a,b}^{\dagger}D_{a,b}\ket{\xi}=\ket{\xi}$ for every
    $\ket{\xi}\in\mathcal K_{a,b}$. In the computational basis,
    \begin{align}
        D_{e,e} & =\Id_{\mathcal L_0},
        \label{eq:mpug_czx_map_ee}                     \\
        D_{e,g} & =\Id_{\mathcal L_1},
        \label{eq:mpug_czx_map_eg}                     \\
        D_{g,e}\ket{0000}=\ket{1100},
                & \quad D_{g,e}\ket{1111}=\ket{0011},
        \label{eq:mpug_czx_map_ge}                     \\
        D_{g,g}\ket{1100}=i\ket{0000},
                & \quad D_{g,g}\ket{0011}=i\ket{1111}.
        \label{eq:mpug_czx_map_gg}
    \end{align}
\end{lemma}

\begin{proof}\leanok
    \uses{thm:mpug_czx_defect_vector_values, def:mpug_czx_prescribed_maps}
    Applying (\ref{eq:mpug_czx_prescribed_map}) to $v(a,b,x)$ and using the
    orthonormality of $v(a,b,0)$ and $v(a,b,1)$ of
    Lemma~\ref{thm:mpug_czx_defect_vector_values} leaves the single term
    $v(e,ab,x)$, which is (\ref{eq:mpug_czx_prescription}). The same
    computation for the adjoint gives
    $D_{a,b}^{\dagger}v(e,ab,x)=v(a,b,x)$, and composing the two identities on
    a linear combination of $v(a,b,0)$ and $v(a,b,1)$ gives the isometry
    property and the inclusion of the images. Substituting
    (\ref{eq:mpug_czx_v_ee})--(\ref{eq:mpug_czx_v_gg}) into
    (\ref{eq:mpug_czx_prescription}) and cancelling the scalars of modulus one
    gives (\ref{eq:mpug_czx_map_ee})--(\ref{eq:mpug_czx_map_gg}).
\end{proof}

\begin{definition}[Prescribed defect maps of the CZX model]
    \label{def:mpug_czx_defect_maps}
    \lean{MPOTensor.CZX.defectMaps,
        MPOTensor.CZX.defectMaps_domain,
        MPOTensor.CZX.defectMaps_prescribed}
    \leanok
    \uses{def:mpug_defect_maps, def:mpug_czx_defect_domains,
        def:mpug_czx_prescribed_maps, thm:mpug_czx_prescribed_action}
    The four domains (\ref{eq:mpug_czx_even_domains}) and
    (\ref{eq:mpug_czx_odd_domains}) together with the four maps
    (\ref{eq:mpug_czx_prescribed_map}) are prescribed defect maps in the sense
    of Definition~\ref{def:mpug_defect_maps} for the group $\Z_2$ and the
    two-site matter space of the CZX model, the isometry condition
    (\ref{eq:mpug_defect_isometry}) being
    Lemma~\ref{thm:mpug_czx_prescribed_action}. That the circuit tuple
    $R^{\mathrm{CZX}}$ lies in the resulting completion class
    (\ref{eq:mpug_completion_class}) is proved unconditionally in
    Lemma~\ref{thm:mpug_czx_tuple_completion} below.
\end{definition}

\section{The displayed CZX tuple as a completion}
\label{sec:mpug_czx_completion}

The four defect domains and four prescribed maps of
Section~\ref{sec:mpug_czx_defect_maps} constrain each fusion operator on one
plane of the two-site matter space. The operators
$(\Id,\Id,w,\tilde\lambda)$ of Definition~\ref{def:mpug_czx_circuit_tuple}
satisfy these constraints, so the displayed tuple belongs to the completion
class of the CZX prescribed defect maps $D^{\mathrm{CZX}}$, and
(\ref{eq:mpug_czx_dimension}) lower-bounds the supremum of the gauge-invariant
dimension over the full class.

\begin{lemma}[The displayed CZX tuple is a completion]
    \label{thm:mpug_czx_tuple_completion}
    \lean{MPOTensor.CZX.circuitTuple_one_mulVec_defectVector,
        MPOTensor.CZX.circuitTuple_gen_one_mulVec_defectVector_zero,
        MPOTensor.CZX.circuitTuple_gen_one_mulVec_defectVector_one,
        MPOTensor.CZX.circuitTuple_gen_gen_mulVec_defectVector_zero,
        MPOTensor.CZX.circuitTuple_gen_gen_mulVec_defectVector_one,
        MPOTensor.CZX.circuitTuple_mulVec_defectVector,
        MPOTensor.CZX.isCompletion_circuitTuple,
        MPOTensor.CZX.circuitTuple_mem_completionClass}
    \leanok
    \uses{def:mpug_czx_defect_maps, def:mpug_czx_circuit_tuple,
        def:mpug_completion_class}
    For every ordered pair $(a,b)$ of labels of $\Z_2$ and every injective
    block $x$, the displayed operator of the pair carries the defect vector of
    the pair to the defect vector of the product label,
    \begin{align}
        U^{\mathrm{CZX}}_{a,b}v(a,b,x) & =v(e,ab,x),
        \label{eq:mpug_czx_tuple_fusion}
    \end{align}
    that is,
    \begin{align}
        \Id\,v(e,b,x)           & =v(e,b,x),
        \label{eq:mpug_czx_tuple_fusion_trivial}  \\
        w\,v(g,e,x)             & =v(e,g,x),
        \label{eq:mpug_czx_tuple_fusion_movement} \\
        \tilde\lambda\,v(g,g,x) & =v(e,e,x).
        \label{eq:mpug_czx_tuple_fusion_modified}
    \end{align}
    Consequently the displayed operator of every pair agrees on
    $\mathcal K_{a,b}$ with the prescribed map $D_{a,b}$, and the displayed
    tuple lies in the full completion class,
    \begin{align}
        R^{\mathrm{CZX}} & \in\mathfrak C(D^{\mathrm{CZX}}).
        \label{eq:mpug_czx_tuple_completion}
    \end{align}
\end{lemma}

\begin{proof}\leanok
    \uses{thm:mpug_czx_phase_tables, thm:mpug_czx_defect_vector_values,
        thm:mpug_czx_defect_domains, thm:mpug_czx_prescribed_action}
    The phase tables (\ref{eq:mpug_czx_w_action}) and
    (\ref{eq:mpug_czx_tilde_lambda_action}), evaluated at the four
    computational basis states carrying the defect vectors of
    Lemma~\ref{thm:mpug_czx_defect_vector_values}, give
    \begin{align}
        w\ket{0000}             & =\ket{1100},
                                & w\ket{1111}
                                & =\ket{0011},
        \label{eq:mpug_czx_w_basis_action}                \\
        \tilde\lambda\ket{1100} & =i\ket{0000},
                                & \tilde\lambda\ket{0011}
                                & =i\ket{1111},
        \label{eq:mpug_czx_tilde_lambda_basis_action}
    \end{align}
    because $e(0000)=e(1111)=0$ and $f(1100)=f(0011)=1$. Substituting
    (\ref{eq:mpug_czx_v_ee})--(\ref{eq:mpug_czx_v_gg}) and cancelling the
    factors $-i$ gives (\ref{eq:mpug_czx_tuple_fusion_movement}) and
    (\ref{eq:mpug_czx_tuple_fusion_modified}); the pairs with left label $e$
    carry the identity operator, so
    (\ref{eq:mpug_czx_tuple_fusion_trivial}) is immediate. The two defect
    vectors of a pair span $\mathcal K_{a,b}$ by (\ref{eq:mpug_czx_domain}),
    and the prescribed map obeys the same identity
    (\ref{eq:mpug_czx_prescription}), so the displayed operator and the
    prescribed map agree on the whole plane, which is
    (\ref{eq:mpug_czx_tuple_completion}).
\end{proof}

\begin{lemma}[Lower bound on the supremum over the CZX completion class]
    \label{thm:mpug_czx_class_maximum}
    \lean{MPOTensor.CZX.bddAbove_range_finrank_gaugeInvariantSubspace,
        MPOTensor.CZX.two_pow_le_iSup_finrank_gaugeInvariantSubspace_completionClass}
    \leanok
    \uses{def:mpug_czx_defect_maps, def:mpug_czx_circuit_tuple,
        def:mpug_completion_class, def:mpug_gauge_invariant_subspace,
        thm:mpug_czx_dimension_of_completion}
    On a periodic chain of $N\geq3$ blocked sites the gauge-invariant
    dimensions of the members of the CZX completion class are bounded above by
    the dimension of the chain space, so their supremum is finite, and
    \begin{align}
        \sup_{R\in\mathfrak C(D^{\mathrm{CZX}})}\dim\mathcal V_N(R) & \geq2^N.
        \label{eq:mpug_czx_class_maximum}
    \end{align}
    Nothing is asserted about the minimum of $\dim\mathcal V_N(R)$ over the
    class, about the exact value of the supremum, or about which completion
    attains it. With the domains and maps entering the class derived
    in~\cite{gap:fbc25_czx_defect_domains}, those extrema are questions about
    this completion class, not questions the source raises;
    \cite[lines~5198--5204]{FrancoRubio2025MPUGauging} asks only whether the
    gauge-invariant subspace of the state-level gauging stays bounded or grows
    with the system size, and nothing is asserted about that growth either.
\end{lemma}

\begin{proof}\leanok
    \uses{thm:mpug_czx_tuple_completion, thm:mpug_czx_dimension_of_completion}
    Every gauge-invariant subspace (\ref{eq:mpug_gauge_invariant_subspace}) is
    a subspace of the chain space, so the dimensions of the members of the
    class are bounded above by the dimension of that space and their supremum
    is finite. By Lemma~\ref{thm:mpug_czx_tuple_completion} the displayed
    tuple lies in the class, and by
    Lemma~\ref{thm:mpug_czx_dimension_of_completion} its gauge-invariant
    subspace has dimension $2^N$, which is therefore at most the supremum.
\end{proof}

\begin{lemma}[The second displayed CZX tuple is a completion]
    \label{thm:mpug_czx_star_tuple_completion}
    \lean{MPOTensor.CZX.matterMatrix_tildeLambdaStar_mulVec_matterKet_zero,
        MPOTensor.CZX.matterMatrix_tildeLambdaStar_mulVec_matterKet_one,
        MPOTensor.CZX.circuitTupleStar_gen_gen_mulVec_defectVector_zero,
        MPOTensor.CZX.circuitTupleStar_gen_gen_mulVec_defectVector_one,
        MPOTensor.CZX.circuitTupleStar_mulVec_defectVector,
        MPOTensor.CZX.isCompletion_circuitTupleStar,
        MPOTensor.CZX.circuitTupleStar_mem_completionClass}
    \leanok
    \uses{def:mpug_czx_defect_maps, def:mpug_czx_star_tuple,
        def:mpug_completion_class}
    The second displayed tuple obeys the same prescription
    (\ref{eq:mpug_czx_tuple_fusion}) as the first,
    \begin{align}
        U^\star_{a,b}v(a,b,x) & =v(e,ab,x),
        \label{eq:mpug_czx_star_tuple_fusion}
    \end{align}
    and therefore lies in the full completion class,
    \begin{align}
        R^\star & \in\mathfrak C(D^{\mathrm{CZX}}).
        \label{eq:mpug_czx_star_tuple_completion}
    \end{align}
\end{lemma}

\begin{proof}\leanok
    \uses{thm:mpug_czx_tuple_completion, thm:mpug_czx_defect_domains,
        def:mpug_czx_star_tuple}
    The pairs $(e,e)$, $(e,g)$, and $(g,e)$ carry the same operators in
    $R^\star$ as in $R^{\mathrm{CZX}}$, so
    Lemma~\ref{thm:mpug_czx_tuple_completion} gives
    (\ref{eq:mpug_czx_star_tuple_fusion}) there. For the pair $(g,g)$ the
    operator is $\tilde\lambda R_\star$, and the defect domain is
    $\mathcal K_{g,g}=\mathcal L_1=\spn_{\C}\{\ket{1100},\ket{0011}\}$ by
    (\ref{eq:mpug_czx_odd_domains}). Since
    $r_\star(1100)=r_\star(0011)=0$, the factor $R_\star$ fixes both spanning
    vectors, so $\tilde\lambda R_\star$ has the same values there as
    $\tilde\lambda$, namely $i\ket{0000}$ and $i\ket{1111}$. The two defect
    vectors of the pair span its domain, which gives
    (\ref{eq:mpug_czx_star_tuple_completion}).
\end{proof}

\begin{lemma}[Two further values of the dimension on the CZX completion
        class]
    \label{thm:mpug_czx_star_class_values}
    \lean{MPOTensor.CZX.finrank_gaugeInvariantSubspace_circuitTupleStar_three,
        MPOTensor.CZX.finrank_gaugeInvariantSubspace_circuitTupleStar_four,
        MPOTensor.CZX.fourteen_le_iSup_finrank_gaugeInvariantSubspace_completionClass,
        MPOTensor.CZX.thirtySix_le_iSup_finrank_gaugeInvariantSubspace_completionClass}
    \leanok
    \uses{def:mpug_czx_defect_maps, def:mpug_completion_class,
        def:mpug_gauge_invariant_subspace, thm:mpug_czx_star_tuple_completion,
        thm:mpug_czx_star_dimension}
    The dimensions $14$ at three sites and $36$ at four sites are values of the
    gauge-invariant dimension on the CZX completion class, so
    \begin{align}
        \sup_{R\in\mathfrak C(D^{\mathrm{CZX}})}\dim\mathcal V_3(R) & \geq14,
        \label{eq:mpug_czx_star_class_three}                                  \\
        \sup_{R\in\mathfrak C(D^{\mathrm{CZX}})}\dim\mathcal V_4(R) & \geq36.
        \label{eq:mpug_czx_star_class_four}
    \end{align}
    Beside the values $8$ and $16$ of the first displayed tuple at the same two
    chain lengths, these exhibit two different values of the dimension on the
    class at each of those lengths. Nothing follows about the minimum, the
    exact supremum, which completion attains it, or the behaviour of any of
    these as $N$ grows.
\end{lemma}

\begin{proof}\leanok
    \uses{thm:mpug_czx_star_tuple_completion, thm:mpug_czx_star_dimension}
    By Lemma~\ref{thm:mpug_czx_star_tuple_completion} the second displayed
    tuple lies in the class, and by
    Lemma~\ref{thm:mpug_czx_star_dimension} its gauge-invariant subspace has
    dimensions $14$ and $36$ at three and four sites. Each value is therefore
    at most the supremum over the class at that chain length, by the same
    boundedness argument as in
    Lemma~\ref{thm:mpug_czx_class_maximum}.
\end{proof}
